\documentclass[11pt,letterpaper,reqno]{amsart}
\makeatletter
\NeedsTeXFormat{LaTeX2e}[2020/10/01]

\RequirePackage[T1]{fontenc}
\RequirePackage{lmodern}
\RequirePackage{microtype}
\RequirePackage[letterpaper,margin=1in,headheight=15pt]{geometry}
\RequirePackage{amsmath,amsthm,amssymb,mathtools}
\RequirePackage{mathrsfs}
\RequirePackage{enumitem}
\RequirePackage{longtable}
\RequirePackage{perpage}
\RequirePackage{xcolor}
\definecolor{manuscriptred}{rgb}{0.6,0.1,0.1}
\definecolor{manuscriptblue}{rgb}{0.1,0.3,0.5}
\definecolor{manuscriptmagenta}{HTML}{60445C}
\RequirePackage[breaklinks=true,colorlinks=true,
  linkcolor=manuscriptred,citecolor=manuscriptblue,urlcolor=manuscriptmagenta,
  filecolor=manuscriptmagenta]{hyperref}
\RequirePackage{cleveref}
\hypersetup{pdfauthor={}}

\setcounter{tocdepth}{1}
\renewcommand{\thefigure}{\Roman{figure}}
\renewcommand{\thetable}{\Roman{table}}
\MakePerPage{footnote}
\renewcommand{\thefootnote}{%
  \ifnum\value{footnote}>9\arabic{footnote}\else\fnsymbol{footnote}\fi}

\theoremstyle{definition}
\newtheorem{ms@theorem}{Theorem}
\renewcommand{\thems@theorem}{\Alph{ms@theorem}}
\newtheorem*{ms@theoremstar}{Theorem}
\newtheorem*{ms@theoremzero}{Theorem 0}
\newtheorem{ms@proposition}{Proposition}[section]
\newtheorem{ms@lemma}[ms@proposition]{Lemma}
\newtheorem{ms@definition}[ms@proposition]{Definition}
\newtheorem{ms@corollary}[ms@proposition]{Corollary}
\newtheorem*{ms@corollarystar}{Corollary}

\newcommand{\ms@majorstart}{%
  \pushQED{\qed}\renewcommand{\qedsymbol}{$\blacksquare\square$}}
\newcommand{\ms@minorstart}{%
  \pushQED{\qed}\renewcommand{\qedsymbol}{$\blacksquare$}}
\newenvironment{theorem}[1][]{\ms@majorstart\ms@theorem[#1]}
  {\popQED\endms@theorem}
\newenvironment{theorem*}[1][]{\ms@majorstart\ms@theoremstar[#1]}
  {\popQED\endms@theoremstar}
\newenvironment{theoremzero}{\ms@majorstart\ms@theoremzero}
  {\popQED\endms@theoremzero}
\newenvironment{proposition}[1][]{\ms@minorstart\ms@proposition[#1]}
  {\popQED\endms@proposition}
\newenvironment{lemma}[1][]{\ms@minorstart\ms@lemma[#1]}
  {\popQED\endms@lemma}
\newenvironment{definition}[1][]{\ms@minorstart\ms@definition[#1]}
  {\popQED\endms@definition}
\newenvironment{corollary}[1][]{\ms@majorstart\ms@corollary[#1]}
  {\popQED\endms@corollary}
\newenvironment{corollary*}[1][]{\ms@majorstart\ms@corollarystar[#1]}
  {\popQED\endms@corollarystar}

\theoremstyle{remark}
\newtheorem{remark}[ms@proposition]{Remark}
\newtheorem*{remark*}{Remark}

\crefname{ms@theorem}{Thm.}{Thms.}
\Crefname{ms@theorem}{Thm.}{Thms.}
\crefname{ms@proposition}{Prop.}{Props.}
\Crefname{ms@proposition}{Prop.}{Props.}
\crefname{ms@lemma}{Lem.}{Lems.}
\Crefname{ms@lemma}{Lem.}{Lems.}
\crefname{ms@definition}{Def.}{Defs.}
\Crefname{ms@definition}{Def.}{Defs.}
\crefname{ms@corollary}{Cor.}{Cors.}
\Crefname{ms@corollary}{Cor.}{Cors.}
\crefname{remark}{Rem.}{Rems.}
\Crefname{remark}{Rem.}{Rems.}

\makeatother


\title[Massive SU(2) gauge theory]{Massive SU(2) gauge theory in three dimensions}
\date{}
\hypersetup{pdftitle={Massive SU(2) gauge theory in three dimensions},pdfauthor={}}
\begin{document}
\begin{abstract}
We study a massive SU(2) gauge theory in three dimensions from
the chamber lattice model with finitely many fundamental fermion
flavours. At sufficiently weak gauge coupling, limits taken first
in the volume and then along a spacing subsequence give centered
correlations of the action density and scalar and pseudoscalar
fermion bilinears. They are covariant under the full Euclidean
group and reflection positive at every plane. They reconstruct
a finite family of Hermitian local Wightman fields, each with
nonzero two-point function. The proof combines the fermion block
expansion with the pure gauge construction, controls the cumulative
mass correction, and derives the insertion and shrinking-plate
bounds needed for rotations and reconstruction.
\end{abstract}
\maketitle
\tableofcontents
\section{Introduction}\label{sec:introduction}

We consider SU(2) gauge fields coupled to $N_f\ge1$ massive
four-component fermions in the fundamental representation.
The lattice regularization uses the balanced chamber Dirac
operator with Wilson parameter $1/2$, defined in
\cref{sec:qcd-lattice-foundation}. At spacing $a=L^{-K}$,
the inverse gauge coupling and bare fermion mass are
$\beta=\mathsf b/a$ and $m_0=a\mathsf m$. The torus measure is
\[
 Z^{-1}e^{-\beta A(U)}|\det D_{U,+}|^{2N_f}\,dU,
 \qquad A(U)=\sum_p(1-\tfrac12\operatorname{tr}U_p).
\]
Here $D_{U,+}$ is the $\tau_3=1$ internal block for one flavour
of the massive chamber operator. Its determinant and that
of the other block are complex conjugates.

\begin{theoremzero}
Fix the number of flavours and a compact positive mass interval.
For sufficiently weak gauge coupling, the chamber model has
subsequential continuum limits of the centered action-density
and scalar and pseudoscalar bilinear correlations on separated
supports, after the infinite-volume limit. Every such limit is
Euclidean covariant and reflection positive at every plane.
It reconstructs a finite family of Hermitian local Wightman
fields, each with a nonzero two-point function.
\end{theoremzero}
The precise statement, including the order of limits, field
family and parameter assumptions, is \cref{C-396-28}.

We retain the action density and a finite family of colour-singlet
fermion bilinears. Let $\tau_0=1,\tau_1,\tau_2,\tau_3$
be the Pauli basis in the internal two-dimensional factor,
and choose Hermitian flavour matrices $t_j$, $1\le j\le N_f^2$,
with $\operatorname{tr}(t_jt_l)=\delta_{jl}$. Define
\[
 (c_0,c_1,c_2,c_3)=(1,i,i,1),\quad
 \Gamma_{a,j}=1_{\rm s}\otimes c_a\tau_a\otimes1_{\rm c}\otimes t_j,
 \quad \mathfrak B=\{\Gamma_{a,j}\},\quad
 \mathcal I=\{0\}\cup\mathfrak B.
\]
The corresponding observables, for real $f\in C_c^1(\mathbb R^3)$, are
\begin{equation}\label{eq:qcd-physical-observables}
 \Phi_{0,K}(f)=a^{-1}\sum_p f(x(p))
                (1-\tfrac12\operatorname{tr}U_p),\qquad
 \Phi_{\Gamma,K}(f)=a\sum_x f(ax)
                      \sum_{v:x_v=x}\bar\psi(v)\Gamma\psi(v).
\end{equation}
In the bilinear sum $x\in\mathbb Z^3$ is a lattice site.
Plaquettes are counted once in each positive coordinate orientation,
and $x(p)$ is their physical center. The basis is reflection real:
the Grassmann reflection sends each observable to its geometrically
reflected test function without an index sign. An improper spatial
isometry instead acts with signs
$\pi_0=1$ and $\pi_{\Gamma_{a,j}}=(+1,-1,+1,-1)_a$.

At fixed $K$, take any weak limit $\nu_K$ of the gauge marginals
as the physical torus sides $2L^m$ tend to infinity; spinors are
antiperiodic in each coordinate. Local fermion expectations in this
state use the Gaussian Berezin functional of $D_U^{-1}$.
Let $S^T_{n,K,\nu_K}$ be the joint cumulants of
\eqref{eq:qcd-physical-observables}. Only tuples with pairwise
disjoint supports are considered. For a limiting family $(S_n)_{n\ge2}$,
its centered moments are
\begin{equation}\label{eq:qcd-centered-moments}
 \sigma(\varnothing)=1,\quad\sigma((\alpha,f))=0,\qquad
 \sigma(X)=\sum_{\substack{\mathcal P\text{ partition of }[n]\\
                       |B|\ge2\ (B\in\mathcal P)}}
                                \prod_{B\in\mathcal P}S_{|B|}(X_B).
\end{equation}

\begin{theorem}[The massive two-colour theory]\label{C-396-28}
Fix $N_f\ge1$ and $0<m_-\le m_+<1/2$. Choose an odd
integer $L\ge13$ and the auxiliary constants in the order and
numerical ranges of \cref{C-396-12}. Then there is
$\mathsf b_0<\infty$ with the following property.
For every $\mathsf b\ge\mathsf b_0$,
$\mathsf m\in[m_-,m_+]$, sequence $K_j\to\infty$ and choices
of the states $\nu_{K_j}$ above, there is a common subsequence
on which
\[
 S^T_{n,K_j,\nu_{K_j}}(X)\longrightarrow S_n(X)
           \qquad(n\ge2)
\]
for every separated tuple $X=((\alpha_i,f_i))_{i=1}^n$.
Every such limit has the following properties.
\begin{enumerate}[label=\textup{(\roman*)}]
\item It is symmetric and multilinear and defines distributions
of locally finite order off the coincidence diagonals.
For supports separated by $d>0$, each in a ball of radius $R$,
\[
 |S_n(X)|\le A(1+d^{-1})^3B(R)^n(n!)^{\varkappa}
                         \prod_i\|f_i\|_{\alpha_i},
\]
where $A,B(R),\varkappa<\infty$ are independent of the positions
of those balls and of the chosen subsequence,
$\|f\|_0=\max(\|f\|_\infty,40M\operatorname{Lip}f)$ and
$\|f\|_\Gamma=\|f\|_\infty$. The fixed constant $M$ belongs
to the construction.
\item Both $S$ and $\sigma$ are translation invariant and
covariant under every $R\in\mathrm O(3)$:
\[
 \sigma(RX)=\pi(X)^{\delta(R)}\sigma(X),\qquad
 \overline{\sigma(X)}=\pi(X)\sigma(X)
 \quad\text{for real }X,
\]
where $\pi(X)=\prod_i\pi_{\alpha_i}$ and
$\delta(R)=0$ for proper rotations and $1$ otherwise;
the same identities hold for $S_n$.
The centered functional is reflection positive at every plane $P$:
\[
 \sum_{i,l}\bar z_i z_l\,
              \sigma(\theta_PX_i\cdot X_l)\ge0
\]
for every finite family of separated tuples on one open side
and every choice of complex $z_i$.
\item There is a separable Hilbert space with a cyclic unit vector
$\Omega$ and $1+4N_f^2$ Hermitian scalar operator-valued tempered
distributions $\varphi_\alpha$ on $\mathbb R^{1+2}$, on a common
invariant dense domain. They are covariant under a strongly
continuous unitary representation of the proper orthochronous
Poincar\'e group fixing $\Omega$. Its translation spectrum lies
in the closed forward cone. Every pair of fields commutes at
spacelike separation. Their Euclidean correlations agree with
$\sigma$ on separated supports. The two-point function of
each $\varphi_\alpha$ is nonzero.
\item The flavour group $\mathrm U(N_f)^2$ acts unitarily,
fixes $\Omega$ and commutes with the Poincar\'e representation.
With $e_\pm=(1\pm\tau_3)/2$ and
$V=e_+\otimes V_1+e_-\otimes V_2$, its action is
$\varphi_\Gamma\mapsto\varphi_{V\Gamma V^{-1}}$ and
$\varphi_0\mapsto\varphi_0$. Space reflection acts unitarily
with the signs $\pi_\alpha$ and exchanges the flavour factors.
Here $\varphi_\Gamma$ is extended real linearly to the reflection-real span of $\mathfrak B$.
\end{enumerate}
The assertions concern centered separated Euclidean correlations.
They include neither convergence of lattice contact terms nor
uniqueness of the vacuum, clustering or a mass gap. The constants
are not uniform as $m_-\downarrow0$.
\end{theorem}

The proof has two model-dependent parts. First, the positive mass
gives a bounded inverse at every lattice gauge field; relative
determinants and periodic-state comparison then give reflection
positivity at coordinate and diagonal mirrors. Second, the
renormalization construction must control the determinant and
the source energies uniformly as $a\downarrow0$. The cumulative
second-order mass identity and its cubic error keep the mass
in its required window. Local source estimates and one common
vacuum energy in the upper and lower partition bounds give
compactness, insertion bounds and shrinking-plate bounds.
These last two bounds permit the finite-family rotation argument
and reconstruction. Nonvanishing follows from the short-distance
two-point laws, not from reflection positivity alone.

We adapt the gauge construction and the general functional results
of \cite{YMM1} to the chamber model defined in
\cref{sec:qcd-lattice-foundation}. The fermion estimates and their
use in the local and history expansions are proved below.
\Cref{def:qcd-gauge-admission} verifies the numerical and structural
hypotheses of the gauge induction; \cref{C-396-12} checks that the
additional matter conditions preserve the required choice of
constants. We use the observables in
\eqref{eq:qcd-physical-observables}. The scalar field indexed
by $0$ is the action density. No gauge potential or coloured
quark field is reconstructed here.

For the estimates below, the physical matrix source is
$J(x)=\sum_{\Gamma\in\mathfrak B}f_\Gamma(x)\Gamma$ and
\[
 Z[w,J]=\Omega\!\left(
     \exp\{\Phi_{0,K}(w)+\sum_\Gamma\Phi_{\Gamma,K}(f_\Gamma)\}
                  \right).
\]
Its fine fermion matrix is $D_U-a_K\hat J$; the hat repeats
the source at the chamber places of its site.

\section{The chamber model and its massive lattice states}
\label{sec:qcd-lattice-foundation}

We use the chamber fermion with Wilson parameter $r=1/2$.
The gauge group is $\mathrm{SU}(2)$, and $N_f\ge1$ is the number
of four-component Dirac flavours in its fundamental representation.
Throughout this section the spacing is one, $\beta\ge0$, and the
lattice mass $m_0$ is positive. The estimates are uniform in the
gauge configuration and the volume; their constants may depend on
$m_0$. Physical coordinates are recovered by $x\mapsto ax$, with
$\beta=\mathsf b/a$ and $m_0=a\mathsf m$.

\subsection{Chambers, spinors and the probability measure}

Let $W_3$ be the group of signed permutation matrices. Its $48$
elements label the chambers at a site of $\mathbb Z^3$.
Put
\[
 A_0=\{0<z_3<z_2<z_1<1/2\},\qquad
 (N_0,N_1,N_2,N_3)=(e_1,e_2-e_1,e_3-e_2,-e_3),
\]
and let $w_s$ be reflection in $N_s^\perp$. A chamber $v=(x,g)$
has alcove $x+gA_0$, site $x_v=x$, and neighbours
\begin{equation}\label{C-391-6}
 v\cdot s=(x+\delta_{s0}ge_1,gw_s),\qquad
 u_s(v)=gN_s/|N_s|,\qquad t_s=3|N_s|.
\end{equation}
Thus $v\mapsto v\cdot s$ is an involution and
$u_s(v\cdot s)=-u_s(v)$. The four neighbours are distinct.
Only $s=0$ changes the site. Define
\[
 U(v,v\cdot s)=1\quad(s=1,2,3),\qquad
 U(v,v\cdot0)=U_{\langle x,x+ge_1\rangle},\qquad
 U_{\langle y,x\rangle}=U_{\langle x,y\rangle}^{-1}.
\]

The fibre is
$\mathbb C^2_{\rm s}\otimes\mathbb C^2_\tau
 \otimes\mathbb C^2_{\rm c}\otimes\mathbb C^{N_f}$,
of dimension $8N_f$. Spin and colour are separate factors.
For the Pauli matrices $\sigma_\mu$ and $\tau_\mu$ put
\[
 \gamma_\mu=\sigma_\mu\otimes\tau_3,\qquad
 \gamma(u)=\sum_\mu u_\mu\gamma_\mu,\qquad
 \gamma_4=1\otimes\tau_1.
\]
All these matrices act trivially on colour and flavour. The massive
operator is
\begin{equation}\label{C-391-10}
 \begin{aligned}
 (D_Uf)(v)&=M_0f(v)-\sum_{s=0}^3h_s(v)U(v,v\cdot s)f(v\cdot s),\\
 h_s(v)&=t_s(\tfrac12-\gamma(u_s(v))),\qquad
 M_0=m_0+\tfrac12\sum_st_s.
 \end{aligned}
\end{equation}
Since $\sum_st_su_s(v)=0$, this is the chamber operator with
bare mass $m_0$. On a finite domain we keep the same diagonal
coefficient and retain only hops whose chambers and gauge link
belong to the domain. A full domain contains every chamber at its
sites and every link and plaquette contained in the region.
On the cubic torus of even side $N\ge4$, the gauge field is
periodic and spinors are antiperiodic in each coordinate.

The variables $\psi_i(v)$ and $\bar\psi_i(v)$ are independent
Grassmann generators. Normalize the Berezin integral to take value
one on $\prod_{v,i}\psi_i(v)\bar\psi_i(v)$. With
\[
 a_p(U)=1-\tfrac12\operatorname{tr}U_p,\qquad
 A(U)=\sum_pa_p(U),\qquad S_F=\bar\psi D_U\psi,
\]
define
\[
 \Omega(F)=\int dU\int d\bar\psi\,d\psi\,
                 e^{-\beta A-S_F}F,\qquad Z=\Omega(1).
\]
The sums and integrals here are over the chosen finite domain or
torus. Haar measures have total mass one.

\begin{proposition}[The determinant and the massive inverse]
\label{C-391-16}
On every finite domain, on an antiperiodic torus, and on the infinite
lattice, $\operatorname{Re}D_U\ge m_0$.
The operator is invertible and $\|D_U^{-1}\|\le m_0^{-1}$.
In finite volume,
\[
 \det D_U=|\det D_{U,+}|^{2N_f}>0,
\]
where $D_{U,+}$ acts on the $\tau_3=1$ block for one flavour.
Consequently $Z>0$ and
\[
 \mu(dU)=Z^{-1}e^{-\beta A(U)}\det D_U\,dU
\]
is a probability measure.
\end{proposition}

\begin{proof}
Let $(R_sf)(v)=U(v,v\cdot s)f(v\cdot s)$, with value zero at
an inactive hop. Then $R_s^*=R_s$ and $\|R_s\|\le1$.
The relation $u_s(v\cdot s)=-u_s(v)$ shows that
$\gamma(u_s)R_s$ is skew-adjoint. Hence
\[
 \operatorname{Re}D_U
       =m_0+\tfrac12\sum_st_s(1-R_s)\ge m_0.
\]
It follows that $\|D_Uf\|\ge m_0\|f\|$. The same estimate
for $D_U^*$ gives a closed, dense range, so $D_U$ is onto even
on the infinite lattice. Its inverse has the stated norm.
The operator commutes with $\tau_3$. Replacing all $\gamma_\mu$
by $-\gamma_\mu$ takes its adjoint, and therefore its two blocks
are $D_{U,+}$ and $D_{U,+}^*$, each repeated $N_f$ times.
Their determinants are nonzero. The Gaussian Berezin formula gives
$Z=\int e^{-\beta A}\det D_U\,dU>0$.
\end{proof}

\subsection{Reflection positivity on finite domains}

A lattice mirror is one of the planes
\[
 x_\mu=c,\quad c\in\tfrac12\mathbb Z,
 \qquad x_\mu\pm x_\nu=c,\quad c\in\mathbb Z,
 \quad \mu<\nu.
\]
Fix a side $H^+$, reflection $\theta$, and the unit normal $n$
pointing into that side. No mirror meets the interior of an alcove.
If the two alcoves of a hop lie on different sides, their common
wall is the mirror, the second chamber is $\theta v$, and
$u_s(v)=-n$ for the chamber on the positive side. A chamber
has at most one such wall. These statements follow directly from
the four inequalities defining $A_0$: the values of a coordinate
in $A_0$ lie in $(0,1/2)$, coordinate differences are nonzero and
have modulus less than $1/2$, and sums of two coordinates lie in
$(0,1)$. The common face of adjacent alcoves then determines the
crossed mirror.

Let $\mathfrak A_+(P)$ contain the generators of chambers on this
side and coefficients depending on gauge links with both endpoints
in $H^+$. Links in the plane are admitted. Define the antilinear,
order-reversing involution
\begin{equation}\label{C-391-12}
 \Theta f=\overline{f\circ\theta},\qquad
 \Theta\psi_i(v)=(\bar\psi(\theta v)\gamma(n))_i,\qquad
 \Theta\bar\psi_i(v)=(\gamma(n)\psi(\theta v))_i.
\end{equation}
The order $F\Theta F$ and the normal into the side of $F$ are
part of this convention.

\begin{theorem}[Finite-domain reflection positivity]
\label{C-391-13}
Let a finite domain be preserved by a lattice reflection. Then
\[
                      \Omega(F\Theta F)\ge0
             \qquad(F\in\mathfrak A_+(P)).
\]
The assertion holds for all nine families, either side, and bounded
Borel coefficients, without a gauge-invariance assumption on $F$.
\end{theorem}

\begin{proof}
\emph{The variables and the positive cone.}
If a gauge link crosses a coordinate mirror through links, write its
variable as $V'^{-1}V$, with $V,V'$ independent Haar variables.
Integrating these two variables gives precisely Haar integration of
the original link. Reflection exchanges $V$ and $V'$. All other
variables split into variables on the two sides and variables in
the plane. Let $J$ be Haar--Berezin integration in these variables.
Let $\mathcal K_+$ be the cone of finite sums $G\Theta G$ with
homogeneous $G$ in the enlarged positive half algebra. It is closed
under multiplication: for homogeneous $G,H$,
\[
 (G\Theta G)(H\Theta H)=(GH)\Theta(GH),
\]
because $H\Theta H$ is even. If $K\ge0$ is a finite matrix and
$b_i$ are odd positive-half elements, diagonalizing $K$ gives
$\sum_{i,j}b_iK_{ij}\Theta b_j\in\mathcal K_+$.

Let $T_+$ be the product $\prod\psi_i(v)\bar\psi_i(v)$ over the
positive-side generators. Reflection maps it to the corresponding
product on the negative side: the two determinants from
$\gamma(n)$ cancel. If $g$ is the coefficient of $T_+$ in $G$,
then only the product of these two top coefficients contributes,
and
\[
 J(G\Theta G)=\int dU_0\,
                   \left|\int dU_+\,g\right|^2\ge0.
\]
Thus $J$ is nonnegative on $\mathcal K_+$.

\emph{The gauge action.}
Plaquettes on either side give $S_G^++\Theta S_G^+$ in
$-\beta A$; plaquettes in the mirror are divided equally.
For a cut plaquette, let $A^p$ be its parallel transport along the
positive half of its boundary. At a diagonal mirror this path joins
the two opposite vertices in the mirror. At a link mirror include
the two positive half-link variables. The reflected path is the
other half of the boundary, with reverse orientation. Hence
\[
 \operatorname{tr}U_p
       =\sum_{a,b}A^p_{ab}\Theta(A^p_{ab}),\qquad
 -\beta A=S_G^++\Theta S_G^+-\beta n_\times+Y,
 \quad
 Y=\frac\beta2\sum_{p\ {\rm cut}}\sum_{a,b}
                              A^p_{ab}\Theta(A^p_{ab})\in\mathcal K_+.
\]
Here $n_\times$ counts cut plaquettes. No plaquette is cut by a
coordinate mirror through sites. The stated decomposition exhausts
the possibilities: along an edge the signed height changes by
$0$ or $\pm1$; a diagonal cut plaquette has heights $0,1,0,-1$,
and a link cut plaquette has two heights $1/2$ and two $-1/2$.

\emph{The fermion action.}
Terms wholly on one side give $S_F^++\Theta S_F^+$ in $-S_F$.
For a crossing pair let $v$ be the chamber on the positive side
and $s$ its crossing wall. If the gauge link is bisected, put
$\chi=V\psi(v)$ and $\bar\chi=\bar\psi(v)V^{-1}$; otherwise
put $\chi=\psi(v)$ and $\bar\chi=\bar\psi(v)$.
The two crossing hops are
\[
 \sum_{i,j}\bar\chi_i\bigl(t_s+\tfrac12t_s\gamma(n)\bigr)_{ij}
                           \Theta\bar\chi_j,
 \qquad
 \sum_{i,j}\chi_j\bigl(t_s-\tfrac12t_s\gamma(n)\bigr)_{ij}
                           \Theta\chi_i.
\]
Indeed $h_s(v)=t_s(1/2+\gamma(n))$ and
$h_s(\theta v)=t_s(1/2-\gamma(n))$; substituting
\eqref{C-391-12} gives these formulas, with one transposition of
odd generators in the second. Both coefficient matrices are
positive, their eigenvalues being $t_s/2$ and $3t_s/2$.
Their sum $X_F$ over crossing pairs belongs to $\mathcal K_+$.
The within-side identity follows also directly from
$\gamma(n)\gamma(u)\gamma(n)=-\gamma(\theta^\circ u)$.
Here $\theta^\circ$ is the linear part of the affine reflection $\theta$.

Put $S^+=S_G^++S_F^+$. All exponents are even, so
\[
 e^{-\beta A-S_F}
       =e^{-\beta n_\times}e^{S^+}\Theta(e^{S^+})e^{X_F}e^Y.
\]
For homogeneous $F$, every term of
$e^{S^+}\Theta(e^{S^+})X_F^kY^lF\Theta F$ belongs to
$\mathcal K_+$. The $X_F$ series is finite, the $Y$ series
converges uniformly in the gauge variables, and the coefficients
$1/(k!l!)$ are nonnegative. Apply $J$ term by term.
For general $F=F_{\rm even}+F_{\rm odd}$ the mixed terms are
odd and integrate to zero. This proves the assertion.
\end{proof}

The same proof gives positivity on the even cubic torus, with
antiperiodic spinors, at its coordinate mirrors. The torus is cut
by the chosen mirror and the parallel mirror at distance $N/2$.
Use a positive slab for the half algebra. Across the first boundary
the calculation above applies. At the far boundary, choose the
representatives so that its reflected chamber is translated by
one period in the normal direction. The outward normal there is
the opposite normal, and the antiperiodic sign from that period
reverses it once more. Thus the far crossing matrices are again
$t_s\pm t_s\gamma(n)/2$, both positive; the gauge half-link
factorisation and the top-monomial calculation are unchanged.
\phantomsection\label{C-391-33}
This proves the antiperiodic coordinate-torus assertion,
including mirrors through sites.
A diagonal plane does not cut this cubic torus into two slabs.

\subsection{Decay, finite-rank perturbations and the specification}

Let $d(v,w)$ be distance in the chamber hopping graph and put
\[
 h_{\max}=\tfrac32(6+6\sqrt2),\qquad
 \kappa_{\mathrm{lat}}=\log\left(1+\frac{m_0}{2h_{\max}}\right),\qquad
 S_{\kappa_{\mathrm{lat}}}=48\left(\frac{1+e^{-\kappa_{\mathrm{lat}}}}{1-e^{-\kappa_{\mathrm{lat}}}}\right)^3.
\]
Since a hop changes its site by zero or one lattice step,
$d(v,w)\ge|x_v-x_w|_1$ and
$\sum_we^{-\kappa_{\mathrm{lat}} d(v,w)}\le S_{\kappa_{\mathrm{lat}}}$.

\begin{lemma}[The propagator at a fixed configuration]
\label{C-389-9}
Writing $Q_U=D_U^{-1}$ on the infinite lattice, one has
\[
 \|Q_U(v,w)\|\le\frac2{m_0}e^{-\kappa_{\mathrm{lat}} d(v,w)}.
\]
Every entry is continuous in $U$ for the product topology and is a
uniform limit of functions of finitely many links. If $U$ has period
$N\mathbb Z^3$, its antiperiodic torus inverse is
\[
 Q_U^N(v,w)=\sum_{n\in\mathbb Z^3}(-1)^{n_1+n_2+n_3}
                                    Q_U(v,w+Nn).
\]
For fixed $v,w$, $Q_U^N(v,w)-Q_U(v,w)=O(e^{-\kappa_{\mathrm{lat}} N})$
uniformly over periodic $U$.
\end{lemma}

\begin{proof}
For $f_R(v)=\min\{d(v,w),R\}$, conjugating $D_U$ by
$e^{\kappa_{\mathrm{lat}} f_R}$ changes only the hopping terms. The absolute
row and column sums of their change are at most
$(e^{\kappa_{\mathrm{lat}}}-1)h_{\max}=m_0/2$. The conjugated operator therefore
has real part at least $m_0/2$ and inverse norm at most $2/m_0$.
Taking its $(v,w)$ entry and then letting $R\to\infty$ proves
the decay estimate. The same argument works on a torus.

The resolvent identity
\[
 Q_U-Q_{U'}=-Q_U(D_U-D_{U'})Q_{U'}
\]
and the decay estimate give an absolutely summable series over the
changed links. There are sixteen chamber hops over a link, each
of norm at most $3(1+r)\|U_b-U'_b\|$, where $r=1/2$.
Consequently the effect of changing links outside an $R$-neighbourhood
of $v$ tends to zero uniformly as $R\to\infty$.
On each fixed neighbourhood the finite sum is continuous in its
link variables. This proves continuity and the local approximation.
For periodic $U$ the proposed image sum converges absolutely by
the decay estimate. Applying $D_U$ gives the antiperiodic delta
function, so uniqueness of the torus inverse proves the formula.
Its terms with $n\ne0$ are bounded by a convergent geometric sum
times $e^{-\kappa_{\mathrm{lat}} N}$, for fixed $v,w$.
\end{proof}

For configurations $W,W'$ agreeing off a finite set of links define
\[
 \delta(W,W')=\det\bigl(D_WD_{W'}^{-1}\bigr)
             =\det\bigl(1+(D_W-D_{W'})Q_{W'}\bigr).
\]
This is a determinant of an identity plus a finite-rank operator.
If $B$ is a finite link set and $V\in\mathrm{SU}(2)^B$, let
$V\vee U$ equal $V$ on $B$ and $U$ elsewhere, and set
$d_B(V\mid U)=\delta(V\vee U,U)$.

\begin{lemma}[The relative determinant]\label{C-386-26}
The relative determinant is positive and multiplicative; for each
fixed finite $B$, the function $d_B(V\mid U)$ is jointly continuous
in $(V,U)$. The cocycle identity is
\[
 \delta(W,W')\delta(W',W'')=\delta(W,W'').
\]
For $\lambda=9/m_0$,
\[
 |\log d_B(V\mid U)|\le128N_f\#B\log(1+\lambda).
\]
Let $R$ be a nonnegative integer. A link is within distance $R$ of
$B$ if one of its endpoints has lattice $\ell^1$ distance at most $R$
from an endpoint of $B$. If $U'=U$ on all these links, then,
uniformly in $V$,
\begin{equation}\label{eq:qcd-determinant-quasilocal}
 |\log d_B(V\mid U')-\log d_B(V\mid U)|
 \le512N_f\#B\lambda^2S_{\kappa_{\mathrm{lat}}} e^{-\kappa_{\mathrm{lat}}(R+1)}.
\end{equation}
The corresponding relative determinants on antiperiodic tori
converge uniformly to $d_B$, over periodic configurations of the
corresponding period, as that period tends to infinity.
\end{lemma}

\begin{proof}
The perturbation $\Delta=D_W-D_{W'}$ is supported on the
sixteen chamber endpoints over each changed link, has rank at
most $128N_f\#B$, and norm at most $9$. The last bound follows
from the matching of chambers across the type-zero walls and
$\|U_b-U'_b\|\le2$. Thus
$|\det(1+\Delta Q_{W'})|\le(1+\lambda)^{128N_f\#B}$.
The inverse is obtained by exchanging $W,W'$, which gives the
lower bound. In the two $\tau_3$ blocks, the finite-rank determinants
are conjugates: $\det(1+AB)=\det(1+BA)$ changes the order in
the adjoint block. Hence $\delta$ is a modulus squared, repeated
$N_f$ times, and is nonzero because $D_WD_{W'}^{-1}$ is
invertible. Multiplicativity is the determinant product identity.
After compressing to the finite set of perturbed endpoints,
continuity follows from \cref{C-389-9}.

For the quantitative assertion, interpolate the link matrices on
$B$ linearly between $U$ and $V\vee U$. These matrices are
contractions, paired on reversed links by adjoint. Each $R_s$
is still a self-adjoint contraction, so the real-part bound and the
decay proof remain valid throughout the interpolation. With
$U_s=(1-s)U+s(V\vee U)$ and $\Delta_B=D_{V\vee U}-D_U$,
\[
 \log d_B(V\mid U)=\int_0^1\operatorname{Tr}(\Delta_BQ_{U_s})\,ds.
\]
Use the same interpolation for $U'$, and put $\Delta'=D_{U'}-D_U$.
Subtracting gives
\[
 -\int_0^1\operatorname{Tr}
                    (\Delta_BQ_{U'_s}\Delta'Q_{U_s})\,ds.
\]
Every nonzero term connects the endpoints over $B$ to a changed
link at distance at least $R+1$. The two kernels contribute
$(2/m_0)^2$ and their exponential decays. Summing the remaining
decay factor costs at most $S_{\kappa_{\mathrm{lat}}}$; the fibre trace and two
hopping perturbations cost at most $32N_f\lambda^2$ for each
of the at most $16\#B$ endpoints. This proves
\eqref{eq:qcd-determinant-quasilocal}.
For the torus comparison apply the same logarithmic identity.
The image formula compares the torus inverse to the inverse in
the periodically changed configuration; the other translates of
$B$ are at distances tending to infinity, so the same resolvent
estimate compares that inverse to $Q_{U_s}$. Both errors tend
to zero uniformly. The determinant bounds just proved are uniform,
so uniform convergence of logarithms also gives that of the ratios.
\end{proof}

Let $A_B$ sum the plaquettes containing a link of $B$, and define
\begin{equation}\label{C-386-10}
 (\gamma_BF)(U)=
 \frac{\int F(V\vee U)e^{-\beta A_B(V\vee U)}d_B(V\mid U)\,dV}
 {\int e^{-\beta A_B(V\vee U)}d_B(V\mid U)\,dV}.
\end{equation}

\begin{lemma}[The specification]\label{C-386-27}
These are probability kernels. They leave functions of the links
off $B$ unchanged, their values do not depend on $U|_B$, and
$\gamma_{B'}\gamma_B=\gamma_{B'}$ when $B\subset B'$.
They map continuous functions to continuous functions; on continuous
local functions their values are uniform limits of continuous local
functions. The set of laws $\nu$ satisfying $\nu\gamma_B=\nu$
for every finite $B$ is convex and weakly compact.
\end{lemma}

\begin{proof}
The numerator has positive continuous weight and its denominator is
positive. The positive relative weight
\[
 \mathrm{rw}(W,W')=
 e^{-\beta\sum_p(a_p(W)-a_p(W'))}\delta(W,W')
\]
is a cocycle. If $U$ is changed only on $B$, the resulting factor
in $\mathrm{rw}(V\vee U,U)$ is independent of $V$ and cancels
in the quotient. This proves independence of $U|_B$; properness
is immediate. For $B\subset B'$, insert the cocycle identity in
the two numerators and denominators and apply Fubini over
$\mathrm{SU}(2)^{B'}=\mathrm{SU}(2)^B
 \times\mathrm{SU}(2)^{B'\setminus B}$ to obtain consistency.
Continuity follows by integration over a compact factor.
For a local $F$, freeze all sufficiently distant links at one.
The plaquette factor is then unchanged near $B$, and
\eqref{eq:qcd-determinant-quasilocal} makes the two normalized
weights uniformly close. Indeed, if $|\log p-\log\widetilde p|
\le\epsilon$, their normalized integrals against $F$ differ
by at most $\|F\|_\infty(e^{2\epsilon}-1)$. This proves the
local approximation. Finally the equations are linear in $\nu$
and weakly closed on continuous local tests, since
$\gamma_BF-F$ is continuous. Such tests determine a measure;
therefore the equations hold for all bounded Borel $F$ as well.
The space of probability measures on the compact metrizable
configuration space is weakly compact.
\end{proof}

\subsection{Limits of tori and local fermion correlations}

For an infinite configuration define the Gaussian Berezin
functional by Wick's rule with covariance $Q_U$:
\[
 \left\langle\psi_{p_1}\cdots\psi_{p_k}
                      \bar\psi_{q_k}\cdots\bar\psi_{q_1}\right\rangle_U
       =\det[Q_U(p_i,q_j)]_{i,j\le k}.
\]
It is zero for unequal numbers of barred and unbarred generators.
For a local element $F$, with continuous link coefficients, put
$h_F(U)=\langle F\rangle_U$ and
$\omega_\nu(F)=\int h_F\,d\nu$.
These formulas also define $\omega_\nu$ for bounded Borel local
coefficients. The functional is on a Grassmann algebra; only its
restriction to gauge-field functions is a probability law.

\begin{proposition}[Torus-limit states]\label{C-389-10}
The periodically extended cubic torus marginals have weak limit
points as their even side tends to infinity. Every such $\nu$
solves \eqref{C-386-10} and is invariant under gauge transformations
and under $\mathbb Z^3\rtimes W_3$. Along each sequence defining
$\nu$,
\[
                 \langle F\rangle_N\longrightarrow\omega_\nu(F)
\]
for every local Grassmann element with bounded Borel local
coefficients. For a monomial $\Psi$ of degree $2k$,
\[
 |\omega_\nu(f\Psi)|\le\|f\|_\infty m_0^{-k}.
\]
At every coordinate lattice mirror, with either side,
\begin{equation}\label{eq:qcd-coordinate-state-rp}
                 \omega_\nu(F\Theta F)\ge0
                     \qquad(F\in\mathfrak A_+(P)).
\end{equation}
In particular the set of solutions of the specification is nonempty.
\end{proposition}

\begin{proof}
The configuration space is a countable product of compact copies
of $\mathrm{SU}(2)$ and is compact metrizable. Its laws are weakly
compact. Gram's inequality and $\|Q_U\|\le m_0^{-1}$ give
$|\langle\Psi\rangle_U|\le m_0^{-k}$, since each column of
the Wick matrix is a subvector of $Q_Ue_q$. The same argument
holds on the torus. By \cref{C-389-9}, $h_F$ is continuous for
continuous local coefficients and
\[
 \sup_{U\;N\text{-periodic}}|\langle F\rangle_U^N-h_F(U)|\longrightarrow0.
\]
Thus weak convergence of the gauge marginals proves convergence
of these local expectations.

For a finite $B$ and continuous local scalar $f$, finite-volume
conditional integration gives exactly
$\mu_N(\gamma_B^Nf)=\mu_N(f)$. The affected plaquettes agree
with those defining $A_B$ once $N$ exceeds the diameter of the
data. By \cref{C-386-26}, $\gamma_B^Nf-\gamma_Bf$ tends to
zero uniformly on the periodic configurations. The continuity of
$\gamma_Bf$ allows passage to the weak limit, proving
$\nu\gamma_Bf=\nu f$. The equality then holds on all bounded
Borel tests because both sides define laws agreeing on the
continuous local functions.

We record why Borel coefficients also converge. For a finite link
set $B$, let $k_B(V,\xi)$ be the normalized positive weight of
\eqref{C-386-10}, and let $k_B^N$ be its torus analogue.
The latter converges uniformly to the former, and $k_B$ is
continuous in $V$ and the infinite gauge configuration $\xi$
on their compact product space.
The $B$-marginal of $\mu_N$ has Haar density
\[
 p_N(V)=\int k_B^N(V,\xi)\,\mu_N(d\xi).
\]
Weak convergence and uniform continuity in $V$ imply uniform
convergence to $p(V)=\int k_B(V,\xi)\,\nu(d\xi)$:
use a finite net in the compact $V$ space to make the weak
convergence uniform. Since $\nu$ solves the specification, $p$
is its $B$-marginal density. Hence the marginals converge in
total variation. This proves convergence on bounded Borel local
scalar functions. For a local Grassmann element $f\Psi$,
$h_\Psi$ is uniformly approximable by continuous functions of
finitely many links. Apply the scalar assertion to $f$ times
each approximation and then remove the approximation. The
torus covariance error is uniform, so it gives the claimed
Grassmann convergence.

For a signed permutation $R$, choose $s_R\in\mathrm{SU}(2)$
with $s_R\sigma(u)s_R^{-1}=(\det R)\sigma(Ru)$, and put
$S_R=s_R\otimes\tau_2^{\delta(R)}$, where $\delta(R)=0$ for
proper $R$ and $1$ otherwise. Then
$S_R\gamma(u)S_R^{-1}=\gamma(Ru)$. Together with the action
$R(x,g)=(Rx,Rg)$ on chambers, this intertwines $D_U$ and
$D_{RU}$. Translations act by reindexing. Thus the cubic torus
marginals are invariant under translations and the signed coordinate
permutations. These maps act continuously on the infinite
configuration space, so the same holds in the limit.
For a gauge transformation and a continuous local test, choose
a periodic gauge transformation agreeing with it at every site
read by that test. Finite-volume gauge invariance then passes
to the weak limit. The lifts of lattice isometries on spinors
intertwine the corresponding Dirac operators, so Wick's rule
gives the same covariance of $\omega_\nu$; a lift's sign cancels
on even elements, and odd elements have zero expectation.

Finally, for a fixed coordinate mirror and local $F$ in its half
algebra, all sufficiently large cubic tori contain $F$ in the
appropriate open slab. Their antiperiodic reflection forms are
nonnegative by the torus argument following \cref{C-391-13}.
The just-proved convergence, applied to $F\Theta F$, gives
\eqref{eq:qcd-coordinate-state-rp}.
\end{proof}



\section{Reflection positivity of the infinite-volume lattice states}
\label{sec:qcd-diagonal-positivity}

We keep $r=1/2$, $m_0>0$ and now assume $\beta>0$.
The states obtained from cubic tori in the preceding section are
reflection positive at the coordinate mirrors. A diagonal plane does
not split these tori. We first prove an assertion about all periodic
solutions of the specification; this permits comparison with states
obtained from tori adapted to a diagonal normal.

Write $\mathcal G$ for the solutions of \eqref{C-386-10}, and
$\mathcal G^{\rm per}$ for those invariant under $N_0\mathbb Z^3$
for some integer $N_0\ge1$. For the proof we use the reflection
\[
 \vartheta_P=(-1)^{\deg}\Theta_P,
\]
where $\Theta_P$ is formed with the normal into the upper side.
Thus $\vartheta_P\psi=-\bar\psi\gamma(n)$ and
$\vartheta_P\bar\psi=-\gamma(n)\psi$, at reflected positions.
For an upper-side element $F$,
\begin{equation}\label{eq:qcd-reflection-order-dictionary}
 \omega_\nu(\vartheta_P(F)F)=\omega_\nu(F\Theta_PF).
\end{equation}
For homogeneous $F$ this follows by interchanging the two factors;
for a general $F$, the mixed-parity terms vanish. The convention
with $\vartheta_P$ is convenient for the crossing Fock space.

\begin{proposition}[All lattice mirrors]\label{C-386-13}
Every $\nu\in\mathcal G^{\rm per}$ is invariant under
$\mathbb Z^3\rtimes W_3$. For every lattice mirror $P$, either
side, and every local half-algebra element $F$ with bounded Borel
link coefficients, and every local element $G$ with bounded Borel
coefficients,
\[
 \omega_\nu(\Theta_PG)=\overline{\omega_\nu(G)},\qquad
 \omega_\nu(F\Theta_PF)\ge0.
\]
The reflection form is Hermitian and vanishes between elements of
different Grassmann parity. In particular these conclusions hold
for every weak limit of the cubic torus marginals.
\end{proposition}

The proof uses two different comparisons. Insertions of unitary
operators at a bond mirror force a reflection identity on its two
neighbouring layers. The field and link equations then propagate
that identity to the whole lattice. This makes the periodic-state
class compact. A second comparison, on tori adapted to a prescribed
mirror, then allows a separation argument within that class.

\subsection{The local equations and a surface comparison}

We identify the chamber $(k,g)$ with the point $x=k+gv_0$,
where $v_0=(v_1,v_2,v_3)=(3/8,1/4,1/8)$ lies in $A_0$,
and write $\pi(x)=k$.
The four neighbours $x^{(i)}$ are numbered in the reverse order
to \eqref{C-391-6}: type $i$ is wall $3-i$. Their weights are
\[
 (w_0,w_1,w_2,w_3)=(6,6\sqrt2,6\sqrt2,6),\qquad
 w(x,x')=\frac{w_i}{2}(r-\gamma(\widehat{x'-x})).
\]
There are eight type-three edges over each gauge link. Only these
edges carry a gauge transporter. In this notation
$D_U=M_0-\sum_i w(x,x^{(i)})U(\pi(x),\pi(x^{(i)}))$.
For an oriented chamber edge put
\[
 \operatorname{hop}_V(x,x')=
 \bar\psi(x)[w(x,x')\otimes V(\pi(x),\pi(x'))\otimes1_{N_f}]\psi(x'),
\]
where $V(\pi(x),\pi(x'))=1$ when the two chambers have the same site.
No result depends on the auxiliary choice of $v_0$.

Every $\omega_\nu$, $\nu\in\mathcal G$, obeys the field equations
\begin{equation}\label{C-386-11}
 \omega_\nu((D_U\psi)_pG)=\omega_\nu(\partial_{\bar\psi_p}G),
 \qquad
 \omega_\nu((\bar\psi D_U)_pG)=-\omega_\nu(\partial_{\psi_p}G).
\end{equation}
Here and below derivatives are left derivatives. For a finite gauge
link set $B$, put $H_B=-\beta A_B$ and
\[
 R_B(V;U)=e^{H_B(V\vee U)}
             \exp\left(\sum_{\text{oriented edges over }B}
                                \operatorname{hop}_{V}(x,x')\right),
 \quad
 \operatorname{Res}_BG=
 \frac{\int R_B(V;U)G(V\vee U)\,dV}{\int R_B(V;U)\,dV}.
\]
The denominator is an even Grassmann element with strictly positive
body. The link equations are
\begin{equation}\label{C-386-12}
                         \omega_\nu(G)=\omega_\nu(\operatorname{Res}_BG).
\end{equation}
For clarity, these equations follow from the specification, not from
an assumed fermionic probability law. Wick's determinant recursion
gives \eqref{C-386-11}. If a finite chamber set $Y$ contains the
endpoints over $B$ and the fields of $G$, let
$C=P_YQ_UP_Y$, $C'=P_YQ_{V\vee U}P_Y$ and
$\Delta=P_Y(D_{V\vee U}-D_U)P_Y$. The resolvent identity gives
$C'^{-1}=C^{-1}+\Delta$; these compressions are invertible since
\[
 \operatorname{Re}\langle f,Cf\rangle
      =\operatorname{Re}\langle D_UQ_Uf,Q_Uf\rangle
      \ge m_0\|Q_Uf\|^2\ge m_0\|f\|^2/\|D_U\|^2.
\]
The finite Gaussian formula consequently gives, for a polynomial
$\xi$ with coefficients independent of the gauge field, the bridge identity
\begin{equation}\label{eq:qcd-gaussian-bridge}
 d_B(V\mid U)\langle \xi\rangle_{V\vee U}
 =\left\langle \xi\exp\left(
       \sum\operatorname{hop}_{V}-\sum\operatorname{hop}_{U}\right)
                                                    \right\rangle_U.
\end{equation}
Apply this coefficientwise to $G(V\vee U)$ and to
$\operatorname{Res}_BG$ under the
numerator of $\gamma_B$. The two numerators agree after multiplication
by the common normalization; $\operatorname{Res}_BG$ is independent
of $U|_B$. Integration against $\nu\gamma_B=\nu$ proves
\eqref{C-386-12}.

A torus is specified by a rank-three period lattice $\Gamma$ and
a sign character $\chi:\Gamma\to\{\pm1\}$. Its systole is
$\min_{g\ne0}|g|_\infty$. The decay and finite-rank determinant
arguments of the preceding section apply with arbitrary such
periods and signs: the image sum uses $\chi(g)$, and the number
of periods of size at most $js$ is at most $(2j+3)^3$ when the
systole is $s$. In particular local covariances and determinant
ratios converge uniformly as $s\to\infty$.

\begin{lemma}[Surface comparison]\label{C-386-30}
Let $\mathcal D$ be a set of site representatives for a torus,
$B=\{\langle y,y+e_j\rangle:y\in\mathcal D,\ j=1,2,3\}$,
and let $s(\mathcal D)$ count links with exactly one endpoint in
$\mathcal D$. Set
\[
 q_0=1+\frac9{2m_0},\qquad
 K(\mathcal D)=(4\beta+192N_f\log q_0)s(\mathcal D).
\]
For nonnegative Borel $f$ depending on $B$,
\[
 e^{-2K(\mathcal D)}\gamma_Bf(U)
     \le\mu_{\mathbb T}(f^{\rm per})
     \le e^{2K(\mathcal D)}\gamma_Bf(U).
\]
Suppose $\Gamma\subset\Gamma'$, $\nu\in\mathcal G$ is
$\Gamma'$-invariant, and $A$ is real continuous and local.
For $\mathbb A=\sum_{v\in\Gamma'/\Gamma}\tau_vA$, let $M$
be the number of terms and $M_{\rm in}$ the number with a
representative supported on $B$. Then
\begin{equation}\label{eq:qcd-lower-tangent}
 \log\mu_{\mathbb T}(e^{t\mathbb A})
 \ge tM\nu(A)-2K(\mathcal D)
            -2|t|\|A\|_\infty(M-M_{\rm in}).
\end{equation}
\end{lemma}

\begin{proof}
Let $Y=\pi^{-1}(\mathcal D)$. Cutting the type-three edges at
its boundary gives the common interior operator $D_Y(V)$.
The torus differs from it by a perturbation of rank at most
$64N_fs(\mathcal D)$. The infinite operator differs from
$D_Y(V)\oplus P_{Y^c}D_UP_{Y^c}$ by a perturbation of rank at
most $128N_fs(\mathcal D)$. Both perturbation norms are at most
$w_3(1+r)/2=9/2$, and the operators before and after cutting have
real part at least $m_0$. Therefore
\[
 (1+\|R\|/m_0)^{-k}
  \le|\det((A+R)A^{-1})|\le(1+\|R\|/m_0)^k
\]
when $\operatorname{rank}R\le k$: use the finite-rank determinant
bound in one direction and $1-R(A+R)^{-1}$ in the other.
Factor the infinite relative determinant through the cut operators.
Its interior factor is $\det D_Y(V)/\det D_Y(U|_B)$; the second
denominator and exterior factors do not depend on $V$.
The torus determinant has the same interior factor, leaving the
two rank bounds above. At most $2s(\mathcal D)$ plaquettes
differ, and $0\le a_p\le2$. Thus the unnormalized torus and
conditional weights differ by a positive factor independent of
$V$ times a factor in $[e^{-K},e^K]$. Normalization proves the
first assertion.

Let $\Phi=t\sum_{\rm in}\tau_vA$. The first assertion and
Jensen's inequality give
$\log\mu_{\mathbb T}(e^\Phi)+2K\ge\gamma_B\Phi(U)$.
Integrate against $\nu$, use its specification and $\Gamma'$
invariance, and bound the omitted terms by
$|t|\|A\|_\infty(M-M_{\rm in})$ on each side. This gives
\eqref{eq:qcd-lower-tangent}.
\end{proof}

We will also use a slowly varying scalar weight on tori. If
$W_j>0$ is continuous and
\begin{equation}\label{eq:qcd-vanishing-tilt}
 \sup_{U=U'\ {\rm off}\ B+\Gamma_j}
                |\log W_j(U)-\log W_j(U')|\longrightarrow0
\end{equation}
for every finite $B$, then every weak limit of
$W_j\mu_{\mathbb T_j}/\mu_{\mathbb T_j}(W_j)$, with systole
tending to infinity, solves the same specification. Indeed its
finite-volume conditional kernel differs uniformly from $\gamma_B$
by the ratio estimate following \cref{C-386-27}, the error being
the sum of \eqref{eq:qcd-vanishing-tilt} and the torus determinant
error. Conditional integration is exact before passage to the limit.
\phantomsection\label{C-386-28}
The image formula also proves convergence of local continuous
Grassmann expectations to $\omega_\nu$ along these weighted laws.

\subsection{Operators inserted at a bond mirror}

Fix a coordinate direction, write its coordinate as $j$, and put
$c=t+1/2$. Write $u_j(b)=U(b^{(j)})$ for horizontal gauge links
and $w_j(z)=U((z,j),(z,j+1))$ for vertical ones. For a finite
horizontal link set $S$ choose a finite nonempty set $\Lambda$
containing its endpoints. Let $X^-(z)$ be the eight chamber sites
over $(z,t)$ at height $t+v_1$, and $X^+(z)=\theta X^-(z)$.
These are exactly the endpoints of the edges crossing the mirror.
Put $X^\pm(\Lambda)=\bigcup_{z\in\Lambda}X^\pm(z)$.
Two-layer elements have continuous coefficients in
$(u_t|_S,w_t|_\Lambda,u_{t+1}|_S)$ and generators at
$X^-(\Lambda)\cup X^+(\Lambda)$.

Let $\mathcal F_\Lambda$ be the Grassmann algebra on the upper
generators, made a Hilbert space by declaring its ordered monomials
orthonormal. An orthonormal change of the one-generator basis
preserves this inner product on every exterior power, by the
determinant formula for its Gram matrix. The endomorphism
$\mathsf O$ induced on this algebra by
\[
 \bar\psi e\longmapsto\bar\psi\,3(1+r\gamma(n))e,
 \qquad e^*\psi\longmapsto(3(1-r\gamma(n))e)^*\psi
\]
is positive and invertible: its one-generator eigenvalues are
$3(1\pm r)>0$. For $\Upsilon\in\operatorname{End}\mathcal F_\Lambda$
define
\begin{equation}\label{C-386-41}
 \operatorname{Ins}[\Upsilon]=
       \sum_{a,a'}\langle a,\Upsilon a'\rangle\vartheta(a')a,
 \qquad \mathrm{Cx}_\Lambda=\operatorname{Ins}[\mathsf O].
\end{equation}
Here $a,a'$ range over an orthonormal monomial basis. This map
is independent of the basis and is a linear bijection onto the
Grassmann algebra of the two sides. It satisfies
\[
 \vartheta\operatorname{Ins}[\Upsilon]
      =\operatorname{Ins}[\Upsilon^*],\qquad
 \mathrm{Cx}_\Lambda
       =\exp\left(\sum_{x\in X^+(\Lambda)}
          \operatorname{hop}_1(x,\theta x)
                   +\operatorname{hop}_1(\theta x,x)\right).
\]
The subscript $1$ means that the crossing transport has been set
to one. To check the second identity, diagonalize $\gamma(n)$.
The crossing terms become
$3(1+\epsilon r)\vartheta(\bar\psi e)\bar\psi e$ and
$3(1-\epsilon r)\vartheta(e^*\psi)e^*\psi$ for
$\gamma(n)e=\epsilon e$. They are commuting square-zero
elements. Expanding their exponential gives the exterior powers
of the same one-generator matrix. On disjoint upper sets, the
product of even insertions is the insertion of their tensor product:
the two signs from interchanging odd factors cancel. These facts
also show $\operatorname{Ins}[\Upsilon]$ is even precisely when
$\Upsilon$ preserves parity.

For the gauge factor let $k_\beta(g)=e^{-\beta(1-\operatorname{tr}g/2)}$.
Write $\chi_l$ for the spin-$l$ character of $\mathrm{SU}(2)$,
$l\in\tfrac12\mathbb N_0$, and $I_n$ for the modified Bessel function.
Its convolution eigenvalues on spin-$l$ matrix coefficients are
\begin{equation}\label{C-386-15}
 \lambda_l=\frac2\beta e^{-\beta}I_{2l+1}(\beta)>0,
 \qquad \lambda_l\le\frac{\beta^{2l}}{(2l+1)!}.
\end{equation}
The Weyl integration formula gives this expression; positivity
also follows from the power series of $I_{2l+1}$.
Consequently $K_S$, convolution with
$k_S(u,u')=\prod_{b\in S}k_\beta(u(b)u'(b)^{-1})$,
maps the trigonometric polynomials bijectively onto themselves.
The series
$q_\beta=\sum_l(2l+1)\lambda_l^{1/2}\chi_l$ converges
uniformly. If $q_u=\prod_bq_\beta(u(b)\,\cdot\,^{-1})$,
then
\begin{equation}\label{C-386-16}
 \langle q_u,q_{u'}\rangle=k_S(u,u'),\qquad \|q_u\|=1.
\end{equation}

On $\mathcal H_{S,\Lambda}=L^2(\mathrm{SU}(2)^S)\otimes\mathcal F_\Lambda$
let $R_g=L_g\otimes\mathsf s_g^{-1}$, where
$(L_gf)(u)=f(g^{-1}.u)$ and $\mathsf s_g$ is the diagonal
gauge action on upper generators. Let $\mathcal W$ be the algebra
of bounded operators commuting with every $R_g$.
For $X\in\mathcal W$ define
\[
 \langle a,\kappa_X(u',u)a'\rangle
       =\langle q_{u'}\otimes a,X(q_u\otimes a')\rangle,
 \quad
 \widehat G_X(u,u')=
 \frac{\operatorname{Ins}[\mathsf O^{1/2}\kappa_X(u',u)\mathsf O^{1/2}]}
                       {k_S(u,u')\mathrm{Cx}_\Lambda}.
\]
The lift of a reduced element is
\begin{equation}\label{eq:qcd-reduced-lift}
 (\mathrm{Lf}\,\widehat Y)(U)=
       \mathsf s^+_{w_t}\widehat Y(u_t,w_t.u_{t+1}),\qquad
                         G_X=\mathrm{Lf}\,\widehat G_X.
\end{equation}
The superscript $+$ means gauge substitution on the upper
generators only. The crossing factor has body one and is invertible.
For a Grassmann polynomial, $|\cdot|_1$ denotes the sum of the
absolute values of its coefficients in the ordered monomial basis.

\begin{lemma}[Insertions and their range]\label{C-386-45}
The element $G_X$ is continuous, local, even and gauge invariant;
$G_1=1$, $\vartheta G_X=G_{X^*}$, and
$\sup_U|G_X(U)|_1\le C_{S,\Lambda}\|X\|$.
Every gauge-invariant two-layer element is a limit in coefficient
supremum norm of elements $G_X$.
There is an antiunitary involution $J$ on $\mathcal H_{S,\Lambda}$
such that $\mathsf JG_X=G_{JXJ}$, where $\mathsf J$ is the
antilinear algebra automorphism induced on generators by
$C_+=(\sigma_2\otimes\tau_1)\otimes\sigma_2$ and on coefficients
by complex conjugation. For every configuration,
$\langle\mathsf JF\rangle_U=\overline{\langle F\rangle_U}$.
\end{lemma}

\begin{proof}
The identities follow from \eqref{C-386-16}, the insertion
identities, and $\kappa_X(u,u')^*=\kappa_{X^*}(u',u)$.
Gauge covariance follows from $L_gq_u=q_{g.u}$, by moving the
gauge action through $X$. The element is even since
$R_{-1}$ acts as parity. For its norm, $k_S\ge e^{-2\beta\#S}$,
$\|\kappa_X\|\le\|X\|$, and all other maps act on a fixed
finite Grassmann algebra; multiplication by the inverse crossing
factor and the lift therefore cost a constant depending only on
the finite pair. Each upper generator is sent by the lift to two
colour components, so this cost is uniform in the gauge field.

A gauge-invariant two-layer element has the form
$\mathrm{Lf}\,\widehat Y$ with reduced $\widehat Y$ invariant
under diagonal gauge substitutions. Indeed set its middle gauge
variables $w_t$ to one by an upper-layer gauge transformation;
the remaining simultaneous transformations on the two layers give
this invariance, and reversing the substitution gives
\eqref{eq:qcd-reduced-lift}. Expand $k_S\widehat Y$ in
Grassmann monomials and approximate each coefficient uniformly
by trigonometric polynomials in $(u,u')$. Gauge averaging preserves
these polynomial spaces and the approximation. To represent an
approximant exactly, write
$m\mathrm{Cx}_\Lambda=\operatorname{Ins}[\Upsilon_m']$
and $\Upsilon_m=\mathsf O^{-1/2}\Upsilon_m'\mathsf O^{-1/2}$.
If $e_\alpha$ is a normalized matrix-coefficient basis,
\[
 \langle q_{u'},
    (\lambda_\alpha\lambda_{\alpha'})^{-1/2}
       |e_\alpha\rangle\langle e_{\alpha'}|q_u\rangle
       =e_\alpha(u')\overline{e_{\alpha'}(u)}.
\]
Thus finite sums of these rank-one operators tensored with
$\Upsilon_m$ represent every approximant. Average the resulting
operator over $R_g$ to put it in $\mathcal W$.
This step uses both $r<1$, for $\mathsf O^{-1/2}$, and
$\beta>0$, for $\lambda_\alpha>0$.

Finally $C_+\overline{D_U}=D_UC_+$ and $C_+^2=1$.
The automorphism $\mathsf J$ fixes the Gaussian action and the
top monomial, hence has the asserted expectation identity.
It commutes with gauge substitutions and with $\vartheta$.
On the upper Fock space it is antiunitary; tensor it with complex
conjugation on $L^2$, where $q_u$ is real. The same insertion
calculation gives $\mathsf JG_X=G_{JXJ}$.
\end{proof}

For a torus $\mathbb T$, write $Z_{\mathbb T}[F]=\Omega_{\mathbb T}(F)$
for its unnormalized Haar--Berezin functional, and let $C(N)$
denote the antiperiodic cubic torus of side $N$.

\begin{lemma}[The operator in the mirror]\label{C-386-46}
On an antiperiodic cubic torus, let finitely many disjoint pairs
$(S_a,\Lambda_a)$ lie in one bond mirror. For even, gauge-invariant
slab elements $F,F'$ there are vectors $\Phi_F,\Phi_{F'}$,
independent of the $X_a$, such that
\begin{equation}\label{C-386-47}
 Z_{\mathbb T}\left[\vartheta(F')\prod_aG_{X_a}F\right]
       =\langle\Phi_F,(\bigotimes_aX_a\otimes1)\Phi_{F'}\rangle.
\end{equation}
Consequently its modulus is at most
$\prod_a\|X_a\|\,Z_{\mathbb T}[\vartheta(F)F]^{1/2}
Z_{\mathbb T}[\vartheta(F')F']^{1/2}$.
\end{lemma}

\begin{proof}
Gauge-fix the vertical crossing links at the two boundary mirrors
to one. Gauge invariance of the integrand and Haar invariance make
this change exact: fix each crossing variable $V$, perform the
gauge transformation $V^{-1}$ at its lower endpoint, and integrate
the remaining lower variables; their Haar measure is unchanged.
The transformation is one on the upper slab.

Let $U_+,U'_+$ be the two slab configurations, with the latter
reflected into the upper slab. The Wilson weight is
\[
 W_+(U_+)W_+(U'_+)
 k_{S_0}(u^{\rm dn},u'^{\rm dn})
 k_{S_0}(u^{\rm up},u'^{\rm up}),
\]
where $S_0$ is the full finite horizontal section. The fermion
weight is $e^{-S_+}\vartheta(e^{-S_+})\operatorname{Ins}[\mathsf O_Y]$,
with $Y$ the upper crossing generators at both boundaries.
At the far boundary the antiperiodic sign changes the apparent
negative crossing coefficient to a positive one, so
$\mathsf O_Y\ge0$. Insertion $G_{X_a}$ cancels its factor
$k_{S_a}\mathrm{Cx}_{\Lambda_a}$ and replaces it by the kernel
of $X_a$ between $q_u$ and the corresponding Fock modes.

For $\xi=Fe^{-S_+}$ define $v_\xi\in\operatorname{Gr}(Y)$ by
$\langle v_\xi,a\rangle=\int_{\rm upper}\xi a$.
The top-monomial normalization and evenness give exactly
\[
 \int\vartheta(\xi')\xi\operatorname{Ins}[\Upsilon]
                =\langle v_\xi,\Upsilon v_{\xi'}\rangle.
\]
There is no sign: products of even insertions on disjoint sets
are tensor products, as checked after \eqref{C-386-41}.
Factor $L^2(\mathrm{SU}(2)^{S_0})$ over the $S_a$ and its
complement, and factor the Fock space over the disjoint mode sets.
In that order set
\[
 \Phi_F=\int dU_+\,W_+(U_+)\,
        q_{u^{\rm dn}}\otimes\mathsf O_Y^{1/2}v_{Fe^{-S_+}}
                                      \otimes q_{u^{\rm up}}.
\]
In these finite-dimensional Fock factors and compact gauge
integrals, Fubini proves \eqref{C-386-47}. Setting all $X_a=1$
identifies $\|\Phi_F\|^2$ with the reflected form, and
Cauchy--Schwarz gives the bound.
\end{proof}

\begin{proposition}[A periodic array of contraction insertions]
\label{C-386-48}
Let $\|X\|\le1$, and suppose a cube of side $2r_0$ contains
the support of $G_X$. If $N'>2r_0$ is even, $i$ is odd,
and $iN'\ge4$,
then on the cubic torus of side $iN'$,
\[
 \left|Z_{\mathbb T}\left[
       \prod_{v\in N'\mathbb Z^3/(iN')\mathbb Z^3}\tau_vG_X\right]
                                                 \right|\le Z_{\mathbb T}[1].
\]
\end{proposition}

\begin{proof}
Group the insertions by their $i$ normal layers, and replace the
operator on a layer by a label in $\{1,X,X^*\}$. Let $F(\sigma)$
be the resulting partition function. The supports of distinct
insertions are disjoint. At layer $k$ use \cref{C-386-46} with
the two boundaries $c+kN'$ and $c+kN'+iN'/2$. Since $i$ is
odd, the second boundary lies between insertion layers. Reflection
and Cauchy--Schwarz give
\[
 |F(\sigma)|^2\le F(R_k^+\sigma)F(R_k^-\sigma),\qquad
                         F(R_k^\pm\sigma)\ge0,
\]
where $R_k^\pm$ replaces label $k$ by $1$, retains one half
of the other labels, and reflects their adjoints into the other
half. On the finite set of label arrays choose a maximizer of
$|F|$. Both reflected arrays also maximize it. A string of $l$
successive labels $1$ can be enlarged, by reflection at its next
boundary, to a string of length $\min(2l+1,i)$. Starting with the
one $1$ created by $R_k^\pm$ gives a constant $1$ maximizer.
Its value is $Z_{\mathbb T}[1]$.
\end{proof}

\subsection{The lower bound for the insertion array}

The following estimate concerns a normalized Gaussian Berezin
functional in a fixed gauge configuration, on any torus. It uses
the mass bound alone.

\begin{lemma}[Separated local elements]\label{C-386-50}
Fix a cube size. There are $\epsilon_0>0$ and $C<\infty$,
depending on that size, $m_0$ and $N_f$ but not on the torus,
the gauge field or the number $R$ of cubes, with the following
property. Let the cubes and all their period translates be disjoint,
and let $V_j$ be even elements supported in them with
$|V_j|_1\le1$. Then
\[
 F(z)=\left\langle e^{z\sum_{j=1}^RV_j}\right\rangle_U^{\mathbb T}
\]
has no zero for $|z|<\epsilon_0$, and its branch $L(0)=0$ of
the logarithm satisfies
\begin{equation}\label{eq:qcd-separated-logarithm}
 \left|L(z)-z\sum_j\langle V_j\rangle_U^{\mathbb T}\right|
                                                     \le C|z|^2R.
\end{equation}
\end{lemma}

\begin{proof}
Write the covariance as $C(p,q)=\langle f_p,g_q\rangle$ with
$\|f_p\|=1$, $\|g_q\|\le m_0^{-1}$. Its block row and
column sums are at most $2S_{\kappa_{\mathrm{lat}}}/m_0$.
For a covariance $K$ and left Grassmann derivatives, put
$\Delta_K=\sum K(p,q)\partial_{\bar\psi_q}\partial_{\psi_p}$.
Wick's rule is $\langle a\rangle_K=[e^{\Delta_K}a]_0$.
In particular differentiation in the covariance inserts
$\Delta_{\dot K}$, and block diagonal covariances factorize
expectations of even elements in distinct blocks.

We use the forest interpolation identity
\cite[Thm.~III.1]{AbdesselamRivasseau1994} for a polynomial $H$ in
edge variables on $n$ vertices:
\begin{equation}\label{C-386-51}
 H(1)=\sum_{\mathcal F}\int_{[0,1]^{\mathcal F}}dt\,
       \left(\prod_{e\in\mathcal F}\partial_{s_e}H\right)
                                      (s^{\mathcal F}(t)).
\end{equation}
The off-diagonal value $s^{\mathcal F}_{ij}$ is the minimum of
the parameters on the forest path from $i$ to $j$, and zero if
there is no path; the diagonal values are one. This matrix is
positive semidefinite: integrate over $u\in[0,1]$ the block
matrices of ones corresponding to components of forest edges
with $t_e\ge u$. Hence $s_{ij}=\langle e_i,e_j\rangle$ for
unit vectors $e_i$. The interpolated covariance is represented
by $e_i\otimes f_p,e_j\otimes g_q$, so Gram's inequality
has the same constant at every interpolation point.

Set $W_j=e^{zV_j}-1$, so $|W_j|_1\le2|z|$ for $|z|\le1$.
Apply \eqref{C-386-51} to
$\langle\prod_{j\in J}W_j\rangle_{C_s}$ for every subset $J$.
Factorization over forest components gives a hard-core polymer
sum with nonempty vertex sets $P$ and activities
\[
 \zeta(P)=\sum_{T\ {\rm tree\ on}\ P}\int dt\,
     \left\langle\prod_{e\in T}\Delta_e\prod_{j\in P}W_j
                                  \right\rangle_{C_{s^T(t)}}.
\]
Let $N_b$ bound the number of generators in one cube and set
$K_G=\max(1,m_0^{-1})^{N_b/2}$.
The sums of the cross-covariance absolute values over a second
cube are bounded by a constant $g_b$, uniformly in the first
cube: one may take $g_b=256N_f^2n_bS_{\kappa_{\mathrm{lat}}}/m_0$ if $n_b$
bounds its chamber sites. Expanding the derivatives and applying
Gram's inequality bounds a tree by the product of its edge sums
times $(2|z|K_G)^{\#P}$. Each differentiation removes specified
generators, and there is no additional factorial at a vertex.
Root at a specified vertex; there are at most $4^{n-1}$ rooted
plane-tree shapes with $n$ vertices, and summing the cube labels
along each tree costs at most $g_b^{n-1}$. Thus
\begin{equation}\label{eq:qcd-polymer-sum}
 \sum_{P\ni j,\ \#P=n}|\zeta(P)|
                      \le(2|z|K_G)^n(4g_b)^{n-1}.
\end{equation}
Choose $\epsilon_0$ so small that
$\sum_{P\ni j}|\zeta(P)|2^{\#P-1}\le1/2$.
For a finite vertex set $I$, let $Z(I)$ be its polymer sum.
Deleting $j$ gives
\[
 \frac{Z(I)}{Z(I\setminus\{j\})}=1+y_j,
 \qquad
 y_j=\sum_{P\subset I,\,j\in P}
             \zeta(P)\frac{Z(I\setminus P)}{Z(I\setminus\{j\})}.
\]
Induction on $\#I$ bounds the quotient in each summand by
$2^{\#P-1}$, proves $|y_j|\le1/2$, and proves nonvanishing.
The terms $\#P\ge2$ in \eqref{eq:qcd-polymer-sum} give
$|y_j-\zeta(\{j\})|\le C_1|z|^2$.
Also $|\zeta(\{j\})-z\langle V_j\rangle|\le C_2|z|^2$.
Order the vertices, multiply these successive ratios, and use
$|\log(1+y)-y|\le|y|^2$ for $|y|\le1/2$.
Summation proves \eqref{eq:qcd-separated-logarithm}.
\end{proof}

\begin{proposition}[The lower tangent for insertions]\label{C-386-52}
Let $\nu\in\mathcal G^{\rm per}$ have period number $N_0$,
and let $D\in\mathcal W$ satisfy $D^*=-D$ and $JDJ=D$.
For an even multiple $N'$ of $N_0$ large enough to separate
the supports, let
$P_i(\epsilon)=\prod_{v\in N'\mathbb Z^3/(iN')\mathbb Z^3}
                         \tau_vG_{e^{\epsilon D}}$.
For all sufficiently small real $\epsilon$, its normalized
fermion expectation is positive at every gauge configuration,
and
\[
 \liminf_{i\to\infty}i^{-3}
       \log\frac{Z_{C(iN')}[P_i(\epsilon)]}{Z_{C(iN')}[1]}
                  \ge\epsilon\omega_\nu(G_D)-C_D\epsilon^2.
\]
The constant and the permitted interval for $\epsilon$ are
independent of $i,N',\nu$ and the configuration.
\end{proposition}

\begin{proof}
Linearity and the norm bound of \cref{C-386-45} show
\[
 G_{e^{\epsilon D}}=1+\epsilon G_D+O_{|\cdot|_1}(\epsilon^2).
\]
It is fixed by $\mathsf J$, since $Je^{\epsilon D}J=e^{\epsilon D}$.
For small $\epsilon$ its finite Grassmann logarithm
$V_\epsilon=\log G_{e^{\epsilon D}}$ exists,
is even and $\mathsf J$-fixed, and satisfies
\[
 |V_\epsilon|_1\le c_2|\epsilon|,\qquad
 |V_\epsilon-\epsilon G_D|_1\le c_3\epsilon^2.
\]
Apply \cref{C-386-50} to its disjoint translates, divided by
$c_2\epsilon$, with parameter $z=c_2\epsilon$.
The expectation is real by the $\mathsf J$ identity, has no
zero on the small real interval, and is one at zero. It is
therefore positive, and its real logarithm gives, uniformly in $U$,
\[
 \langle P_i(\epsilon)\rangle_U^{C(iN')}
 \ge\exp\left(\epsilon\sum_v
                  \langle\tau_vG_D\rangle_U^{C(iN')}
                                     -C_D\epsilon^2i^3\right).
\]
The function $h_{G_D}$ is real continuous and uniformly
approximable by real local functions. The image formula replaces
each torus expectation by its infinite-volume value with an error
$O(e^{-\kappa_{\mathrm{lat}} iN'})$ uniform in its position. Approximate
$h_{G_D}$ by a real local $A$ within $\delta$ and integrate
the preceding inequality against the torus gauge law.
Apply \eqref{eq:qcd-lower-tangent} to the array of $A$ on the
cube $\{0,\ldots,iN'-1\}^3$. Its surface is $6(iN')^2$;
the number of array terms meeting a fixed-width boundary layer
is $O((iN')^2)$. Divide by $i^3$ and let $i\to\infty$ at
fixed $N'$, then $\delta\to0$. The surface, boundary and image
errors vanish, and $\nu(h_{G_D})=\omega_\nu(G_D)$.
\end{proof}

\begin{theorem}[The two-layer principle]\label{C-386-54}
For every periodic solution and every continuous two-layer
element $Y$ at a bond mirror,
\[
                         \omega_\nu(\vartheta Y)=\overline{\omega_\nu(Y)}.
\]
\end{theorem}

\begin{proof}
For $D$ as in \cref{C-386-52}, $e^{\epsilon D}$ is unitary.
Apply the upper bound \cref{C-386-48} for odd $i$, and the
lower tangent along the same indices. Their quotient is positive,
so
$0\ge\epsilon\omega_\nu(G_D)-C_D\epsilon^2$ for both signs
of small real $\epsilon$. Thus $\omega_\nu(G_D)=0$.
Let $\phi(X)=\omega_\nu(G_X)$. It is linear and
$\phi(JXJ)=\overline{\phi(X)}$. For $JXJ=X$, its value is
real, and $(X-X^*)/2$ is an allowed $D$, giving
$\phi(X^*)=\phi(X)$. Decompose a general $X$ into its two
$J$-real parts to conclude
$\phi(X^*)=\overline{\phi(X)}$.

For a gauge-invariant $Y$, the range assertion of
\cref{C-386-45} now proves the result by uniform coefficient
approximation. The expectation is continuous in this norm by the
mass bound on its finitely many generators. For general $Y$,
average over gauge transformations on the finite sites of the
two layers which it reads. Gibbs states are gauge invariant:
resampling the links incident to one site, Haar substitution and
the relative-determinant cocycle give this assertion first for
one-site transformations, then for arbitrary finite sets.
The averaged element is gauge invariant and remains two-layer;
the expectation and its reflected expectation are unchanged.
This proves the identity for $Y$.
\end{proof}

\subsection{Two layers determine the chain}

Fix the same coordinate direction and $c=t_0+1/2$. For each
horizontal site let $\widehat w_{t_0}=1$ and
$\widehat w_{j+1}(z)=\widehat w_j(z)w_j(z)$.
Temporal gauge replaces all vertical gauge links by one and
horizontal layers by
\[
 \widetilde u_j(b)=\widehat w_j(b_-)u_j(b)\widehat w_j(b_+)^{-1}.
\]
Each such coordinate and each $\widehat w_j$ reads finitely many
links. Pull a local element of the new variables back by this
substitution and the corresponding gauge substitution on its
generators; denote the map by $\mathrm{tg}^*$.
It is a unital homomorphism and sends a continuous local element
to a continuous local element.

The fermion sites above a gauge plane have six thin layers,
ordered by their normal displacements
$-v_1,-v_2,-v_3,v_3,v_2,v_1$. Write them $X_{j,k}$,
$1\le k\le6$, or $\mathbb L_{6j+k}$. Every site has one
neighbour in each adjacent thin layer and two in its own thin
layer. The upward edge types are $2,1,0,1,2,3$ as $k$ runs
from $1$ to $6$, and the downward types are $3,2,1,0,1,2$.
In temporal gauge these upward and downward transports are one.
At layers $k=1,6$ every transport is one; at each other layer
the only nontrivial transport reads one horizontal coordinate
of plane $j$. These statements follow by applying the four
wall reflections to the signed coordinates of $gv_0$.
In particular the endpoints of the fermion edges over a
horizontal link in plane $j$ belong to its layers $2,3,4,5$.

The chain algebra consists of continuous functions of finitely
many $\widetilde u_j(b)$ with local Grassmann polynomials.
In it the field equations are \eqref{C-386-11} with the
temporal-gauge operator $\widetilde D$. The link equation at
horizontal $S$ in plane $j$ resamples $\widetilde u_j|_S$
with the even Grassmann kernel
\begin{equation}\label{eq:qcd-chain-kernel}
 \widetilde k_{S,j}(v)=
 k_S(\widetilde u_{j-1}|_S,v)e^{-\beta V_S(v\vee\widetilde u_j)}
 k_S(v,\widetilde u_{j+1}|_S)e^{-S_F^{S,j}(v)}.
\end{equation}
Here $V_S$ sums the horizontal plaquettes containing $S$, and
$S_F^{S,j}$ is minus the sum of the oriented fermion hops over
those links. Denote its normalized integration by
$\widetilde\gamma_{S,j}$. Its denominator is invertible, having
body at least $e^{-8\beta\#S}$.

The functional $\widehat\omega=\omega_\nu\circ\mathrm{tg}^*$
obeys these chain equations and is bounded on each fixed finite
Grassmann sector. Indeed temporal gauge intertwines $D_U$ and
$\widetilde D$ by the colour matrices $\widehat w_j$.
The chain rule for left derivatives gives the same intertwining
for the field equations. Resampling a horizontal link does not
change $\widehat w_j$; Haar substitution
$v(b)=\widehat w_j(b_-)V(b^{(j)})\widehat w_j(b_+)^{-1}$
therefore intertwines its link equation with
\eqref{eq:qcd-chain-kernel}. The sector bounds follow from
the fixed-degree Wick bound and bounded colour substitutions.

\begin{proposition}[Uniqueness from two layers]\label{C-386-58}
Two locally bounded functionals on the chain algebra satisfying
the field and link equations agree everywhere if they agree on
gauge planes $t_0,t_0+1$ and thin layers
$X_{t_0,6}\cup X_{t_0+1,1}$.
\end{proposition}

\begin{proof}
Subtract the functionals and call the difference $\delta$.
It obeys the homogeneous chain equations. There are two
propagation steps.

\emph{Adding a fermion site.} Suppose the null algebra includes
a neighbour $x'$ of a new site $x$, the three other neighbours
of $x'$, and the coordinates read by their transports.
The hopping matrix at the missing edge is invertible,
\[
 w(x',x)^{-1}=\frac2{w_i(r^2-1)}(r+\gamma(\widehat{x-x'})).
\]
Solve the field equation at $x'$ for $\psi(x)$, multiplying
by the inverse hopping and the inverse transporter; do the same
on the right for $\bar\psi(x)$. After multiplication by a
monomial and application of $\delta$, the field equations
replace $\widetilde D\psi$ and $\bar\psi\widetilde D$ by
left derivatives. Each term has one fewer generator at $x$.
Induction on that number proves nullity after adjoining $x$.

\emph{Adding a gauge plane.} Let $\epsilon=\pm1$.
Suppose the null algebra contains $\widetilde u_j|_S$,
$\widetilde u_{j-\epsilon}|_S$, all horizontal plaquette
coordinates in plane $j$ needed by $V_S$, and the endpoints
of the fermion hops over $S$ in that plane; the new coordinates
are $\widetilde u_{j+\epsilon}|_S$. Factor
\[
 \widetilde k_{S,j}(v)
       =k_S(v,\widetilde u_{j+\epsilon}|_S)\Phi_\epsilon(v),
\]
where $\Phi_\epsilon$ is even, invertible and does not read
the new plane. For $\chi$ in the old null algebra and free of
$\widetilde u_j|_S$,
resampling
$G=\varphi(\widetilde u_j|_S)\chi\Phi_\epsilon^{-1}$
gives
\[
 \widetilde\gamma_{S,j}G
       =Z_{S,j}^{-1}(K_S\varphi)
                     (\widetilde u_{j+\epsilon}|_S)\chi.
\]
Its value under $\delta$ is zero by the link equation.
Since $K_S$ maps trigonometric polynomials onto themselves,
Peter--Weyl density \cite[Thm.~37.5(ii)]{Etingof} and local boundedness imply
$\delta(Z_{S,j}^{-1}H)=0$ for every $H$ in the new coordinates
and the old coordinates other than $\widetilde u_j|_S$.
For an arbitrary element of the enlarged algebra, its resampling
is of this form. Its value under $\delta$ is therefore zero.

To check exhaustion, write $[p,q;a,b]$ for gauge planes
$p,\ldots,q$ and thin layers $\mathbb L_a,\ldots,\mathbb L_b$.
A null interval extends by a thin layer at its upper end if
that end has $k=1$ or $6$, or if its gauge plane lies in $[p,q]$;
the same holds at its lower end. This is precisely the first
step and the thin-layer geometry above. It extends by a gauge
plane above $q$ once its thin layers $2,\ldots,5$ are included,
and similarly below $p$, by the second step. Starting from
$[t_0,t_0+1;6t_0+6,6t_0+7]$, obtain
\[
 [t_0,t_0+1;6(t_0-1)+5,6(t_0+2)+2].
\]
From
$[t_0-s,t_0+1+s;6(t_0-s-1)+5,6(t_0+s+2)+2]$
add the two outside gauge planes, then six thin layers on
each side, to obtain the interval for $s+1$.
These intervals exhaust every local element.
\end{proof}

\begin{proposition}[Bond-reflection invariance]\label{C-386-59}
Every periodic Gibbs law is invariant under every coordinate
bond reflection. Moreover
\[
 \mathcal G^{\rm per}
       =\{\nu\in\mathcal G:\nu\text{ is }2\mathbb Z^3\text{-invariant}\},
\]
so this set is nonempty, convex and weakly compact.
\end{proposition}

\begin{proof}
Let $\nu_1=\nu$ and $\nu_2=\theta_*\nu$ for a bond mirror.
The reflected law is again periodic Gibbs, by reflection covariance
of the specification. Pull both functionals back by the same
temporal gauge. On the two-layer algebra in
\cref{C-386-58}, its pullback is an element $Y$ of the physical
two-layer algebra. Reflection covariance of the Gaussian functional
and \cref{C-386-54} give
\[
 (\omega_{\nu_2}\circ\mathrm{tg}^*)(G)
     =\overline{\omega_\nu(\vartheta Y)}
     =\omega_\nu(Y)
     =(\omega_{\nu_1}\circ\mathrm{tg}^*)(G).
\]
Both pulled-back functionals obey the chain equations, so
\cref{C-386-58} identifies them everywhere. For a gauge-invariant
local scalar $F$, evaluate it on temporal-gauge configurations
to get $f$ with $\mathrm{tg}^*f=F$; hence $\nu_1(F)=\nu_2(F)$.
Both laws are gauge invariant, so averaging any local scalar
over its finite gauge variables identifies them on all local
continuous functions. Thus $\nu=\theta_*\nu$.
The product of reflections in consecutive bond mirrors is a
translation by two. Every periodic solution is therefore
$2\mathbb Z^3$-invariant; the converse is immediate. The right
side is a weakly closed convex subset of the compact Gibbs
class of \cref{C-386-27}, and contains a cubic-torus limit.
\end{proof}

\subsection{Adapted tori and comparison states}

Fix a mirror $P=\{n\cdot x=c\}$, with
$n\in\{e_\mu,e_\mu\pm e_\nu\}$, and put $q=|n|^2$.
Take orthogonal integer vectors $f_0=n,f_1,f_2$: for a coordinate
normal use the remaining axes; for $n=e_\mu\mp e_\nu$ use
$f_1=e_\mu\pm e_\nu$ and the remaining axis as $f_2$.
For fixed $N'$ and integer $N\ge2$, put
\[
 a=2NN',\quad h=qN',\quad
 \Gamma=\bigoplus_{j=0}^2a\mathbb Zf_j,\quad
 \Gamma'=\bigoplus_{j=0}^2N'\mathbb Zf_j,\quad
                         \chi(af_j)=-1.
\]
The torus has systole $a$ and normal height period $qa$; the
parallel mirrors $c_j=c+jh$ split it into $2N$ slabs.
Write $g_1=1$, $g_{j+1}=\theta_{n,c_j}g_j$.
Then $g_j$ takes the first slab to the $j$th and
\[
 \theta_{n,c_m}g_j=g_{2m+1-j},\quad
 g_{j+2N}=\tau_{af_0}g_j.
\]
For odd $j$, $g_j=\tau_{(j-1)N'f_0}$; for even $j$ this
translation is $g_j\mathrm{fl}$, where
$\mathrm{fl}=\theta_{n,c+h/2}$ reflects the first slab.

The finite-domain positivity proof in the preceding section
applies to these tori and slabs. Represent the generators in one
half-period slab; at the far boundary, the normal reverses and
the character $\chi(af_0)=-1$ reverses the crossing sign.
Both boundaries therefore have positive crossing matrices.
The gauge plaquette and Haar factorizations are the same as in
\cref{C-391-13}. In the present order convention this gives
$Z_{\mathbb T}[\vartheta(F)F]\ge0$ for every local slab
element $F$, including odd ones.

Write $\Gamma_\parallel=\Gamma\cap n^\perp$ and
$\Gamma'_\parallel=\Gamma'\cap n^\perp$.
For a real local $A$ supported in the first slab define
\[
 \mathbb A=\sum_{v\in\Gamma'/\Gamma}\tau_vA,\qquad
 \mathbb A^{\rm f}=\sum_{j=1}^{2N}
       \sum_{v\in\Gamma'_\parallel/\Gamma_\parallel}g_j\tau_vA,
 \qquad M=(2N)^3.
\]
The folded law $\mu_N^{\rm f}=e^{t\mathbb A^{\rm f}}
\mu_{\mathbb T}/\mu_{\mathbb T}(e^{t\mathbb A^{\rm f}})$ satisfies
\begin{equation}\label{eq:qcd-fold-comparison}
 \mu_{\mathbb T}(e^{t\mathbb A})
      \le\mu_{\mathbb T}(e^{t\mathbb A^{\rm f}}),\qquad
 \log\mu_{\mathbb T}(e^{t\mathbb A^{\rm f}})
                                      \le tM\mu_N^{\rm f}(A).
\end{equation}
Here are the finite-array details behind the first inequality.
For slab functions $f_j$, write
\[
 \mathrm{cy}(f_1,\ldots,f_{2N})=\mu_{\mathbb T}(\prod_jg_jf_j).
\]
Cauchy--Schwarz at each wall bounds its squared modulus by the
two arrays obtained by copying either half into the other. Their
values are nonnegative. Normalize each $f$ by
$\|f\|=\mathrm{cy}(f,\ldots,f)^{1/(2N)}$ when this is positive.
Among the finite arrays with entries taken from a fixed list,
maximize the modulus after this normalization. Both half-reflections
of a maximizer remain maximizers. Choose one with a longest
constant cyclic string. Reflecting at its end doubles that string
if its length is at most $N$, and makes it constant otherwise.
Thus the maximum is a constant array and equals one. It follows that
\[
 |\mathrm{cy}(f_1,\ldots,f_{2N})|\le\prod_j\|f_j\|.
\]
For the present application take
$f=\exp(t\sum_v\tau_vA)$ and $f'=\mathrm{fl}f$;
their norms are strictly positive and equal, since translation
by $N'f_0$ carries one full fold to the other.
The unfurled array alternates $f,f'$, and the folded array is
constant. This proves the first inequality in
\eqref{eq:qcd-fold-comparison}.

The folded law is invariant under all walls $c_j$ and under
$\Gamma^{\rm f}=\Gamma'_\parallel\oplus2N'\mathbb Zf_0$:
the displayed identities for $g_j$ permute its terms. Each term
therefore has expectation $\mu_N^{\rm f}(A)$. Finally Jensen's
inequality applied to
$\mu_N^{\rm f}(e^{-t\mathbb A^{\rm f}})
  =\mu_{\mathbb T}(e^{t\mathbb A^{\rm f}})^{-1}$
gives the second inequality.
The folded law is reflection positive, with fermions, at $P$:
$e^{t\mathbb A^{\rm f}}=\Phi\vartheta\Phi$, where $\Phi$
is the exponential of the terms in one half-period slab.
It is a scalar, so the tilted form is the original positive
form applied to $\Phi F$.

\begin{lemma}[Comparison states]\label{C-386-34}
For $\nu\in\mathcal G^{\rm per}$ and real continuous local
$A$ supported in the upper half of $P$, there is
$\nu_0\in\mathcal G^{\rm per}$ such that
\[
 \nu(A)\le\nu_0(A),\qquad
 \omega_{\nu_0}(\vartheta(F)F)\ge0
\]
for every continuous local upper-half element $F$.
\end{lemma}

\begin{proof}
Choose $N'$ a multiple of a period number of $\nu$, large
enough that the support of $A$ lies in the first slab above $P$.
In the preceding construction put $t=a^{-1/2}$ and choose
site representatives
\[
 \mathcal D=\{y\in\mathbb Z^3:
             0\le f_j\cdot y<a|f_j|^2,\ j=0,1,2\}.
\]
There are $qa^3$ sites. Each fixed value of $f_j\cdot y$
has at most $qa^2$ representatives; therefore
$s(\mathcal D)\le36qa^2$. If the support of $A$ fits in a
cube of side $r_A$, the translates not wholly inside this
domain number at most $9r_Aqa^2$. Combine
\eqref{eq:qcd-lower-tangent} and \eqref{eq:qcd-fold-comparison}:
\[
 tM\nu(A)-2K(\mathcal D)
      -2t\|A\|_\infty(M-M_{\rm in})
       \le tM\mu_N^{\rm f}(A).
\]
Here $M=(a/N')^3$. Dividing by $tM$ makes the errors
$O(a^{-1/2})+O(a^{-1})$. Let $\nu_0$ be a weak limit of
the folded laws. A link belongs to at most $2\#B_A$ folded
support sets, since each odd or even family is an array of
distinct translations. Hence a change on a fixed finite $B$
changes $t\mathbb A^{\rm f}$ by at most
$4t\|A\|_\infty\#B_A\#B\to0$.
The slowly varying weight argument after
\eqref{eq:qcd-vanishing-tilt} proves $\nu_0\in\mathcal G$ and
convergence of its local continuous Grassmann expectations.
The group $\Gamma^{\rm f}$ is independent of $N$ and contains
$4N'\mathbb Z^3$, so $\nu_0$ is periodic. Continuity of $A$
gives $\nu(A)\le\nu_0(A)$. The finite folded forms are positive,
and their local continuous Grassmann convergence proves the
second assertion.
\end{proof}

\subsection{Separation and bounded Borel coefficients}

\begin{proof}[Proof of \cref{C-386-13}]
Fix $P$ and let $\mathcal R$ be the periodic Gibbs laws whose
forms $\omega_\nu(\vartheta(F)F)$ are nonnegative for every
continuous upper-half element. It is a nonempty compact convex
set: each condition is affine and weakly closed because the
fixed-configuration Wick expectation is continuous; nonemptiness
follows from \cref{C-386-34}; compactness follows from
\cref{C-386-59}.

If $\nu\in\mathcal G^{\rm per}\setminus\mathcal R$, there is
a real continuous scalar $f$ and $\delta>0$ with
$\nu(f)\ge\nu'(f)+3\delta$ for all $\nu'\in\mathcal R$.
Indeed finitely many continuous integrals separate the point
from the compact convex set: otherwise the finite-intersection
property would put the point in it. In that finite-dimensional
space the nearest-point separating functional gives $f$.
Approximate $f$ uniformly by a real local $A_0$ within
$\delta$, and translate it by $2kn$ into the upper half of
$P$. All periodic Gibbs laws are $2\mathbb Z^3$-invariant by
\cref{C-386-59}, so the translated $A$ still satisfies
$\nu(A)\ge\nu'(A)+\delta$ on $\mathcal R$.
But \cref{C-386-34} supplies $\nu_0\in\mathcal R$ with
$\nu(A)\le\nu_0(A)$, a contradiction. Thus every periodic
solution is positive at $P$.

Scalar reflection positivity now gives reflection invariance.
The form $s(F',F)=\nu(\overline{F'\circ\theta}F)$ on
continuous upper-half scalar functions is Hermitian, so
$\nu(F)=\nu(F\circ\theta)$ there. Translate an arbitrary
local scalar function far enough into that half by $2kn$,
and use $\theta\tau_{2kn}=\tau_{-2kn}\theta$ and translation
invariance to get the same equality for every local function.
The reflections in all nine mirror families generate
$\mathbb Z^3\rtimes W_3$, proving the full lattice invariance.
Reflection covariance of the Dirac inverse and Wick's rule give
$\omega_\nu(\vartheta G)=\overline{\omega_\nu(G)}$.
The same identity holds with $\Theta_P$, since the two differ
only on odd elements, whose values vanish.

For bounded Borel coefficients, write $F=\sum_jf_j\xi_j$
with finitely many half-side links and monomials. Approximate
$f_j$ by continuous $f_{j,k}$ in $L^2$ of their finite-link
marginal of $\nu$. Reflection invariance gives the same
convergence for $f_{j,k}\circ\theta$. Each Wick kernel
$h_{\vartheta(\xi_i)\xi_j}$ is bounded. Cauchy--Schwarz
therefore shows convergence of every term of the reflected
form. This extends positivity and Hermitian symmetry. Terms
between different parities are odd and vanish.
On the lower half, apply the upper form to $\vartheta F$;
interchanging homogeneous odd factors produces its parity sign.
Changing the normal in $\Theta_P$ produces precisely the same
sign. Thus \eqref{eq:qcd-reflection-order-dictionary} gives
$\omega_\nu(F\Theta_PF)\ge0$ with the normal into either
chosen side. Cubic-torus limits are periodic Gibbs laws by
\cref{C-389-10}, so the last assertion follows.
\end{proof}

\section{The pure gauge input and local expansions}\label{sec:qcd-gauge-input}

We use the finite gauge construction of \cite{YMM1} with the following
conventions and hypotheses. Fix its odd blocking integer $L\ge13$,
an integer $r_c\ge2$, an integer $\mathfrak p_0\ge6r_c+1$,
and a power $M=L^{m'}$. Fix also the positive fermion block
penalty $b_F$, subject to the bounds used below. These choices
precede the cutoff and volume. Here $K\ge1$ and the fine spacing is $a_K=L^{-K}$,
the physical torus has side $\ell=2L^m$, and level $j$ has spacing
$\lambda_j=L^{j-K}$ and $n_j=\ell^3L^{3(K-j)}$ sites.
The calibration and running coefficients are
\[
 \mathsf t_j=L^{K-j}\mathsf b,\qquad
 x_0=\beta=L^K\mathsf b,\qquad
 x_{j+1}=L^{-1}x_j-b_{j+1},\quad |b_{j+1}|\le B_\sharp.
\]
The last bound implies $|x_j-\mathsf t_j|\le\bar B$ for a
fixed $\bar B$. Require $x_j\ge\gamma^{-2}$ and
$MR(x_K)\le L^m$, where
\[
 R(x)=\min\{L^q:q\in\mathbb Z_{\ge0},\ L^q\ge(\log x)^{r_c}\},\quad
 \epsilon(x)=A_0x^{-1/2}(\log x)^{\mathfrak p_0},\quad
 \delta(x)=A_1x^{-1/2}(\log x)^{\mathfrak p_0}.
\]
For the continuum construction fix $0<g\le\gamma$ and take
$\mathsf b$ in the closure of
$(g^{-2}+\bar B,L(g^{-2}+B_\sharp)-\bar B]$.
Then $g^{-2}\le x_K\le L(g^{-2}+B_\sharp)$.
Write $X_g=L(g^{-2}+B_\sharp)$. We take the allowed physical
sides to be the cofinal set $\mathfrak L(g)=\{2L^m:m\ge m_g\}$, where
\[
 m_g=1+m'+\max\left\{
 \left\lceil\log_L R(X_g+\bar B)\right\rceil,
 \left\lceil\log_L C_{\rm s}\right\rceil\right\}.
\]
Here $C_{\rm s}\ge2$ is the fixed chart-room constant of
\cite[Cor.~69.4 and its preceding torus-room calculation]{YMM1}.
This is the room for the larger local charts as well as the
finite induction. Omitting finitely many sides does not alter
the class of infinite-volume torus limits.
The fermion mass is $m_0=L^{-K}\mathsf m$, with
$\mathsf m\in[m_-,m_+]$ and $0<m_-\le m_+<1/2$.

We use the threshold and growth functions
\[
 \mathfrak b(t)=t^{1/2}\delta(t),\qquad
 \phi(t)=t^{-1/2}(1+\log t)^{\mathfrak p_0},\qquad
 \widehat\phi(t)=(1+\log t)^{\mathfrak p_0},
\]
\[
 \phi_\sharp(t)=t^{-1/2}(1+\log t)^{3\mathfrak p_0},\qquad
 \psi_\sharp(t)=t^{-1/2}(1+\log t)^{9\mathfrak p_0},\qquad
 \mathfrak m_j=L^{j-K}\mathsf m.
\]

The blocks of one step are the cubes of $L^3$ sites centered
at the next lattice. The gauge map $\mathcal M$ is the global
real block average of \cite[Equations~(37)--(38)]{YMM1}.
On its small-curvature domain it is the symmetrized contour
mean: for a coarse bond $c$ its fine contours start at $c_-$,
follow the straight fine path from $x$ to $x+L\hat\mu$, and end at
$c_+$, with $x\in B(c_-)$. The symmetrized path transports
give matrices $W_x$ near one, and
\[
 \mathcal M(U)(c)=
 \exp\left[L^{-3}\sum_{x\in B(c_-)}\log W_x\right]U[c].
\]
The continuation is used only on the logarithm branch specified
there. Outside the small-curvature domain use the Borel real
extension of \cite[Prop.~4.5 and its preceding definition]{YMM1}.
This fixes the global map as well as its analytic branch;
it does not choose a branch separately in each history.

At each level let $\mathcal P_j$ be the partition into cubes of side
$M$ in level-$j$ units. A polymer $X\in\mathcal D_j$ is a
nonempty wall-connected union of these cubes. Its size $d_j(X)$
is the least number of edges in a tree formed from edges of the
cubes of $X$ and meeting every closed cube of $X$. A single
corner is allowed as a tree. We use the entropy and joining bounds,
in particular
\[
 \sum_{X\ni\square}e^{-c_Ed_j(X)}\le C_E,
 \qquad
 d_j(X_1\cup X_2)\le d_j(X_1)+d_j(X_2)+c_0
\]
for intersecting polymers, with the fixed geometric constant
$c_0$. The closed enlargement $X^\sharp$ has the reach fixed
by that construction. Parameters on cubes outside this
enlargement localize a kernel; a finite difference in a cube
set $\mathcal K$ is assigned to $X^\sharp\cup\bigcup\mathcal K$.

In level-$j$ units, a regular pair $(\mathbf U,\mathbf J)$ is
represented by a triple $(U,A',\mathbf J)$, with $U$ real and
$\mathbf U=e^{iL^{-j}A'}U$ complex. It belongs to
the analytic core $\mathcal U_j^\circ(X;b_0,b_1)$ of
\cite[Def.~49.1]{YMM1}: its curvature and current have the
indicated small bounds, its local real gauge potentials have
the stated bounds, and every required reduced block tower
exists on the prescribed logarithm branch and satisfies the
critical-background bounds at every earlier level.
The core is not defined by curvature alone. Two representatives
of the same pair have the same $(\mathbf U,\mathbf J)$;
their real reference configurations may differ. At radius
$a\le\alpha_0$ use
$(b_0,b_1)=(a,(\alpha_1/\alpha_0)a)$ at all levels.
In a restriction from level $j'$ to level $j\le j'$, put
$\lambda=L^{-(j'-j)}$. The forms $A'$ and $\mathbf J$ scale by
$\lambda$ and $\lambda^3$, respectively, while $\mathbf U$ is
unchanged, as in \cite[Equation~(757)]{YMM1}.

For a regular coarse field $V$, $U_j(V)$ denotes the admitted
small-curvature critical gauge orbit of Wilson's action on
the fibre $\mathcal M^jU=V$. At a representative of this
orbit the shift system $\mathcal A(B;P)$ is the
holomorphic coordinate change for the gauge fluctuation:
the fine field read by a step is
$U'=U\exp\mathcal A(B;P)$, with the associated site gauge
factor. Each term of a kernel expansion is assigned to the
polymer hull containing its path and parameter cubes.
The shift is bounded by
$C_{\mathcal A}\delta(\mathsf t_j)$ on the fluctuation
polydisc and has the exponential tails across
the layers. We retain the site gauge factor in the
fermion averages; gauge covariance carries it to their
generators.

The following results of \cite{YMM1} are the external pure gauge
input. The verification in \cref{C-396-12}, using
\cref{def:qcd-gauge-admission}, supplies the hypotheses of its
finite-renormalization theorem (D) on the sequences used here.
It gives the critical backgrounds, shift systems, histories,
coordinate maps, positive gauge-fixing weights, retained
forms and pivot determinants; its local terms are
holomorphic on the cores, local and gauge covariant, with
exponential polymer bounds. Its history decomposition is
linear in a bounded Borel density, coefficientwise for a
finite-dimensional algebra-valued density.
We use \cite[Thm.~V and Def.~52.1]{YMM1}: holomorphic local Gaussian
kernels and interaction, the real-reference bounds, the
gauge covariance and independence of the triple representing
a given pair $(\mathbf U,\mathbf J)$, the polymer
bounds and its stated numerical inequalities. The output
has rate $\kappa_+=(1+4\beta_s)\kappa$. The quadratic
term caused by changing the gauge Gaussian coefficient,
denoted by q in the expansion, is kept as a separate scalar family.
We also use \cite[Prop.~51.16]{YMM1}: at fixed radius the change of an old array gains
the size of the shift, and at radius $a$ its bound is
\[
 \frac65K_\Lambda\frac{C_{\mathcal A}\delta(\mathsf t_k)}a
 \sum_{i\le k}L^{-(k-i)}\epsilon_i\,
 e^{-(1-2\delta_w)\kappa d_k(Y)}.
\]
Linearity permits explicit normalizer families and small
fluctuation families to be localized separately.
These assertions include their original existence, domain
and numerical hypotheses; the constants are chosen in
their admitted order. The matter-specific uses of them
are verified below.

\subsection{The centered conditional integral}

On a history, some gauge variables remain outside the integration
region. The conditional Gaussian therefore has a nonzero mean.
The expansion must be estimated about that mean: a bound obtained
by taking the absolute value of an uncentered density would introduce
an extensive error depending on the external field.

We recall the conditional-integral estimate of
\cite[Prop.~75.1, Lem.~75.2 and Prop.~75.3]{YMM1}.
Let $(C,\Gamma,\mathbf m)$
be a holomorphic Gaussian system satisfying its mean, error,
geometry and tested-bond conditions. Suppose $F(B;n)$ has the
required holomorphy in $B$ on the prescribed tube and is
holomorphic in the center $n$ on its prescribed domain for
each real $B$. If
\[
 |F(B;n)|\le F_0\exp\{\rho_+\|B-n\|^2\},
 \qquad \rho_+=(LM)^{-3},
\]
then the centered conditional integral satisfies
\begin{equation}\label{eq:qcd-centered-contour}
 |J^{\mathbf m}[F(\cdot;\mathbf m)]|
 \le F_0 e^{c_P^+N_\Lambda(X)}e^{-\gamma_P\#e^+},
 \qquad c_P^+=c_P+C_{\rm mn}^3/\lambda_r.
\end{equation}
Here $N_\Lambda(X)$ counts the integration cubes and $e^+$ is
the failed-bond set of the conditional integral. A finite
difference in a set $\mathcal K$ of parameter cubes gains
$e^{-(\bar\kappa-1)\#\mathcal K}$. The estimate uses the
actual complex Gaussian system at the point of evaluation.
In the contour proof its centered argument is $\beta'-i\eta$,
with $\|\eta\|^2\le16d_*^{-2}\epsilon_m^2\|\beta'\|^2$.
Thus the bound contains no factor depending on the size of
the conditional mean.

\begin{lemma}[Algebra-valued conditional integration]
\label{lem:qcd-centered-interface}
The estimate \eqref{eq:qcd-centered-contour}, including its
finite-difference version, holds for a holomorphic integrand
with values in a finite-dimensional normed algebra, with its
norm in place of the absolute value. A subsequent Gaussian
Berezin integration preserves the pointwise norm estimate
whenever it is a contraction for the fibre norm on the whole
parameter domain. The joined finite-difference and component
statements additionally require the holomorphy, locality and
component restrictions specified in \cref{C-399-41}.
\end{lemma}

\begin{proof}
Apply the scalar conditional-integral theorem to
$\lambda(F)$ for each norm-one linear functional $\lambda$.
Its growth bound is the same as that of $F$, and the contour
integral commutes with $\lambda$. Taking the supremum gives
the norm bound, with no factor for the dimension of the algebra.
The same argument applies before taking finite differences.
The pointwise Berezin extension follows from its uniform
contraction bound; the parameter-dependent joined expansion
is treated in \cref{C-399-41}. In the even Grassmann algebra, products
commute and the fibre norm is submultiplicative; hence the
scalar Mayer majorants also bound the corresponding products.
\end{proof}

This functional estimate is only one input to a history step.
The gauge domain, the fermion covariance, the origin assigned
to each quadratic term, and the identities of the actual density
must also be checked. We make those checks in
\cref{C-399-32,C-399-41,C-399-43,C-399-44}.

\section{Fermion block integration}\label{sec:qcd-block-integration}

We now express a fermion block integral as a Schur complement
and a determinant. Coercivity gives the inverse needed in this
identity; the fibre norm controls its Grassmann coefficients,
and the normalized block average fixes the Gaussian block kernel.

Throughout this section the fermion is the balanced chamber operator
with $r=\frac12$, and $N_f\ge1$ is fixed.  Its fibre dimension is
$\nu=8N_f$.  The maps $\iota\varphi(x,g)=48^{-1/2}\varphi(x)$,
$\Pi_0=\iota\iota^*$ and $\Pi_1=1-\Pi_0$ distinguish the chamber mean
from the other components.  In this section $D[U]$ is massless:
it is the foundational operator $D_U-m_0$.
Write $\Delta[U]=\operatorname{Re}D[U]$ and $\theta=b_F>0$.
For an oriented site bond $b=\langle x,y\rangle$ set
$\nabla_U\varphi(b)=U_{\langle x,y\rangle}\varphi(y)-\varphi(x)$;
its squared norm is summed over both orientations of each bond.
Restricted gradient norms use both orientations of bonds with both
endpoints in the indicated set.
The lattice operator satisfies
\[
 \Delta\ge0,\quad D^*D\le c_{\mathrm{disc}}\Delta,
 \quad c_{\mathrm{disc}}=9(6+6\sqrt2),\qquad
 \langle f,\Delta f\rangle\ge
 \frac{\lambda_\rho}{2}\|\Pi_1f\|^2+
 c_\Delta\|\nabla_U\iota^*f\|^2,
 \quad c_\Delta=\frac{\lambda_\rho}{96}.
\]
Here $\lambda_\rho>0$ is the first nonzero eigenvalue of the
within-site chamber Laplacian.  Every constant below is independent
of the cutoff, volume and history unless a dependence is displayed.

Here are proofs of these elementary bounds.  Write
$D=\sum_sh_s(1-R_s)$, where $R_s$ is a self-adjoint
involution, $h_s=t_s(r-\gamma_s)$ and
$\Delta=\sum_s\rho_s(1-R_s)$, $\rho_s=rt_s>0$.
Thus $\|(1-R_s)f\|^2=2\langle f,(1-R_s)f\rangle$.
Weighted Cauchy--Schwarz gives
\[
 \|Df\|^2\le
 2\Bigl(\sum_s\frac{t_s^2(1+r)^2}{\rho_s}\Bigr)
                    \langle f,\Delta f\rangle
 =c_{\rm disc}\langle f,\Delta f\rangle.
\]
For the balanced chamber weights
$\sum_st_s=6+6\sqrt2$ and $r=1/2$, giving the displayed
constant.  In particular $\|D\|\le c_{\rm disc}$ and
$0\le\Delta\le c_{\rm disc}$.
The three within-site reflections generate the chamber group.
Their graph Laplacian has constants as its kernel, and
its least positive eigenvalue is $\lambda_\rho$; hence
$\sum_{s\ge1}\rho_s(1-R_s)\ge\lambda_\rho\Pi_1$.
For $f=\iota\varphi+f'$, $f'=\Pi_1f$, the remaining hop gives
\[
 \langle f,(1-R_0)f\rangle
 \ge\tfrac14\|(1-R_0)\iota\varphi\|^2
                      -\tfrac12\|(1-R_0)f'\|^2
 \ge\tfrac1{24}\|\nabla_U\varphi\|^2-2\|f'\|^2.
\]
Each of the six signed coordinate directions occurs eight times
among the 48 chambers. With both bond orientations in the norm, this gives
$\|(1-R_0)\iota\varphi\|^2=\tfrac16\|\nabla_U\varphi\|^2$
and hence the coefficient $1/24$. This hop form is also
nonnegative.  Take its fraction
$\tau=\lambda_\rho/(4\rho_0)\le1$ and add the within-site
bound to obtain the stated coercivity.
Finally, for every real $m\ge0$,
\[
 (D+m)^*(D+m)\le(c_{\rm disc}+m)(\Delta+m),
\]
because the difference after using $D^*D\le c_{\rm disc}\Delta$
is $m(c_{\rm disc}-\Delta)\ge0$.  The massive operator
therefore belongs to the disc of diameter $c_{\rm disc}+m$
and has positive definite real part when $m>0$.

\subsection{The fibre norm and Gaussian integration}

Let $P$ be a finite set of places, with independent barred and unbarred
copies of the fibre at each place.  For positive weights $h(s)$, equal
on the two copies, define
\[
 \|F\|_h=\inf\left\{\sum_k|a_k|\prod_lh(s_{kl}):
 F=\sum_ka_ke_{k1}\cdots e_{kn_k},\quad
 e_{kl}\text{ a unit vector in a fibre at }s_{kl}\right\}.
\]

\begin{lemma}[The fibre norm]\label{C-387-18}
The even Grassmann algebra with this norm is a commutative unital
Banach algebra.  The norm is unchanged by fibrewise unitary
substitutions.  If $T_{s's}$ are the blocks of a linear substitution
and $\sum_{s'}h'(s')\|T_{s's}\|\le h(s)$, then
$\|F\circ T\|_{h'}\le\|F\|_h$.
For Berezin integration over a subset $P_0$,
\[
 \|\mathrm{Ber}_{P_0}F\|_h\le
     \prod_{s\in P_0}h(s)^{-2\nu}\|F\|_h.
\]
If $\mu_C$ is Gaussian Grassmann integration with covariance $C$ on
$P_0$ and $h\ge\|C\|^{1/2}$ there, then
$\|\mu_CF\|_h\le\|F\|_h$.  More generally,
\[
 \left\|\int F(\Psi+\Psi')\,d\mu_C(\Psi')\right\|_h
 \le\|F\|_{h+\|C\|^{1/2}\mathbf1_{P_0}}.
\]
\end{lemma}

\begin{proof}
Concatenating representations proves submultiplicativity.  If
$|F|_h$ denotes the coefficient norm in an orthonormal basis, then
$\|F\|_h\le|F|_h\le\|F\|_{\sqrt\nu h}$: the second inequality
follows by expanding unit vectors, whose coefficient sums are at most
$\sqrt\nu$.  In particular the infimum defines a norm.  The even
algebra is commutative and finite dimensional, hence complete.
A fibrewise unitary carries the defining unit vectors to unit vectors;
applying the inverse gives equality of the norms.
For a unit vector $e$ at $s$, write $Te=\sum_{s'}T_{s's}e$ and apply
the assumed weighted row bound.  Multiplication and the defining
infimum give the substitution estimate.

Berezin integration of a product of unit vectors vanishes unless
there are $2\nu$ factors at each integrated place.  Its top coefficient
is a determinant of unit columns, with modulus at most one by
Hadamard's inequality.  This removes exactly the weights displayed
in the estimate.
Gaussian integration contracts $m$ unbarred and $m$ barred factors
to $\det[\langle a_i,Cb_j\rangle]$.  Its modulus is at most
$\|C\|^m$, by the exterior-product form of Gram's inequality;
unbalanced products have zero integral.  The integrated weights pay
this factor.  Finally, expand each factor of
$F(\Psi+\Psi')$ into its two terms.  Summing the bounds for the
subsets taken at $\Psi'$ replaces each weight $h(s)$ by
$h(s)+\|C\|^{1/2}\mathbf1_{P_0}(s)$.
\end{proof}

We normalize the Berezin integral so that
$\mathrm{Ber}\,e^{-\bar\Psi A\Psi}=\det A$.  At invertible $A$,
Gaussian moments are the minors of $A^{-1}$.  Integration in disjoint
sets of barred--unbarred pairs commutes; each such pair is an even
operation.  Thus all the following identities hold coefficientwise
also for measures with Grassmann values.  In finite dimension this
last assertion follows by applying scalar integration to each
coefficient.

\subsection{Averages defined at every gauge field}

For a block $B(y)$ and its contour families $\mathcal G(y,x)$ put
\[
 \begin{aligned}
 S(V;y,x)&=\frac1{\#\mathcal G(y,x)}\sum_{\Gamma\in\mathcal G(y,x)}V(\Gamma),\\
 n(V;y)&=L^{-3}\sum_{x\in B(y)}\det S(V;y,x),\qquad
 q(V;y,x)=n(V;y)^{-1/2}S(V;y,x).
 \end{aligned}
\]
\[
 (Q(V)\varphi)(y)=L^{-3/2}\sum_{x\in B(y)}q(V;y,x)\varphi(x).
\]
The contours run through the coordinate directions in all orders;
on the coordinate lines through the block centre there is one contour.

\begin{lemma}[The normalized block average]\label{C-387-24}
At every real gauge field,
\[
 (3L-2)L^{-3}\le n(V;y)\le1,\qquad Q(V)Q(V)^*=1.
\]
The average is block local, Borel, gauge covariant and covariant under
the lattice isometries preserving the block partition.  It preserves
the spinor boundary signs.  For
$R_k=Q(\mathcal M^{k-1}U)\cdots Q(U)\iota^*$ one has
$R_kR_k^*=1$ and $R_k\Pi_1=0$.
\end{lemma}

\begin{proof}
A real linear combination of $\mathrm{SU}(2)$ matrices is a real
quaternion.  Hence $SS^*=(\det S)1$ and $0\le\det S\le1$ for the
average of unit quaternions.  At the $3L-2$ sites on the three
coordinate lines, $S$ is a single transporter and its determinant is
one.  This proves both bounds on $n$.  Since the blocks are disjoint,
\[
 (QQ^*)(y,y')=\delta_{yy'}L^{-3}n(V;y)^{-1}
                      \sum_x\det S(V;y,x)=\delta_{yy'}.
\]
Path transporters are gauge covariant, and an isometry permutes the
contours.  The formulas are local and Borel; translating a whole
block by a period gives the same formula with the prescribed spinor
sign.  The assertions for $R_k$ follow by composition and
$\iota^*\iota=1$, $\iota^*\Pi_1=0$.
\end{proof}

\begin{lemma}[The first fluctuation at an arbitrary field]\label{C-387-29}
On a block $B(y)$,
\[
 \|\varphi\|_B^2\le2|Q(V)\varphi(y)|^2+
             12L^3(L-1)\sum_{b\subset B}|\nabla_V\varphi(b)|^2.
\]
Here the explicit sums over $b\subset B$ count each bond once.
Consequently the first fluctuation operator, on any union of
blocks and with zero continuation off that union, is
$\Gamma_0=(D[V]+m_0)_{AA}+\theta q_1^*q_1$,
where $A$ is the set of chambers over the union and
$q_1=L^{-1/2}Q(V)\iota^*$. It satisfies
\[
 \operatorname{Re}\Gamma_0\ge m_0+c_{\mathrm{fl}},\qquad
 c_{\mathrm{fl}}=\min\left\{\frac\theta{2L},
                    \frac{c_\Delta}{6L^3(L-1)}\right\}>0.
\]
These assertions hold for every real $\mathrm{SU}(2)$ configuration.
Let $c_{\mathrm{hop}}=9/2$ be the bound on the site-changing
hopping coefficient, and put
$\delta_0=c_{\mathrm{fl}}/[4e(c_{\mathrm{hop}}+6\theta)]$.
For every set $Y$ of chambers, the compression $\Gamma_{0,Y}$ obeys
\[
 |\Gamma_{0,Y}^{-1}(v,v')|
 \le\frac2{c_{\mathrm{fl}}}
       e^{-\delta_0|x_v-x_{v'}|_1}.
\]
The same estimate holds on the mirror tori used below.
\end{lemma}

\begin{proof}
Let $\delta^2=\sum_{b\subset B}|\nabla_V\varphi(b)|^2$ and
$c_L^2=3(L-1)/2$.  Telescoping along a contour gives
$|V(\Gamma)\varphi(x)-\varphi(y)|\le c_L\delta$.
Average this identity over the contours.  Since $n\le1$,
\[
 |Q\varphi(y)|\ge L^{3/2}(|\varphi(y)|-c_L\delta),\qquad
 |\varphi(x)|\le|\varphi(y)|+c_L\delta.
\]
Hence
$\|\varphi\|_B^2\le
L^3(L^{-3/2}|Q\varphi(y)|+2c_L\delta)^2$,
which is the first inequality.  Write $f=\iota\varphi+f'$.
The chamber coercivity has coefficient $\lambda_\rho/2$ on
$\|f'\|^2$ and $2c_\Delta$ on the internal, once-counted bond
differences of $\varphi$. Since $q_1f=L^{-1/2}Q\varphi$,
summing the block inequality gives
\[
 \|f\|^2\le c_{\mathrm{fl}}^{-1}
          (\langle f,\Delta f\rangle+\theta\|q_1f\|^2).
\]
Adding $m_0\|f\|^2$ proves the fluctuation floor.

For a real weight $\eta$ changing by at most $\delta_0$ across
a bond, twisting $D$ changes its norm by at most
$c_{\mathrm{hop}}(e^{\delta_0}-1)$. Write
$q_1^*q_1=L^{-1}\iota P\iota^*$, where $P$ is the sum of
the block projections. The variation of $\eta$ from a block
centre is at most $3L\delta_0/2\le1$, so
\[
 \|e^\eta Pe^{-\eta}-P\|\le e^{3L\delta_0}-1
                         \le6eL\delta_0.
\]
The two changes together cost at most $c_{\mathrm{fl}}/2$.
Compression commutes with multiplication by $e^\eta$; hence
the twisted compressed operator has real part at least
$m_0+c_{\mathrm{fl}}/2$ and inverse norm at most
$2/c_{\mathrm{fl}}$. Take
$\eta(x)=\delta_0\operatorname{dist}_1(x,x_{v'})$ to obtain
the kernel bound. The argument is local to blocks and walls,
so it also applies on every mirror torus.
\end{proof}

Let $q_1=L^{-1/2}Q(U)\iota^*$ and
$q_{j+1}=L^{-1/2}Q(\mathcal M^jU)$ for $j\ge1$.
They satisfy $q_{j+1}q_{j+1}^*=L^{-1}$.
The normalized kernel of one step is
\[
 \mathrm{Ker}_j=\theta^{-\nu n_{j+1}}
 \exp\{-\theta\langle\bar\Psi_{j+1}-\bar\Psi_jq_{j+1}^*,
                         \Psi_{j+1}-q_{j+1}\Psi_j\rangle\}.
\]
Its integral in the new variables is one: translate these variables
by $q_{j+1}\Psi_j$ and use the Gaussian determinant.  Gauge covariance
is immediate from that of $q_{j+1}$ and the contragredient barred
action.  Thus the transformation
$\mathbf T_F\nu=\mathcal M_*(\mathrm{Ber}_j(\mathrm{Ker}_j\nu))$
commutes with full Berezin integration.  This exact identity is made
before the small- and large-field tests.

\begin{lemma}[Composition of Gaussian block kernels]\label{C-387-30}
Two normalized block kernels with parameters $(b,Q)$ and $(b',Q')$,
where $Q'Q'^*=c'1$, compose to the normalized kernel with average
$Q'Q$ and parameter $b''$ given by
\[
 (b'')^{-1}=(b')^{-1}+c'b^{-1}.
\]
In particular the parameter after $k$ steps is
\[
 \theta_k=\frac{b_F(1-L^{-1})}{1-L^{-k}},\qquad
 \mathfrak a_k=L^{-k}\theta_k.
\]
\end{lemma}

\begin{proof}
Translate the intermediate variables by $Q\Psi$.  For
$w=\Psi''-Q'Q\Psi$ their Gaussian integration leaves the quadratic
form
\[
 b'\bar w\{1-b'Q'(b+b'Q'^*Q')^{-1}Q'^*\}w
       =\frac{bb'}{b+b'c'}\bar ww.
\]
The identity
$Q'(b+b'Q'^*Q')^{-1}Q'^*=c'/(b+b'c')$
also gives the determinant
$b^{\nu(n'-n'')}(b+b'c')^{\nu n''}$, which converts the two
normalizations into $(b'')^{-\nu n''}$.
Iterating with $c'=L^{-1}$ gives the displayed geometric sum.
\end{proof}

Set
\[
 c_b=\frac{b_F(1-L^{-1})}{b_F(1-L^{-1})+1/2},\qquad
 c_{\mathrm{tw}}=\max\{c_{\mathrm{disc}}+1,b_F\}.
\]
The latter is a uniform tower bound on the mass range used here.
For a site matrix $C$ write
\[
 \mathsf S_1(C\psi)=\sum_v\bar\psi(v)C(x_v)\psi(v),\qquad
 \mathsf S_1(\psi)=\mathsf S_1(1\psi).
\]

\subsection{The Schur complement and its determinant}

For $c>0$ let $\mathsf C(c)$ denote the operator disc
$\{X:X^*X\le c\operatorname{Re}X\}$.
At a step split the old variables into stored variables $S$ and
integrated variables $A$.  If their quadratic form is $\mathbb D_j$,
let
\[
 \Gamma_j=(\mathbb D_j)_{AA}+\theta q_{j+1}^*q_{j+1}.
\]
Write $q=q_{j+1}$ and $\Psi^\bullet=(\Psi_S,\Psi')$. The two
background substitutions on the integrated variables are
\[
 \Psi_A^{\rm bg}=\Gamma_j^{-1}
       (\theta q^*\Psi'-(\mathbb D_j)_{AS}\Psi_S),\qquad
 \bar\Psi_A^{\rm bg}=
       (\theta\bar\Psi' q-\bar\Psi_S(\mathbb D_j)_{SA})\Gamma_j^{-1}.
\]
The stored variables are unchanged. The second substitution is
the row twin of the first. On complex gauge fields the symbol
$q^*$ in these formulas is continued by its holomorphic row
extension $q^+$, which agrees with the adjoint at real fields.
With $W=\Psi_A-\Psi_A^{\rm bg}$ and
$\bar W=\bar\Psi_A-\bar\Psi_A^{\rm bg}$, completing the square gives
\[
 \bar\Psi\mathbb D_j\Psi+
 \theta(\bar\Psi'-\bar\Psi_Aq^*)(\Psi'-q\Psi_A)
 =\bar\Psi^\bullet\mathbb D_{j+1}\Psi^\bullet+
                          \bar W\Gamma_jW.
\]
The row twin above is not the adjoint of the unbarred background map.

\begin{lemma}[Exact factorization]\label{C-397-12}
The tower is formed at every real gauge configuration when $m_0>0$.
Writing $D_0=D[U]+m_0$, one has
\[
 \det D_0=\left(\prod_{i<j}\theta^{-\dim H_{i+1}}
                      \det\Gamma_i\right)\det\mathbb D_j.
\]
If $X^*X\le c\operatorname{Re}X$ and
$\operatorname{Re}X>0$, then $\operatorname{Re}(X^{-1})\ge c^{-1}$.
For $D_0=D[U]+m_0$, it holds with $c=c_{\rm disc}+m_0$.
Writing $c_*:=\max\{c,\theta\}$, the tower forms obey
$\|\mathbb D_j\|\le c_*$, whereas
$\|\Gamma_j\|\le c_*+\theta$.
More precisely, on the accumulated stored variables and the
current variables, write $S_i$ for the one-particle space first
stored at level $i$ and $H_j$ for the current level-$j$ space. Then
\[
 \operatorname{Re}\mathbb D_j^{-1}
 \ge c^{-1}1_{S_0}\oplus
        \theta^{-1}1_{S_1\oplus\cdots\oplus S_{j-1}\oplus H_j}.
\]
\end{lemma}

\begin{proof}
The matrix before eliminating $\Psi_A$ is
$\mathsf T^*(\mathbb D_j\oplus\theta)\mathsf T$ with
$\mathsf T(\Psi,\Psi')=(\Psi,\Psi'-q\Psi_A)$ and
$\det\mathsf T=1$.  Its Schur complement is $\mathbb D_{j+1}$;
therefore
\[
 \det\mathbb D_j=\theta^{-\dim H_{j+1}}
                         \det\Gamma_j\det\mathbb D_{j+1},
 \qquad
 \mathbb D_{j+1}^{-1}=\mathsf q\mathbb D_j^{-1}\mathsf q^*
                                      +(0\oplus\theta^{-1}),
\]
with $\mathsf q$ the identity on all previously stored variables
and $q$ on the integrated variables.  The inverse inequality is the identity
\[
 \operatorname{Re}(X^{-1})-c^{-1}
  =X^{-*}(\operatorname{Re}X-c^{-1}X^*X)X^{-1}.
\]
Suppose the displayed block floor holds at level $j$.
On every stored block $\mathsf q$ is the identity, so its
floor is preserved.  On the new variables the transported
positive term is supplemented by $\theta^{-1}$ times the
identity.  This proves the floor at level $j+1$ even though
$qq^*=L^{-1}$; no isometry of $q$ is being assumed.
The uniform lower bound is consequently $c_*^{-1}$.
The inverse identity makes this equivalent to
$X^*X\le c_*\operatorname{Re}X$, or
$\|X-c_*/2\|\le c_*/2$, for $X=\mathbb D_j$.
It gives $\|\mathbb D_j\|\le c_*$.
A compression preserves this disc because
$\|PXPf\|^2\le\|XPf\|^2$.
The penalty $\theta q^*q$ is nonnegative with norm
$\theta/L\le\theta$; its disc diameter is at most
$\theta$.  Discs add, so
$\Gamma_j\in\mathsf C(c_*+\theta)$.
Moreover the compressed real part of $\mathbb D_j$
is positive definite, so $\Gamma_j$ is invertible.
The Schur complement has positive real part: evaluate the
original positive-real-part form on its background vector.
This supplies the next invertible form and closes the
induction.  The determinant identity then iterates as claimed.
\end{proof}

On the history without large-field regions the closed form is
\[
 \mathcal K_k=D[U]+m_0+\mathfrak a_kR_k^*R_k,
 \qquad
 \mathbb D_k=L^k\mathsf D^{(k)},\qquad
 \mathsf D^{(k)}=\mathfrak a_k-
       \mathfrak a_k^2R_k\mathcal K_k^{-1}R_k^*.
\]

\section{Fluctuation operators in a regular gauge background}\label{sec:qcd-regular-operators}

The block integral requires estimates on its covariance that remain
uniform as the lattice spacing decreases.  The small mass does not
provide such estimates.  The block penalties control the light modes,
while the chamber gap controls the heavy modes.  We first combine these
two estimates, then localize the inverse.  The localization will also
let us integrate over a polymer using only its configuration.

For a history with regions $\Omega_1\supset\cdots\supset\Omega_k$,
put $X_n^i=\Omega_n\setminus\Omega_{n+1}$ for $n<i$ and
$X_i^i=\Omega_i$.  These sets partition $\Omega_1$.  On $X_n^i$
let $\mathsf L_i=L^n$, and let it be one outside $\Omega_1$.
The separation and regularity conditions used below are
\[
 \operatorname{dist}_\infty(\Omega_{n+1},\Omega_n^c)>RL^n,
 \qquad
 |\partial U(p)-1|\le\epsilon L^{-2n}
 \quad\text{if a corner of }p\text{ lies in }X_n^{j+1}.
\]
There is no regularity requirement outside $\Omega_1$.
Here $\partial U(p)=U_p$ is the oriented plaquette product.
Define
\[
 \mathsf s_i(x)=\min_y\{\mathsf L_i(y)+|x-y|_1/R\},\qquad
 \mathcal K_i=(D[U]+m_0)_{\Omega_1\Omega_1}+\mathsf W_i,
 \qquad
 \mathsf W_i=\sum_{n=1}^i\sum_{B\subset X_n^i}
                    \theta_nL^{-n}\iota P_B\iota^*.
\]
Here $P_B=Q_B^*Q_B$, and $Q_B$ is the row of the iterated site
average on the block $B$ of side $L^n$.
Equivalently set $T_n^i=1_{X_n^i}q_n\cdots q_1$;
then $\mathsf W_i=\sum_n(T_n^i)^+\theta_nT_n^i$,
where $+$ denotes the row twin on the complex class and
the adjoint at real fields.  At real fields its norm squared on a block
is $L^{-n}$.  At the top integrated places of step $j$
write $T_A=1_{A_j}q_j\cdots q_1$ and
$c_j=\theta_j+\theta q_{j+1}^+q_{j+1}$.  The block
inverse identity gives
\[
 \Gamma_j^{-1}=c_j^{-1}+c_j^{-1}\theta_jT_A
           \mathcal K_{j+1}^{-1}T_A^+\theta_jc_j^{-1}.
\]
One way to verify it is to eliminate the bare field
first or the last coarse field first in the same
three-block Gaussian matrix.  The bare Schur form
after the second elimination is
$\mathcal K_j-T_A^+\theta_jc_j^{-1}\theta_jT_A$,
which is $\mathcal K_{j+1}$.  The corresponding block
of its inverse is the displayed formula.  At $j=0$,
$\Gamma_0^{-1}=\mathcal K_1^{-1}$.
Put $\mathcal E_U(g)=\|\Pi_1g\|^2+\|\nabla_U\iota^*g\|^2$.
We use $c_G=c_\Delta$ in the bound
$\langle f,\Delta f\rangle\ge c_G\mathcal E_U(f)$.
The separation gives
$1\le\mathsf s_i\le\mathsf L_i\le L\mathsf s_i$ and
$|\mathsf s_i(x)-\mathsf s_i(x')|\le R^{-1}$ on a bond.
The distance $d_{h,i}$ is the shortest-path distance
whose bond $\langle x,x'\rangle$ has length
$1/\max\{\mathsf s_i(x),\mathsf s_i(x')\}$.
A weight has rate $\delta$ if its difference across
every bond is at most $\delta$ times this length.

\begin{lemma}[Regular averages]\label{C-397-23}
Assume that the admitted pure-gauge block-average estimate applies in
a block $B$ of side $L^n$, so that
\[
 |\partial(\mathcal M^iU)(p)-1|\le2\epsilon L^{2(i-n)}
 \quad(0\le i<n),\qquad
 \epsilon\le\min\{\varepsilon_\flat,(400L^2)^{-1}\}.
\]
Then
\[
 \|(1-P_B)\varphi\|_B^2\le c_PL^{2n}\|\nabla_U\varphi\|_B^2
                         +\epsilon_P\|\varphi\|_B^2,
 \quad c_P=\frac{e^{1/(L-1)}L^2}{2(L^2-1)},
 \quad\epsilon_P=50(64L^2)^2L^{-2}\epsilon^2.
\]
\end{lemma}

\begin{proof}
Write $V_i=\mathcal M^iU$,
$\bar\omega_i=2\epsilon L^{2(i-n)}$ and
$r_i=e^{9L^2\bar\omega_i/8}-1$.
Rearranging the steps of two contours bounds the difference of their
transporters by $r_i$.  For unit quaternions,
\[
 \left|\frac1m\sum_ag_a\right|^2
 =1-\frac1{2m^2}\sum_{a,b}|g_a-g_b|^2.
\]
Consequently $q(V_i;y,x)=w(y,x)p(y,x)$, with
$p(y,x)\in\mathrm{SU}(2)$, $|w-1|\le r_i^2$ and
$L^{-3}\sum_xw(y,x)^2=1$.
The same contour rearrangement, and the formula for the admitted
gauge block average on the double block of a bond, give
\[
 |p(y,c_-)V_i(c)-p(y,c_+)|\le c_T\bar\omega_i,
\]
\[
 |V_{i+1}(c)q(V_i;c_+,z')-
                 q(V_i;c_-,z)V_i([z,z'])|\le c_T\bar\omega_i,
 \qquad c_T=64L^2.
\]
For the second inequality the two paths have rearranged step
sequences of length at most $(5L-3)/2$; their difference is at most
$e^{25L^2\bar\omega_i/8}-1$, and
$\bar\omega_i\le(200L^2)^{-1}$ bounds the resulting exponentials
linearly.  These comparisons use only the block or double block in
question, including bonds crossing the faces of $B$.

Let $Q_i=Q(V_{i-1})\cdots Q(V_0)$ and
$\mathsf P_i=Q_i^*Q_i$.  These are nested orthogonal projections,
so $1-P_B=\sum_{i<n}(\mathsf P_i-\mathsf P_{i+1})$ orthogonally.
For $u=Q_i\varphi$ on one block put $g(z)=p(y,z)u(z)$ and
$\tilde w=L^{-3/2}w(y,\cdot)$.
Since $\|\tilde w\|=1$,
\[
 \|g\|^2-|\langle\tilde w,g\rangle|^2
 \le2\|g-\bar g\|^2+2r_i^4\|g\|^2.
\]
The cube Poincar\'e inequality and the first transport comparison
therefore yield
\[
 \|(\mathsf P_i-\mathsf P_{i+1})\varphi\|^2
 \le\tfrac12L^2\|\nabla_{V_i}Q_i\varphi\|^2+
        (3L^2c_T^2\bar\omega_i^2+2r_i^4)\|\varphi\|^2.
\]
The second transport comparison, Cauchy--Schwarz along the straight
paths of length $L$, and the fact that each fine bond belongs to
exactly $L$ of these paths give, for every $\eta_i>0$,
\[
 \|\nabla_{V_{i+1}}Q(V_i)u\|^2
 \le(1+\eta_i)L^2\|\nabla_{V_i}u\|^2+
          6(1+\eta_i^{-1})c_T^2\bar\omega_i^2\|u\|^2.
\]
Take $\eta_i=L^{-(n-i)}$. The product $\prod_{i<n}(1+\eta_i)$ is at most
$e^{1/(L-1)}$, and iteration gives
\[
 \|\nabla_{V_i}Q_i\varphi\|^2
 \le e^{1/(L-1)}L^{2i}\|\nabla_U\varphi\|^2+
       30c_T^2\epsilon^2L^{3(i-n)-3}\|\varphi\|^2.
\]
Sum over $i<n$.  The gradient sum is at most
$L^{2n}/(L^2-1)$ and the remaining geometric sums are bounded by
$\epsilon_P$.  This proves the assertion.
\end{proof}

\begin{lemma}[Localized ellipticity]\label{C-397-26}
There is $c_e>0$, depending only on the chamber operator, such that
$\|D[1]g\|^2\ge c_e\mathcal E_1(g)$ for finitely supported $g$.
Let $\mathsf s\ge1$ have bond differences at most $R^{-1}$, with
$R\ge6R'$ and $R'\ge8$.  Suppose the cube of side
$2R'\mathsf s(x)$ around every $x$ in the support of $g$ has
plaquette bound $\epsilon'\mathsf s(x)^{-2}$.  Then
\[
 \|D[U]g\|^2\ge\frac{c_e}{16}\mathcal E_U(g)
 -C_{\mathrm{pu}}(1+c_e)c_{\mathrm{hop}}^2
        (R'^{-2}+R'^2\epsilon'^2)\|\mathsf s^{-1}g\|^2.
\]
\end{lemma}

\begin{proof}
The free chamber symbol has heavy block of real part at least
$\lambda_\rho$, off-diagonal light--heavy blocks of size
$O(|p|)$, and light Schur complement $i\gamma(p)+O(|p|^2)$.
The two block equations therefore bound
$|p|\,|\iota^*\hat g|+|\Pi_1\hat g|$ by a constant times
$|\hat D(p)\hat g|$ near zero.  Away from zero the coercive
real part bounds the inverse uniformly.  Fourier transformation
gives the stated free inequality.

Choose the partition of unity on the covering lattice with
\[
 \rho(c)=R'\mathsf s(c)/8,\quad
 \beta_c(x)=\beta((x-c)/\rho(c)),\quad
 Z=\sum_c\beta_c^2,\quad\chi_c=\beta_cZ^{-1/2},
\]
where $\beta$ is one on the half unit cube, vanishes outside the
unit cube and has Lipschitz constant two.  Slow variation of
$\mathsf s$ gives
\[
 \sum_c\chi_c^2=1,\qquad
 \sum_c|\chi_c(x')-\chi_c(x)|^2
                    \le C(R'\mathsf s(x))^{-2}.
\]
Indeed, the centres with nonzero $\beta_c(x)$ lie within
$2R'\mathsf s(x)/15$ of $x$, while those within
$R'\mathsf s(x)/17$ give $\beta_c(x)=1$; these two counts and
the difference formula for $\beta_cZ^{-1/2}$ give the bound.
The family is covariant under the periods, so it can be summed
over centres modulo the periods even if a localizing cube exceeds
the torus side.

Only wall-zero hops change the site, and hence
$\|[D,\chi]f\|\le c_{\mathrm{hop}}
\||\nabla\chi|^\sharp f\|$.
Thus
\[
 \|Dg\|^2\ge\tfrac12\sum_c\|D\chi_cg\|^2
                 -Cc_{\mathrm{hop}}^2R'^{-2}
                                  \|\mathsf s^{-1}g\|^2.
\]
On the support cube of $\chi_c$, axial gauge represents each link
as a product of at most $2R'\mathsf s(c)$ plaquette boundaries,
so $|U^u-1|\le4R'\epsilon'\mathsf s(c)^{-1}$.
Apply the free inequality to $u\chi_cg$, continued by zero, and
compare $D[U^u]$ with $D[1]$.  Comparing the corresponding gradient
energies and summing the partition gives
\[
 \sum_c\mathcal E_U(\chi_cg)
 \ge\tfrac12\mathcal E_U(g)-CR'^{-2}\|\mathsf s^{-1}g\|^2.
\]
The two comparisons each use
$\|a+b\|^2\ge\frac12\|a\|^2-\|b\|^2$.
Together with the first localization inequality they give the
factor $c_e/16$ and the displayed error.
\end{proof}

\begin{proposition}[Weighted energy and fluctuation accretivity]
\label{C-397-28}
There are $R_0,C_S<\infty$ and $\epsilon_*>0$, depending only on
$L,b_F$ and the fixed chamber operator, such that if $R\ge R_0$
and $0<\epsilon\le\epsilon_*$ in the stated separation and
regularity conditions, then for $1\le i\le j+1$,
\[
 \|\mathsf s_i^{-1/2}f\|\le
 C_S\|\mathsf s_i^{1/2}\mathcal K_if\|,
 \qquad
 \|\mathsf s_i^{-1/2}\mathcal K_i^{-1}
                         \mathsf s_i^{-1/2}\|\le C_S.
\]
The same holds for $\mathcal K_i^*$.  Consequently
$\operatorname{Re}\Gamma_j\ge2/C_D$ for a uniform $C_D$.
Every compression of $\Gamma_j$ has an inverse with exponentially
decaying kernel in the units of level $j$.
\end{proposition}

\begin{proof}
Fix $i$, suppress its subscript and continue $f$ by zero outside
$\Omega_1$.  Write
$N=\|\mathsf s^{-1/2}f\|$, $E=\|\mathsf s^{1/2}\mathcal Kf\|$
and $g=\mathsf s^{1/2}f$.  The real part gives
\[
 \langle f,\Delta f\rangle+\langle f,\mathsf Wf\rangle
                         +m_0\|f\|^2\le NE.
\]
The penalty blocks have
$\|\mathsf s^{1/2}\mathsf Wf\|^2
\le\theta\langle f,\mathsf Wf\rangle$.
The commutator with $\mathsf s^{1/2}$ costs at most
$2c_{\mathrm{hop}}R^{-1}N$.
Across the boundary of $\Omega_1$ the scale is at most
$1+R^{-1}$ and the real part pays for the crossing differences,
without a regularity condition on the exterior bonds.  These facts
give
\[
 \|Dg\|^2\le3E^2+(3\theta+C_{\mathrm{bh}})NE
                           +12c_{\mathrm{hop}}^2R^{-2}N^2,
 \qquad C_{\mathrm{bh}}=12(1+r)^2/r=54.
\]
The constants are independent of $m_0$: accretivity permits dropping the positive
mass in this inequality.

Let $d_1=2R'L$ and choose $\zeta=0$ within $d_1$ of
$\Omega_1^c$, $\zeta=1$ beyond $2d_1$ and
$|\nabla\zeta|\le d_1^{-1}$.  Separation puts this boundary zone
in $X_1$, where $\mathsf L=L$.  In the zone the unweighted real
part already controls the gradient.  On the support of $\zeta$,
the regularity condition allows \cref{C-397-26} with
$\epsilon'=L^2\epsilon_*$.  It follows that the weighted gradient
$Y=\sum_b\mathsf s(b_-)|\nabla_U\iota^*f(b)|^2$, with both
orientations included in this sum, satisfies
\[
 Y\le\frac{64}{c_e}\|Dg\|^2+C_0
       (R^{-2}+R'^{-2}+R'^2L^4\epsilon_*^2)N^2
                          +2Lc_G^{-1}NE.
\]
Here $c_G>0$ is the coercivity constant in the real part, and
$C_0$ depends on that constant, $c_e$ and $c_{\mathrm{hop}}$.
Apply \cref{C-397-23} block by block and use the penalty on the
averaged component and the heavy-component gap.  This gives
\[
 N^2\le(2/\lambda_\rho+2L/\theta)NE
                             +L^2c_PY+L\epsilon_PN^2.
\]
Combining the last three inequalities yields
$(1-\vartheta)N^2\le A E^2+BNE$.
Choose $R'$ first so its $R'^{-2}$ error is at most $1/8$;
then impose the separation floor on $R_0$ so its $R^{-2}$
error is at most $1/8$; finally decrease $\epsilon_*$ so its
$R'^2L^4\epsilon_*^2$ and $L\epsilon_P$ errors are each at
most $1/8$.  Thus $\vartheta\le1/2$ and
$N\le(2B+\sqrt{2A})E$.  The adjoint replaces $\gamma$ by
$-\gamma$ and has the same proof.

The closed Schur formula is
\[
 \Gamma_j^{-1}=c_j^{-1}+c_j^{-1}\theta_jT_A
               \mathcal K_{j+1}^{-1}T_A^*\theta_jc_j^{-1},
 \qquad c_j=\theta_j+\theta q_{j+1}^*q_{j+1}.
\]
It gives $\|\Gamma_j^{-1}\|\le2/\theta+LC_S$.
At $j=0$ the all-field block Poincar\'e inequality gives
$\operatorname{Re}\Gamma_0\ge m_0+c_{\mathrm{fl}}$.
The disc inequality then gives
\[
 \operatorname{Re}\Gamma_j\ge
 \{(c_{\mathrm{tw}}+\theta)\|\Gamma_j^{-1}\|^2\}^{-1}.
\]
Choose $C_D$ to make this at least $2/C_D$.

For the decay, conjugate the fine operator by $e^\eta$, with
$|\eta(x)-\eta(x')|\le\delta/\max\{\mathsf s(x),\mathsf s(x')\}$.
The weighted norm of its change is at most
$3e(c_{\mathrm{hop}}+\theta)L\delta$.
The Neumann series therefore preserves the weighted inverse bound
for a fixed small $\delta>0$.  Lift a Lipschitz weight on level
$j$ to the fine lattice by
$\tilde\eta(x)=\min_y\{\eta(y)+\alpha L^{-j}|x-y|_1\}$.
Across a block of level $j$ it differs from its centre value by
at most $3\alpha/2$; the Schur formula consequently gives
$\|e^\eta\Gamma_je^{-\eta}-\Gamma_j\|\le C_\Gamma\alpha$.
Choose $C_\Gamma\alpha\le1/C_D$.
Compression commutes with the weights and preserves the real-part
floor.  The inverse of every weighted compression has norm at most
$C_D$; taking $\eta(y)=\alpha|y-y'|_1$ proves its kernel bound.
\end{proof}

\subsection{Complex fields, sources, and local kernels}

For the complex class at depth $j$, give sites outside $\Omega_1$
depth zero and sites in $X_n^{j+1}$ depth $n$.
For a bare bond $b$, let $n(b)$ be the smaller of its endpoint
depths. We set $\mathsf U(b)=e^{iX(b)}U(b)$ on positive bonds,
$\mathsf U(\bar b)=\mathsf U(b)^{-1}$, with
$X(b)\in\mathfrak{sl}(2,\mathbb C)$, and require
$|X(b)|\le\alpha L^{-n(b)}$.

\begin{lemma}[Complete fermion averages at complex fields]
\label{lem:qcd-complex-block-averages}
There are $a_Q,\sigma_Q>0$ and $C_Q<\infty$, depending only
on $L$, with the following property. Let $1\le n\le K$ and
let $B=B^n(y)$ be a block
of side $L^n$, let $U$ be real on its inner bonds, and suppose
its plaquettes satisfy $|U_p-1|\le aL^{-2n}$, $a\le a_Q$.
Let $\mathsf U=e^{iX}U$ there, with
$L^n\sup|X|\le\sigma\le\sigma_Q$. Put
$V_i=\mathcal M^iU$, $V_i^{\rm c}=\mathcal M^i\mathsf U$ for
$0\le i<n$.
The normalized one-step kernels and their row twins are
holomorphic on this common domain, with
\[
 q_i^{\rm c}=n(V_i^{\rm c})^{-1/2}S(V_i^{\rm c}),\qquad
 q_i^+=n(V_i^{\rm c})^{-1/2}
             [(\operatorname{tr}S(V_i^{\rm c}))1-S(V_i^{\rm c})].
\]
Their complete products $q_n^{\rm c}(y,x)$, along the unique
nested chain from $x$ to $y$, and the reversed products
$(q_n^{\rm c})^+(x,y)$ satisfy
\[
 \sup_{x\in B}|q_n^{\rm c}(y,x)-q_n(U;y,x)|
 +\sup_{x\in B}|(q_n^{\rm c})^+(x,y)-q_n(U;y,x)^*|
 \le C_Q\sigma.
\]
For the rows normalized by $L^{-3n/2}$, the same comparison
holds in operator norm, and $Q_n^{\rm c}(Q_n^{\rm c})^+=1$.
Consequently the chamber average $R_n^{\rm c}=Q_n^{\rm c}\iota^*$
has $R_n^{\rm c}(R_n^{\rm c})^+=1$, with
$(R_n^{\rm c})^+=\iota(Q_n^{\rm c})^+$.
The kernels read only the bonds of $B$, preserve the prescribed
boundary signs, and have the gauge and lattice covariance of
the real averages.

For a regular history with $\epsilon\le a_Q$,
$\alpha\le\sigma_Q$ and $C_Q\alpha\le1/32$, the transported
comparison holds as well. If $m\le n$, $c=\langle y,y'\rangle$
is a level-$m$ bond with both endpoints in $\Omega_n$,
$x\in B^m(y)$ and $x'=x+L^m\hat c$, then both the real and
complex transported comparisons hold:
\[
 V_m(c)q_m(y',x')=q_m(y,x)U([x,x'])+E,\qquad
 |E|\le c_T\epsilon L^{2(m-n)},
\]
\[
 V_m^{\rm c}(c)q_m^{\rm c}(y',x')
 =q_m^{\rm c}(y,x)\mathsf U([x,x'])+E^{\rm c},\qquad
 |E^{\rm c}|\le c_T\epsilon L^{2(m-n)}
                         +C_Q\alpha L^{m-n+1}.
\]
This includes bonds crossing the faces of the level-$n$
blocks. At a bond from $\Omega_n$ to its adjacent lower
layer, with $m\le n-1$, the same statement holds with
$n-1$ in place of $n$.
\end{lemma}

\begin{proof}
We first estimate the iterated gauge fields and then the products
defining the fermion averages. By
\cite[equation (1063), p.~473]{YMM1}, on the double block,
after reducing fixed radii, the iterates are holomorphic and
\[
 V_i^{\rm c}=e^{iX_i}V_i,\quad
 \xi_i:=\sup|X_i|\le C_{\rm av}\sigma L^{i-n},\qquad
 |(V_i)_p-1|\le\omega_i:=2aL^{2(i-n)}.
\]
Write $c_{\rm ia},c_\sharp$ for its curvature and perturbation
radii, and choose
\[
 a_Q=\min\{\varepsilon_\flat,c_{\rm ia}/2,(400L^2)^{-1}\}.
\]
To use that double-block statement on $B$, reflect its data
into the neighbouring block, put $U=1$ and $X=0$ on the
crossing bonds, and observe that every crossing plaquette is
one. The average in $B$ is independent of this extension by
\cite[Lem.~14.2]{YMM1}.

The contour rearrangement estimate gives
$r_i=e^{9L^2\omega_i/8}-1$. The real quaternion identity
used in \cref{C-397-23} gives
\[
 n(V_i)\ge1-r_i^2/2\ge1/2,\qquad |q(V_i)|\le1+r_i^2.
\]
A one-step contour has at most $3(L-1)/2$ bonds. Telescoping
its transports gives $|S(V_i^{\rm c})-S(V_i)|\le CL\xi_i$.
The determinant formula for two-by-two matrices consequently
gives $|n(V_i^{\rm c})-n(V_i)|\le C_0L\xi_i$, where
$C_0\ge1$ also bounds the square-root and kernel changes on
$\{z:|z-r|\le1/4,\ 1/2\le r\le1\}$. This constant depends
only on $L$. Take $C_Q$ larger than $32C_0LC_{\rm av}$ and
than the fixed constants in the transport comparison below,
and define
\[
 \sigma_Q=\min\{c_\sharp/2,(60LC_{\rm av})^{-1},
                   (4C_0LC_{\rm av})^{-1},(32C_Q)^{-1}\}.
\]
Thus $3L\xi_i\le1/20$ and the determinant change is at most
$1/4$. All these values lie in the right
half-plane, so use its single square-root branch, positive
at real fields. The displayed formulas for $q_i^{\rm c}$
and $q_i^+$ are holomorphic and differ from their real
values by at most $CL\xi_i$.

For every two-by-two matrix $S$,
$S[(\operatorname{tr}S)1-S]=(\det S)1$. Thus
\[
 L^{-3}\sum_{x\in B(z)}q_i^{\rm c}(z,x)q_i^+(x,z)=1.
\]
Iterating this identity gives the complete normalization.
Moreover
\[
 \sum_{i<n}r_i^2\le13a^2,\qquad
 \sum_{i<n}L\xi_i\le C\sigma,
\]
with constants independent of $n$. The products of all
factor norms are bounded by
$\exp\sum_{i<n}(r_i^2+CL\xi_i)\le2$.
Telescoping the complete product now proves the asserted
$C_Q\sigma$ comparison; the reversed row product has the
same bound. The normalized operator comparison follows by
Cauchy--Schwarz over the $L^{3n}$ sites. Locality, covariance
and boundary signs pass through these same formulas.

For the history assertion, inner bare bonds of $\Omega_n$
have $n(b)\ge n$. Apply the block result at level $m$ with
$a=\epsilon L^{2(m-n)}$ and
$\sigma=\alpha L^{m-n}$. Here is the iteration of the real
one-step comparison in the proof of \cref{C-397-23}.
Writing $E_m$ for its complete defect, product substitution
gives
\[
 |E_m|\le(1+r_{m-1}^2)L\sup|E_{m-1}|
                                  +2c_T\omega_{m-1}.
\]
The real product bounds and
$\sum_{i<m}L^{m-1-i}\omega_i
 \le2\epsilon L^{2(m-n)}/[L(L-1)]$
give $|E_m|\le c_T\epsilon L^{2(m-n)}$ with the same
$c_T$, since the real product is at most $21/20$ and
$L\ge13$. This iteration
uses all the bonds of the double block, including those
crossing faces. Write $q_m$ for the complete real kernels.
Subtracting its transport identity from the complex one gives
\[
 \begin{split}
 E^{\rm c}-E={}&(V_m^{\rm c}(c)-V_m(c))q_m^{\rm c}(y',x')
       +V_m(c)(q_m^{\rm c}(y',x')-q_m(y',x'))\\
 &-(q_m^{\rm c}(y,x)-q_m(y,x))\mathsf U([x,x'])
       -q_m(y,x)(\mathsf U([x,x'])-U([x,x'])).
 \end{split}
\]
The pure double-block bound also estimates $V_m^{\rm c}(c)$
when $c$ crosses a block face or $m=n$. The straight path has
$L^m$ bonds, so its complex transport changes by at most
$C\alpha L^{m-n}$. Each of the four terms above is therefore
bounded by $C_Q\alpha L^{m-n}$, with fixed product bounds.
The sum is at most $5C_Q\alpha L^{m-n}$, and $L\ge13$
gives the stated weaker bound with $L^{m-n+1}$.
For a bond into the lower layer, both blocks
and their double block lie in $\Omega_{n-1}$; the same
argument uses that region. No product of unrestricted
one-step Lipschitz constants is used.
\end{proof}

Reduce $\epsilon_*$, if necessary, so that $\epsilon_*\le a_Q/2$.

\begin{lemma}[Stability]\label{C-397-29}
The preceding inverse estimates persist jointly on a sufficiently
small complex gauge class and at every complex matrix source
$m_0-a_K\hat J$, where $J$ acts on spin, the internal doublet
and flavour and is the identity on colour, satisfying
\[
 \|J\|_{[j]}:=\sup_x\|L^{j-K}J(x)\|\le\varrho_\sharp.
\]
The radii $\alpha_1^*>0$ and $\varrho_\sharp>0$ depend on the
earlier operator constants, not on $m_0$, the volume, the cutoff or
$N_f$.  In particular no condition on the real part of the source
mass is imposed.  Every fluctuation compression is invertible and
holomorphic there, with the same exponential kernel estimate after
enlarging its constant by a fixed factor.
\end{lemma}

\begin{proof}
For $\widetilde{\mathcal K}=\mathcal K+E$, the weighted inverse
identity and a Neumann series apply if
\[
 \|\mathsf s^{1/2}E\mathsf s^{1/2}\|\le(2C_S)^{-1}.
\]
They give $\|\mathsf s^{-1/2}\widetilde{\mathcal K}^{-1}
\mathsf s^{-1/2}\|\le2C_S$.  A twist of rate $\delta$ changes
the fine operator by at most
$3e(c_{\rm hop}+\theta)L\delta$ in this norm.  A complex gauge
change costs at most
$8(1+\theta)(c_{\rm hop}+C_Q)\alpha$; the comparison of
\cref{lem:qcd-complex-block-averages} is used for the block penalty.  The source costs
at most $L\|J\|_{[j]}$.  Thus one can first choose
\[
 \delta_*\le[24eC_S(c_{\rm hop}+\theta)L]^{-1},\qquad
 \bar\alpha\le\min\{\sigma_Q,(32C_Q)^{-1},
       [64C_S(1+\theta)(c_{\rm hop}+C_Q)]^{-1}\},
\]
and then lower the source radius.  These choices apply on the local
mirror tori as well as on the original history.

Put $C_{\mathrm{inv}}=\max\{2/\theta+2LC_S,2/c_{\rm fl}\}$.
The Schur formula bounds the fluctuation inverse by $C_{\mathrm{inv}}$.
Choose $C_D$ to dominate $2(c_{\rm tw}+\theta)C_{\mathrm{inv}}^2$
and $4/c_{\rm fl}$, choose the fluctuation twist rate
$0<\alpha_D\le\delta_*/4$ after $C_D$, and increase $R_0$ until
\[
 \alpha_DR_0\ge\max\{2\log L,32\},\qquad
 R_0\ge1+16/\delta_*.
\]
Increasing $R_0$ does not increase $C_S$.  The source radius may be
taken below
\[
 \min\{(8L^2C_S)^{-1},(4C''_\Gamma C_D)^{-1},
 c_{\rm fl}/4,[1+9C_S^2(c_{\rm hop}+\theta)^2]^{-1}\},
 \qquad C''_\Gamma=16\theta^2C_S^2.
\]
Finally choose $\alpha_1^*\le\bar\alpha$ so that the change of
the fluctuation form has norm at most $(4C_D)^{-1}$.
The real-part floor of the real form and these perturbation bounds
give the inverse bounds for every compression.  Inserting exponential
weights gives their kernel bounds.  In finite volume all entries
are holomorphic matrix functions; the locally uniform bounds give
the same assertion for the infinite-volume kernels.
\end{proof}

We record the free fact needed both in localization and in the
mass calculation.  For the balanced chamber operator, with $t=3$
and $\rho_0=3/2$,
\[
 \iota^*\widehat D(p)\iota
 =i\sum_\mu\gamma_\mu\sin p_\mu+
       \tfrac12\sum_\mu(1-\cos p_\mu),\qquad
 \Pi_1\operatorname{Re}\widehat D(p)\Pi_1
       \ge\lambda_\rho\Pi_1.
\]
The light--heavy couplings are $O(|p|)$.  Indeed the balanced
within-site hops cancel on a chamber-independent field, leaving
only $(\rho_0-3\gamma(e))\nabla_e\varphi$; the same argument for
$D^*$ bounds both couplings.  The heavy block has inverse norm
at most $\lambda_\rho^{-1}$.  Its Schur complement on the light
space is therefore
\[
 m+i\gamma\cdot p+O(|p|^2),
\]
uniformly for $m\ge0$ in a bounded interval.  Near zero its inverse
differs from $(m+i\gamma\cdot p)^{-1}$ by at most
$C|p|^2/(m^2+|p|^2)$.  To see this, write the complement as
$m+i\gamma\cdot p+E$, $\|E\|\le C|p|^2$, and apply the
resolvent identity when $C|p|^2\le\frac12(m^2+|p|^2)^{1/2}$.
Away from zero, the real part has a positive gap: equality in
$\sum_s\rho_s(1-\widehat R_s(p))\ge0$ forces the vector to be
constant in the chamber variable and $e^{ip\cdot e}=1$ for all
coordinate directions.  Compactness supplies the remaining bound.
This proves the free ellipticity used above and excludes additional
zero modes.  Differentiating the Schur formula gives, for each
multi-index $\zeta$,
\[
 \|\partial_p^\zeta(\widehat D(p)+m)^{-1}\|
       \le C_\zeta(m+|p|)^{-1-|\zeta|}.
\]
The singular light contribution to the heavy output has one
additional factor $|p|$, since the coupling is a difference.
The heavy-heavy inverse also has a regular, uniformly bounded part.

\begin{lemma}[Local regularity]\label{C-397-32}
There are $\epsilon_{\rm lr},C_{\rm lr}>0$ with the following property.
For a site $x_0$ put $Q_r=\{x:|x-x_0|_\infty\le r\}$.
Let $l\ge4$, let $U$ be real on the bonds meeting $Q_{2l}$,
and suppose the curvature at every plaquette meeting $Q_{2l}$
is at most $\epsilon_{\rm lr}l^{-2}$.
Let $F,s$ be defined on the places over $Q_{2l}$ and its neighbours,
with $(D[U]+\mu)F=s$ over $Q_{2l}$ and $0\le\mu\le1$.
Write $\|F\|_2=\|F\|_{\ell^2(Q_{2l})}$ and
$\|s\|_\infty=\sup_{Q_{2l}}|s|$. At the centre,
\[
 |F|\le C_{\rm lr}(l^{-3/2}\|F\|_2+l\|s\|_\infty),
\]
\[
 |\nabla_U\iota^*F|+|\Pi_1F|
 \le C_{\rm lr}\{l^{-5/2}\|F\|_2+
                   (1+\log l)\|s\|_\infty\}.
\]
Put $\varphi=\iota^*F$. For $x\in Q_{l/4}$,
$x'=x+\rho e$, where $e$ is a coordinate unit vector and
$1\le\rho\le l/4$, the straight-path transport satisfies
\[
 |U([x,x'])\varphi(x')-\varphi(x)|
 \le C_{\rm lr}\rho\left\{
 l^{-5/2}\|F\|_2+
 (1+\log(l/\rho))\|s\|_\infty+
 \epsilon_{\rm lr}l^{-1}\sup_{Q_{3l/2}}|F|\right\}.
\]
\end{lemma}

\begin{proof}
The free symbol estimates just proved, a dyadic partition around
$p=0$, and integration by parts on each dyadic annulus give
\[
 |G_\mu(x)|\le C(1+|x|)^{-2}e^{-\mu|x|/C},\qquad
 |\nabla G_\mu(x)|+|\Pi_1G_\mu(x)|
       \le C(1+|x|)^{-3}e^{-\mu|x|/C}.
\]
For completeness, on an annulus $|p|\asymp2^{-n}$ the symbol
has size $O(2^n)$ and its $q$th derivative size $O(2^{(q+1)n})$.
Its Fourier coefficient is bounded by
$C_q2^{-2n}(1+2^{-n}|x|)^{-q}$.  Summing with $q>3$ gives the
first bound.  A difference or a heavy projection removes $2^n$,
giving the second.  A contour shift of size $c\mu$ gives the
exponential factor.  Subtracting the kernels at distance $\rho$
and splitting at $|x|=\rho$ gives the asserted logarithmic modulus.

Choose a rooted axial gauge on $Q_{2l}$; summing plaquette
curvatures over its coordinate surfaces gives
$|U^u(b)-1|\le C\epsilon_{\rm lr}l^{-1}$, and at the root
the sharper bound proportional to its distance from the root.
For a cutoff $\chi$, equal to one on $Q_l$ and with
$|\nabla\chi|\le C/l$, extend by zero and write
\[
 \chi uF=G_\mu\{\chi us+[D[1],\chi]uF+
                 \chi(D[1]-D[U^u])uF\}.
\]
The source term uses respectively
$\sum_{|x|\le Cl}(1+|x|)^{-2}\le Cl$ and
$\sum_{|x|\le Cl}(1+|x|)^{-3}\le C(1+\log l)$.
The commutator is on the cutoff shell; Cauchy--Schwarz with the
free kernel gives $Cl^{-3/2}\|F\|_2$, or
$Cl^{-5/2}\|F\|_2$ after a difference or heavy projection.
The field term is absorbed in the weighted local supremum
$\sup_x\operatorname{dist}(x,\partial Q_{2l})^{3/2}|F(x)|$
by taking $\epsilon_{\rm lr}$ small.  For the differentiated
estimate use the axial gauge rooted at the point being estimated;
the additional distance in $U^u-1$ gives precisely the same
absorption.  The kernel difference calculation gives the modulus
of continuity.  
\end{proof}

\begin{proposition}[Local walks and background maps]\label{C-397-40}
Choose a power $M_w\ge L$ of $L$ after the operator and local-regularity
constants, and assume $R\ge LM_w$.  There are local walks
$\mathcal W_\omega$, $\omega=(\square_0,\ldots,\square_n)$,
such that
\[
 \mathcal K_i^{-1}=\sum_\omega\mathcal W_\omega,
 \quad
 \|\mathsf s^{-1/2}\mathcal W_\omega\mathsf s^{-1/2}\|
       \le2LC_Sq_w^n,
 \quad q_w=C_q/M_w,
\]
and their kernels and gauge/source dependence are contained in
$X_\omega=\bigcup_lW_{\square_l}$.  With
$|g|_{(a)}=\sup_x\mathsf s(x)^a|g(x)|$,
\[
 |\mathcal W_\omega g|_{(1)}
       \le C_{\rm sup}q_{\rm sup}^n
                         |1_{W_{\square_n}}g|_{(2)},
 \qquad q_{\rm sup}=C_q^{\rm sup}/M_w.
\]
The constants are chosen so that
$C_{\rm nb}\max\{q_w,q_{\rm sup}\}\le e^{-3}$.
The same bounds hold for the row twins and jointly on the complex
gauge class and the source ball.  Every covariance obtained by
switching off walk cubes has a uniform Gram constant $\Gamma_F$
on the closed decoupling polydisc, factorizes over the remaining
components, and agrees with the original covariance when all
parameters are one.
For the Gaussian integrations we take the mean of the right and left
localized inverses.  This covariance, denoted by $S_j^{\rm sym}$ below,
also preserves conjugation and charge conjugation at every value of
the decoupling parameters.

For the background map to level $n$ in a history ending at $k$,
write $\mathcal H_{(n)}=\mathcal H^{\rm loc}_{(n)}+
\sum_\omega\mathcal H_{(n);\omega}$.  Its local part reads one
block and has norm at most $2L^{-2(n_P-n)}$ on a column of level
$n_P$.  The walk part obeys
\[
 \sum_P|\mathcal H_{(n);\omega}(x;P)|
       \le C_H\frac{L^n}{\mathsf s_k(x)}q_{\rm sup}^{|\omega|}.
\]
For neighbouring level-$n$ sites $x,x'$, let $U_n(x,x')$
be their level-$n$ link transporter. Then
\[
 \sum_P|\mathcal H_{(n);\omega}(x;P)-
       U_n(x,x')\mathcal H_{(n);\omega}(x';P)|
 \le C_H\Bigl(\frac{L^n}{\mathsf s_k(x)}\Bigr)^2
 \Bigl(1+\log_+\frac{\mathsf s_k(x)}{L^n}\Bigr)
 q_{\rm sup}^{|\omega|}.
\]
At level zero the heavy part has this latter gain.  In particular,
for every $0<\alpha<1$ the second bound is at most
$C(\alpha)(L^n/\mathsf s_k(x))^{1+\alpha}q_{\rm sup}^{|\omega|}$.
On the history without large fields $\mathsf s_k=L^k$.
All these estimates are uniform in $k$, $K$, the volume and $N_f$.
\end{proposition}

\begin{proof}
At the first step $\mathcal K_1=\Gamma_0$ and
$1\le\mathsf s_1\le L$. The preceding all-field bound
also admits a matrix source of norm at most
$\varrho_\sharp\le c_{\mathrm{fl}}/4$: the twisted real
part loses at most $c_{\mathrm{fl}}/4$, so the kernel
prefactor becomes $4/c_{\mathrm{fl}}$. For a weight of
rate at most $\min\{\delta_*,\delta_0\}/8$, the twisted
kernel is bounded by
$(4/c_{\mathrm{fl}})e^{-7\delta_0|x_v-x_{v'}|_1/8}$.
Summing over the $\nu_0=48$ chambers at each site gives
\[
 |e^\eta\mathcal K_1^{-1}e^{-\eta}g|_{(1)}
 \le\frac{4L\nu_0}{c_{\mathrm{fl}}}
       \left(\coth\frac{7\delta_0}{16}\right)^3|g|_{(2)}.
\]
Include this number in $C_{\rm sup}$ before choosing the
walk constants. It applies to the full and mirror inverses
at every real configuration. The regular complex domains
are treated by the perturbation bounds below.

Partition the regular layers into cubes $\square$ of side
$M_wL^{n(\square)}$ and choose $\sum h_\square^2=1$, with
$h_\square$ supported in the concentric cube $W_\square$ of
triple side and $|\nabla h_\square|\le
CL/(M_wL^{n(\square)})$.  Separation allows the history and the
configuration on $W_\square$ to be reflected to a mirror torus.
The reflected layers retain separation and regularity and their
smoothed scales differ by at most $L$; hence its local inverse
$\mathrm{Sl}_\square$ has the weighted bound $2LC_S$.
The right parametrix
$\sum h_\square\mathrm{Sl}_\square h_\square$ has error made
of $[\mathcal K_i,h_\square]\mathrm{Sl}_\square h_\square$.
The finite-range hops and the block penalties give
\[
 \|\mathsf s^{1/2}[\mathcal K_i,h_\square]
                         \mathsf s^{1/2}\|\le C/M_w.
\]
Indeed the penalty kernel has row sum at most two and its
commutator costs its row sum times
$\operatorname{osc}_Bh_\square\le3L^{n(B)}|\nabla h_\square|$.
The error norm is below one after $M_w$ is chosen.  Its Neumann
series gives exactly the stated walks and the factor $q_w$ per
step.  Only consecutive overlapping cubes occur, with at most
$C_{\rm nb}$ continuations.

Here is the supremum estimate needed for the legs.  Split $g$ into
the blocks $B$ of the layers.  The twisted energy estimate gives
\[
 \|\mathsf s^{-1/2}e^{\delta_*d(\cdot,y_B)}
                 \mathcal K_i^{-1}1_Bg\|
       \le C L^2|g|_{(2)}.
\]
At $l=\varepsilon_0\mathsf s(x)\ge4$, keep the real operator
$D[U]+m_0$ in the local regularity estimate and write
\[
 (D[U]+m_0)F=1_Bg-\mathsf W_iF+YF,\qquad
 Y=a_K\hat J-(D[\mathsf U]-D[U]).
\]
On this cube the source and complex-hop bounds give
$\|Y\|\le(L\varrho_\sharp+3c_{\rm hop}\alpha_1^*)/\mathsf s(x)$.
In the weighted supremum
$\sup_x\mathsf s(x)e^{\delta_*d(x,y_B)/2}|F(x)|$, the penalty,
mass and complex-field terms have coefficients at most
$1/18$, $5/32$ and $1/96$, respectively; their sum is at most
$1/4$.  Choose $\varepsilon_0$ first and then lower
$\epsilon_*$ and $\alpha_1^*$ to secure these inequalities.
Where $l<4$, use the weighted energy bound directly.
It follows that
\[
 |\mathcal K_i^{-1}1_Bg(x)|\le
 C\mathsf s(x)^{-1}e^{-\delta_*d(x,y_B)/2}|g|_{(2)}.
\]
The number of layer places within distance $t$ is bounded by
$C(1+te^{t/R})^3$.  The floor $R\ge1+16/\delta_*$ and the
distance $R$ between successive complete layers make its sum
against the displayed exponential finite, uniformly in the
history.  This proves the supremum bound for the full inverse
and for the local mirror inverse, independently of $M_w$.
Moreover
\[
 \|[\mathcal K_i,h_\square]\|_{(2)\leftarrow(1)}
        \le C(c_{\rm hop}+\theta)L^2/M_w.
\]
Multiplying these bounds along a walk gives $q_{\rm sup}$ per
step. For $\omega=(\square_0,\ldots,\square_m)$, write
$\mathcal W_\omega g=h_{\square_0}F$, where
$F=(\mathcal K^{\rm mir}_{\square_0})^{-1}g_1$.
Here $g_1=h_{\square_0}g$ if $m=0$, and
$g_1=-h_{\square_0}[\mathcal K_i,h_{\square_1}]F_1$ otherwise,
where $\mathcal W_{(\square_1,\ldots,\square_m)}g
=h_{\square_1}F_1$. The preceding bounds give
$|g_1|_{(2)}+|F|_{(1)}\le Cq_{\rm sup}^{|\omega|}G$, with
$G=|1_{W_{\square_m}}g|_{(2)}$.
Apply the real local estimate to $F$ in the first cube.
Put $s_* =\mathsf s_i(x)$ and $l=\varepsilon_0s_*$ for a fixed small
$\varepsilon_0>0$.
The energy and supremum bounds just proved give
\[
 \|F\|_{\ell^2(Q_{2l})}\le Cq_{\rm sup}^{|\omega|}s_*^{1/2}G,
 \qquad \sup_{Q_{3l/2}}|F|\le Cq_{\rm sup}^{|\omega|}s_*^{-1}G.
\]
In the real local equation the forcing is $g_1-\mathsf W_iF+YF$;
its supremum on $Q_{2l}$ is at most
$Cq_{\rm sup}^{|\omega|}s_*^{-2}G$.
Thus all three terms in the modulus of \cref{C-397-32}, including
the curvature term, are bounded by
$Cq_{\rm sup}^{|\omega|}\rho s_*^{-2}(1+\log(s_*/\rho))G$.
Multiplication by $h=h_{\square_0}$ preserves this bound: the extra
term $(h(x')-h(x))U([x,x'])F(x')$ has norm at most
$C\rho s_*^{-2}|F|_{(1)}$, by the gradient bound for $h$.
For a level-$n$ difference take $\rho=L^n$. The normalized
average contributes $L^n$, giving the stated squared scale ratio
and logarithm. If $4L^n>\varepsilon_0s_*$, or $l<4$, use the size
bound instead: $t=L^n/s_*$ is then bounded below by a fixed positive
constant, and $t\le Ct^2$. The complete-average comparison pays
the transport defect, including bonds crossing layer faces.
The heavy-component estimate in \cref{C-397-32}, with the same forcing
and energy bounds, gives
\[
 |\Pi_1\mathcal W_\omega g(x)|
 \le Cq_{\rm sup}^{|\omega|}\mathsf s_i(x)^{-2}
              (1+\log\mathsf s_i(x))G.
\]
For $l<4$ this follows from the size bound, since $\mathsf s_i(x)$
is then bounded.
All these estimates retain the same $q_{\rm sup}^{|\omega|}$;
no general complex-source extension of the real lemma is used.

For the background maps take $i=k$. The background equations have the form
\[
 \psi=\mathcal K_k^{-1}S\Psi,
 \quad S\Psi=\sum_m(T_m^k)^+\theta_m\Psi_m-D_{A_0S_0}\Psi_{S_0},
 \qquad\|S\|_{(2)\leftarrow(0)}\le C.
\]
At level $n\ge1$, eliminating the block equations gives
\[
 \mathcal H_{(n);\omega}=
 \left(q_n\cdots q_1-\theta_n^{-1}
 \sum_{m>n}(q_m\cdots q_{n+1})^+
                     1_{X_m^k}\theta_mT_m^k\right)
                 \mathcal W_\omega S.
\]
The leftover term is $\mathcal H^{\rm loc}_{(n)}$.
In the per-chamber supremum norm, $\iota^*$ has norm $\sqrt{48}$:
it is $48^{-1/2}$ times the sum over the chambers at a site.
Thus each average is at most $2\sqrt{48}L^n$ times the supremum;
$L^{m-1}\le\mathsf s\le L^m$ on its layer gives the size bound.
Write $\psi_\omega=\mathcal W_\omega S\Psi$ and
$\varphi_\omega=\iota^*\psi_\omega$.
For neighbouring blocks the real complete-average defect is at most
$c_T\epsilon_*L^{2(n-m)}$. On the complex class,
\cref{lem:qcd-complex-block-averages} adds
$C_Q\alpha_1^*L^{n-m+1}$. Multiplying by the normalized average
factor $L^n$ and the size bound
$|\varphi_\omega|\le2\sqrt{48}|\psi_\omega|_{(1)}/\mathsf s_k(x)$
bounds this contribution by
\[
 C(c_T\epsilon_*L^2+C_Q\alpha_1^*L^2)
 \left(\frac{L^n}{\mathsf s_k(x)}\right)^2
 |\psi_\omega|_{(1)}.
\]
For $\rho\le\varepsilon_0\mathsf s_k(x)/4$, the complex and
real transports along a straight path of length $\rho$ differ by at most
$C\alpha_1^*L\rho/\mathsf s_k(x)$. At $\rho=L^n$ this
difference, multiplied by the same size bound and average factor,
has the same squared gain. Changing the level-$n$ link transporter
costs at most $CC_Q\alpha_1^*L\,L^n/\mathsf s_k(x)$ and hence
also has this gain after multiplication by the background size.
Together with the logarithmic difference bound in the first cube,
these estimates give the stated bound, with the same walk factor,
including across faces between layers. The local term has size
$O(L^{-2(m-n)})$, so
its possible jump has the same gain.  At level zero apply the
heavy-component bound above to $g=S\Psi$; the bound on $S$ gives
the asserted gain.

At zero source all kernels are trivial in flavour.  The row sum
of the norms of their matrix entries is bounded by $\nu'=8$
times the norm on vector-valued bounded functions: expand a row
vector in an orthonormal basis and choose each column vector
with its maximizing phase.  At a source use instead
\[
 \mathcal K(J)^{-1}=
 [1-\mathcal K(0)^{-1}a_K\hat J]^{-1}\mathcal K(0)^{-1}
\]
in the row norm
$\|A\|_{\rm row}=\sup_x\mathsf s(x)
\sum_y\|A(x,y)\|\mathsf s(y)^{-2}$.
Since $\|a_K\hat J\|_{(2)\leftarrow(1)}\le L\|J\|_{[j]}$,
lower the source radius so the Neumann ratio is at most $1/2$.
For a difference or heavy component of a background row, put the
zero-source inverse in the first factor of the resolvent identity.
Its gain is retained; only the remaining factors cost two.  Thus
the column sums have no extra factor $N_f$.

To file walks into polymers of side $M\ge3M_w$, let $c_0$ satisfy
$d(X\cup Y)\le d(X)+d(Y)+c_0$ for meeting polymers and
$d(Y)\le c_0$ in a cube of side $2M$.  Divide a walk into
segments of at most $\lfloor M/(3M_w)\rfloor$ cubes.  Each
segment lies in a cube of side $M$.  Consequently
\[
 d(Y_\omega)\le c_0+6c_0(|\omega|+1)M_w/M.
\]
Put $q=\max\{q_w,q_{\rm sup}\}$. With $c_fM\ge C_{\rm dc}(\kappa+2)$ and
$c_f=(3LM_w)^{-1}$, choose $C_{\rm dc}$ so that
$6c_0(\kappa+2)M_w/M\le1/8$. For each fixed first cube
$\square$, counting continuations gives
\[
 \sum_{\substack{\omega:\square_0=\square\\Y_\omega=Y}}q^{|\omega|}
 \le22e^{(\kappa+2)c_0}e^{-(\kappa+2)d(Y)},
\]
with the same bound for a fixed last cube. For a fixed place
$x$ of a polymer $X$, sum instead over walks whose first cube
is one of the at most $C_{\rm nb}$ cubes for which
$h_{\square_0}$ does not vanish on the layer block containing
$x$. Then
\[
 \sum_\omega e^{|\omega|+\kappa[d(X\cup Y_\omega)-d(X)]_+}
                 q^{|\omega|}
 \le A_\kappa:=2C_{\rm nb}e^{(\kappa+2)c_0}.
\]
The same bound holds with the admissible last cubes of $x$.
The offset $c_0$ is essential for walks of no step.

For decoupling, let $Z_\circ$ be a union of fine cubes, enlarge it
by $D_{\rm ex}$ to $Z_\circ^\sharp$, and put $\bar s(\Delta)=1$
on its big cubes and $\bar s(\Delta)=s(\Delta)$ elsewhere.
Here $|s(\Delta)|\le e^{\bar\kappa}$.  Weaken a walk by
$\bar s(\omega)=\prod_{\Delta\cap X_\omega\ne\varnothing}
\bar s(\Delta)$, each cube counted once.  Define
\[
 \Gamma_j^{-1}(Z_\circ,s)=\chi_{Z_\circ}
 \left[c_j^{-1}+c_j^{-1}\theta_jT_A
       \Bigl(\sum_\omega\bar s(\omega)\mathcal W_\omega\Bigr)
                        T_A^+\theta_jc_j^{-1}\right]\chi_{Z_\circ};
\]
at $j=0$ use the restricted weakened fine inverse directly.
A chain of $n+1$ cubes meets at most $8+CnM_w/M$ big cubes.
If it starts and ends in $Z_\circ$ and leaves $Z_\circ^\sharp$,
then $n+2\ge2D_{\rm ex}/(3LM_w)$, in level-$j$ units.
Indeed a walk cube has side at most $3LM_w$ in those units,
and the walk goes outside the enlargement and returns.
Therefore
\[
 |\bar s(\omega)|(C_{\rm nb}q_w)^n
 \le e^{8\bar\kappa}
        (e^{C\bar\kappa M_w/M}e^{-3})^n
 \le e^{8\bar\kappa}e^{-2n}.
\]
The floor $c_f\min\{M,D_{\rm ex}\}\ge
C_{\rm dc}\bar\kappa$, with $C_{\rm dc}$ sufficiently large,
ensures both inequalities. In particular
$D_{\rm ex}\ge15LM_w\bar\kappa$ gives
$n\ge10\bar\kappa-2\ge8\bar\kappa$, so the last bound is
at most $e^{-n}$ for walks which leave the enlargement.
For the other walks every weakening factor is one, and the
same bound follows from $C_{\rm nb}q_w\le e^{-3}$.
The series converges
uniformly on the larger polydisc $|s(\Delta)|\le
e^{\bar\kappa+1}$. On the complex class, the block-average
bounds and $C_Q\alpha_1^*\le1/32$ give
\[
 \|c_j^{-1}\theta_jT_A\mathsf s^{1/2}\|^2,
 \ \|\mathsf s^{1/2}T_A^+\theta_jc_j^{-1}\|^2\le\tfrac54L,
 \qquad \|c_j^{-1}\|\le\tfrac65\theta^{-1}.
\]
Schur's test over the first and last cubes bounds the weighted
norm of the walk sum by
\[
 2LC_S\sum_{n\ge0}e^{-n}\le16LC_S/5.
\]
At $j=0$ its plain norm is at most $16/(5c_{\rm fl})$. Thus
\[
 \|\Gamma_j^{-1}(Z_\circ,s)\|\le\Gamma_F^2,
 \qquad \Gamma_F^2=2/\theta+16C_{\rm nb}L^2C_S
                                      +10C_{\rm nb}/c_{\rm fl}.
\]
A walk contributing to the entry at $(y,y')$ satisfies
$3M_w(n+1)+2\ge|y-y'|$ in level-$(j+1)$ units.
The endpoint offsets are absorbed by the geometric sum:
\[
 \sum_{\substack{n\ge0\\3M_w(n+1)+2\ge|y-y'|}}e^{-n}
 \le\frac{e^{1+2/(3M_w)}}{1-e^{-1}}e^{-Lc_f|y-y'|}
 \le5e^{-c_f|y-y'|}.
\]
The entry prefactor is at most $\tfrac52C_{\rm nb}L^2C_S$
for $j\ge1$ and $2C_{\rm nb}/c_{\rm fl}$ for $j=0$.
The block term $c_j^{-1}$ joins only places with $|y-y'|<1$;
its bound is at most $(2/\theta)e^{-c_f|y-y'|}$.
These estimates give entry decay with the same $\Gamma_F^2$.
Only walks leaving the
enlargement change when $s$ changes from one, so their sum is
at most $\Gamma_F^2e^{8\bar\kappa}e^{-c_fD_{\rm ex}}$.
Turning a parameter to zero removes every walk crossing that
cube.  Connectedness of $X_\omega$ then makes the covariance a
direct sum on the surviving components.  At parameters one it
is the original inverse.  The left parametrix proves the
row-twin statements.  Local holomorphy and uniform convergence
prove these assertions.

The right inverse alone need not preserve the real structure at
intermediate parameters.  To retain that structure, form the left
parametrix as well. On functions supported where all neighbours
and blocks remain inside $W_\square$, the full operator and
the mirror operator agree. Thus the corresponding columns of
$\mathrm{Sl}_\square\mathcal K_{j+1}$ are identity columns.
Consequently
\[
 \Bigl(\sum_\square h_\square\mathrm{Sl}_\square h_\square\Bigr)
       \mathcal K_{j+1}=1-\mathfrak N^{\rm L},\qquad
 \mathfrak N^{\rm L}=-\sum_\square
       h_\square\mathrm{Sl}_\square[h_\square,\mathcal K_{j+1}].
\]
Expanding its inverse gives
\[
 \mathcal W^{\rm L}_\omega=
 \Bigl[\prod_{i=n}^{1}
       (-h_{\square_i}\mathrm{Sl}_{\square_i}
                       [h_{\square_i},\mathcal K_{j+1}])\Bigr]
       h_{\square_0}\mathrm{Sl}_{\square_0}h_{\square_0}.
\]
The order in this product runs from $n$ on the left to $1$ on
the right.  These are the same factors as in a right walk, in
reverse order.  Their supports are still $X_\omega$, and their
weighted bounds are the same.  The supremum and difference estimates
start at the last cube instead of the first; the argument above
therefore gives the same constants.  Denote by
$\Gamma_{j,\mathrm L}^{-1}(Z_\circ,s)$ the covariance obtained
from these walks.

Write $\mathsf U^{\mathfrak c}(b)=(\mathsf U(b)^*)^{-1}$,
$J^{\mathfrak c}=\gamma_4J^\dagger\gamma_4$ and
$s^{\mathfrak c}(\Delta)=\overline{s(\Delta)}$.  Define
\begin{equation}\label{eq:qcd-symmetric-covariance}
 S_j^{\rm sym}(Z_\circ,s;\mathsf U,J)
 =\frac12\left[\Gamma_j^{-1}(Z_\circ,s;\mathsf U,J)
 +\gamma_4\Gamma_j^{-1}
       (Z_\circ,s^{\mathfrak c};\mathsf U^{\mathfrak c},
                                   J^{\mathfrak c})^*\gamma_4\right].
\end{equation}
Its second summand is $\Gamma_{j,\mathrm L}^{-1}$.
Indeed, at a real configuration and a source fixed by the
conjugation, $Y\mapsto\gamma_4Y^*\gamma_4$ fixes the local
operators and reverses their products.  It exchanges $T_A,T_A^+$
and the two commutators, hence carries each right walk to its left
walk.  At a fixed real factor $U$, write $\mathsf U=e^{iX}U$.
Both sides of this identity are holomorphic in $(s,X,J)$ on the
connected product of the parameter polydisc, the convex complex
gauge class and the source ball.  They agree on its real points,
an open subset of a real form of this space; the holomorphic
identity theorem proves equality throughout the domain and
continuity gives it on the closed class.

The transpose identity for the chamber operator gives the analogous
assertion for charge conjugation. Put
$\mathcal C=(\sigma_2\otimes1_\tau)\otimes\epsilon
\otimes1_{N_f}$, with $\epsilon=i\sigma_2$ on colour;
all transposes are in the standard tensor-product bases.
Define the linear source transform by
$\widehat{J^{\mathsf c}}=\mathcal C\widehat J^{\mathsf T}
\mathcal C^{-1}$. For physical sources in the span of
$\mathfrak B$, which are scalar on spin, this reduces to
transpose on the internal doublet and flavour. This differs from the antilinear
$J^{\mathfrak c}=\gamma_4J^\dagger\gamma_4$ above.
On colour the transpose identity uses
$\epsilon A^{\mathsf T}\epsilon^{-1}=(\operatorname{tr}A)1-A$,
which exchanges $q$ and $q^+$.  Thus charge conjugation reverses
the walk factors and exchanges the two parametrices, replacing
the source by $J^{\mathsf c}$. Their mean satisfies
\[
 \gamma_4S_j^{\rm sym}(Z_\circ,s^{\mathfrak c};
          \mathsf U^{\mathfrak c},J^{\mathfrak c})^*\gamma_4
      =S_j^{\rm sym}(Z_\circ,s;\mathsf U,J),
\]
\[
 \mathcal C S_j^{\rm sym}(Z_\circ,s;\mathsf U,J)^{\mathsf T}
                 \mathcal C^{-1}
      =S_j^{\rm sym}(Z_\circ,s;\mathsf U,J^{\mathsf c}).
\]
Both summands have norm at most $\Gamma_F^2$ and the same
component restrictions; so does their mean.  Each is the whole
covariance when $s=1$.  Holomorphy of the second summand in
\eqref{eq:qcd-symmetric-covariance} follows because it is the
conjugate of a holomorphic function at conjugate arguments.

If $Z_\circ$ extends outside $\Omega_{j+1}$, first intersect
its fine cubes with $\Omega_{j+1}$ and continue its parameters
as one on $Z_\circ^\sharp$.  The compression to the integrated
places is unchanged by this intersection.  The completed
parameters and hence the walk sums are the same; if the
intersection is empty they act on the zero space.  This extends
the construction and its component identities to these regions.
Parameters are supplied on all big cubes.  A walk can leave
$\Omega_{j+1}$ through an earlier layer, so restricting the
parameter set to the top layer would not give the stated
component decomposition.
\end{proof}

\subsection{The local fermion normalizer}

\begin{lemma}[Rows reaching stored layers]\label{C-397-33}
Choose the weight constant $C_H$ at least twice the enlarged
whole-background row constant and above the fine-resolvent
constant below. On the regularity and source domains of
\cref{C-397-29}, with $j=k-1$, one has
\begin{equation}\label{eq:qcd-fine-resolvent-kernel}
 |\mathcal K_k(\mathsf n)^{-1}(v,v')|
 \le C_H(1+|x_v-x_{v'}|_1)^{-2}
                   e^{-\alpha_Dd_{h,k}(x_v,x_{v'})},
 \qquad \mathsf n=-a_KJ.
\end{equation}
The same bound holds jointly with the complex gauge class.
If a row $x$ of level $n$
lies in the layer $\mathcal L_m$, $n\le m$, then at a real
regular configuration
\[
 \sum_P|\mathcal H_{(n)}(x;P)|
            \le\tfrac12 C_HL^{-(m-n)}.
\]
The sum includes stored columns of every level.  The row
twins obey the same estimate.  At the admitted complex
configurations and matrix sources the bound holds with
$C_H$ in place of $C_H/2$.  The localized walk rows,
their neighbour differences and their sums with the
weights $e^{|\omega|}$ have the constants $C_H$ and
$C_H(\alpha)$ of \cref{C-397-40} on these domains as
well.  On the empty history this gives
$\sum_y|\mathsf H_{n,k}(x,y)|\le C_HL^{-(k-n)}$ and
the difference bound $C_H(\alpha)L^{-(1+\alpha)(k-n)}$.
\end{lemma}
\begin{proof}
\emph{The fine resolvent.}
We give the point-source argument, since a weighted operator
bound alone does not give the power in
\eqref{eq:qcd-fine-resolvent-kernel}. Retain
$\alpha_D\le\delta_*/4$ and the separation floors chosen in
\cref{C-397-29}. Write $\mathsf s=\mathsf s_k$, $d=d_{h,k}$,
and fix a place $v'$ and a unit fibre vector $\xi$. Put
\[
 x'=x_{v'},\qquad s'=\mathsf s(x'),\qquad
 F=\mathcal K_k(\mathsf n)^{-1}\delta_{v'}\xi,
 \qquad \eta(x)=\delta_*d(x,x').
\]
All constants in this argument depend only on the previously
chosen operator constants. They are independent of the history,
volume, cutoff and flavour multiplicity. Let
$t=|x-x'|_1$ and $\rho=|x-x'|_\infty$. The Lipschitz bound
on $\mathsf s$ gives
\[
 t\le(Rs'+1)(e^{d/R}-1)\le2s'de^{d/R},\qquad
 \max\{\mathsf s(x)/s',s'/\mathsf s(x)\}\le e^{2d/R}.
\]
Consequently, if $\delta R\ge4$,
\begin{equation}\label{eq:qcd-metric-to-bare-decay}
 e^{-\delta d}\le C_\delta s'^2(1+t)^{-2},\qquad
 C_\delta=4\sup_{u\ge0}(1+u)^2e^{-\delta u/2}.
\end{equation}

Choose the fixed $\varepsilon_0>0$ in the local regularity
argument small enough for its absorptions below, and reduce
$\epsilon_*$ so that
$\epsilon_*L^2\varepsilon_0^2\le\epsilon_{\rm lr}$.
If
$\varepsilon_0s'<64$, the twisted inverse bound gives
\[
 |F(x)|\le2C_S(\mathsf s(x)s')^{1/2}e^{-\eta(x)}
       \le2C_Ss'e^{-15\delta_*d/16}.
\]
Apply \eqref{eq:qcd-metric-to-bare-decay} with
$\delta=\delta_*/2$. The remaining exponential has rate
at least $\alpha_D$, and $s'<64/\varepsilon_0$ bounds
the remaining prefactor. This proves the claim in this case.

Suppose now that $l=\varepsilon_0s'\ge64$.
The separation of the layers puts $Q_{2l}(x')$ inside
$\Omega_1$, with its required neighbouring bonds, and gives
\[
 |U_p-1|\le\epsilon_*L^2s'^{-2},\qquad
 \tfrac34s'\le\mathsf s\le\tfrac54s'
\]
on this cube and on the layer blocks meeting it. The lengths
of these blocks lie between $3s'/4$ and $5Ls'/4$.
Use the axial gauge rooted at $x'$, extended by the identity
outside $Q_{2l}(x')$, so that $u(x')=1$ and
\[
 |U^u(b)-1|\le\epsilon_*L^2s'^{-2}|z-x'|_1
\]
for bonds at $z$ in the cube. After this unitary gauge change
the operator is
\[
 \mathcal K=D[U^u]+m_0+\mathsf W-\mathsf Y,
 \qquad
 \mathsf Y=a_K\widehat J-(D[\mathsf U^u]-D[U^u]).
\]
The second term in $\mathsf Y$ is zero at real links.
The penalty kernel on a block of level $n$ is bounded by
$5\theta L^{-4n}$, by the complete-average estimates.
The finite-range coefficients of $\mathsf Y$ near $x'$ are
bounded by
$\upsilon=(L\varrho_\sharp+3c_{\rm hop}\alpha_1^*)/s'$.

\emph{A local approximation with two terms.}
Let $G=(D[1]+m_0)^{-1}$ be the free kernel on $\mathbb Z^3$
from \cref{C-397-32}, read in the lifted cutoff cube.
Choose cutoffs $\chi,\chi'$ with gradients
at most $4/l$, where $\chi=1$ on $Q_l(x')$ and vanishes
off $Q_{3l/2}(x')$, and $\chi'=1$ on $Q_{3l/2+1}(x')$
and vanishes off $Q_{2l-1}(x')$. Put
\[
 \widetilde g=\mathsf Y\chi G\delta_{v'}\xi,
 \qquad F^{\rm p}=\chi G\delta_{v'}\xi+\chi'G\widetilde g.
\]
The support of $\widetilde g$ lies where $\chi'=1$.
Direct multiplication gives
$\mathcal KF^{\rm p}=\delta_{v'}\xi+g+g'$, where
\[
 \begin{split}
 g={}&[D[U^u],\chi]G\delta_{v'}\xi
       +\chi(D[U^u]-D[1])G\delta_{v'}\xi
       +\mathsf W\chi G\delta_{v'}\xi,\\
 g'={}&[D[U^u],\chi']G\widetilde g
       +\chi'(D[U^u]-D[1])G\widetilde g
       +\mathsf W\chi'G\widetilde g
       -\mathsf Y\chi'G\widetilde g.
 \end{split}
\]
In particular the second term in $F^{\rm p}$ cancels the
first occurrence of $-\mathsf Y\chi G\delta_{v'}\xi$.
This cancellation is needed for the weighted norm of the error.

The free bound $|G(z)|\le C(1+|z|)^{-2}$ gives the following
bounds for the three terms of $g$, in their displayed order:
\[
 Cl^{-3}\ \text{on the cutoff shell},\qquad
 C\epsilon_*L^2s'^{-2}(1+|z-x'|)^{-1},\qquad
 C\theta\varepsilon_0s'^{-3}.
\]
The last term is supported on the blocks meeting the cutoff;
their union has at most $C(Ls')^3$ sites. For it, sum
$5\theta L^{-4n}(1+|z-x'|)^{-2}$ over the cutoff cube.
For the second approximation use the elementary convolution
bound in three dimensions,
\[
 \sum_y(1+|z-y|)^{-2}(1+|y-x'|)^{-2}
       \le C(1+|z-x'|)^{-1}.
\]
Thus $|G\widetilde g(z)|\le C\upsilon(1+|z-x'|)^{-1}$.
The four terms of $g'$ are bounded respectively by
\[
 Cs'^{-1}l^{-2},\qquad C\epsilon_*L^2s'^{-3},\qquad
 C\theta\varepsilon_0^2s'^{-3},\qquad
 Cs'^{-2}(1+|z-x'|)^{-1},
\]
on the second shell, the cutoff cube, its meeting blocks,
and the neighbouring cutoff cube. Here and below fixed
powers of $\varepsilon_0,L$ are included in $C$.
On all these supports $\mathsf s\asymp s'$ and $\eta$ is
bounded. Squaring and summing the displayed bounds gives
\[
 \|\mathsf s^{1/2}e^\eta(g+g')\|_2\le C s'^{-1}.
\]
For example, the square of the last term's weighted norm is bounded by
\[
 Cs'\,s'^{-4}
 \sum_{|z-x'|_\infty\le2l}(1+|z-x'|_\infty)^{-2}
 =O(s'^{-2}).
\]
In contrast, the uncancelled $\widetilde g$ would have
weighted norm of order $\upsilon s'^{1/2}$.
Also $|F^{\rm p}(x)|\le C(1+\rho)^{-2}$ on $\rho\le2l$
and $F^{\rm p}=0$ outside this cube.

\emph{The remainder.}
Set $F^{\rm r}=uF-F^{\rm p}=-\mathcal K^{-1}(g+g')$.
The twisted inverse estimate yields
\[
 \|\mathsf s^{-1/2}e^\eta F^{\rm r}\|_2
       \le C_{\rm par}s'^{-1}.
\]
Its equation is
\[
 (D[U^u]+m_0)F^{\rm r}
       =-(g+g')-\mathsf WF^{\rm r}+\mathsf YF^{\rm r}.
\]
On a block of level $n$, Cauchy--Schwarz and the penalty
kernel bound give
\[
 |\mathsf WF^{\rm r}(x)|
 \le C\theta L^{-5n/2}\|F^{\rm r}\|_{\ell^2(B)}
 \le C\theta\mathsf s(x)^{-2}e^{-\eta(x)}s'^{-1}.
\]
Define
\[
 v(x)=
 \begin{cases}
 (1+\rho)^{-2},&\rho\le l,\\
 (1+l)^{-2}(s'/\mathsf s(x))e^{-\eta(x)},&\rho>l,
 \end{cases}
 \qquad A=\sup_x|F^{\rm r}(x)|/v(x).
\]
This supremum is finite on the finite torus. When $\rho\le32$,
or when $\rho>l$ and $\varepsilon_0\mathsf s(x)<64$,
the weighted $\ell^2$ bound already gives
$|F^{\rm r}(x)|\le Cv(x)$.
At every remaining site apply \cref{C-397-32} to the real
operator $D[U^u]+m_0$ on a cube of half-side
\[
 l_x=\rho/8\quad(\rho\le l),\qquad
 l_x=\varepsilon_0\mathsf s(x)/16\quad(\rho>l).
\]
These lengths are at least four. The enlarged cube stays at
distance at least $7\rho/10$ from $x'$, and $v\le3v(x)$
there. The same separation and curvature bounds as above
verify the local regularity hypotheses.

For $\rho\le l$, the local $\ell^2$ contribution is bounded by
$C\rho^{-2}(\rho/s')^{1/2}$. The penalty contribution is
$C l_xs'^{-3}$. The curvature, shell and remaining source
errors cost at most
$C(\epsilon_*L^2s'^{-2}+l^{-2}+s'^{-2})$ when their
supports meet the cube. Each is at most $Cv(x)$.
For $\rho>l$, the first two contributions are bounded by
\[
 C\bigl(l_x^{-3/2}\mathsf s(x)^{1/2}
              +l_x\theta\mathsf s(x)^{-2}\bigr)s'^{-1}e^{-\eta(x)}
 \le C\mathsf s(x)^{-1}s'^{-1}e^{-\eta(x)}\le Cv(x).
\]
The errors $g+g'$ away from their centre are $O(s'^{-3})$.
They contribute only within distance $3Ls'$ of $x'$, where
$\mathsf s(x)\asymp s'$ and $\eta$ is bounded, so their
contribution has the same bound.

It remains to absorb $\mathsf YF^{\rm r}$. If $n'$ is the
least level of a block meeting the local cube, its contribution
is at most
\[
 3C_{\rm lr}l_x
 (L^{1-k}\varrho_\sharp+3c_{\rm hop}\alpha_1^*L^{1-n'})A v(x).
\]
Here $l_x\le\varepsilon_0L^k/8$ and
$l_xL^{-n'}\le\varepsilon_0/6$. Choose
$C_{\rm lr}\varepsilon_0\le1/8$,
$L\varrho_\sharp\le1/8$ and
$c_{\rm hop}\alpha_1^*L\le1/64$, reducing the existing
radii if necessary. The coefficient is at most $1/2$.
Thus $A\le C+A/2$, proving $|F^{\rm r}(x)|\le Cv(x)$.

For $\rho\le2l$, $d(x,x')\le8\varepsilon_0$, so the bounds
on $F^{\rm p},F^{\rm r}$ give the asserted kernel estimate.
For $\rho>2l$ the approximation vanishes and
\[
 |F(x)|\le Cl^{-2}e^{-7\delta_*d(x,x')/8}.
\]
Apply \eqref{eq:qcd-metric-to-bare-decay} once more, with
$\delta=\delta_*/2$. Its factor $s'^2$ cancels $l^{-2}$
up to the fixed $\varepsilon_0^{-2}$, and the remaining
rate exceeds $\alpha_D$. This proves
\eqref{eq:qcd-fine-resolvent-kernel}. The argument used the
complex perturbation and the matrix source together; it
does not require a positive real part for their sum.

\emph{The background rows.}
For a column $P$ of level $n_P\ge1$, its fine background
is $\mathcal K_k^{-1}(T^k_{n_P})^*\theta_{n_P}\delta_{y_P}$.
The source has norm $L^{-n_P/2}$, is supported in its
block and has $\mathsf s\le L^{n_P}$ there.  The weighted
inverse estimate therefore bounds
$\|\mathsf s^{-1/2}e^{\delta_*d(\cdot,y_P)}\psi\|$
uniformly.  For a fine stored column the source is
$-D_{A_0S_0}\delta_P$ near the fine layer, and has the
same bound.  The local regularity argument used for
\cref{C-397-40}, with the penalty and mass terms
absorbed in the weighted supremum, gives
$|\psi(x)|\le C\mathsf s(x)^{-1}e^{-\delta_*d(x,y_P)/2}$.
Passing through the level-$n$ averages yields
$C_H^\circ L^{-(n_P-n)}e^{-\alpha_Dd(x,y_P)}$ for
$n\le n_P$.  For $n>n_P$ the preliminary factor is
$L^{n-n_P}$; each intervening layer has distance at
least $R$, and $\alpha_DR\ge2\log L$ pays this factor
from half the exponential.  The number of places in
a distance ball of radius $t$ is bounded by
$C(1+te^{t/R})^3$, and the separation floor makes its
sum against the remaining exponential uniformly finite.
Thus the full row sum is bounded by
$(1+C_\Sigma L^{3/2})C_H^\circ L^{-(m-n)}$.
Choose $C_H\ge2(1+C_\Sigma L^{3/2})C_H^\circ$.

The neighbour and heavy-component bounds apply the
same local regularity estimate to the column equation.
Its logarithmic modulus is bounded by a power
$\alpha<1$; the transported-average defect has the
same gain.  Summation over columns uses the preceding
distance count.  The complex and source perturbation
may double these whole-kernel constants.  The constants
for the local walk rows were already formed on the
complex source domain in \cref{C-397-40}; choose $C_H$
and $C_H(\alpha)$ to include them from the outset.
Neither the geometric walk ratio nor $A_\kappa$ is
increased.  This proves the assertions.
\end{proof}

In particular a formal background of level $n<m$ in a
stored layer, given weight at least
$\frac12C_HL^{-(m-n)}h_F$, can be read through the
whole real background at generator weight $h_F$ with
norm at most one.  This reserve will be used when a
history is finally evaluated.

\begin{proposition}[Closed-walk expansion]\label{C-397-42}
The difference $\log\det\Gamma_j(U;J)-\log\det\Gamma_j(1;J)$
is a sum of local, strictly gauge-invariant scalar terms, vanishing
at $U=1$. Its individual array members obey
\[
 \sup_P|N^{\rm F}(X;P)|\le C_NN_f e^{-(\kappa+1)d(X)}.
\]
This is a polymerwise bound, not the family norm obtained
by summing $e^{\kappa d(X)}|N^{\rm F}(X)|$ through a cube.
They are holomorphic on the same
complex gauge and source domains and are covariant under the
isometries preserving the partitions. At $J=0$ they are real at
real configurations. The flat-field part at a nonconstant $J$ is
a local source-dependent sum; at $J=0$ it is a vacuum number per
block with a thermodynamic limit. The bounds are uniform in the
volume, cutoff and history.
\end{proposition}

\begin{proof}
For $j\ge1$ let $c_j=\theta_j+\theta q^*q$ and
$P=Lq^+q$, an idempotent.  The Schur determinant identity gives
\[
 \det(\Gamma_j+y)=\det(c_j+y)\det\mathcal K_y/\det\mathcal K_j,
 \quad
 \mathcal K_y=\mathcal K_j-T_A^+(\theta_j-p_y)T_A,
\]
\[
 p_y=\frac{y}{1+y/\theta_j}(1-P)+
       \frac{y+\theta/L}{1+(y+\theta/L)/\theta_j}P,
 \qquad N_y=\partial_y\mathcal K_y
       =T_A^+\theta_j(c_j+y)^{-2}\theta_jT_A.
\]
Here $\mathcal K_0=\mathcal K_{j+1}$ and
$\mathcal K_\infty=\mathcal K_j$.  Since
$p_y^2\le\theta p_y$ and $p_y\ge p_0$, the weighted energy
proof applies uniformly for $y\ge0$, with $LC_S$ in place of
$C_S$.  Integrate the logarithmic derivative and subtract the
flat configuration:
\[
 \log\det\Gamma_j(U)-\log\det\Gamma_j(1)
 =-\int_0^\infty\!\{\operatorname{Tr}(\mathcal K_y(U)^{-1}N_y(U))
       -\operatorname{Tr}(\mathcal K_y(1)^{-1}N_y(1))\}\,dy.
\]
The block-local determinant of $c_j+y$ is configuration independent:
its eigenspace dimensions are fixed by $P^2=P$.
Use the local walk expansion for $\mathcal K_y^{-1}$, and file
each trace walk at the polymer of its visited cubes.  The trace
closes the walk and is gauge invariant.  Before the integral,
its rank bound is $\nu(3LM_w)^3$, its two $T_A$ factors cost
$L^{-j}$, and the weighted inverse costs $CL^{j+1}$.
The remaining integral is bounded by
$\int_0^\infty\|\partial_yp_y\|\,dy\le C\theta$.
Thus the coefficient bound is $C\nu q_w^{|\omega|}$.
The walk count and the sum over trace origins give the
individual array bound at rate $\kappa+1$; one spare unit
pays the number of origins in its polymer.  Summing this
bound through a cube without the $e^{\kappa d(X)}$ weight
is allowed by the actual polymer count and the large
floor on $\kappa$.  No weighted family bound is inferred
from a single spare unit of the rate.

At $j=0$ use
\[
 \log\det\mathcal K_1=
 \int_0^\infty\operatorname{Tr}\bigl((1+y)^{-1}
                         -(\mathcal K_1+y)^{-1}\bigr)\,dy.
\]
After subtracting the flat configuration, the one-cube term is
$O((1+y)^{-2})$; a walk of $n\ge1$ steps has $n+1$
resolvents and is integrable.  The same closed-walk count applies.
The source is read only in visited cubes, with the same bound
when both configurations have the same source. At $J=0$, conjugation by
$\gamma_4$ exchanges a real kernel with its adjoint; averaging
the left and right trace-walk expansions makes every filed
coefficient real.  Gauge covariance, lattice covariance and
holomorphy follow term by term.

At the flat field and $J=0$ the fluctuation symbol has real part bounded
below and is analytic in a fixed momentum strip.  Its principal
logarithm has modulus at most $C\nu$ per block.  The Fourier
coefficients therefore decay exponentially. Let $\delta_F>0$ be
a common lower bound for this strip width and the Fourier decay
rate, depending only on the earlier operator constants. The antiperiodic
momentum sum differs from its Brillouin-zone integral by
$C\nu e^{-\delta_F\ell L^{K-j}}$, by summing its nonzero Fourier
coefficients at the period multiples.  This gives the bounded
vacuum number, its thermodynamic limit, and the same uniformity.
\end{proof}

The following local form of the construction is needed before any
cluster expansion.  Its input is a configuration on a polymer;
extending that configuration to a regular field on the whole torus
is not part of the construction.

\begin{lemma}[Operators formed on a polymer]\label{C-397-47}
Assume the separation, regularity and complex-field bounds above,
with $R\ge\max\{R_0,LM_w\}$, $M\ge3M_w$ and
$c_f\min\{M,D_{\rm ex}\}\ge C_{\rm dc}\bar\kappa$.
Let $X$ be a union of fine cubes.  Suppose $\mathsf U=e^{i\Xi}U$
is given on its inner bonds, the real field $U$ satisfies the
regularity bounds at its inner plaquettes, and
$|\Xi(b)|\le\alpha_1^*L^{-n(b)}$ there.  The source is given
on its sites and satisfies
$\sup_{x\in X}\|L^{j-K}J(x)\|\le\varrho_\sharp$.
At $j=0$ an arbitrary real configuration is also allowed.
Then all walks whose supports lie in $X$, their background-map
terms, and the closed-walk normalizer on $X$ are defined from
these data and satisfy the preceding bounds with the same constants.

If $Z_\circ^\sharp\subset X$ and parameters vanish off the big
cubes contained in $X$, the two weakened covariances and
$S_j^{\rm sym}(Z_\circ,s)$ are defined, have norm at most
$\Gamma_F^2$, and obey their component and symmetry identities.
The comparison with parameters one now means parameters
$\mathbf1_X$, zero off $X$.
The normalizer can be pointed at its trace places $p$ so that
\[
 \sum_pN_p(X,\mathsf U;J)=N^{\rm F}(X,\mathsf U;J),\qquad
 |N_p(X,\mathsf U;J)|\le C_NN_fe^{-(\kappa+1)d(X)}.
\]
The same bound holds for the sum over places sent to one site
of the next level.  These objects are holomorphic at a fixed
real factor and depend on $\mathsf U$, not on its representation
as $e^{i\Xi}U$.  If the data extend to a global configuration,
the local and global constructions agree on these terms.
The identities with the whole inverse and whole determinant
remain assertions on a torus.
\end{lemma}
\begin{proof}
A local inverse reads the configuration on $W_\square$ only.
Reflect it to its mirror torus and put the identity on the
links crossing a mirror.  A crossing plaquette is the identity;
every other plaquette is copied from $W_\square$.  On a reversed
bond the copied non-unitary part is
$-U(b)^{-1}\Xi(b)U(b)$, with the same norm.  The copied regions
and source therefore satisfy the hypotheses of the weighted
inverse estimate.  This constructs the mirror inverse without
using any bonds outside $X$.

Fine cubes are unions of the blocks of level $j+1$, and all
lower blocks nest in them.  Thus a block meeting a retained
walk lies in $X$.  The averages, penalties, commutators and
finite-range factors adjacent to the walk read only inner
bonds and sites of $X$.  Their bounds are precisely those used
in the proof of \cref{C-397-40}.  The products have the same
bounds, and the restricted walks are a subset of the walks
already counted.  The row-difference argument takes place
on the mirror torus of its first cube; it also uses no global
extension.

When the parameters vanish off $X$, every walk leaving $X$
has coefficient zero.  The weakened covariance is therefore
the convergent sum of the walks just constructed.  Schur's
test and the component argument give its bound and restriction
identities.  The left walks reverse the same local factors,
so the proof of \eqref{eq:qcd-symmetric-covariance} applies
on $X$ as well.

For the normalizer, file the trace walks at $X$ and take the
trace of the diagonal block at $p$ before integrating in $y$:
\[
 N_p(X,\mathsf U;J)=-\int_0^\infty\sum_{\omega:\widehat X_\omega=X}
 \left\{\operatorname{tr}_p
 [T_A\mathcal W_\omega(\mathsf U,y;J)T_A^+\partial_yp_y]
                   -(\mathsf U\to1)\right\}\,dy.
\]
At level zero use the diagonal trace of the fine resolvent
walk itself.  Cyclicity of the trace proves the sum rule.
A diagonal block has trace bounded by $\nu$ times its operator
norm.  The places in one next-level block are fewer than
those in a threefold walk cube, which is the rank count in
\cref{C-397-42}.  The same integral and walk count therefore
give both pointed bounds.  Reversing the trace walk proves
reality at real fields and zero source.

All formulas are holomorphic products and uniformly convergent
sums of functions of the links.  The principal square root
in the normalized average is fixed by the common right
half-plane.  Consequently two allowed decompositions of the
same complex links give the same averages and local inverses,
and hence the same objects.  Restricting global data leaves
each factor and the set of retained walks unchanged.  This
proves the last assertion as an equality of formulas.
\end{proof}

On the history without large fields, a pair in a local
analyticity core has first member $e^{iL^{-(j+1)}A'}U$.
Its curvature and complex displacement bounds put it in
\cref{C-397-47}.  Thus the local objects are defined on the
whole core and are functions of its pair.  This is the
domain assertion used below when coefficients are evaluated
at complex, locally defined backgrounds.

\subsection{The early bound for the gauge coupling coefficient}

The normalizer action coefficient must be bounded before
$M$ and $\kappa$ are chosen.  The following finite-matrix
form of the four-propagator argument gives that
bound without using the later polymer norm.

\begin{proposition}[The fermion contribution to the coupling]
\label{C-397-43}
On the history without large fields, let $c_k^{\rm F}$
be the action coefficient of the field-dependent
$\log\det\mathcal K_k$, in level-$k$ units, and put
$c_0^{\rm F}=0$.  For $0\le\mathfrak m_k\le1$,
\[
 |c_k^{\rm F}|\le N_fc_2^{\rm F},\qquad
 b_{k+1}^{\rm F}=c_{k+1}^{\rm F}-L^{-1}c_k^{\rm F},
 \qquad |b_{k+1}^{\rm F}|
                 \le(1+L^{-1})N_fc_2^{\rm F},
\]
where $c_2^{\rm F}$ depends only on the fixed chamber
operator, $L$ and $b_F$.  It is independent of
$k,K,M,\kappa$ and the mass in this range.
The coefficients are real and their dependence on the
top mass is uniformly Lipschitz on this interval.
\end{proposition}

\begin{proof}
Work first at one flavour, in level-$k$ units with
$h=L^{-k}$ and one coarse block as Bloch cell.  The
operator is
\[
 K=D_h+m+W,\qquad D_h=h^{-1}D,
 \quad W=\theta_kR^+R,\quad0\le m\le1,
 \qquad G=K^{-1}.
\]
The codomain of $R$ has dimension eight, including
colour; chamber multiplicity enlarges its domain
and not its rank.  Also $\|W\|\le b_F$ at real
fields.  The scaled energy inverse estimate gives
$\|G\|\le C$.  Its proof discards nonnegative mass
terms and thus applies throughout the stated mass
interval (indeed for every nonnegative mass for
this operator bound).

We first justify all jets of the penalty that enter
the polarization.  Write a fine site as $y+z$, with
$y$ its coarse centre and $z$ in a fixed unit cube.
Strip the coarse Bloch phase $e^{i\xi\cdot y}$,
and for an external wave also strip
$e^{ip\cdot y}$.  On the cell its gauge amplitude
is then $T e^{ip\cdot z_b}$, with bond midpoint
$z_b$ of uniformly bounded size.  A fixed complex
strip of external and internal momenta has a
uniform amplitude bound.  At the flat field apply
the complete-average bound stated before \cref{C-397-29} to
\[
 X_b=h\{t_1A_{1,b}(p_1)+t_2A_{2,b}(p_2)\}.
\]
A fixed small complex disc in $t_1,t_2$ is inside
its domain, independently of $k$.  The actual
complete normalized kernel bound and Cauchy's
formula give its first two gauge jets, with up
to two needed external or internal momentum jets,
uniformly in $k$.  The lattice-normalized row of
$R$ has entries $h^{3/2}q_k^c(y,x)\iota^*$.
Thus its squared row norm, and the corresponding
jet norms, are bounded by $Ch^3\#B=C$.
The remaining phase $e^{i\xi\cdot z}$ has bounded
derivatives as well.  This is why absolute cell
coordinates cannot be used in the jet estimate.

Every gauge or momentum derivative of
$W=\theta_kR^+R$ needed here is a finite sum of
products of two such maps.  Its rank is at most
a fixed multiple of eight and its operator and
trace norms are uniformly bounded.  On complex
fields one differentiates the holomorphic row
twin $R^+$, not an antiholomorphic adjoint.
Any polarization term containing one of these
penalty jets has a uniformly bounded trace:
use its trace norm and the operator norms of
all remaining factors.  The latter remain
bounded also after momentum differentiation,
by the resolvent identity below.

For the terms without a penalty jet a uniform
operator norm is insufficient, because the cell
dimension grows as $h^{-3}$.  Put
$G_0=(D_h+m+1)^{-1}$.  In a cell fibre $\xi$ the
fine modes are $q_n=\xi+2\pi n$.  Free ellipticity
gives
\[
 \|G_0(\xi)\|_4^4
 \le dC_{\rm el}^4\sum_n(1+|\xi+2\pi n|)^{-4}
 \le C,
\]
where $d$ is the fixed one-flavour fine-site
dimension, including chamber multiplicity, and
$\|\cdot\|_4$ is the fourth Schatten norm.
The finite alias sum is bounded by the infinite
three-dimensional sum, uniformly in $h$ and $\xi$.
The resolvent identity is
\[
 G=G_0-G_0(W-1)G,
 \qquad \|G\|_4\le\|G_0\|_4
                     [1+\|W-1\|\|G\|]\le C.
\]
This is the constant-penalty resolvent expansion
in the norm appropriate for the
four-propagator count.

Only elementary matrix trace inequalities are
needed.  Hilbert--Schmidt Cauchy--Schwarz gives
\[
 \|AB\|_2^2=\operatorname{Tr}(A^*ABB^*)
          \le\|A\|_4^2\|B\|_4^2.
\]
Apply it to two adjacent pairs and once more
to their trace to obtain
$|\operatorname{Tr}(A_1A_2A_3A_4)|
\le\prod_i\|A_i\|_4$.  The ideal inequality
$\|ABC\|_4\le\|A\|\|B\|_4\|C\|$
allows any bounded vertices between the four
resolvents.

The exact second field variation of $\log\det K$
is $\operatorname{Tr}(GK_{A_1A_2})-
\operatorname{Tr}(GK_{A_1}GK_{A_2})$.
Keep its penalty vertices and penalty seagull;
they were bounded by rank above.  The mixed
bare seagull vanishes, because a gauge-dependent
hop uses one coordinate direction.  In the
symmetric momentum assignment the remaining
bare bubble is
\[
 \operatorname{Tr}\{G(\xi-p/2)V_1(\xi)
                           G(\xi+p/2)V_2(\xi)\},
\]
where the bare vertices have bounded norm and
are independent of the external momentum.
More precisely, let $\pi_{\mu\nu}(p)$ be the infinite-volume,
volume-normalized second field variation of $\log\det K$ along
the waves $e^{ip\cdot b_-}T$ and $e^{-ip\cdot b_-}T$ in
directions $\mu$ and $\nu$, with $\tfrac12\operatorname{tr}T^2=1$.
The action-coefficient convention of \cite[Section~43]{YMM1} is
$c_k^{\rm F}=-\partial_{p_1}\partial_{p_2}\pi_{12}(0)$.
The midpoint phases used above give the same second derivative,
since the polarization and its first jet vanish at zero.
For internal derivatives,
\[
 \partial_iG=-GK_iG,\qquad
 \partial_{12}G=GK_1GK_2G+GK_2GK_1G-GW_{12}G.
\]
Here $K_i$ is bounded and $K_{12}=W_{12}$ has
bounded rank, since $D_{h,12}=0$.
Every resulting term is therefore either a
four-resolvent trace bounded by the preceding
inequality or a bounded-rank trace already
treated.  Integration over the fixed coarse
Brillouin zone gives $|c_k^{\rm F}|\le c_2^{\rm F}$.
No polymer constant has entered.

The same argument holds on a fixed complex
disc about each real $m\in[0,1]$.  The scaled
inverse Neumann estimate supplies a common
disc radius; the constant-mass free resolvent
retains its fourth-norm bound after that
radius is decreased.  Cauchy's estimate now
applies to the uniformly bounded coefficient
itself, proving the Lipschitz assertion.
Holomorphy of matrix entries alone would not
have sufficed for it.  The $\gamma_4$ real
structure gives reality, and spectator
flavours multiply every trace by $N_f$.

Finally the exact determinant identity gives
$\log\det\Gamma_k=\log\det c_k+
\log\det\mathcal K_{k+1}-\log\det\mathcal K_k$.
The first determinant is field independent.
Changing from level-$k$ units to level-$k+1$
units multiplies the old action coefficient
by $L^{-1}$, giving the claimed telescope.
At $k=0$, $\Gamma_0=\mathcal K_1$, so the
correct convention is $c_0^{\rm F}=0$;
the unblocked bare determinant, which has
an infrared-divergent coefficient at zero
mass, does not enter this bound.
\end{proof}

\section{The small-field fermion step}

We first describe the gauge ingredients required by a local step.
The gauge chart has a holomorphic shift system
$h_Y(s;P,B)$ on regular complex pairs $P$, with size
$c_h\sup|B|$, supported in its wall components, compatible with
restriction, and equal to the actual shift at $s=1$ on real
background pairs.  It has a holomorphic site factor
$\mathfrak r\in\mathrm{SL}(2,\mathbb C)$ with
$|\mathfrak r-1|\le C_{\mathfrak r}\sup|B|$.  The backgrounds
compose up to the residual gauge group.  Their real curvatures
are at most $\bar C_{\rm reg}\epsilon(\gamma^{-2})L^{-2n}$;
we require this to be at most $\epsilon_*L^{-2n}$, and require
$\alpha_0L^2\le\epsilon_*$, $L\alpha_1\le\alpha_1^*$.
These bounds put a depth-$k+1$ background into the class of
\cref{C-397-47} at depth $k+1$.  Its chart field need only be
regular at depth $k$; for its depth-$k+1$ operators we use it
as a complex perturbation of the more regular background.
The distinction is needed for the vertex estimate.
The pure gauge conditions follow from
\cref{def:qcd-gauge-admission,C-396-12}; the following argument checks the
data and the matter arrays to which the single-step results apply.

\begin{lemma}[The tested block fields]\label{C-388-38}
Let the calibration frame and canonical auxiliary data be those
of the admitted gauge construction.  Assume
$\epsilon(\gamma^{-2})\le\varepsilon_\flat$ and
$2\epsilon(\gamma^{-2})\le\varepsilon_{\mathcal W}\alpha_0$,
where $\varepsilon_{\mathcal W}$ is the radius for the chart
sum rule.  Every block field $W$ passing the density test at
level $k+1$ lies in that sum-rule domain.  Its chart, axial
background, pivot inverses and fluctuation operator are defined,
and its background pairs lie in the local analyticity cores.
The set $\mathcal V^\chi_{k+1}$ of these fields is open.

For every finite-dimensional vector-valued function $F$ for
which the tested integrand is bounded and Borel, the scalar
chart identity holds coordinate by coordinate:
\[
 \mathbf T\bigl(\mathfrak G_k[\mathsf t_k]
 \chi^{\rm sf}_{\mathsf t_k}
 \mathbf1_{\{\mathcal M(V)\in\mathcal V^\chi_{k+1}\}}
 F(V)\,dV\bigr)
   =\chi_{k+1}(T_\ell)\,\mathsf I[F](W)\,dW.
\]
The right side is defined as zero off $\mathcal V^\chi_{k+1}$.
This includes the finite Grassmann algebra of the step.
\end{lemma}
\begin{proof}
The calibrations decrease and have zero increments.  The
admitted floor on $\log\gamma^{-2}$ gives
$R(Lx)\le LR(x)$ for $x\ge\gamma^{-2}$, where $R$ is the
least power of $L$ above $(\log x)^{r_c}$.  The prescribed torus
side satisfies $L^m\ge LM R(L(g^{-2}+\bar B))$.  Therefore
\[
 M R(\mathsf t_i)\le L^{m+K-i-1}.
\]
This verifies the frame geometry at every depth.  Monotonicity
of the test sizes on the final coupling range gives their
admitted bounds in that frame.

For each plaquette, the density-test construction assigns a
density cube, independently of the configuration, whose test
implies $|W_p-1|<2\epsilon(\mathsf t_{k+1})$.
All these tests pass, so
\[
 \max_p|W_p-1|<2\epsilon(\mathsf t_{k+1})
                         \le\varepsilon_{\mathcal W}\alpha_0.
\]
The admitted chart theorem at this radius supplies its
invertible pivots, axial background and operator $\mathsf B_W$;
the latter vanishes on block gauge directions.  The canonical
auxiliary data supply the interaction sum rule at this field,
including the rule that localizes an arbitrary admitted bare
array.  The background bound places each restricted pair in
the claimed core.  Strictness of the tests makes
$\mathcal V^\chi_{k+1}$ open.

Put $F'=\chi^{\rm sf}_{\mathsf t_k}
\mathbf1_{\{\mathcal M(V)\in\mathcal V^\chi_{k+1}\}}F$.
On the chart over a tested $W$, its configurations $V_W(B)$
satisfy $\mathcal M(V_W(B))=W$ and the small-field test;
hence $F'(V_W(B))=F(V_W(B))$.  Apply the bounded Borel
chart identity to each coordinate of $F'$.  Its pushforward
is carried by $\mathcal V^\chi_{k+1}$, and multiplication
by the inverse-image indicator becomes multiplication by
$\chi_{k+1}(T_\ell)$ after pushforward.  This proves the
measure identity with the stated support.
\end{proof}

\begin{proposition}[The objects of a local step]\label{C-388-34}
For the gauge construction of \cref{C-396-12}, under the displayed operator and
radius bounds, the scalar conditioned Gaussian, the chart
shift, the site factor, the local gauge vertices and
$S_k^{\rm sym}$ form a local step on each admitted pair
domain.  They are holomorphic on its slices, obey component
restrictions at zero parameters, and agree with the actual
step at parameters one on a torus.  Their joined integral
has the scalar absolute-density majorant after the integrand
is replaced by its fibre norm.  At a tested block field its
sum rule and its coefficientwise measure identity hold
pointwise as in \cref{C-388-38}.
\end{proposition}
\begin{proof}
The chart and its scalar conditioned integral are supplied
by the admitted one-step construction, independently of an
old array.  The shift bound, regularity of the background
and $L\alpha_1\le\alpha_1^*$ put every local first member
in \cref{C-397-47}.  That lemma constructs the covariance
on the polymer itself, with Gram constant $\Gamma_F$;
the same restricted data give its component identities.
For an even integrand $F$, the Gaussian contraction bound
gives $\|\mu_{S_k^{\rm sym}}F\|\le\|F\|$.
Apply the scalar majorant through each norm-one linear
functional on the finite output algebra and take the
supremum.  Holomorphy follows from the finite Gaussian
moment formulas and the locally uniform majorant.
At parameters one the covariance is the whole inverse,
and the chart identity of \cref{C-388-38} applies.
An old scalar array is admitted separately by its norm;
none of these assertions uses the pure gauge array
induction for the matter array.
\end{proof}

For a family $F=(F(X))$ of local terms put
\[
 \|F\|_{k;a}=\sup_\square\sum_{X\supset\square}
 e^{\kappa d_k(X)}\sup_{P,B}\|F(X;P,B)\|_h.
\]
The supremum is over the regular pairs of first radius $a$ and
second radius $(\alpha_1/\alpha_0)a$, and the fluctuation radius
of the step.  The radius system is constant over the levels.
It is admissible when
$0<a\le\alpha_0$ and $(1+\beta_s)\delta_k\le s^\star a$.
The background legs of level $n\le k$ have weights
$h^{(k)}(n)=C_HL^{-(k-n)}h_F$; the fluctuation legs have weight
$\Gamma_F$.  Choose
\[
 h_F\ge2A_\kappa\Gamma_F L/(L-1),\qquad
 c_\psi=2A_\kappa\Gamma_F/h_F.
\]
After a step, a leg of level $n$ splits into a formal new
background of weight $C_HL^{-(k+1-n)}h_F$ and a localized
fluctuation image of weight
$A_\kappa C_HL^{-(k-n)}\Gamma_F$.  Thus the split fits the old
weight.  The stronger displayed floor leaves room also at weight
$2\Gamma_F$ on the fluctuation.

\begin{lemma}[The vertex at one fine field]\label{C-388-14}
At the background $\mathsf U$ of level $k+1$, let
$\mathsf U'=\mathsf U\exp\mathcal A(B;\mathsf U)$ be the
fine field of the chart.  Rotate the integrated generators by
the site factor and shift them by their background.  The exponent
of the quadratic form and block kernel becomes
\[
 \bar\Psi_{k+1}\mathbb D_{k+1}(\mathsf U)\Psi_{k+1}
       +\bar W\Gamma_k(\mathsf U)W-\mathcal V_k^{\rm F}.
\]
The potential has a local decomposition into bilinear terms,
holomorphic in the pair and in $B$, vanishing at $B=0$, with
\[
 \|\mathcal V_k^{\rm F}\|_{k;a}
       \le C_Vh_F^2\sup|B|.
\]
The constant is uniform in $a$, $k$, $K$ and the volume.
Moreover its normalized fermion integral is exactly
\[
 \mu_{\Gamma_k(\mathsf U)^{-1}}e^{\mathcal V_k^{\rm F}}
 =\frac{\det\Gamma_k(\mathsf U')}{\det\Gamma_k(\mathsf U)}
 \exp\{-\bar\Psi_{k+1}
 [\mathfrak r^{-1}\mathbb D_{k+1}(\mathsf U')\mathfrak r
             -\mathbb D_{k+1}(\mathsf U)]\Psi_{k+1}\}.
\]
The first and second derivatives of the site factor belong to
this identity and are included in the average vertices.
\end{lemma}

\begin{proof}
Gauge covariance carries the quadratic form at the chart field
to that at $\mathsf U'$, with $q$ replaced by
$\mathfrak r^{-1}q$ and its row twin by $q^+\mathfrak r$.
The rotation of the integrated generators has Berezin Jacobian
$\prod\det\mathfrak r\det\mathfrak r^{-1}=1$.
The rotation of the outer generators remains in the block kernel.
At $B=0$ the two forms coincide.  Completing the square at
$\mathsf U$ leaves no linear fluctuation term; define the vertex
as the negative of the difference of the two forms.

One can estimate this difference without differentiating a long
product of Schur complements.  Let $\mathbb M$ be the one-shot
matrix on $(\psi,\Psi_k,\Psi_{k+1})$, consisting of $D+m_0$
and the two block kernels.  Its $\psi$ block is $d=\mathcal K_k$.
For the perturbed matrix $\mathbb M'$ write
$\delta\mathbb M=\mathbb M'-\mathbb M$ and
$\mathbb G=d^{-1}\oplus0$.  The Schur identity is
\[
 \mathfrak Q'-\mathfrak Q
   =\langle\bar\Phi^{\rm bg},\delta\mathbb M
                            \Phi^{{\rm bg}\prime}\rangle,
 \qquad
 \Phi^{{\rm bg}\prime}
     =\sum_{m\ge0}(-\mathbb G\delta\mathbb M)^m\Phi^{\rm bg}.
\]
Indeed multiplying a block matrix by its background column
annuls its $\psi$ component; subtracting the two such products
gives the first identity.  The two background equations give
the second by a Neumann series.

The Dirac part of a vertex at a bond has size
$CL^{-k}\sup|B|$.  A cube contains $O(L^{3k})$ bare places,
and each fine leg has weight $CL^{-k}h_F$.  Its bilinear norm
is therefore $C\nu h_F^2\sup|B|$ since
$L^{3k-k-2k}=1$.  For an average vertex the product rule gives
\[
 \|\delta r_k\|+\|\delta q\|
 \le C\sup|B|\sum_{n<k}L^{n-k}+C\sup|B|
 \le C_{\mathfrak r}\sup|B|.
\]
Each such vertex stands beside a residual
$\Psi_k-r_k\psi=\theta_k^{-1}\mathbb D_k\Psi_k$ or
$\Psi_{k+1}-q\Psi_k$, of bounded weight in level-$k$ units.
It has the same estimate.  The normalized averages, not their
unnormalized contour sums, enter this calculation.

Replace every inverse in the series by its local walks.
Their sum with the weight for positive polymer growth is
$C_K=A_\kappa C_{\rm sup}$.  Thus a link costs
$C_G\sup|B|$, where $C_G$ includes the finite-range vertex,
the at most two blocks it reads, and the recounting of its cubes.
Choose $\beta_0$ after $C_G$ so that $C_G\beta_0\le1/2$ and
the shifted fields remain in the complex class.  The series then
converges locally uniformly for $\sup|B|\le\beta_0$.

Apply the admitted shift system to a chain and take its finite
differences in the cube parameters.  A difference at a cube costs
$(e^{\bar\kappa}-1)^{-1}\le e^{-(\bar\kappa-1)}$ by Cauchy's
formula.  It vanishes unless its wall component meets the chain.
File it at the polymer containing the vertices, walks, transport
paths and parameter cubes.  If $C_g^n$ bounds the number of
connected sets of $n$ cubes through a given cube, require
\[
 x:=C_ge^{(\kappa+2)c_0-(\bar\kappa-1)}\le e^{-8}.
\]
The sum over parameter components costs at most
$\prod_{c\in R}(1+\sum_{n\ge1}x^n)\le e^{|R|/8}$, where
$R$ is the set of occupied chain cubes.  Its exponential is paid
by the geometric factor per link and the weights
$e^{|\omega|}$ in the walk sum.  This bounds the sum of chain
norms with rate $\kappa+2$ for each outer vertex.
For a fixed cube $\square$, the two spare units pay the outer
vertex location, since
\[
 d_k(Y)\ge\operatorname{dist}(c,\square)/M,\qquad
 \sum_c e^{-2\operatorname{dist}(c,\square)/M}\le37.
\]
This proves the summed family bound at rate $\kappa$.
Schwarz's lemma on $t\mapsto tB$ gives its factor $\sup|B|$.
The localization sums to the original vertex at real chart
fields because the finite differences telescope.  Both sides
are holomorphic on the connected fluctuation polydisc, so the
identity holds throughout it.  This also proves covariance of
the total and its independence of the chosen filing.

Finally integrate the completed quadratic form at the perturbed
field.  Its integrated block is exactly $\Gamma_k(\mathsf U')$:
$\mathfrak r\mathfrak r^{-1}$ cancels there.  Its outer Schur
complement is $\mathfrak r^{-1}\mathbb D_{k+1}(\mathsf U')
\mathfrak r$.  Dividing by the unperturbed determinant proves
the stated exact identity.  At $k=0$ no fine inverse is eliminated;
the same argument is its finite local expression and the series
is empty.
\end{proof}

\subsection{The joined cluster expansion}

In this subsection the conditional mean is zero.  The history
version, with its stored field and conditional mean, is proved
in \cref{C-399-41,C-399-32}.
Here are the scalar gauge inputs required for this expansion.
They are the precise inputs to be matched to the admitted pure
construction.  The conditioned integral $\mathfrak J(Z,s)$ is
linear, normalized at one, holomorphic for
$|s(\Delta)|<e^{\bar\kappa}$, agrees with the restricted Gaussian
integral at $s=1$, and factorizes on wall components when the
parameters off them vanish.  Its positive absolute-density
majorant satisfies, with the scaled fluctuation $a=\mathsf t_k^{1/2}B$,
\[
 |\mathfrak J(Z,s)f|\le\mathfrak J^{\rm abs}(Z,s)|f|,
 \quad
 \mathfrak J^{\rm abs}(Z,s)F
 \le F_0e^{c_{\rm Gauss}(\hat\varepsilon+\rho)N(Z)}
 \quad(0\le F(B)\le F_0e^{\rho \mathsf t_k\|B\|^2}).
\]
Here $c_{\rm Gauss}$ is the fixed Gaussian volume constant, and
$N(Z)\le18(LM)^3\#Z$ counts real fluctuation variables.
This is an exponential volume bound, not a uniform operator
bound in $Z$.  Finite differences in a set $\mathcal K$ of
parameters cost $(e^{\bar\kappa}-1)^{-\#\mathcal K}$.
Test defects lie in $[0,1]$ and a marked large-field bond costs
$p=e^{-\gamma_{\rm tail}\mathfrak b(\mathsf t_k)^2/2}$.

We use the scalar Mayer and cluster count for at most four
origins, each of norm at most $c_{\rm cl}=w_*$ at its input
rate, with output rate $\kappa_+=(1+4\beta_s)\kappa$.
There are two ways a scalar origin can enter this count.
An origin of sort $\Sigma_{\rm v}$ is given on the tested
fluctuation domain, and each vertex inserts the tests on its
support.  An origin of sort $\Sigma_{\rm q}$ is given at all
real fluctuations.  The quadratic correction is given on the
larger domain, but in the step it must be used with the tests
of the first sort.
The scalar count reads a potential only through its norm,
measurability and quadratic growth; it requires neither scalar
gauge invariance nor vanishing at $B=0$.  Localization reads
the pair on the resulting polymer and commutes with cube
isometries.  These scalar facts, including their rate and geometry conditions,
are the external cluster input verified in \cref{C-396-12}.

Here is the norm used for one origin, at a fixed admitted pair:
\begin{equation}\label{eq:qcd-origin-norm}
 \|\mathbb V\|_{\star,r}
 =\sup_{Y,a}
 \frac{e^{r d_k(Y)}\|\mathbb V(Y;a)\|}
 {1+M^{-3}\displaystyle\sum_{b\in\mathfrak X^{\circ+}\langle Y\rangle}|a(b)|^2}.
\end{equation}
The set $\mathfrak X^{\circ+}\langle Y\rangle$ consists of the
positively oriented nonpivot fluctuation bonds read by $Y$.
The variable is $a=\mathsf t_k^{1/2}B$: the supremum is over
the tested domain for sort $\Sigma_{\rm v}$ and all real
fluctuations for sort $\Sigma_{\rm q}$. The numerator uses
absolute value for a scalar origin and the fibre norm for an
even Grassmann-valued origin. All bounds are uniform in the
admitted pair. This is the quadratic-growth norm of the scalar
cluster construction in \cite[Section 47.2, equation (704)]{YMM1}; the label
$\Sigma_{\rm q}$ specifies the domain, not polynomial degree.
Its input rate is
$\lambda_V=(\mathsf C_{\rm ex}/L)\lambda_1+2\lambda_{\rm s}$.
The admitted rate inequalities give
$\lambda_V\le(1-2\delta_w)\kappa\le\kappa-1$.
Consequently the summed potential norm in \cref{C-388-19}
dominates $\|\mathbb V\|_{\star,\lambda_V}$ on the same domain.

\begin{lemma}[The quadratic origin and Gaussian normalization]
\label{C-388-37}
Put $\tau=x_k/\mathsf t_k$, with $|\tau-1|<\sigma_\star$.
In the scaled fluctuation $a=\mathsf t_k^{1/2}B$, let
\[
 \mathcal V^{\rm q}(\square;a)
       =-\tfrac12(1-\tau)(a,\widehat G_\square a),
 \qquad \|\mathcal V^{\rm q}\|_{\star,\lambda_V}
       \le |1-\tau|K_{\mathsf G}\le w_*/2.
\]
It vanishes on polymers other than single fine cubes.  Use
its restriction $\mathcal V^{\rm q}|_{\rm v}$ as a tested
origin in the step, and use the unrestricted family only
for its untested Gaussian logarithm
$G_\tau(X)=E^\circ[\{\mathcal V^{\rm q}\}](X)$.
For at most three additional tested origins of sizes $w_o$,
\[
 \sum_X\left(E[\{\mathcal V^{\rm q}|_{\rm v}\}
                           \cup\underline{\mathbb V}](X)-G_\tau(X)\right)
      =\log\mathsf E_k\exp\sum_{o,Y}\mathbb V^o(Y),
\]
where $\mathsf E_k$ divides the tested integral by the
untested integral of $e^{\mathcal V^{\rm q}}$.  Polymer by
polymer the difference is bounded by
\[
 \left(K_\Pi\sum_{o\ne{\rm q}}w_o+3K_P^\Pi p_*\right)
                           e^{-\kappa_+d(X)},\qquad
 p_* =\mathsf C_{\rm ex}((L+2)M)^3e^{4\lambda_1}p.
\]
These assertions hold also for even Grassmann-valued
additional origins, using their fibre norm.
\end{lemma}
\begin{proof}
On the set where all tests pass, the restricted and
unrestricted quadratic potentials coincide.  The sum of
their cube gauge-fixing operators is the operator of the
chart at a tested field, by \cref{C-388-38}.  The finite
Mayer identities and the convergent connected expansion
therefore give the displayed identity of totals, with
$G_\tau$ equal to the logarithm of the divisor.

For the local estimate split the difference into three
parts.  First remove the additional origins while keeping
the tested quadratic origin.  Every remaining term is
marked by one of the removed origins, and the marked
cluster bound is $K_\Pi\sum_{o\ne{\rm q}}w_o$.
Second change the quadratic origin from the tested to
the unrestricted sort, at the same finite threshold.
Third remove the threshold from that unrestricted
quadratic system.  The last part is marked by a failed
bond and costs $K_P^\Pi p_*$.

To estimate the second part, expand each test attached to
a quadratic vertex by
\[
 \mathsf{in}_{E'}=\sum_{E''\subset E'}
                  (-1)^{\#E''^+}\mathsf{out}_{E''}.
\]
This identity may be used at every real fluctuation because
the quadratic potential is given there.  The summand
$E''=\emptyset$ is the unrestricted Mayer summand on the
same region.  The difference consists of two sums: terms
with a nonempty added failed-bond set, and unrestricted
terms whose original failed-bond set meets the support
of a quadratic vertex.  Both have a marked failed bond.
An added bond has an endpoint in a vertex support; its
enlargement contains a big cube of the original connected
region and has tree size zero.  Give it weight
$pe^{4\lambda_1}$.  At most $6((L+2)M)^3$ such bonds
can meet a fixed big cube.  The correspondence between
the vertex set, the original failed bonds and the added
failed bonds is injective: the latter are exactly those
in the tested vertex supports.  The scalar connected
count therefore applies to the enlarged set of marks.
Each of these two sums costs $K_P^\Pi p_*$, and its
absolute majorant also justifies the refinement of the
series.  Adding the three estimates proves the bound.
For Grassmann-valued origins, contraction and duality as
in \cref{C-388-34} give the same scalar majorant.
\end{proof}

\begin{lemma}[Grassmann-valued cluster expansion]\label{C-388-19}
For a local potential define
\[
 \|V\|_r=\sup_{Z,P,B}\sup_{\square\subset Z}
       \sum_{\square\subset Y\subset Z}
          e^{r d_k(Y)}\|V(Y;P|_Y,B)\|.
\]
Here $Z$ ranges over the level-$k$ polymers, including the
torus, $P$ over its admitted pairs of radius $a$, and $B$
over the fluctuation domain. All terms in a sum use the
restriction of the same $P,B$. The fibre norm has weights
$h^{(k+1)}$ on the backgrounds and $h_W$ on the fluctuations.
The level family norm, with the supremum inside the polymer
sum and rate $\kappa$, dominates $\|V\|_{\kappa-1}$.
Let $\mathcal V=\mathcal V^{\rm g}+\mathcal V^{\natural}$ be
even, local and holomorphic on the pair radii, with
\[
 |\mathcal V^{\rm g}(Y;B)|\le\varepsilon_V
       (1+M^{-3}\sum_{b\subset Y}\mathsf t_k|B(b)|^2)e^{-\lambda_Vd_k(Y)},
  \qquad\|\mathcal V^{\natural}\|_{\kappa-1}
                      \le\varepsilon_{\natural},
 \quad \varepsilon_V+\varepsilon_{\natural}\le c_{\rm cl}.
\]
The joined gauge and fermion integral is the exponential of a
sum of local even terms.  After subtracting the $q$-only gauge
part its family $\widetilde E$ satisfies
\[
 \begin{aligned}
 \|\widetilde E(X)\|&\le K'_M
 (\varepsilon_V+\varepsilon_{\natural}+p)e^{-\kappa_+d(X)},\\
 \|\widetilde E\|_{k+1;a}
 &\le K_M(\varepsilon_V+\varepsilon_{\natural}+p),
 \qquad K_M=4096K'_M,\\
 K'_M&=\max\{K_\Pi,3K_P^\Pi\mathsf C_{\rm ex}
                         ((L+2)M)^3e^{4\lambda_1}\}.
 \end{aligned}
\]
The constants do not depend on the radius $a$ or on the
dimension of the Grassmann algebra.  The logarithm is formed
in its even commutative algebra.
\end{lemma}

\begin{proof}
For the symmetrized decoupled covariance
$S(Z,s)=S_k^{\rm sym}(Z,s)$ set
\[
 \mathfrak J_F(Z,s)F=
 \mathfrak J(Z,s)[B\mapsto\mu_{S(Z,s)}F(B)].
\]
Its Gram bound and Lem.~\ref{C-387-18} give
$\|\mu_SF\|\le\|F\|$ on the closed decoupling polydisc.
For every linear functional $\ell$ of norm at most one,
\[
 |\ell(\mathfrak J_FF)|
 \le\mathfrak J^{\rm abs}\|F\|.
\]
Taking the supremum over $\ell$ gives this bound in the fibre
norm.  Its coordinates are holomorphic in $s$, by the Gaussian
moment formulas and the scalar quadratic-growth majorant.
The gauge functional factorizes over wall components and the
fermion covariance is a direct sum there; Gaussian integration
of even elements therefore factorizes as well.

Use the same Mayer indices as in the scalar construction.
Submultiplicativity gives
$\|e^G-1\|\le e^{\|G\|}-1$.  At every real $B$, the norm
of a Mayer product is bounded by the corresponding product
with potentials replaced by their norms.  The positive
conditioning majorant bounds its integral; scalar Cauchy
estimates, applied through $\ell$, bound each parameter
bracket.  Thus each local activity is bounded by the
scalar majorant.  Scalar test defects multiply its norm by
$1-\chi$ and have the same tail bound.

To justify the logarithm, let $z(\gamma)$ be the resulting
commuting activities and let their scalar majorant satisfy
\[
 \sum_{\gamma'\nsim\gamma}\|z(\gamma')\|e^{a(\gamma')}
                                      \le a(\gamma).
\]
The scalar connected-graph bound then bounds absolutely the
series
\[
 S=\sum_{n\ge1}\frac1{n!}\sum_{\gamma_1,\ldots,\gamma_n}
 \rho^T(\gamma_1,\ldots,\gamma_n)\prod_i z(\gamma_i).
\]
The exponential identity for commuting formal variables is
coefficientwise the identity $\Xi=e^S$ for the gas.  Absolute
convergence makes it an identity in the Banach algebra;
$\Xi^{-1}=e^{-S}$.  No assertion that a Grassmann-valued
quantity is a positive number is used.  Regroup the series
by its union polymer to obtain the local logarithm.
The local logarithm to be retained is
\[
 \widetilde E(X)=E[\{\mathcal V^{\rm q}|_{\rm v},
            \mathcal V^{\rm g},\mathcal V^\natural\}](X)-G_\tau(X).
\]
Apply \cref{C-388-37}; in addition to the terms marked by
the two nonquadratic origins, it includes the change of
sort and the removal of the threshold from the quadratic
origin.  Its bound has no additive $|1-\tau|$ term.

Finally, polymers of size $n$ through a cube number at most
$2048\cdot9216^n$.  The floor
$4\beta_s\kappa\ge2\log18432+1$ in particular gives
\[
 \sum_{X\ni\square}e^{-(\kappa_+-\kappa)d(X)}\le4096.
\]
This proves the family bound.  Uniform majorants give
holomorphy.  All operations use only data on the polymer, so
they preserve locality and covariance.  Since the radii enter
only through the input bounds, outputs at two radii agree on
their common domain.
\end{proof}

\begin{lemma}[First variation of the expansion]\label{C-388-39}
Use the setting and norms of \cref{C-388-19}, with
$\varepsilon_V+\varepsilon_\natural<c_{\rm cl}$.
Let $f=(f(Y))$ be even, local and holomorphic on the same pair
and fluctuation domains, vanish at $(B,\bar W,W)=0$, and have
finite potential norm $\|f\|=\|f\|_{\kappa-1}$. At each pair,
its cluster-origin norm satisfies
$\|f\|_{\star,\lambda_V}\le\|f\|$. On the disc
$|\lambda|\|f\|<c_{\rm cl}-\varepsilon_V-\varepsilon_\natural$,
the activity is holomorphic and
\[
 \|\widetilde E[\mathcal V+\lambda f](X)
                   -\widetilde E[\mathcal V](X)\|
 \le K_\Pi|\lambda|\|f\|_{\star,\lambda_V}
                       e^{-\kappa_+d_{k+1}(X)}.
\]
Its derivative $E'[\mathcal V;f]$ is linear in $f$ and satisfies
\[
 \|E'[\mathcal V;f](X)\|
 \le K_\Pi\|f\|_{\star,\lambda_V}e^{-\kappa_+d_{k+1}(X)},
 \qquad
 \sup_\square\sum_{X\supset\square}e^{\kappa d_{k+1}(X)}
       \|E'[\mathcal V;f](X)\|\le K_M\|f\|.
\]
The same difference and derivative bounds hold when $\lambda f$
is added as a fourth tested origin to any admitted system of
at most three tested origins, for
$|\lambda|\|f\|_{\star,\lambda_V}<w_*$. The derivatives are
holomorphic in any additional parameters for which the data
are holomorphic with uniform bounds on an open set.
\end{lemma}
\begin{proof}
Keep $\lambda f$ as a separate origin. Mark every vertex belonging
to it. The terms of the expansion without marked vertices are
exactly those of the original system, with the same regions,
tests and conditional integrals. Absolute convergence therefore
identifies the difference of the expansions with the sum of
terms containing a marked vertex. The marked-origin estimate
of \cite[Thm. R and its proof]{YMM1}, transferred to the
even Grassmann algebra by the contraction and duality argument
of \cref{C-388-34}, bounds that sum by
\[
  K_\Pi|\lambda|\|f\|_{\star,\lambda_V}e^{-\kappa_+d(X)}.
\]
This argument reads only the sizes of the at most four origins.
In particular it applies to the system consisting of the tested
quadratic origin, $\mathcal V^{\rm g}$, $\mathcal V^\natural$
and $\lambda f$.

On the stated disc the last two origins can be joined into
$\mathcal V^\natural+\lambda f$. The connected sum is unchanged
by joining origins of the same sort, and the subtracted Gaussian
family is independent of $\lambda$. This proves the first bound.
Divide by $|\lambda|$ and let $\lambda\to0$. Joint holomorphy
in two insertion parameters proves linearity of the derivative.
The polymer count in \cref{C-388-19} gives
\[
 \sum_{X\supset\square}e^{-(\kappa_+-\kappa)d_{k+1}(X)}
 \le4096,\qquad4096K_\Pi\le K_M,
\]
and proves the family bound. A Cauchy integral on a smaller
insertion circle proves holomorphy in the additional parameters.
\end{proof}

\subsection{Relevant extraction and old activities}

The relevant coefficient is extracted from a translation-covariant
total.  Its individual polymer terms must nevertheless retain
gauge invariance: the estimate of an old term will compare it
with the trivial pair in a local gauge.

\begin{lemma}[Joint gauge invariance]\label{C-388-36}
At every level the local potentials, the pointed activities
and their scalar parts are invariant under the joint action
on the pair and the fermion generators.  For an $\mathrm{SU}(2)$
site map $u$, writing $\alpha_u$ for the induced automorphism,
\[
 F(X;P^u)=\alpha_uF(X;P).
\]
This holds polymer by polymer and, for a pointed family,
pointer by pointer, whenever both pairs belong to its domain.
\end{lemma}
\begin{proof}
The block kernel and its completed quadratic form are gauge
covariant.  Their difference in \cref{C-388-14} is therefore
jointly invariant.  The local inverses, walks and parameter
brackets are covariant on their supports.  The symmetric
covariance preserves this covariance, and Gaussian integration
commutes with $\alpha_u$.  Hence each Mayer term and each
connected sum has the displayed property.  In the normalizer
the covariance is under a trace, giving scalar invariance.
The operation $n^{-1}N_p$ used for pointing commutes with
$\alpha_u$, since it counts generators at a place.  The
action is extracted in its plaquette shares, each invariant,
and the mass in its local bilinear, also invariant.  These
observations prove the assertion inductively, starting with
the vertex and the normalizers.  Taking degree zero gives
strict scalar invariance of each term.
\end{proof}

\begin{lemma}[The local mass coefficient]\label{C-388-21}
Suppose a local even family $E$ has norm at most $\varepsilon$
and its total at the trivial pair is translation covariant,
cubic covariant, and invariant under colour,
$\mathrm U(N_f)\times\mathrm U(N_f)$ and charge conjugation.
After reading its legs at their backgrounds, its quadratic
kernel $\kappa(u)$ satisfies
\[
 \sum_u\|\kappa(u)\|_{\rm tr}e^{c_M|u|}
                  \le C^\psi_{am}h_F^{-2}\varepsilon,
 \quad \sum_u\kappa(u)=\mathsf{cf}_m(E)1,
 \quad \sum_u u_\mu\kappa(u)=\mathsf{cf}_z(E)\gamma_\mu,
\]
where $c_M=\min\{\alpha_D/4,\kappa/(c'M)\}$.
In particular
$|\mathsf{cf}_m(E)|,|\mathsf{cf}_z(E)|
 \le C^\psi_{am}(\nu h_F^2)^{-1}\varepsilon$.
The coefficients are additive, holomorphic, and real at real
data with the $\gamma_4$ real structure.
\end{lemma}

\begin{proof}
The degree-two fibre norm is the weighted trace norm of its
matrix kernel: a singular-value decomposition gives one
inequality, and testing against matrices of operator norm at
most one gives the other.  For a term $E(X)$, substitution
of the background kernels consequently gives a kernel
$\Sigma_X(y,y')$ with
\[
 \sum_{y,y'}\|\Sigma_X(y,y')\|_{\rm tr}
 e^{\alpha_D(d(y,X)+d(y',X))/2}
                      \le C h_F^{-2}\|E(X)\|.
\]
The factors $C_HL^{-(j-n)}$ are precisely the leg weights and
cancel.  Sum over one cell of the translation-covariant total.
The inequality
$|y-y'|\le d(y,X)+\operatorname{diam}X+d(y',X)$ and
$\operatorname{diam}X\le c'M(1+d_j(X))/2$ pays its exponential
with the background and polymer rates.  The remaining sums
are bounded uniformly, proving the kernel bound.

The common flavour phase rules out $\Psi\Psi$ and
$\bar\Psi\bar\Psi$.  Colour invariance makes the kernel the
identity in colour.  Independent flavour rotations on the
two $\tau_3$ eigenspaces give
$a_+(u)\otimes e_++a_-(u)\otimes e_-$, with identity in flavour.
The binary octahedral action on spin has a scalar part and one
vector part.  Its zeroth moment is therefore
$\mu_+e_++\mu_-e_-$ and its first moments are
$a_+\sigma_\mu e_++a_-\sigma_\mu e_-$.
Inversion has lift $1\otimes\tau_2$, exchanges the eigenspaces
and reverses $u$; hence $\mu_+=\mu_-$ and $a_+=-a_-$.
This is the asserted form.  The $\gamma_4$ conjugation fixes
the scalar zeroth moment and the first moment because
$\gamma_4\gamma_\mu\gamma_4=-\gamma_\mu$; their coefficients
are real.  The trace norm of $m1$ is $\nu|m|$.
The coefficient formulas and local uniform convergence prove
the remaining assertions.
\end{proof}

At the trivial field put $\eta_j=\theta_j/(\theta_j+\mathfrak m_j)$
and $\mathfrak s_j=\eta_j^2$.
The background of a constant block field $c$ is
$L^{-j}\eta_j\iota c$.  Thus the zeroth moment of
$L^{-j}\mathsf S_1(\psi^{\rm bg})$ is $\eta_j^2$.
At small block momentum the light background symbol is
$\lambda(p)=\theta_j(\theta_j+\mathfrak m_j+i\gamma\cdot p)^{-1}
+O(|p|^2)$.  The heavy components and nonzero folded momenta
contribute only $O(|p|^2)$.  Since the barred map is the row
twin, its symbol is also $\lambda(p)$, not $\lambda(p)^*$.
The first moment of the mass form is consequently
$-2\eta_j^2\gamma_\mu/(\theta_j+\mathfrak m_j)$; that of the
bare kinetic form is $\eta_j^2\gamma_\mu$.  The zeroth moment
of the latter is zero.  If $\eta_j\ge c_b>0$, removing a mass
coefficient
\[
 \mu=-\mathsf{cf}_m(E)/\eta_j^2
\]
makes the remaining quadratic zeroth moment vanish, with
norm cost at most
\[
 K_\mu=1+c_{\mathsf S}C^\psi_{am}c_b^{-2}/\nu,
 \qquad c_{\mathsf S}=48M^3\nu C_H^2.
\]
Indeed a cube contains $48M^3L^{3j}$ bare places, each mass
bilinear has norm at most $\nu(C_HL^{-j}h_F)^2$, and its
coefficient is $L^{-j}\mu$; the powers of $L$ cancel.
We do not remove a wave-function coefficient.

\begin{lemma}[Reading an old activity]\label{C-388-24}
Let a level-$j$ family $F$ be holomorphic on the whole pair core,
depend on both members only through their pair, and be jointly
gauge invariant there. Suppose its Grassmann degree-zero part
$F_0(X;P)$ vanishes for every polymer $X$ and every pair $P$ in
its core. Suppose also that it is covariant under its cube
translations, has norm at most $\varepsilon_j$ and has vanishing
quadratic zeroth moment at the trivial pair. With
$n=k-j$, its local reading at level $k$ has norm at most
\[
 \omega_n\varepsilon_j,
 \qquad\omega_n=C_7L^{n/4}.
\]
Its degree at least four and the gauge-dependent part of its
quadratic term have norm at most $C_7L^{-n}\varepsilon_j$.
The constants are uniform on the constant pair-radius systems.
For matrix sources the same conclusion holds whenever the
quadratic kernel admits cell units of total mass zero and
uniform first absolute moment as in the proof below.
\end{lemma}

\begin{proof}
There is no scalar term by hypothesis. A leg read $n$ levels
later costs $L^{-n}$, whereas a cube
contains $L^{3n}$ old cells.  Thus degree $2p\ge4$ costs
$L^{-(2p-3)n}$.  Take $n_0$ with
$2(2A_\kappa C_H)^2L^{-n_0}\le1$.  At $n\ge n_0$ expand
each leg first to level $j$ and then to level $k$ by its local
walks.  Each stage costs at most $2A_\kappa C_H$; the extra
factors beyond the fourth are bounded by a geometric factor
$2^{-(2p-4)}$, which also pays polynomials in the number of
legs.  The finitely many lags below $n_0$ are absorbed in $C_7$.

The quadratic comparison needs two powers of $L^{-n}$ from
the gauge field.  Fix a polymer $X$ at a cell centre $z_c$,
and fix a local or walk term in each of its two legs.  Put
\[
 D=M(1+d_j(X)),\qquad
 D'=D+9M_w\sum_i(|\omega'_i|+1).
\]
Join each column to $z_c$ through its walk cubes and a tree
of cubes of $X$.  Such paths have length at most $C_tD'$;
one can take $C_t=9$ for the tree paths.  Hold the transported
column vectors at $z_c$ fixed and call the resulting
coefficient $\mathcal E(P)$.  Both the coefficient of $F(X)$
and the first-stage walks vary with the pair.  By
\cref{C-397-47}, $\mathcal E$ is holomorphic on the open
set $\mathcal O$ where the restricted pair lies in its
core and the first member on the walk support has a small
logarithm.  It is invariant under real site gauges fixing
$z_c$ and has a uniform bound $B_{\mathcal E}$ there.

This invariance extends along complex gauge paths staying
in $\mathcal O$.  The derivative at zero of
$\mathcal E(P^{e^{\epsilon\xi}})$ is complex linear in
$\xi$.  It vanishes for the real gauge Lie algebra and
hence for its complexification.  Along a differentiable
path of complex gauges fixing $z_c$, the derivative is
therefore zero whenever the transformed pair stays in
$\mathcal O$.

Let $s_Q=\min\{M,\mathsf oLM\}$, using the cube size allowed
by the regular-pair core.  If
$2(C_tD'+M+2)>s_QL^n$, the trivial bound $2B_{\mathcal E}$
is at most
$2B_{\mathcal E}(8C_t/s_Q)^2L^{-2n}D'^2$.
Otherwise a cube $Q$ about $z_c$, of side at most $s_Q$
in level-$k$ units, contains the coefficient bonds, the
walk supports and their needed neighbours.  Its regular
real gauge, normalized to be one at $z_c$, gives in
level-$j$ units, with $\eta_j=L^{-j}$,
\[
 P=(e^{i\eta_j A_Q},J_Q),\qquad
 |A_Q|\le C_{\rm cube}L^{-n},\quad
 |\nabla^{\eta_j}A_Q|\le C_{\rm cube}L^{-2n},\quad
 |J_Q|\le C_{\rm cube}L^{-3n}.
\]
These bounds follow by taking the logarithm of the three
factors supplied by the shift, the complex displacement
and the real regular gauge.  The logarithm is Lipschitz
on the admitted ball.  Converting a covariant difference
to a plain difference adds a product of connection and
displacement sizes, which has the required $L^{-2n}$.

Let $\rho_A$ be the logarithmic radius whose pairs lie in
$\mathcal O$.  The radial disc through $(A_Q,J_Q)$ first
gives $2B_{\mathcal E}C_{\rm cube}L^{-n}/\rho_A$ by
Schwarz's lemma; if its radius is below one use
$2B_{\mathcal E}$.  For large $n$ remove the connection
at the centre.  Set
$A^c_\nu=A_Q(\langle z_c,z_c+\eta_je_\nu\rangle)$ and
\[
 g_v(x)=h_3h_2h_1,\qquad
 h_\nu=\exp(iv(x-z_c)_\nu A^c_\nu),\qquad0\le v\le1.
\]
The matrices need not commute, and this order is kept.
Writing $g^>_\mu,g^<_\mu$ for the products before and
after $h_\mu$, the connection on a $\mu$-bond becomes
\[
 A_{Q,v}=\operatorname{Ad}(g^>_\mu h_\mu)
 \mathsf S(\operatorname{Ad}(g^<_\mu)A_Q,-vA^c_\mu),
 \qquad e^{i\eta_j\mathsf S(A,B)}=e^{i\eta_jA}e^{i\eta_jB}.
\]
Each product of the $h_\nu^{\pm1}$ is bounded by
$e^{3C_{\rm cube}M/2}$.  The Lipschitz estimate for
$\mathsf S$ gives a constant $C_{\rm lg}$ such that
\[
 |A_{Q,v}|\le C_{\rm lg}L^{-n},\quad
 |\nabla^{\eta_j}A_{Q,v}|\le C_{\rm lg}L^{-2n},\quad
 |\operatorname{Ad}(g_v)J_Q|\le C_{\rm lg}L^{-3n}.
\]
In the difference estimate, differentiating a conjugating
factor contributes the product of two $L^{-n}$ factors.
At $v=1$, $\mathsf S(A^c_\mu,-A^c_\mu)=0$; subtracting
the centre value gives
$|A_{Q,1}(b)|\le C_{\rm lg}L^{-2n}|x-z_c|_1$.
Choose $n_{\rm lg}$ so
$C_{\rm lg}L^{-n_{\rm lg}}\le\rho_A$.
For $n\ge n_{\rm lg}$ every transformed pair remains
in $\mathcal O$.  On the bonds read by the term,
$|x-z_c|_1\le6C_tD'$, so the endpoint admits a radial
disc of radius at least
$\rho_A/(6C_{\rm lg}C_tL^{-2n}D')$.
The gauge invariance along the path and Schwarz's lemma
on this disc give
\[
 \|\mathcal E(P)-\mathcal E(1,0)\|
       \le2B_{\mathcal E}C_{\rm gf}L^{-2n}D'^2.
\]
The trivial bound covers radii below one.  The finitely
many intermediate lags use the first radial estimate
and $L^{-n}\le L^{n_{\rm lg}}L^{-2n}$.

To sum these bounds, fix each walk simultaneously at all rows of
$X$. A walk term vanishes at rows for which its first cube is
not an allowed first cube; for a barred leg use the last cube
of its row twin. The rows in a level-$j$ block have the same
allowed first cubes, at most $C_{\rm nb}$ of them. Consequently
the number of possible first cubes, and likewise of last cubes,
is at most
\[
 N_X=C_{\rm nb}M^3\#X
       \le8C_{\rm nb}M^3(1+d_j(X)).
\]
For a fixed first cube the positive weighted walk sum is at
most $A_\kappa$. Including the one local term, the sum of the
suprema of the row bounds for one leg at level $n'\le j$ is
therefore at most
\[
 2C_HL^{-(j-n')}(1+N_XA_\kappa)
 \le2C_HL^{-(j-n')}(1+8C_{\rm nb}M^3A_\kappa)(1+d_j(X)).
\]
Apply this to the two legs in $B_{\mathcal E}$. Absolute
convergence and bilinearity identify their sum with the
quadratic activity. The additional factor $(1+d_j(X))^2$
must be retained when summing the bounds of its individual
holomorphic terms. Together with the path factor $D^2$, it
costs at most
\[
 M^2\sup_{d\ge0}(1+d)^4e^{-7\kappa d/8}<\infty.
\]
This number and $(1+8C_{\rm nb}M^3A_\kappa)^2$ are included
in $C_7$ and the later old-activity constants, before choosing
the fibre weight. Higher-degree terms are evaluated at one
pair and retain the ordinary row walk count.

The two legs contribute $L^{-2n}$, the gauge comparison
$L^{-2n}$, and the cell count $L^{3n}$, giving
$CL^{-n}\varepsilon_j$ per new cube.
Enlarge the filing polymer by the at most eight level-$k$
cubes meeting $Q$.  They lie in a cube of side $2M$ and
cost the fixed factor $e^{\kappa c_0}$.  The polynomials
in $D'$ are paid by the old polymer rate and the walk
weights $e^{|\omega'|}$.  The complex transports are
paid by the same rates under
$2C_t\alpha_1M\le\kappa$ and
$18C_t\alpha_1M_w\le1$.

It remains to estimate the trivial-pair quadratic form.  Write
its kernel as units $\beta_u(y,y')$ filed at a cell $c(u)$,
with
\[
 \sum_{u:c(u)=c}\sum_{y,y'}\|\beta_u(y,y')\|
 (1+|y-z_c|+|y'-z_c|)\le Ch_F^{-2}\varepsilon_j,
 \qquad \sum_{u:c(u)=c}\sum_{y,y'}\beta_u(y,y')=0.
\]
At zero sources these are the polymer kernels; translation
covariance makes the last sum equal to the vanishing mass
$M^3\mathsf{cf}_m(E)1$.
At sources the stated cell-unit condition is precisely what
is used.  For each unit, carry both legs along its paths to
$z_c$ and write
$\Phi(y)=\Phi(z_c)+[\Phi(y)-\Phi(z_c)]$.
The reference value of a first-stage leg is its row sum
at the trivial configuration times $\Phi(z_c)$.  Its
remainder is a sum of transported differences.  Their
paths use straight bare transports, whereas the kernel
estimate uses the iterated block field.  At neighbours
$x,x'$ one must include both terms in
\[
 \Phi(x)-\mathsf U([x,x'])\Phi(x')
 =\Phi(x)-(\mathcal M^j\mathsf U)(x,x')\Phi(x')
 +[(\mathcal M^j\mathsf U)(x,x')-\mathsf U([x,x'])]\Phi(x').
\]
The complete-average comparison bounds the last transport
defect by $2(c_T\epsilon_*L^{-2n}+C_Q\alpha_1L^{1-n})$.
Here the block-centre kernel differs from one by at most
$C_Q\alpha_1L^{-n}+14\epsilon_*^2L^{-4n}$, so that
comparison of the normalized kernels also compares the
two transports.  Multiplication by the leg size
$2C_HL^{-n}$ gives the bound
$4C_H(c_T\epsilon_*+C_QL\alpha_1)L^{-(1+\alpha)n}$
for every $\alpha<1$.  Thus the extra term has the
required H\"older gain.  The terms without differences
cancel cell by cell before localization.  A term with a difference gains
$L^{-3n/4}$ by the background estimate with $\alpha=3/4$;
the other leg costs $L^{-n}$.  Its cell bound is
$C\varepsilon_jL^{-11n/4}$, hence after $L^{3n}$ cells it
is $C\varepsilon_jL^{n/4}$.  The first absolute moment
pays the path length.  Moving a diagonal unit to the cell of
its site changes a new cube sum only by units near its
boundary; their number $O(L^{2n})$ and leg cost $L^{-2n}$
give $C\varepsilon_j$.

Localize these differences by the walk bounds with their
difference gain.  For completeness, the exact hull inequality
for the tree size is $d_{j+1}(X^{\rm up})\le L^{-1}d_j(X)$,
where $X^{\rm up}$ is the union of coarse cubes containing a
fine cube of $X$. In coordinates with fine corners in
$\mathbb Z^3$ and coarse corners in $L\mathbb Z^3$, round
$Lq+t$, $0\le t<L$, to $Lq$ if $t\le\sigma$ and to
$L(q+1)$ otherwise, choosing $\sigma\in\{0,\ldots,L-1\}$
in each coordinate. Every lattice point of a closed coarse cube is
mapped to one of its corners. Thus the image of the tree is connected,
uses edges of coarse cubes in $X^{\rm up}$, and meets each of
them: a vertex meeting a fine closed cube maps to a corner of
its containing coarse cube. A coordinate edge survives for
exactly one cut position. The mean number of image edges is
at most $L^{-1}$ times the original count; some choice does
no worse, and a spanning tree has no more edges.
Iterating gives $d_k(X^{\rm up})\le L^{-n}d_j(X)$.
Add the walk polymers one at a time and use the positive
parts of their size increments; size need not be monotone
under enlargement.  Their weighted sums cost $A_\kappa$ per
leg.  The spare old rate
$\kappa(1-L^{-n})d_j(X)$ pays all remaining powers of the
diameter, and $e^{|\omega|}$ pays the path-length polynomials
and the recounting of cubes reached by a walk.  This proves
the asserted local family norm.

On a torus the paths include the antiperiodic seam.  A
nonwrapping piece and its lift have the same cell number,
because the seam signs at the two reference legs cancel.
The torus cell sum minus its infinite-lattice value is
therefore the sum of wrapping pieces minus that of
infinite-lattice pieces that do not fit.  A piece in
either sum has $d_j(Z_p)\ge\ell_j/(2M)-2$.
File a wrapping piece at its coarse hull and a nonfitting
piece at the image of its coarse hull; both wrap.
Their two remaining legs and the count of cells give
\[
 C\varepsilon_jL^n\sum_{d\ge\ell_j/(2M)-2}(1+d)^3
       e^{-\kappa(1-L^{-n})d}\le C\varepsilon_j,
\]
since $\ell_j/(2M)\ge L^n$.  On a nonwrapping output
polymer every retained term is the image of its lifted
term, so the local constructions agree there.
The argument at matrix sources uses the stated cell
units and their first absolute moments.  On a smaller
constant-radius system, $C_{\rm cube},C_{\rm lg}$ and
$\rho_A$ scale by at most $a/\alpha_0$; their ratios and
the lag threshold are no larger.  This proves uniformity
on those systems.
\end{proof}

The potential of a complete step consists of the gauge vertex
and the variations of the old scalar terms, together with
\[
 \mathcal V^{\natural}=\mathcal V_k^{\rm F}
       +\mathcal V^{\rm s}+\mathcal V^{\rm old}_{\ge2}.
\]
All terms vanish at $(B,\bar W,W)=0$, although the fermion
part need not vanish at $B=0$.  At the fixed pair radii,
with $2a_{\rm an}=\mathsf c_{\rm os}\alpha_0$, put
$\mathsf s_a=c_h\delta_k/a_{\rm an}$. Apply \cref{C-388-24}
to $E_{\ge2}^{(j)}$, whose degree-zero part vanishes identically.
The preceding bounds, Cauchy's estimate for an old coefficient in
the shift, and
the fluctuation-leg room give
\[
 \|\mathcal V_k^{\rm F}\|\le C_Vh_F^2\delta_k,
 \quad \|\mathcal V^{\rm s}\|
       \le(c_\psi+C_G\delta_k)c_{\mathsf S}h_F^2|s_k|,
\]
\[
 \|\mathcal V^{\rm old}_{\ge2}\|
 \le(c_\psi+\mathsf s_a+C_G\delta_k)
                         \sum_{j\le k}\omega_{k+1-j}\varepsilon_j.
\]
To obtain the factor $c_\psi$, expand a homogeneous product
under $a+b$.  If $h_a+h'_b\le h$ and $h_b\le h'_b$, its part
with at least one $b$ has norm at most $(h_b/h'_b)\|F\|_h$;
this is the convexity inequality
$(h_a+h_b)^n-h_a^n\le(h_b/h'_b)
[(h_a+h'_b)^n-h_a^n]$.  Apply it to the localized $W$ images.
The pure scalar localization gives
$|\mathcal V^{\rm old}_0|\le\mathsf s_aK_{\rm old}$ with
$K_{\rm old}\le C(E_0+C_NN_f+B_\sharp+1)$.
At radius $a$, replace $\mathsf s_a$ for an old fermion
coefficient by $(\alpha_0/a)\mathsf s_a$; an explicit old
scalar family keeps its fixed-radius bound, while its
fluctuation family has factor $C_{\mathcal A}\delta_k/a$.
Admissibility gives $(\alpha_0/a)\mathsf s_a\le1/9$.
Thus these estimates also hold on each admissible constant
radius system.  They supply the input sizes for
Lem.~\ref{C-388-19}. Their simultaneous use in the matter
induction is verified in \cref{C-396-12}.

\section{The cumulative mass term}

We must control the mass terms accumulated over all preceding
scales. The three-square identity first gives the sign of each
Gaussian slice. We then bound their sum and identify it with
the entire second-order contribution of the interacting step.
Finally we estimate the cubic remainder and close the forward
mass recursion. These are separate parts of the argument;
positivity of the slices alone does not bound their sum.

\begin{lemma}[The three-square identity]\label{C-398-4}
Let $m\ge0$, $a=\mathfrak a_k$, $R=R_k(U)$, $P=R^*R$,
$\eta=a/(a+m)$ and
$\varphi=a(D[U]+m+aP)^{-1}R^*\chi$, whenever the inverse exists.
Then
\[
 \operatorname{Re}\langle\chi,\mathsf D^{(k)}[U;m]\chi\rangle
 -\frac{am}{a+m}\|\chi\|^2
 =(a+m)\|R\varphi-\eta\chi\|^2+
     m\|(1-P)\varphi\|^2+\langle\varphi,\Delta\varphi\rangle.
\]
The right side is nonnegative and vanishes at the flat field on
constant chamber-independent fields. Consequently, for constant
chamber-independent $\chi$, the coefficient of order two in a
real gauge fluctuation about the flat field, and its mean under
every centered Gaussian slice, are nonnegative.
\end{lemma}

\begin{proof}
The defining equation is
$(D+m)\varphi=aR^*(\chi-R\varphi)$ and
$\mathsf D^{(k)}\chi=a(\chi-R\varphi)$.  Therefore
\[
 \operatorname{Re}\langle\chi,\mathsf D^{(k)}\chi\rangle
 =a\|\chi-R\varphi\|^2+
             m\|\varphi\|^2+\langle\varphi,\Delta\varphi\rangle.
\]
Use $RR^*=1$ and
$\|\varphi\|^2=\|R\varphi\|^2+\|(1-P)\varphi\|^2$.
Completing the remaining square uses
\[
 a\|\chi-u\|^2+m\|u\|^2
 =(a+m)\|u-\eta\chi\|^2+\frac{am}{a+m}\|\chi\|^2.
\]
At the flat field $D[1]R^*\chi=0$ for constant $\chi$, so
$\varphi=\eta R^*\chi$ and all three terms vanish.  Each square
vanishes at zero fluctuation; its quadratic coefficient is a sum of
squares of first derivatives.  Gaussian averaging preserves this
nonnegativity.
\end{proof}

For a bare mass $m$, put
\[
 \mathsf r_k(U;m)[\chi]=\eta^{-2}
 \left(\operatorname{Re}\langle\chi,
             \mathsf D^{(k)}[U;m]\chi\rangle
            -\frac{\mathfrak a_km}{\mathfrak a_k+m}\|\chi\|^2\right),
 \qquad\eta=\mathfrak a_k/(\mathfrak a_k+m).
\]
Let $\Phi_k(A;m)$ be its coefficient of $t^2$ at $U=e^{itA}$
on constant block fields, per block site and with normalized
trace over their fibre.  If $\gamma_{k,i}$ is the unit-curvature
slice defined below, set
\[
 \sigma_i(k;m)=\int\Phi_k(A;L^{-k}m)\,\gamma_{k,i}(dA),
\]
so the argument of $\sigma_i$ is in level-$k$ mass units.
For the label mass $\mathsf m$ let
$\mathfrak m_k=L^{k-K}\mathsf m$ and set
\[
 F(k;m)=\sum_{i<k}\sigma_i(k;m),\qquad
 G_k(x)=\sum_{i<k}\frac{L^{k-i}}{x_i}
                         \sigma_i(k;\mathfrak m_k),\qquad
 \mu_k^{\mathrm{1L}}(x)=G_k(x)-LG_{k-1}(x),\quad G_0=0.
\]
The boundedness of $F$, and the fact that the entire second order of
the interacting step is $\mu_k^{\mathrm{1L}}$, are separate estimates;
the three-square identity alone does not prove them.

\subsection{A bound for the sum of the gauge slices}

For chamber-wall types $j\in\{0,1,2,3\}$ put
$\varrho_j=rt_j$ and $\nabla_j[U]=1-R_j[U]$.
Thus $\langle f,\Delta f\rangle
=\tfrac12\sum_j\varrho_j\|\nabla_j[U]f\|^2$;
this chamber difference is distinct from the site-bond difference.
Write $\mathbf Q_k^s$ for the composite straight gauge average and
$Q'_k$ for the composite site mean.  The linearized gauge block
map has the form
\[
 \mathbf Q_k^{\rm lin}=\mathbf Q_k^s-d\bar n_k,
 \qquad \mathbf Q_k^{\rm lin}d\lambda
                  =d(\lambda|_{L^k\mathbb Z^3}).
\]
This follows directly by linearizing the contour averages: their
first variation on a constant spinor is its contour mean
$i\bar n_k(A)$.  Put $T_k=\ker\mathbf Q_k^{\rm lin}$ and
$\mathfrak G_k=\{d\lambda:\lambda|_{L^k\mathbb Z^3}=0\}$.
The gauge Gaussian slice $i$ has curvature covariance
$\mathcal C_{(i+1)}=P_{i+1}-P_i$, where $P_i$ projects onto
the closed curvatures whose level-$i$ plaquette average vanishes.
These are the linearized, zero-background gauge inputs used here.
They concern the pure Gaussian construction, not matter estimates.

The quadratic coefficient of the three squares can be written
explicitly. Write $P_k^{\rm F}=R_k^*R_k$ for the fermion
projection, distinct from the curvature projection $P_k$.
At the flat field let $\chi_c=R_k^*c$,
$\kappa_1=\partial_tR_k(e^{itA})^*c|_{t=0}$,
$V_A=\partial_tD[e^{itA}]|_{t=0}$, and
$\nabla_j^{(1)}=\partial_t\nabla_j[e^{itA}]|_{t=0}$.
Choose the representative with $\bar n_k(A)=0$ and set
$\zeta=\mathcal K_k[1;m]^{-1}(V_A\chi_c-m\kappa_1)$.
Differentiating the three squares, whose constant terms vanish,
gives
\[
 \Phi_k(A;m)=\left\langle
 (\mathfrak a_k+m)|R_k\zeta|^2+
 m|(1-P_k^{\rm F})\zeta+\kappa_1|^2+
 \sum_j\frac{\varrho_j}2
        |\nabla_j^{(1)}\chi_c-\nabla_j\zeta|^2
 \right\rangle_c.
\]
Here $\langle\cdot\rangle_c$ is normalized trace over constant
block spinors per block.  For the proof, the exact identity
$\varphi_t=\eta(1-\mathcal K_t^{-1}D_t)R_t^*c$ gives
$R_t\varphi_t-\eta c=-t\eta R\zeta+O(t^2)$,
$(1-P_t^{\rm F})\varphi_t=-t\eta((1-P_k^{\rm F})\zeta+\kappa_1)+O(t^2)$,
and $\nabla_j[U_t]\varphi_t=
t\eta(\nabla_j^{(1)}\chi_c-\nabla_j\zeta)+O(t^2)$.
Here $P_t^{\rm F}=R_k(U_t)^*R_k(U_t)$.
Dividing the quadratic coefficient by $\eta^2$ proves the formula.
Gauge covariance makes it invariant under $\mathfrak G_k$.
Indeed replacing $A$ by $A+d\lambda$ is the first-order gauge
transform by $e^{-it\lambda}$; the nonnegative block form is
zero at $t=0$, so its quadratic coefficient depends on this
tangent vector alone.  Its colour trace is the sum of the same
quadratic form in the three colour components.

In a Bloch fibre $\xi\ne0$, $d$ identifies
$T_k/\mathfrak G_k$ with the closed curvatures of vanishing
level-$k$ average.  If two representatives have the same curl,
their difference is $d\lambda$; the linearized constraint
forces $\lambda$ to be constant on the coarse sites, and this
constant can be subtracted.  Conversely choose $\lambda$ with
$\lambda|_{L^k\mathbb Z^3}=\bar n_k(B)$ for
$B\in\ker\mathbf Q_k^s$ to put $B+d\lambda$ in $T_k$.
The slice spaces are nested and
\[
 \sum_{i<k}\mathcal C_{(i+1)}=P_k.
\]
Consequently $F(k;m)$ is the Gaussian mean of $\Phi_k$ on
$T_k/\mathfrak G_k$ with density $e^{-|dA|^2/2}$.
This identity is essential: the sum is bounded as one Gaussian
mean, rather than by $k$ separate bounds.

\begin{lemma}[The constrained covariance]\label{C-398-11}
Use level-$k$ units, $h=L^{-k}$, fine modes
$q_n=\xi+2\pi n$, $\hat q_{n\mu}=2h^{-1}\sin(hq_{n\mu}/2)$,
and $k_n=\sum_\mu\hat q_{n\mu}^2$.
At unit spacing put $k_1(t)=4\sum_\mu\sin^2(t_\mu/2)$.
There is a representative $A''$ of the quotient satisfying
$\mathbf Q_k^sA''=0$ and $(1-\widetilde P_k)d^*A''=0$,
where $\widetilde P_k$ is the projection onto block constants.
For a nonnegative quadratic observable invariant under
$A\mapsto A+d\lambda$, $Q'_k\lambda=0$, its mean under
this representative equals its mean under the covariance
\[
 \widehat C=\widehat{\mathbb K}^{-1/2}
 \{(1-\Pi)+\Upsilon^{-1}|(1-\Pi)w\rangle
                          \langle(1-\Pi)w|\}
                   \widehat{\mathbb K}^{-1/2}
 \le\Upsilon_0^{-1}\widehat{\mathbb K}^{-1},
 \quad\Upsilon_0=(2/\pi)^6,
\]
in each colour, where $\widehat{\mathbb K}$ acts as $k_n$.
Here $\Pi$ projects orthogonally onto the three constraint
vectors $e_\mu$, $w$ is the unit vector representing the block
part of $d^*A$ in the coordinates
$a=\widehat{\mathbb K}^{1/2}A$, and $\Upsilon=|\Pi w|^2$.
These vectors are constructed in the proof.
\end{lemma}

\begin{proof}
Choose first a gauge with $\bar n_k=0$.  A further gauge
change with $Q'_k\lambda=0$ uniquely solves
$(1-\widetilde P_k)d^*d\lambda=
 -(1-\widetilde P_k)d^*A$:
its homogeneous equation has energy $|d\lambda|^2=0$ and
block mean zero, so it has no kernel in the finite Bloch fibre.
For $a=\widehat{\mathbb K}^{1/2}A$, the three straight-average
constraints are $a\perp e_\mu$, with orthonormal vectors
$e_\mu$ supported on their coordinate components.  The block
part of $d^*A$ is $\langle w,a\rangle$ for a unit vector $w$.
The gauge-fixed form is therefore
$|a|^2-|\langle(1-\Pi)w,a\rangle|^2$ on
$\operatorname{ran}(1-\Pi)$, where
$\Pi=\sum_\mu|e_\mu\rangle\langle e_\mu|$.
Inverting this rank-one perturbation gives the displayed
covariance with $\Upsilon=|\Pi w|^2$.

If $v_n$ is the symbol of the block mean, put
$S=\sum_nv_n^2/k_n$ and
$S_\mu=\sum_nv_n^2/(\hat q_{n\mu}^2k_n)$.
The straight-average identity $\mathbf Q_k^sd=dQ'_k$
gives $|\langle e_\mu,w\rangle|^2=S^2/S_\mu$.
Since $\hat q_{n\mu}^2\ge\hat q_{0\mu}^2$,
\[
 \Upsilon=\sum_\mu S^2/S_\mu\ge Sk_0\ge v_0^2
 =\prod_\mu\left[\frac{\sin(\xi_\mu/2)}
                       {h^{-1}\sin(h\xi_\mu/2)}\right]^2
 \ge(2/\pi)^6.
\]
The inverse perturbation is at most $\Upsilon^{-1}$, proving
the covariance bound.  Every field in the straight-average
constraint space splits uniquely as $A''+d\lambda$ with
$Q'_k\lambda=0$.  Its form splits into $|dA''|^2$ plus the
independent gauge energy
$|(1-\widetilde P_k)d^*d\lambda|^2$.  Thus an invariant
observable has the same mean under the two laws.
\end{proof}

To apply this lemma to $\Phi_k$, replace $\kappa_1$ by
$(1-P_k^{\rm F})\kappa_1$ in its second square and in $\zeta$; call
the resulting expression $\Phi_k^{\rm nat}$.  It agrees with
$\Phi_k$ when $\bar n_k=0$.  Under $Q'_k\lambda=0$, its
$\zeta$ changes by $i\lambda\chi_c$ and
$(1-P_k^{\rm F})\kappa_1$ by $-i\lambda\chi_c$.
The first square is unchanged because $R_k\lambda\chi_c=0$;
the second and third are unchanged by cancellation.
It therefore has the covariance representation just proved.

\begin{proposition}[Bounded cumulative loop]\label{C-398-4-bound}
There is $B_F<\infty$, depending only on $L$, $b_F$ and the
fixed chamber operator, such that $0\le F(k;m)\le B_F$ for
all $k\ge1$ and $0\le m\le1$ in level-$k$ units.
\end{proposition}

\begin{proof}
In a Bloch fibre put $d_n=h^{-1}\widehat D(hq_n)$,
$g_n(m)=(d_n+m)^{-1}$ and $R=(v_n\iota^*)_n$.
The one-shot Schur formula gives
\[
 D_k(\xi;m)=
 [\theta_k^{-1}+\sum_nv_n^2\iota^*g_n(m)\iota]^{-1}.
\]
The free estimates give
$\|\widehat D(p)^{-1}\|\le c_1/|p|$ and
$\|\operatorname{Re}\widehat D(p)^{-1}\|\le c_{\rm re}^{-1}$.
For the latter, the free real part is bounded above by a
constant times $|p|^2|\Pi_0f|^2+|\Pi_1f|^2$ near zero:
the within-site terms annihilate $\Pi_0$ and the remaining
hop is a difference.  The free ellipticity estimate bounds
this by $C|\widehat D(p)f|^2$.  Multiplying by the inverse
gives the assertion; compactness treats the rest of the zone.
Also $\|\Pi_1g_n(m)\|\le Ch$, by the heavy block inverse
and its difference coupling.  The gauge vertex matrices are
uniformly bounded.

First take $m=0$.  Accretivity gives $\|D_k\|\le\theta_k$;
near zero the light block is $i\gamma\cdot\xi+O(|\xi|^2)$,
$v_0^2\ge\Upsilon_0$, and the remaining modes give a bounded
sum.  Hence $\|D_k\|\le C|\xi|$ there.  With $X=RD^{-1}$,
$R\mathcal K^{-1}=\theta_k^{-1}D_kX$.
The Gaussian mean of the first square is
\[
 3\theta_k^{-1}\int\frac{d^3\xi}{(2\pi)^3}
 \operatorname{Tr}(D_kX\widehat V\widehat C
                    \widehat V^*X^*D_k^*).
\]
Its trace before the two $D_k$ factors is bounded by
$C\sum_nv_n^2|q_n|^{-2}k_n^{-1}\le C|\xi|^{-4}$.
Thus its integrand is at most $C|\xi|^{-2}$, integrable
uniformly in $k$ in three dimensions.

For the third square write
$\zeta=\zeta^0-\zeta^1$, with
$\zeta_n^0=g_nV_n\chi_c$ and
$\zeta^1=gR^*D_kX V\chi_c$.
The gradient energy of $g_nf$ is at most
$hc_{\rm re}^{-1}|f|^2$ because
\[
 h^{-1}\sum_j\frac{\varrho_j}2|\nabla_jg_nf|^2
       =\langle f,\operatorname{Re}g_nf\rangle.
\]
The block correction therefore costs at most $Ch$ times the
already bounded first square.  For the main part, bound
$\widehat C$ by $\Upsilon_0^{-1}\widehat{\mathbb K}^{-1}$
and unfold the fine modes.  The substitution $q=t/h$ gives
the same finite integral for every $h$:
\[
 \int_{[-\pi,\pi]^3}\frac{d^3t}{(2\pi)^3k_1(t)}
 \sum_j\frac{\varrho_j}2
  \|\widehat\nabla_j^{(1)}-
             \widehat\nabla_j\widehat D^{-1}\widehat V\|_{\rm HS}^2
 \le C(1+c_{\rm re}^{-1})\int\frac{d^3t}{k_1(t)}<\infty.
\]
The inequality uses
$(D^{-1})^*\Delta D^{-1}=\operatorname{Re}D^{-1}$.
Using only $\|D^{-1}\|$ here would lose this bound.

We next control the average variation in the mass-dependent
second square.  The exact contour recursion is
\[
 \bar n_{j+1}(A)=\bar n_j(A)|_{L^{j+1}\mathbb Z^3}
                        +\mathsf k(\mathbf Q_j^sA),
 \quad
 (1-P_k^{\rm F})\kappa_1=-i\sum_{j<k}\mathfrak n_A^{(j)}\chi_c,
\]
where $\mathfrak n_A^{(j)}$ is the contour mean of the
straight average $\mathbf Q_j^sA$ in its own step, minus its
block mean.  To prove the recursion insert
$\mathbf Q_j^{\rm lin}=\mathbf Q_j^s-d\bar n_j$ and
$\mathsf k(d\vartheta)=Q'\vartheta-
\vartheta|_{L^{j+1}\mathbb Z^3}$.  In the first variation
of the composite fermion average the remaining $\bar n_j$
terms telescope, giving the second identity.
Each $\mathfrak n_A^{(j)}$ has block mean zero.

Under the covariance bound, a link of $\mathbf Q_j^sA$ has
variance at most $C_QL^j$ in bare units.  Indeed its multiplier
obeys
\[
 |\hat f_j(t)|\le L^{-3j}
       \prod_\nu\min(L^j,\pi/|t_\nu|)
                    \min(L^j,\pi/|t_\mu|),
\]
and its variance is at most
$\Upsilon_0^{-1}\int|\hat f_j(t)|^2/k_1(t)\,d^3t$.
Rescaling $t=L^{-j}s$ bounds this by $C_QL^j$; the remaining
integral of
$|s|^{-2}\prod_\nu\min(1,\pi^2/s_\nu^2)
\min(1,\pi^2/s_\mu^2)$ is finite.
A contour uses at most $3L$ links.  Minkowski's inequality
therefore bounds the standard deviation of the sum by
$CL\sum_{j<k}L^{j/2}\le C_LL^{k/2}$.
The bare mean square is at most $C_LL^k$, or $C_L$ in
level-$k$ units.  This controls
$\mathbb E|(1-P_k^{\rm F})\kappa_1|^2$ uniformly.

Finally $\|\mathcal K(m)^{-1}\|\le C_{\mathcal K}$ in each
fibre, uniformly in $m\ge0$ and $k$.  The Woodbury formula
proves it: nonzero folded modes have a uniform gap, while on
the zero folded light mode
\[
 \mathcal K(0)^{-1}_{00}
 =S_0^{-1}D_k[\theta_k^{-1}+
                 \sum_{n\ne0}v_n^2S_n^{-1}],
\]
where $S_n$ is the light Schur complement of $d_n$ at mass zero.
The product $S_0^{-1}D_k$ is bounded near $\xi=0$. The heavy blocks
are bounded as above.  Since $\mathcal K(0)^{-1}$ is accretive,
$(1+m\mathcal K(0)^{-1})^{-1}$ has norm at most one.
The same cancellation and the covariance bound give
$\mathbb E|\zeta_0|^2\le C$: the nonzero-mode sum is bounded
by $C\sum_{n\ne0}|q_n|^{-2}k_n^{-1}+C|\xi|^{-2}$ and the
zero-mode part by $C|\xi|^{-2}$, all integrable.
The resolvent identity gives
\[
 \zeta(m)-\zeta_0=-m\mathcal K(m)^{-1}
                         [\zeta_0+(1-P_k^{\rm F})\kappa_1].
\]
The energy identity for $u=\mathcal K(m)^{-1}f$ bounds its
gradient and block energy by
$\operatorname{Re}\langle u,f\rangle\le C_{\mathcal K}|f|^2$.
Thus the first and third squares at $m$ are at most twice
those at zero plus a bounded multiple of
$m^2\mathbb E|\zeta_0+(1-P_k^{\rm F})\kappa_1|^2$; the second is
at most $Cm\mathbb E(|\zeta_0|^2+|(1-P_k^{\rm F})\kappa_1|^2)$.
All are bounded for $0\le m\le1$, proving the result.
\end{proof}

\subsection{The whole second order}

For a translation-covariant total $P$ at the trivial pair, let
$\|P\|_j$ be the infimum of the norms of its local
cube-covariant representatives.  This quotient norm is used
only for totals; a particular decomposition can have a larger
norm.  In particular the cancellations below are identities
of totals, not identities polymer by polymer.

The infinite-lattice slices and these norms are compatible
with the thermodynamic limit.  Here is the required argument.
In each nonzero Bloch fibre the closedness and vanishing
plaquette-average conditions are kernels of finite matrices
of trigonometric polynomials.  Their orthogonal projections
are bounded by one and real analytic outside the common zero
set of their maximal minors, a closed null set.  The gauge
representative constructed above gives a slice covariance
$\widehat C_A(\xi)$ bounded by
$C_{[i]}(1+|\xi|^{-2})$, with $i$ fixed; its fibres are the
same on every torus containing that momentum.  Its zero
fibre on a torus is bounded independently of the side.
For a coefficient of a quadratic gauge variation of a fermion
form, exponential locality of the flat kernels bounds its
bond kernel by $Ce^{-c(|b-y|+|b'-y|)}$.  Its mean is the
pairing of this kernel with the slice covariance.  On
$|\xi|\ge\delta$ the Riemann sums converge, while the sum
and integral over $0<|\xi|<\delta$ are at most $C\delta$;
the zero fibre contributes $O(N^{-3})$.  Dominated convergence
then gives the same mean on the infinite lattice.
Gauge-invariant quadratic forms read only the curl law:
two covariances with the same curl image differ by symmetric
terms with one gauge-direction factor, which the polarized
quadratic form annihilates.  These observations justify
interchanging the gauge and fermion integrations below.

If torus totals with representatives of norm at most $A$
converge coefficientwise, choose a diagonal subsequence of
their representatives on every fixed polymer.  Each lies
in a bounded finite-dimensional Grassmann space.  Fatou's
lemma gives a limiting representative of norm at most $A$;
the omitted wrapping polymers are exponentially small.
Thus $\|P\|_j\le A$ in the limit.  This proves the norm
passages used below without selecting an unbounded
representative of the quotient.

Let $\sigma_*>0$ be the relative coupling radius of the pure
one-step expansion in \cite[equation~(890) and Thm.~V]{YMM1},
and put $r_x=\tfrac14\min\{\tfrac12,\sigma_*\}$.
About a real admissible coupling sequence $x^\circ$, the full
coupling domain is the product of discs
$|x_i-x_i^\circ|<r_xx_i^\circ$; the inner domain has radii
$\tfrac23r_xx_i^\circ$. The sharp mass windows remain separate
hypotheses.

For real $x_i>0$, let $A_i$ be the centered gauge slice in bare
lattice units with covariance $L^{-i}x_i^{-1}$ times that of
$\gamma_{k,i}$, and let $\mathbb E_i$ denote its Gaussian mean.
The slices are independent, and the same slice is used for
every $k>i$. At complex $x_i$, $\mathbb E_i$ on the quadratic
polynomials below is its continuation linear in $x_i^{-1}$.
Let $\delta_A$ and $\delta_A^2$ denote the first and second
homogeneous Taylor terms, so $\delta_A^2$ includes the factor
$1/2$.  At the trivial configuration put
\[
 \mathscr T_k=\sum_{i<k}
 \left\{-\mathbb E_i\delta_{A_i}^2
            \langle\bar\Psi,\mathbb D_k\Psi\rangle
 +\tfrac12\mathbb E_i
       (\delta_{A_i}\langle\bar\Psi,\mathbb D_k\Psi\rangle)^2
 \right\}.
\]
It has degrees two and four.  Its last slice is
$\mathscr A_k$.  Write $\mathcal J_k$ for Gaussian fermion
integration under
$\Psi_k=\mathsf H_k\Psi_{k+1}+W$ and its row twin.

\begin{lemma}[Transport of the second order]\label{C-398-20}
As polynomials of degrees two and four,
\[
 \mathscr T_{k+1}=(\mathcal J_k\mathscr T_k)_{\ge2}
                                      +\mathscr A_{k+1}.
\]
\end{lemma}

\begin{proof}
At fixed block average the Schur formula gives
\[
 \mathbb D_{k+1}[\mathbb D_k+X]
 =\mathbb D_{k+1}+\bar{\mathsf H}_kX\mathsf H_k
       -\bar{\mathsf H}_kX\Gamma_k^{-1}X\mathsf H_k+O(X^3).
\]
For an earlier slice $i<k$, its level-$k$ block field has no
linear variation because $A_i\in T_{i+1}\subset T_k$.
Its quadratic variation takes values in $\mathfrak{su}(2)$;
its mean is an invariant colour vector and hence zero.
Thus $q$ may be held fixed inside these means.  The displayed
formula gives the transported quadratic Taylor term and the
additional term
$\mathbb E_i\bar{\mathsf H}_k(\delta_A\mathbb D_k)
\Gamma_k^{-1}(\delta_A\mathbb D_k)\mathsf H_k$.

On the other side Wick's rule applied to
$\tfrac12\langle\bar\Psi,N\Psi\rangle^2$ gives its
uncontracted quartic term and exactly this bilinear term.
There are two cross contractions, each with coefficient
$1/2$.  The self-contraction is
$-\operatorname{Tr}(N\Gamma_k^{-1})
\langle\bar\Psi',\bar{\mathsf H}_kN\mathsf H_k\Psi'\rangle$;
it vanishes since the linear vertex has traceless colour
matrices and the flat covariance is the identity in colour.
The remaining fully contracted term is constant and is
discarded.  This proves the formula slice by slice.
In finite volume Fubini is algebraic; the exponential fermion
kernels and the integrable slice covariance prove absolute
convergence at each infinite-lattice coefficient.  The new
slice is $\mathscr A_{k+1}$ by definition.
\end{proof}

Define recursively, with all old terms read at the current
level,
\[
 \widetilde T^{(j)}=\mathscr A_j+
 \left(\mathcal J_{j-1}\sum_{l<j}T^{(l)}
       -\sum_{l<j}T^{(l)}\circ(\mathsf H_{j-1},
                                     \bar{\mathsf H}_{j-1})\right)_{\ge2},
\]
\[
 \mu_j^{(2)}=-\mathsf{cf}_m(\widetilde T^{(j)})/\eta_j^2,
 \quad T^{(j)}=\widetilde T^{(j)}+
                 L^{-j}\mu_j^{(2)}\mathsf S_1(\psi^{\rm bg}),
 \quad G_k^{(2)}=\sum_{j\le k}L^{k-j}\mu_j^{(2)}.
\]
These expressions include the contractions of earlier
four-fermion interactions; retaining only the last slice
would give a different mass.

\begin{proposition}[The sum rule]\label{C-398-24}
Each $T^{(j)}$ has quadratic zeroth moment zero at every
later level and
$\|T^{(j)}\|_j\le C_Th_F^4/\mathsf t_{j-1}$.  Moreover
\[
 \sum_{j\le k}T^{(j)}-L^{-k}G_k^{(2)}
                       \mathsf S_1(\psi^{\rm bg})
                   =(\mathscr T_k)_{\ge2},
 \qquad G_k^{(2)}=G_k(x),\quad
                  \mu_k^{(2)}=\mu_k^{\rm1L}(x).
\]
\end{proposition}

\begin{proof}
All terms have the symmetries used in
Lem.~\ref{C-388-21}.  The Gaussian curvature slice is
invariant under the isometries preserving its level, since
its orthogonal projection is defined by centered plaquette
averages.  The fermion tower and its averages have the same
covariance.  Thus the zeroth moments are scalar.  Their
subtraction gives zero at level $j$.  The flat map
$\mathsf H_j$ carries constants to constants, multiplying
by
$\theta L^{-2}/[\mathfrak d_j+\theta/L]>0$, where
$\mathfrak d_j=\theta_j\mathfrak m_j/(\theta_j+\mathfrak m_j)$.
It therefore preserves a vanishing zeroth moment at every
later level.

The identity holds at $k=1$.  If it holds at $k$, apply
Lem.~\ref{C-398-20}.  Gaussian integration carries the
bare mass form to the next bare background mass form plus
a constant.  In particular its coefficient becomes
$L^{-k}=L\,L^{-(k+1)}$; no new normalization occurs.
The recursive definition of $\widetilde T^{(k+1)}$ then
gives the identity at $k+1$ with
$G_{k+1}^{(2)}=LG_k^{(2)}+\mu_{k+1}^{(2)}$.

For the norm, Wick's rule for the new gauge slice and
Cauchy's estimate on the vertex give
$\|\mathscr A_j\|_j\le4C'C_V^2h_F^4/\mathsf t_{j-1}$.
Contractions of earlier quartics are the coefficient linear
in the old-activity origin of the joined expansion.  Its
marked count and the leg-room estimate give
\[
 K_Mc_\psi\sum_{l<j}C_7L^{(j-l)/4}\|T^{(l)}\|_l.
\]
This is a local family bound, obtained with the decoupled
covariances; using the whole unlocalized covariance to
integrate each polymer would not give such a bound.
Put
$S_2^\sharp=\sum_{n\ge1}L^{-3n/4}(1+n\log L)^{9\mathfrak p_0}$
and require
$K_\mu K_Mc_\psi C_7S_2^\sharp\le1/4$.
Since $\mathsf t_{l-1}=L^{j-l}\mathsf t_{j-1}$, extraction and induction
close with $C_T=8K_\mu C'C_V^2$.
The torus-to-infinite norm argument above supplies a
representative with the same bound.

Finally take quadratic zeroth moments in the sum rule.
The left side is $-\eta_k^2G_k^{(2)}$.  In the right side
the quartics have no quadratic moment.  The mean Hermitian
quadratic form on constants is, by the three-square identity,
$\eta_k^2\sum_{i<k}L^{k-i}x_i^{-1}
\sigma_i(k;\mathfrak m_k)=\eta_k^2G_k(x)$, with the minus
sign in $\mathscr T_k$. The factor $L^{k-i}/x_i$ combines the
slice covariance multiplier $L^{-i}/x_i$ with the factor $L^k$
in $\mathbb D_k$. The scalar moment equals this
Hermitian moment at real coefficients; linearity in
$1/x_i$ extends it to the complex coupling discs.
Hence $G_k^{(2)}=G_k(x)$ and taking successive increments
proves the mass identity.
\end{proof}

\subsection{The third-order remainder}

The marked scalar expansion also gives a local comparison
of tested and untested gauge second moments at real
coefficients.  For bonds $b,b'$ it writes their difference
as $\sum_Xc_X(b,b')$, $X$ containing both bonds, with
\[
 |c_X(b,b')|\le K_{\rm cp}x_k^{-1}
 e^{-\gamma_{\rm tail}\mathfrak b(\mathsf t_k)^2/2}
                                  e^{-\kappa_+d(X)}.
\]
To prove this, put $a=\mathsf t_k^{1/2}B$ and insert two linear sources
$\sigma_1\mathsf t_k^{-1/2}a^a(b)+
\sigma_2\mathsf t_k^{-1/2}a^{a'}(b')$ in the $q$ origin.
These are $\sigma_1B^a(b)+\sigma_2B^{a'}(b')$ in unscaled units.
The inequality
$|a(b)|\le\tfrac12M^{3/2}(1+M^{-3}|a(b)|^2)$
puts this origin below its allowed norm on the bidisc
$|\sigma_i|\le\rho_0=c_{\rm cl}\mathsf t_k^{1/2}M^{-3/2}$,
after assigning half its allowance to the original $q$
potential.  The tested logarithm minus the untested
Gaussian logarithm is exactly the part marked by test
defects, of norm at most $K_Pp_*e^{-\kappa_+d(X)}$ on
that bidisc.  Its $\sigma_1\sigma_2$ coefficient is the
second-moment difference because both first moments
vanish by evenness.  Cauchy's estimate costs
$\rho_0^{-2}\le CM^3c_{\rm cl}^{-2}x_k^{-1}$, giving
the stated $K_{\rm cp}$.  Locality makes this coefficient
zero unless $X$ contains both bonds.  This proof uses
the admitted marked scalar count and quadratic-growth
majorant; a scalar probability tail alone would not
give the local comparison.

Put $\phi(t)=t^{-1/2}(1+\log t)^{\mathfrak p_0}$ and
retain the sharper mass window
$|s_i|\le2\bar B_M/\mathsf t_i$ throughout the following step.
Deform its potential by
\[
 \mathcal V(z)=z_3\mathcal V^0(B)+
 \mathcal V_k^{\rm F}(z_1B)+z_2\mathcal V^{\rm old}(z_1B)
                                  +z_4\mathcal V^{\rm s}(z_1B).
\]
Here the fluctuation argument is in unscaled units.
The weight of $n\in\mathbb N_0^4$ is
$|n|_w=n_1+2n_2+n_3+2n_4$.
With $c_\diamond=c_{\rm cl}/4$ choose the radii
\[
 r_1=\min\{c_\diamond/(C_Vh_F^2\delta_k),
                                  \beta_0/(2\delta_k)\},
 \quad r_2=\frac{c_\diamond}
 {c'_\psi\sum_{j\le k}C_7L^{(k+1-j)/4}\varepsilon_j},
\]
\[
 r_3=c_\diamond/(K_0^\sharp\delta_k),\qquad
 r_4=c_\diamond/(c'_\psi c_{\mathsf S}h_F^2|s_k|),
\]
with infinite radius when its denominator is zero.
The number $c'_\psi$ bounds the displayed leg and shift
factors on the enlarged fluctuation disc; it is fixed before
the small coupling ceiling.  The scalar potential norm in
$r_3$ admits its quadratic Gaussian growth.  Replacing it
by the supremum over the fluctuation domain would lose
two powers of the threshold logarithm.

\begin{proposition}[Second-order activities and cubic mass error]
\label{C-398-28}
There are constants $C_2,C_3,C'_\rho$ such that the
small-field step preserves
\[
 \|E_{\ge2}^{(j)}\|_j\le C_2\phi(\mathsf t_{j-1})^2,
 \qquad E_{\ge2}^{(j)}=T^{(j)}+R^{(j)},\quad
 \|R^{(j)}\|_j\le C_3\phi(\mathsf t_{j-1})^3
\]
at the trivial pair, the latter as a norm of totals.
Its extracted mass satisfies
\[
 \mu_{k+1}=\mu_{k+1}^{\rm1L}(x)+\rho_{k+1},
 \qquad |\rho_{k+1}|\le C'_\rho\phi(\mathsf t_k)^3.
\]
The estimates hold at real coupling sequences in their
window and at complex sequences in the fixed inner
coupling discs, under the same sharp mass windows.
\end{proposition}

\begin{proof}
On $|z_i|\le r_i$, the four potential pieces have sizes
at most $c_\diamond$ each.  Their local logarithm is
holomorphic with norm at most $2K_Mc_{\rm cl}$, including
the test tail.  If $\widetilde E(z)=\sum_na_nz^n$,
Cauchy's estimate gives
\[
 \|a_n\|_{k+1}\le2K_Mc_{\rm cl}\prod_ir_i^{-n_i}.
\]
Here $r_1^{-1},r_3^{-1}\le C_r\phi(\mathsf t_k)$,
$r_2^{-1}\le C_rC_2\phi(\mathsf t_k)^2$ and
$r_4^{-1}\le C_r\phi(\mathsf t_k)^2$.  The last estimate uses
the sharp mass window.  Choose the coupling ceiling so
that every finite $r_i\ge2$.

There is no fermion coefficient of weight zero or one.
The coefficient linear in $z_1$ vanishes: the first Taylor
coefficient of $\mathcal V_k^{\rm F}(z_1B)$ is linear in $B$,
and the Gaussian, tests and $q$ weight are even.  Terms depending only on $z_3$ are
scalar.  At $z_1=0$ the gauge and fermion integrals
factorize, so mixed totals with a positive $z_3$ power
and an old-activity or mass power vanish.  The coefficient
linear in $z_4$ is scalar: its single-$W$ terms integrate
to zero and its $\bar WW$ term is constant.
The coefficient linear in $z_2$, as a total, is exactly
\[
 \mathcal J_k\sum_{j\le k}E_{\ge2}^{(j)}
       -\sum_{j\le k}E_{\ge2}^{(j)}
                         \circ(\mathsf H_k,\bar{\mathsf H}_k).
\]
The marked expansion of this origin, with Cauchy's estimate
in $z_2$, gives the local bound
\[
 K_Mc_\psi\sum_{j\le k}C_7L^{(k+1-j)/4}\varepsilon_j.
\]
Its normalization cancels the gauge integral, so no test
tail is present in this particular coefficient.

We justify this identity also on the infinite lattice,
where the old family is defined.  For a single old term
$v(Y)$ let $L_{X,Y}(v)$ be the coefficient linear in its
marked origin in the localized logarithm.  On a torus,
linearity of the cluster coefficients and the exact
Mayer resummation give
\[
 \sum_XL^T_{X,Y}(v)=\mu_{\Gamma_T^{-1}}v(Y).
\]
The term $v$ reads no gauge fluctuation, so its scalar
normalization cancels.  A summand with output diameter
at least $R$ has exponentially small norm, uniformly
in the torus.  Let its side increase.  Nonwrapping
local terms agree with their infinite-lattice terms,
while the tail tends to zero.  The finitely many
Gaussian moments in the right side converge with
their covariance entries.  This proves
$\sum_XL_{X,Y}(v)=\mu_{\Gamma^{-1}}v(Y)$ on the
infinite lattice.  To sum over $Y$, fix one external
coefficient.  The exponential bound on the output
support makes the sum over $X,Y$ absolutely convergent;
the sums may therefore be exchanged.  Terms without
a difference of external legs cancel through the
zero cell sums from \cref{C-388-24}.  This proves the
identity for all the finite-norm old families used here,
including their cell-unit representation at sources.

The displayed identity is an identity of totals.
The expression $\mu_{\Gamma^{-1}}v(Y)$ itself reads
the gauge configuration through the whole covariance
and is not a local family.  The local representative
is $\sum_YL_{X,Y}(v)$, formed with the decoupled
covariances.  The remainder estimate below uses the
quotient norm of totals at exactly this point.

At the trivial pair the coefficient $z_1z_3$ also vanishes:
the expectation of one adjoint colour vector multiplied
by a colour-invariant scalar is an invariant vector and
is zero.  This follows from colour invariance.  By the exact identity in
Lem.~\ref{C-388-14}, the $z_1^2$ coefficient is
$\mathscr A_{k+1}$, apart from the tested-Gaussian
comparison.  The linear determinant vertex has zero
colour trace, and the mean of a second-order
$\mathfrak{su}(2)$ field inserted through a first
derivative is an invariant colour vector; neither adds
a term.

To justify the untested measure used here when $x_k\ne \mathsf t_k$,
write its density as $e^{-x_kH-\mathsf t_kG}$ and the comparison
Gaussian as $e^{-x_kH-x_kG}$, with $H$ the curvature
energy and $G$ the gauge-fixing energy.  At the unit
coefficient the known curvature covariance is the
orthogonal curvature projection.  Gaussian elimination
in the gauge directions then shows $G_{\min}(dA)=0$:
otherwise the induced curvature form would contain this
nonzero extra nonnegative term.  Thus both Gaussians
have the same curl law.  Their quadratic fermion
observable is invariant under the residual linear gauge
directions, so the covariance argument above makes
their means equal.  Only the tests change the result.
The local second-moment comparison bounds this
difference by a total of norm
$C_\chi \mathsf t_k^{-1}e^{-\gamma_{\rm tail}\mathfrak b(\mathsf t_k)^2/2}$.
Indeed a polymer has at most $8(1+d(X))$ cubes, so it
holds at most $C(LM)^6(1+d(X))^2$ ordered bond pairs.
The exponential output rate pays these squares and
the leg distances; summing through a cube uses
$\sum_{n\ge0}(n+1)^22^{-n}=12$.
The thermodynamic norm passage above preserves the bound.
Equivalently one obtains the comparison directly by
using the unrestricted origin
$\mathcal V^{\rm q}+\zeta\mathcal V_1+
\zeta^2\mathcal V_2$, where $\mathcal V_1,\mathcal V_2$
are the first two fermion vertices.  In the difference
of its tested and untested logarithms every term has
a failed-bond mark.  Cauchy's estimate in $\zeta$
selects its second coefficient and gives the same
bound in the quotient norm of totals.

For complex $x_k$, regard this difference as a holomorphic
function with values in the quotient norm of totals.
On the full coupling disc it is bounded by
$C\phi(\mathsf t_k)^2$, and on its real diameter by the last
exponentially small bound.  Apply the scalar two-constants
theorem through bounded linear functionals on that
quotient. In the coordinate $(x_k-x_k^\circ)/(r_xx_k^\circ)$,
the inner disc has radius $2/3$; there its norm is
at most
\[
 C_\chi\phi(\mathsf t_k)^2
     e^{-c_{\rm hm}\gamma_{\rm tail}\mathfrak b(\mathsf t_k)^2/2},
\]
where $c_{\rm hm}>0$ is the minimum harmonic measure of
the real diameter on that compact inner disc.  A final
coupling ceiling makes this at most
$C_\chi\phi(\mathsf t_k)^3$.

The Cauchy radii now bound all remaining coefficients
of weight at least three by
\[
 \sum_{|n|_w\ge3}\|a_n\|_{k+1}
                         \le(C_0+C_1C_2)\phi(\mathsf t_k)^3.
\]
This follows by summing the four geometric series:
each weight-one inverse radius costs $C_r\phi$ and
each weight-two inverse radius costs $C_r\phi^2$;
the exceptional coefficient linear in $z_2$ is kept
separately.  The ceiling $C_r^2C_2\phi\le1$ bounds
terms with two or more old-origin factors uniformly
by a constant plus a constant times $C_2$.

For the full family of fermion terms, before the
trivial-pair split, weight two contributes at most
$C_{(2)}\phi^2$.  After mass extraction its bound is
\[
 K_\mu[C_{(2)}+(C_0+C_1C_2)\phi]\phi^2+
 K_\mu K_Mc_\psi\sum_{j\le k}C_7L^{(k+1-j)/4}
                                    \varepsilon_j.
\]
The geometric calibration relation bounds the last
sum by $C_7S_2^\sharp C_2\phi(\mathsf t_k)^2$; the defining
logarithmic factor in $S_2^\sharp$ pays the ratio of the
calibration logarithms.
Require $K_\mu K_Mc_\psi C_7S_2^\sharp\le1/4$,
take $C_2=4K_\mu(C_{(2)}+C_0)$, and then choose
$\gamma$ so $K_\mu C_1\phi(\gamma^{-2})\le1/4$.
This closes the second-order bound.

At the trivial pair subtract the exact $T$ sum rule.
The remaining pre-extraction total is the tested
$z_1^2$ error, the coefficient linear in the old
remainders, and the coefficients of weight at least
three.  The same calculation, using their cubic
induction bounds, closes with
\[
 C_3=4K_\mu(C_\chi+C_0+C_1C_2).
\]
Extract its zeroth moment.  Since $\eta_{k+1}\ge c_b$,
Lem.~\ref{C-388-21} gives
$|\rho_{k+1}|\le
c_b^{-2}C^\psi_{am}h_F^{-2}C_3\phi(\mathsf t_k)^3$;
discarding the additional factor $\nu^{-1}\le1$
only weakens the bound.  This proves the proposition.
The squared mass shift has weight four precisely
because $|s_k|=O(\mathsf t_k^{-1})$; the argument does not
assert a cubic error in a wider $O(\phi(\mathsf t_k))$ window.
\end{proof}

\begin{proposition}[Closing the forward mass flow]\label{C-396-33}
Put $\overline B=LB_\sharp/(L-1)$ and let
$x_0=L^K\mathsf b$, $\mathsf b\ge\gamma^{-2}+\overline B$.
Suppose the small-field steps give
\[
 x_{k+1}=L^{-1}x_k-b_{k+1},\quad |b_{k+1}|\le B_\sharp,
 \qquad s_{k+1}=Ls_k+\mu_{k+1},\quad s_0=0,
\]
and, whenever the preceding mass windows hold,
\[
 \mu_j=\mu_j^{\mathrm{1L}}(x)+\rho_j,\qquad
 |\rho_j|\le C_\rho\phi_\sharp(x_{j-1})^3,
 \quad \phi_\sharp(t)=t^{-1/2}(1+\log t)^{3\mathfrak p_0},
 \qquad 0\le F(k;m)\le B_F.
\]
For the final-level assertion suppose also the calibration
bound already given by \cref{C-398-28},
$|\rho_j|\le C'_\rho\phi_\sharp(\mathsf t_{j-1})^3$, with
$\mathsf t_i=L^{K-i}\mathsf b$. For a sufficiently small $\gamma$, chosen
after these constants,
\[
 |s_k|\le\frac{3B_F+1}{x_k}\quad(0\le k\le K).
\]
At the final level,
\[
 \left|s_K-\frac{F(K;\mathsf m)}{\mathsf b}\right|
 \le C_s\mathsf b^{-3/2}(1+\log\mathsf b)^{9\mathfrak p_0}.
\]
\end{proposition}

\begin{proof}
The coupling recursion gives
$|x_i-L^{K-i}\mathsf b|\le\overline B=LB_\sharp/(L-1)$.
Choose $\gamma$ so that all $x_i$ exceed $4\overline B$.
Then, for $i\le k$,
$\frac12L^{k-i}x_k\le x_i\le2L^{k-i}x_k$.
The one-loop increments telescope exactly:
\[
 s_k=\sum_{j\le k}L^{k-j}\mu_j
       =G_k(x)+\sum_{j\le k}L^{k-j}\rho_j,
 \qquad |G_k(x)|\le\frac{2B_F}{x_k}.
\]
The last inequality uses nonnegativity of the individual
$\sigma_i$ and the bound on their sum.
Assume the asserted mass window at levels below $k$.
The remainder estimate is then applicable.  A geometric sum gives
\[
 \sum_{j\le k}L^{k-j}\phi_\sharp(x_{j-1})^3
 \le S_3x_k^{-3/2}(1+\log x_k)^{9\mathfrak p_0},
\]
where
\[
 S_3=\sup_{x\ge2}\sum_{n\ge0}
 L^n(\tfrac12L^{n+1})^{-3/2}
 \left(\frac{1+\log(2L^{n+1}x)}{1+\log x}\right)^{9\mathfrak p_0}
 <\infty.
\]
Choose $\gamma^{-2}\ge e^{18\mathfrak p_0}$ and
$C_\rho S_3\gamma(1+2\log\gamma^{-1})^{9\mathfrak p_0}\le\frac14$.
The sum of the remainders is at most $1/(4x_k)$.
Thus $|s_k|\le(2B_F+\frac14)/x_k\le(3B_F+1)/x_k$,
closing the induction.

At $k=K$,
\[
 \left|G_K(x)-\frac{F(K;\mathsf m)}{\mathsf b}\right|
 \le 2\overline B B_F\mathsf b^{-2},
\]
because
$|L^{K-i}/x_i-\mathsf b^{-1}|\le2\overline B\mathsf b^{-2}$.
Along the calibrations $\mathsf t_{j-1}=L^{K-j+1}\mathsf b$ the
third-order remainders sum to at most
$C'_\rho S_3^\sharp\mathsf b^{-3/2}
(1+\log\mathsf b)^{9\mathfrak p_0}$, with
\[
 S_3^\sharp=\sum_{n\ge0}L^{-(n+3)/2}
                    (1+(n+1)\log L)^{9\mathfrak p_0}<\infty.
\]
Adding the two bounds gives the result with
$C_s=2\overline BB_F+C'_\rho S_3^\sharp$.
\end{proof}

\subsection{Scalar extraction, pointing, and the order of constants}

A scalar family used here is strictly gauge invariant term
by term.  Its first derivative at the trivial configuration
is zero because it would be an invariant adjoint colour
vector.  Differentiating gauge invariance once more shows
that its Hessian annihilates exact one-forms in either slot.
For a term on $X$, choose a rooted radial gauge in its
bounding box.  Each gauge potential is then a sum of curls
over a surface of area at most $C\operatorname{diam}(X)^2$.
Thus its Hessian is $T_X[dA,dA]$, with norm at most a
polynomial in $\operatorname{diam}(X)$ times the norm of
the term, by Cauchy's estimate.  The polymer decay pays
this polynomial.  The total has a bounded continuous
curvature symbol, so its potential symbol begins at
order two.  Colour and cubic invariance make that
quadratic part a multiple of the Wilson-action Hessian.
This argument excludes a first-order Chern--Simons
density because the individual terms are strictly
invariant, rather than invariant up to a difference.
It is the scalar classification needed for
taking out the gauge-action coefficient; parity is
not used for it.

The scalar local extraction has the admitted coefficient
bound $|\mathsf{cf}(F)|\le C_{am}r_\alpha^{-2}\|F\|$ on
the usable pair radius.  Removing its constant and its
multiple of the action therefore costs
$2+C_{am}r_\alpha^{-2}$.  Applying this to the joined
expansion gives
\[
 \|F_0^{(k+1)}\|_{k+1}\le
       \tfrac12C_w\phi(\mathsf t_k)+\theta_sh_F^2|s_k|,
 \qquad |b_{k+1}^{\rm r}|\le
       C_b\phi(\mathsf t_k)+\theta_sh_F^2|s_k|.
\]
Indeed the input sizes are
$(K_0+C_Vh_F^2)\delta_k$ for the gauge and Dirac vertices,
$c'_\psi c_{\mathsf S}h_F^2|s_k|$ for the mass variation,
$3\sum_{j\le k}\omega_{k+1-j}\varepsilon_j$ for old
fermion terms, and the test tail.  Their last sum is at
most $3C_7S_\omega C_2\phi(\mathsf t_k)^2$, where
$S_\omega=\sum_{n\ge1}L^{-n/4}(1+n\log L)^{3\mathfrak p_0}$.
Multiplying by $K_M$ and the extraction factor proves
these bounds after $C_w,C_b,\theta_s$ are chosen.
On a constant smaller radius the coefficient bound
costs $(\alpha_0/a)^2$ and the action costs
$(a/\alpha_0)^2$; they cancel.  Hence the same extraction
factor works on all the stated radius systems.
The normalizer part is bounded before the polymer
constants by Prop.~\ref{C-397-43}; the admitted
pure normalizer increment is bounded by
$(1+L^{-1})c_{2,\mathrm{loop}}$.
Choose
$B_\sharp\ge(1+L^{-1})(c_{2,\mathrm{loop}}+
N_fc_2^{\rm F})+1$.
The remaining increment is at most one once
$\gamma$ is small, by the displayed scalar
extraction estimate and the sharp mass window.
This supplies $|b_{k+1}|\le B_\sharp$ in
Prop.~\ref{C-396-33}, subject to the
admitted pure step inputs throughout.

Pointing the new fermion terms can be done without
dividing them equally among arbitrary polymer cubes.
For a homogeneous degree-$n$ term, $n\ge2$, let
$N_p$ count its legs at place $p$ and put
$F_p=n^{-1}N_pF$.  Then $\sum_pF_p=F$ and the norm
of the sum over any subset of places is at most
$\|F\|$.  Along the potential path
$\mathcal V_t=\mathcal V_0+t\mathcal V^{\rm F}$ define
\[
 E^\bullet(X,p)=\int_0^1
       E'[\mathcal V_t;(\mathcal V^{\rm F})_p](X)\,dt.
\]
Summing the pointers gives
$E[\mathcal V_0+\mathcal V^{\rm F}](X)-
E[\mathcal V_0](X)$.  The marked cluster majorant
bounds this derivative, so the preceding estimates
hold with the pointer summed; send $p$ to its block
site of the new level.  Scalar remainders use the
gauge bond pointing, and closed fermion
normalizers use their diagonal trace pointer.  The
field-independent wrapping loop terms use the
equal thirds on the three wrapping rows.
Quadratic and cubic coefficient extraction commutes
with this linear pointing operation.

The scalar pointing must retain the same quadratic
normalization as \cref{C-388-37}.  Let
$\mathbb V^{\rm int}=\mathbb V^{\rm J}+\mathbb V^{\rm G}
 +\tau\mathbb V^{\rm A}$ be the gauge vertices and
$\mathbb V'$ the old scalar variation.  Their complex
polydisc sizes satisfy
\[
 n^{\rm int}\le\tfrac72K_P\delta_k\le\tfrac7{16}w_*,
 \qquad n'\le K_\Lambda\mathsf s_aN_\star\le\tfrac14w_*.
\]
Write $\widehat g_b$ for the bond pointing of the
quadratic gauge-fixing family.  Point its tested and
untested logarithms separately, integrating the
insertion $-(1-\tau)\widehat g_b$ along the path
$\sigma\mathcal V^{\rm q}$, $0\le\sigma\le1$.
Call the two terms $G^{\rm v}(X,b)$ and $G_\tau(X,b)$.
Let $E^\chi[\emptyset](X,b)$ be the test-only pointing.
Finally integrate the marked derivative along the
gauge-vertex and old-variation path to form
$E^{\bullet,\rm v}(X,b)$, keeping the tested quadratic
origin fixed on this path.  Their sum is
\[
 E^\chi[\emptyset](X,b)+G^{\rm v}(X,b)-G_\tau(X,b)
                                      +E^{\bullet,\rm v}(X,b).
\]
Summing over $b$ gives exactly the degree-zero local
logarithm, with $G_\tau$ subtracted.  The four path
integrals give the tests alone, the change to the
tested quadratic system, the untested quadratic
divisor, and the remaining potential; their
intermediate endpoints cancel.

The sort-change proof in \cref{C-388-37}, differentiated
by adding an insertion of size $K_{\mathsf G}$ on a
disc of radius $w_*/K_{\mathsf G}$, gives
\[
 |G^{\rm v}-G_\tau|(X,b)
 \le3|1-\tau|K_{\mathsf G}w_*^{-1}K_P^\Pi p_*
                         e^{-\kappa_+d(X)}.
\]
The other path derivative has size at most
$9(n^{\rm int}+n')$, so its pointed cluster term is
bounded by $9K_\Pi(n^{\rm int}+n')$ times the same
exponential.  Send bonds to sites by the tent weights
of the chart.  Their sum at a site costs at most
$3L^3$.  The resulting pointing constant is
\[
 K_{\mathbf M}^{\rm v}=3L^3\max\left\{
 K_\Pi K_{\mathsf G},
 (1+\tfrac32\mathsf C_{\rm ex}((L+2)M)^3)K_P^\Pi e^{4\lambda_1},
 32K_\Pi K_P,\ 81K_\Pi K_\Lambda\right\}.
\]
It replaces $K_{\mathbf M}$ in the scalar extraction
bounds.  The Lipschitz estimate in the fourth family
splits its change into a change of insertion and a
change of system.  These have constants $9K_\Pi$
and at most $72K_\Pi$, respectively, when
$n^{\rm int}+n''\le w_*$ and
$\|\mathbb V'-\mathbb V''\|\le w_*/2$.
Restrictions and joint gauge transformations commute
with every operation, giving local pointed terms
on the whole pair cores.

The array supplied to the next step has a
degree-zero bound in a supremum over polymers,
pointers and configurations, rather than the summed
fermion-family norm.  At each new level its explicit
families give
\begin{equation}\label{eq:qcd-degree-zero-array}
 \begin{split}
 &\sup_{X,z}e^{(\kappa+1)d_{k+1}(X)}
      \sup_P|E^{(k+1)}_0(X;P;z)|\\
 &\quad\le E_{\mathbf N}+C_{\mathsf n}
 +C(C_N+M^3)N_f\\
 &\qquad+C(1+L^{-1})(c_{2,\mathrm{loop}}+N_fc_2^{\rm F})
 +CK_{\mathbf G}\bar B\gamma^2
 +C_w\phi(\gamma^{-2})\le E_0 .
 \end{split}
\end{equation}
The first two terms are the gauge normalizer and its
retained one-site constant, the next two are the
closed fermion loops and the extracted action
coefficient, and the last two are the Gaussian
correction and the fluctuation family.
For a wrapping row of $N/M$ cubes the free determinant
error is at most $C\nu NM^2e^{-\delta_FN}$.
The floor $\delta_FM\ge\kappa+2$ bounds it by
$C\nu M^3$ after multiplication by
$e^{(\kappa+1)N/M}$; its pointed shares obey the
same bound.  The extracted plaquette shares have
tree size zero.  The pointed fluctuation estimate
has rate $\kappa_+\ge\kappa+1$.
This proves \eqref{eq:qcd-degree-zero-array} anew
at each level.

On the complex coupling disc replace the correction
bound by $CK_{\mathbf G}(\bar B\gamma^2+r_x)$ and
the fluctuation bound by its complex mass-window
bound.  At a complex fine mass the closed-loop
action coefficient is bounded by the Cauchy estimate
$C C_{am}r_\alpha^{-2}C_NN_f$.
Field-independent wrapping numbers drop out of every
background variation and pass outside the linear
block map.  The common choice of $E_0$ covers these
versions as well.  A single spare unit of polymer
decay does not imply a summed family bound: only
the scalar background variation, of second order
in curvature, has the summed bound $K_{\rm old}$
used above.

The choices used in these proofs have the following
order.  First fix the chamber operator, $N_f$, $L$
and $b_F$.  Form the free and weighted energy constants,
then the complex and source radii, the local supremum
constants, $M_w$ and the walk ratios.
Choose the admitted pure geometry and polymer constants
subject to their conditions and the additional displayed
floors on $R$, $M$, $\bar\kappa$ and $\kappa$;
then form $A_\kappa$.
The constants $\Gamma_F$, $K_M$, $c_{\mathsf S}$,
$C^\psi_{am}$, $K_\mu$ and $C_7$ are then fixed before
$h_F$; homogeneity in $h_F$ cancels from each of them.
The cumulative bound $B_F$ also precedes $h_F$, and
sets $\bar B_M=3B_F+1$.
Raise $h_F$ to meet the leg room and
\[
 h_F\ge8K_\mu K_MA_\kappa\Gamma_F C_7S_2^\sharp.
\]
Next form $C_r,C_{(2)},C_0,C_1,C_\chi$, then $C_2$,
$C_T$, $C_3$, $C'_\rho$ and the forward-flow constant.
Finally lower $\gamma$ to meet the finitely many
smallness conditions, including
\[
 \log\gamma^{-2}\ge18\mathfrak p_0,
 \quad C_r^2C_2\phi(\gamma^{-2})\le1,
 \quad K_\mu C_1\phi(\gamma^{-2})\le\tfrac14,
\]
\[
 \sup_{t\ge\gamma^{-2}}
   e^{-c_{\rm hm}\gamma_{\rm tail}\mathfrak b(t)^2/4}/\phi(t)
 \le1,
 \quad C'_\rho S_3^\sharp\psi_\sharp(\gamma^{-2})\le\tfrac12.
\]
For fixed earlier admissible constants these last
ceilings hold for sufficiently small $\gamma$.
Prop.~\ref{C-396-12} verifies a simultaneous assignment
with the pure, history and source conditions.


\section{Histories and the common vacuum energy}\label{sec:qcd-histories}

Let $T_\ell=(\mathbb R/\ell\mathbb Z)^3$ be the physical torus.
We use the admitted pure gauge histories in the calibration frame
$\mathsf t_j=L^{K-j}\mathsf b$. Their tests, coordinates, stored
gauge factors and profiles $\varphi_j$ are unchanged. The running
coefficients satisfy $|x_j-\mathsf t_j|\le\bar B$ and
$|s_j|\le\bar B_M/x_j$.
The weight of the calibration frame is the constant
$\mathsf{bw}=\beta=\mathsf t_0$. The running coefficients
enter the exponent and the Gaussian covariance $x_j^{-1}C_j$;
they do not change the tests, the coordinates or the block kernels.
The histories and their free-coordinate measures are those of
\cite[Section~44 and Thm.~D]{YMM1}. Write $\mathcal S=\mathcal S(h)$
for the discrete label of $h$, distinct from its free evaluation
coordinates. The integration regions
satisfy $\Omega_j\supset\Lambda_j$, and $Z_j=\Lambda_j^c$.
Put $\Omega_0=T_\ell$,
$\mathcal L_j=\Omega_j\setminus\Omega_{j+1}$ for $j<k$,
and $\mathcal L_k=\Omega_k$.
For a region $E$, $\#_j(E)$ counts its level-$j$ sites, or
equivalently its level-$j$ unit cells. The places $\mathfrak P_h$
consist of the 48 chambers over each site of $\mathcal L_0$
and one coarse place per level-$j$ site of $\mathcal L_j$ for $j\ge1$.
The fermions on the layers with $j<k$ are stored; those on the
last region are current. The profiles $\varphi_j$ are the fixed
cutoffs of these histories, equal to one on the prescribed inner
region and zero on $Z_j$. Write
\[
 n_j=\ell^3L^{3(K-j)},\quad \nu=8N_f,\quad
 \Sigma_{\mathcal L}(h)=\sum_{j=0}^k\#_j(\mathcal L_j),\quad
 \Sigma_Z(h)=\sum_{j=1}^k\#_j(Z_j),\quad
 \Sigma'(h)=\#\mathfrak P_h.
\]
There are $48$ chamber places per fine site, and
$\Sigma_{\mathcal L}\le\Sigma'\le48\Sigma_{\mathcal L}$.
The pure geometry gives
$\Sigma_{\mathcal L}\le n_k+L^3\Sigma_Z$.
The layer separation is $R_\flat=8LM$ in the units of the lower
of two consecutive levels.

Write $\varepsilon_j^{\rm leg}$ for the small fermion-origin
size at step $j$, and retain $n_j$ for the number of sites.
The extension to a history has two tasks. The gauge expansion
must be applied to the actual coefficients and conditional
Gaussian of that history. In addition, a fermion term created
at an earlier level must remain small when its background
fields are read at the present level. We verify these tasks
separately. The verification in \cref{C-396-12} supplies the
gauge charts, conditional structures and coordinate identities
for both sequences. Bounds on the matter coefficients are
established below.

The exponent of a history has the form
\begin{equation}\label{eq:qcd-history-format}
 -A(\beta_h^F,U)-\langle\bar\Psi,\mathbb D_h\Psi\rangle
 -M_h+E_h^{\rm reg}+E_h^{\rm bd}-\mathcal E^F_{k,\mathcal S},
 \qquad
 \beta_h^F=\beta-\sum_{j\le k}L^jb_j\varphi_j.
\end{equation}
For a scalar plaquette weight $q$, the weighted action here is
$A(q,U)=\sum_pq(p)a_p(U)$, with the history profiles transferred
to plaquettes by the convention of the cited gauge construction.
In particular,
\[
 -A(\beta_h^F,U)=-\beta A(U)
                    +\sum_{j\le k}L^jb_jA(\varphi_j,U).
\]
The last sum is bounded by a constant times
$\Sigma_{\mathcal L}(h)$ in the scalar estimate below and
is included in its layer constant. This permits the history
integration estimate to use the fixed frame weight $\beta$.
Here $s_h=\sum_{j\le k}L^{-j}\mu_j\varphi_j$ is the mass
profile and $M_h=\mathsf S_1(s_h\psi^{\rm bg})$ is its
quadratic mass member. On $\Omega_j\setminus\Lambda_j$ its coefficient is
$L^{-(j-1)}s_{j-1}$: the entire output on that collar is retained
as a boundary term, with no local part extracted there.

\subsection{The conditional expansion with fermion contractions}

The conditional mean near a collar changes with the decoupling
parameters. Consequently the local expansion must contract the
fermions after each corresponding conditional gauge integral.
Using a fixed whole-volume mean in every term would lose the
component identities on which the expansion rests.

\begin{lemma}[Local contractions in the centered expansion]
\label{C-399-41}
Consider an admitted centered gauge frame at step $k+1$,
including the first-step frame, with all the class-size,
external-field and Gaussian bounds required by
\cite[Prop.~76.2 and Thm.~AF]{YMM1}.
Let its values lie in the finite even Grassmann algebra with
integrated and external generators, with integrated weights
at least $\Gamma_F$. Suppose the contraction maps
$\mathscr L(Z,s;P)$ have the following properties:
they fix external elements, are linear over them and have
norm at most one; their matrix entries are holomorphic in
the decoupling parameters on a neighborhood of the closed
unit polydisc and on the pair slices, and are Borel in the
real data; they read the pair alone, independently of the
integration field, external field and stored coordinates;
they factor over wall-components when the separating
parameters vanish; their restrictions agree on subdomains;
they intertwine gauge transformations and conjugation; and,
at parameters one on the torus, they restrict one whole
contraction $\mathscr L_*$. Then the centered expansion
followed by these contractions has the scalar frame's
majorants, with the fibre norm in place of absolute value.
In particular its cluster family satisfies
\[
 \|\widetilde E(X)\|
 \le b_{\rm cs}^{\mathfrak F}e^{-\kappa''d_{k+1}(X)}
 \le\tfrac12e^{-\kappa''d_{k+1}(X)},\qquad
 \kappa''=\kappa_++1+c_E.
\]
The part passing through a specified class of ordinary
origins of total size $\sigma$ is bounded by
$K^{\rm lb}\sigma e^{-\kappa''d_{k+1}(X)}$, where
\[
 K^{\rm lb}
 =eC_{\rm cs}(LM)^3
  \exp\{C_{\rm cs}(\bar c_P^++\bar r_H^++\kappa'')\}.
\]
Take the larger of the first-step and positive-level constants
in this formula. It is fixed before the fibre weight.
The conclusions hold on the admitted external-field box
with complex room $r_{\rm bd}=C_{\rm xt}\mathfrak b(\mathsf t_k)$.
They include locality on $X$, joint pair and field holomorphy,
Borel dependence, gauge covariance, conjugation, component
factorization, and the exponential identity for the whole
integral. Each centered member uses the conditional mean
of its own decoupled structure. Inside $\Lambda_{k+1}$
this mean and the external field vanish.
\end{lemma}
\begin{proof}
For a Mayer index with region $Z$, form the scalar conditional
integral coefficient by coefficient, with its conditional mean
$\mathbf m(Z^{\rm I},s^{\rm en})$, and then apply
$\mathscr L(Z,s;P)$. Denote the resulting function by $f_\iota(s)$.
Its finite parameter brackets are the terms of the joined
expansion. Use all big cubes as guards, so the parameter set
outside $Z^\sharp$ is the full set of cubes of the ambient
domain. At parameters supported in a region, both the
conditional structure and $\mathscr L$ split over its
wall-components. The bracket identities and the cancellation
of disconnected terms therefore hold in the even algebra,
which is commutative.

To estimate a term, apply a linear functional of norm one
on the external algebra. The composite of this functional
with $\mathscr L$ has norm at most one. The scalar contour
estimate of \cref{lem:qcd-centered-interface}, at the
parameter-dependent center, bounds it by the same majorant.
Taking the supremum over the functionals gives the fibre
norm bound. Summation over the rooted indices and then over
connected clusters is exactly the scalar majorant summation.
The version with a marked origin gives its factor
$e\max\{K_{\rm or},K_{\rm or}^+\}\le K^{\rm lb}$; this proves
the selective bound, rather than a bound by the size of
all scalar origins. Projections onto a fixed degree have
norm at most one and obey the same estimate.

All finite brackets use values of the members and
contractions on their own domains. Their restriction,
holomorphy and symmetry properties therefore pass to the
terms, and uniform absolute convergence passes them to
the sums. At parameters one the conditional structure is
the whole structure and its center is the whole conditional
mean. Finite Grassmann expansion commutes with the
coefficientwise gauge integral; the exponential identity
follows by the same finite-volume resummation. On an
interior domain there are no external bonds, and every
cube meets the integration region. Its indices are then
the one-lattice indices, with center zero. Equality with
a second frame requires equality of its actual members
and contractions; locality alone does not assert that equality.
\end{proof}

For our contractions the covariance is the symmetric weakened
covariance of \cref{C-397-40}, restricted to the integrated
fermion places. Parameters outside $\Omega_{k+1}$ are retained.
The Gram bound gives the contraction norm, and
\cref{C-397-47} gives the intrinsic domain and restriction
properties. The symmetry in \eqref{eq:qcd-symmetric-covariance}
is needed here: Wick's rule intertwines Grassmann conjugation
with $S\mapsto\gamma_4S^\dagger\gamma_4$. The unmodified
one-sided parametrix would not by itself give that identity
at weakened parameters.

\begin{lemma}[Moving an external field to its stored value]
\label{C-399-42}
On the admitted pair and external-field boxes, let
$F_X(a;P;q)$ be a local family with values in the finite
external even algebra. Assume that it reads the field inside
$X$ alone, is jointly holomorphic in the pair and field
variables, has its stated external legs in $X$, and satisfies
$\|F_X\|\le b e^{-r d(X)}$ there. The coordinate change from
the field before the strip to the stored field produces local
rows whose sum is
\[
 \sum_X F_X(a_{\rm before};P_X;q_X)
       -\sum_X F_X(a_{\rm st};P_X;q_X).
\]
Each row is evaluated at its own pair and stored-data
restriction. Their norm bound is
$\eta_{\rm sr}b e^{-(r-1-c_E)d}$, where
$\eta_{\rm sr}=2^{14}e^{-\check R_{k+1}}<1/8$ under the
chosen floor $\check R_{k+1}\ge12$ and the strip conditions.
They retain holomorphy, locality and the stated symmetries,
are supported on made boundary domains, and vanish for the
empty history.
\end{lemma}
\begin{proof}
The strip construction is a finite difference of values of
$F_X$ along the admitted external-field substitution. Its
Cauchy contours use the complex room
$C_{\rm xt}\mathfrak b(\mathsf t_k)$; the real stored field
occupies only the prescribed half of the doubled box.
Apply a norm-one functional to the external algebra.
Linearity identifies the resulting row with the scalar
transfer row. Its contour estimate and the sum over
domains give the factor $2^{14}e^{-\check R_{k+1}}$ and
the loss $1+c_E$ in the rate. Taking the supremum gives
the norm estimate. The scalar sum rule is a linear
identity, hence holds coefficient by coefficient. Its
restriction identities prove the support and locality
claims. This argument concerns this external-field row;
it is not an assertion that every coordinate identity is
supplied by a single transfer estimate.
\end{proof}

\subsection{The running coupling in the calibration frame}

The tests and coordinates are fixed by $\mathsf t_k$, whereas
the action has coefficient $x_k$. The difference changes
both the integrated quadratic form and the factor left on
the collar. We first keep track of this difference exactly.

\begin{lemma}[The quadratic form of the actual step]
\label{C-399-33}
Put $\tau=x_k/\mathsf t_k>0$, with $|\tau-1|\le1/4$.
In the fluctuation scaled by $\mathsf t_k^{1/2}$ write
$A_1=G+A^{\rm act}$ for the calibration quadratic form,
where $G$ is its block gauge-fixing form. Let $I$ and $E$
be the integrated and stored collar bonds. Then
$G=G_{II}\oplus G_{EE}$, and the actual form is
\[
 A_\tau=G+\tau A^{\rm act}=\tau A_1+(1-\tau)G.
\]
Use the conditional Gaussian of $\tau A_1$, of covariance
$x_k^{-1}C_k$ in the unscaled variable, and put the correction
$-\tfrac12(1-\tau)(a,G_{II}a)$ among the tested ordinary
origins. Its mean is
$-K_1a'$, where
$K_1=(A_1)_{II}^{-1}(A_1)_{IE}$ and $a'$ is the collar
field before the strip. The calibration formed factor on
that field is multiplied by
\begin{equation}\label{eq:qcd-collar-mismatch}
 e^{R(a')},\qquad
 R(b')=-\tfrac12(\tau-1)
       \langle b',(S_1-G_{EE})b'\rangle,\quad
 S_1=(A_1)_{EE}-(A_1)_{EI}(A_1)_{II}^{-1}(A_1)_{IE}.
\end{equation}
At stored coordinates retain $R(a'_{\rm st})$ as a scalar
boundary term and file $R(a')-R(a'_{\rm st})$ in its strip
rows. The tests, geometric histories and stored gauge
factors remain those of the calibration frame.
\end{lemma}
\begin{proof}
The action coefficient is one number on the whole current
region. Thus $\tau$ multiplies all three action blocks,
including the mixed and collar blocks. The gauge-fixing
form is a sum of terms inside individual coarse blocks.
Every such block lies wholly in the integration region
or wholly in its collar; hence its mixed block vanishes.
Complete the square with $\tau A_1$:
$a=\widetilde a-K_1a'$. The factor $\tau$ cancels from
$K_1$, so every test and vertex reads this same shifted
field, and the q-origin is evaluated there as well.
The Gaussian leaves
\[
 \exp\{-\tfrac12\tau\langle a',S_1a'\rangle
             -\tfrac12(1-\tau)\langle a',G_{EE}a'\rangle\}.
\]
Factoring out the calibration Schur form proves
\eqref{eq:qcd-collar-mismatch}. The retained and rejected
localizations of $-\frac12(b',S_1b')$ are identities in
$b'$. Their strip sum rules, together with the transfer
rule for the local forms $\frac12(b',G_{EE}b')$, give
exactly the difference of $R$ at the two collar fields.
No collar gauge-fixing term is inserted a second time
among the integrated q-origins.
\end{proof}

\begin{lemma}[Bounds after real calibration]
\label{C-399-44}
Under the admitted calibration-frame conditions, the
conditional structure used in \cref{C-399-33} is obtained
by replacing its covariance and connecting operator by
$\tau^{-1}C$ and $\tau^{1/2}\Gamma$. In the variable
$a'=x_k^{1/2}B$, external fields and thresholds are
multiplied by $\tau^{1/2}$. The centered quadratic-growth
norm increases by at most $4/3$, and the Gaussian test
exponent becomes $\tau\gamma_P$. The other geometric
frame constants remain those of the calibration frame.
The q-origin is confined to the integrated fine cubes.
Its untested all-real-fluctuation expansion gives
$\mathbf G_\tau$; in the actual step it has the tested
ordinary sort of \cref{C-388-37}.
\end{lemma}
\begin{proof}
The change of variables is real and positive. Substitution
in the covariance, thresholds and conditional mean gives
the stated transformations. The only increase in a
quadratic-growth denominator is bounded by
$\max(\tau,\tau^{-1})\le4/3$; the Gaussian exponent scales
by $\tau$. The q-origin returns precisely the integrated
gauge-fixing difference in \cref{C-399-33}. Its restriction
to integrated cubes follows from the separate collar block
there. Applying \cref{C-388-37} distinguishes its two
sorts and supplies the comparison with $\mathbf G_\tau$.
These changes of variables prove a statement for real
$\tau$; complex-parameter holomorphy is obtained on the
stated coefficient domains, not by treating this real
Gaussian substitution as a complex change of variables.
\end{proof}

\subsection{The scalar entries and their local variations}

The pure frame accepts a local array and boundary entries
with specified bounds. We must supply the actual entries
of the matter construction. In particular the coupling
extraction has already been removed from the regular
activities, and must not be inserted again as a counterterm.

\begin{definition}[The scalar list of a history]\label{C-399-46}
At level $k$, use calibration $\mathsf t_k$ and zero
counterterms in the local array. Its regular entries
are the pointed degree-zero activities with their vacuum
and action parts removed. Its made boundary entries are
the scalar parts of earlier generated terms. Its formed
factor family is the unchanged calibration family.
If $X(p;j)$ is the pointer cell assigned to a plaquette at
level $j$, set
\[
 w(p)=\sum_{j\le k}L^{j-k}b_j
            1_{\{X(p;j)\not\subset\Lambda_j\}}.
\]
With the localizations of the Jacobian, gauge-fixing,
action, regular array, old boundary and formed factors
denoted respectively by
$V^J,V^G,V^A,V^E,V^B,V^\Phi$, the two scalar classes are
\[
 V_{\rm rg}^F=V^J+V^G+\tau V^A+V^E-V^w,\qquad
 V_{\rm bd}^F=V^B+V^\Phi.
\]
Here $V^w$ is the variation of $L^kA(w;U)$ under the
shift. The q-origin of \cref{C-399-33} is a separate class.
\end{definition}

\begin{lemma}[Admission and sum rules for the scalar list]
\label{C-399-47}
Suppose the earlier history bounds hold and the actual
array satisfies \eqref{eq:qcd-degree-zero-array}.
Take $B\ge8M^3E_0$ and the unchanged formed-factor
bound
$B_{\rm fr}^{\flat}=K_{\rm ch}C_{\rm xt}^2+
12M^3C_z a_{\rm nf}$.
Then the scalar list has the required local holomorphic
entries, pair ties and Borel dependence, with array bound
$E_0$, made-entry bound $B$ and edge bound $8M^3E_0$.
Moreover $|w|\le\bar B$, and $w$ vanishes on every plaquette
whose pointer cell is wholly in the current integration
region. Its nonzero plaquettes have a corner in $Z_k$.
The sum of its class potentials is the variation of the
actual degree-zero exponent, each part counted once.
To state their sizes, put
\[
 g_k=\mathsf t_k^{-1/2},\quad y_k=\log\mathsf t_k,\quad
 \mathfrak b_k=\mathfrak b(\mathsf t_k),\quad
 \alpha_{0,k}=C_0g_ky_k^q,\quad q>\mathfrak p_0.
\]
Here $\alpha_{0,k}$ is the gauge frame's analytic radius.
With $\delta_k=\delta(\mathsf t_k)$, write
\[
 \vartheta_{0,k}=\frac{\delta_k}{\alpha_{0,k}},\qquad
 \mathsf s_k^\sharp[E]=\frac{\delta_k}{\alpha_0}
       (1+E+\alpha_0^2B_\sharp)+M^3g_k\mathfrak b_k^3.
\]
These are the quantities in
\cite[Prop.~58.26, equations~(1318)--(1319)]{YMM1},
evaluated at the calibration coupling. The tier bounds are
\[
 \|V_{\rm rg}^F\|_{r_V}
 \le\max(1,\tau)C_{\rm eff}
       (\mathsf s_k^\sharp[E_0]+\vartheta_{0,k}g_k^2y_k^{2q}),
 \qquad
 \|V_{\rm bd}^F\|_{r_V}
 \le C_{\rm eff}(B+B_{\rm fr}^{\flat})\vartheta_{0,k}.
\]
Consequently the effective threshold
$y_k\ge Y_{\rm eff}^F(\epsilon)$ makes these bounds at
most $\max(1,\tau)\epsilon$ and $\epsilon$.
\end{lemma}
\begin{proof}
The regular entries satisfy the array bound by their
construction, with their prescribed pointing. Restriction
to an edge gives at most $8M^3$ shares; their norms, rather
than a second sum over arbitrary pointers, give the edge
bound. The made entries obey their bounds on their full
filing domains. The admitted frame's localizations apply
to an arbitrary list with precisely these bounds and pair
ties; no induction theorem for an unspecified matter
array is used.

For the coupling term use
$L^kx_k=\beta-\sum_{j\le k}L^jb_j$. Before smoothing the
pointer shares its identity is
\[
 -\beta A(U)+\sum_{j\le k}L^jb_j
       \sum_{X(p;j)\subset\Lambda_j}a_p(U)
       =-L^kx_kA(U)-L^kA(w;U).
\]
The profiles and their rim shares restore this identity
because their weights add to
$1_{\{X(p;j)\subset\Lambda_j\}}$: in the interior the
profile is one, and on the rim its complement is
distributed among shares with sum one. This explains why
$-V^w$ replaces the counterterm class. The geometric
series in the bounds on $b_j$ gives $|w|\le\bar B$.
Every excluded pointer has a corner in $Z_k$; the
separation of $Z_k$ from $\Omega_{k+1}$ gives
$d_k(Y)\ge7L\check R_{k+1}-2$ whenever its variation
on $Y$ is nonzero.

The local sum rules now give respectively the change of
the pointed regular array, of $L^kA(w;U)$, of the old
boundary entries at unchanged stored data, and of the
formed factors. On an interior domain the last three
boundary contributions vanish, and $V^E$ is the flat
array variation. Applying the admitted localization
bounds to this verified list gives the two displayed
estimates. Their constants depend on the three entry
ceilings; the definition of $Y_{\rm eff}^F$ is the
threshold that makes their right sides at most the
stated sizes.
\end{proof}

\begin{lemma}[Variation of an earlier boundary coefficient]
\label{C-399-48}
Let $g_{j,X}$ be an earlier made or edge family with
values in a finite normed space. Suppose it depends on
the pair, is Borel in the real and stored data, is
holomorphic on the admitted slices of its filing domain, and
satisfies $\|g_{j,X}\|\le b^\bullet e^{-\kappa d_j(X)}$.
Apply the ordinary or pinned coefficient localization of the
history frame, keeping the prescribed stored data fixed.
The resulting coefficients are local, analytic on the tiers
and Borel there. Their sum is the difference of the values
at the shifted and unshifted configurations. Summed over
the earlier made and edge families, their bound at each
output polymer $Y$ is
\[
 C^{\rm bd}\vartheta_{0,k}
       (b^{\rm mk}+b^{\rm ed})e^{-r_Vd_k(Y)}.
\]
They vanish unless $d_k(Y)\ge7L\check R_{k+1}-2$.
For a jointly gauge-invariant made term with external
legs, depending on the pair, the coefficient at a
pinned substitution need not itself be a function of the
pair. Its product with the inverse joint action of the
levelling function, formed at the same parameters, is a
function of the pair and has the joint covariance and tie
of the original term.
\end{lemma}
\begin{proof}
The coefficient localization is a finite bracket, with
coefficients $\pm1$, of values at substituted triples.
The substituted triples do not depend on $g_{j,X}$.
Consequently localization commutes with every norm-one
linear functional $\ell$. The scalar localization estimate
uses the analyticity and bound of $\ell\circ g_{j,X}$;
it does not require scalar gauge invariance. Its Cauchy
radius contributes
$C_{\rm sb}C_{\mathcal A}a^\bullet
e^{C_{\rm sb}\bar\kappa}$. The domain-count constant
$C_B(0)$ includes this factor $C_{\rm sb}$. The made and edge radii satisfy
$a^{\rm mk}\le c_-^{-1}\vartheta_{0,k}$ and
$a^{\rm ed}\le\vartheta_{0,k}$. Thus the choice
\[
 C^{\rm bd}\ge
 C_{\mathcal A}e^{C_{\rm sb}\bar\kappa}C_B(0)(1+c_-^{-1})
\]
gives the stated bound after taking the supremum over
$\ell$. The bracket sum rule gives the difference of the
two endpoint values, at unchanged stored data.

For the assertion about the pair, write
\[
 \Pi_{j,X}(s,\lambda;\mathfrak u)
 =\theta_{u_{s,\lambda}}^{-1}
  g_{j,X}\bigl(T_{X|h'}(s,\lambda;\mathfrak u);q|_X\bigr),
\]
where $\theta_v$ turns the nonstored generators by $v$ and
leaves the stored generators fixed. The coefficient and
$u_{s,\lambda}$ use the same parameters. At two triples
with one pair, let $u$ and $\widetilde u$ be their
levelling functions. The substituted coefficients are
related by $\widetilde g=\theta_{\widetilde u u^{-1}}g$.
For real turns this follows from joint covariance: the
connecting gauge transformation is one at the tops and
stored sites, so it preserves the stored data. Holomorphy
in the turn parameters gives the identity at the pinned
turn. Hence
$\theta_{\widetilde u}^{-1}\widetilde g
=\theta_u^{-1}g$. This proves pair dependence of $\Pi$;
the same covariance identity gives its tie.

At fixed real configuration and stored data, normalization
to a tame point is linear in the two form coordinates.
It preserves the substituted coefficient and the levelling
function. At the normalized point these are locally
holomorphic in the form coordinates, the fluctuation and
the interpolation parameters. Their composition with
normalization is therefore analytic on the tiers.
The maps are Borel at fixed parameters, and the finite
brackets preserve Borel dependence. These assertions
concern the admitted filing domains and tiers.
\end{proof}

\begin{lemma}[Local bound for the retained collar correction]
\label{C-399-49}
The term $R$ in \eqref{eq:qcd-collar-mismatch} has a local
family supported on domains meeting stored collar bonds:
\[
 R_{k+1}(X)
 =(\tau-1)\{K_{\rm kept}(X)+B_{\rm cross}(X)
                              +Q_X^{\rm gf}\},
 \qquad
 Q_X^{\rm gf}=\tfrac12\sum_{\square\subset X}
                  (a'_{\rm st},\widehat G_\square a'_{\rm st}).
\]
The last sum is assigned to a single big collar cube and
is zero on other domains. On the full admitted complex
pair and stored-field domain,
\[
 \begin{gathered}
 |R_{k+1}(X)|\le r_ke^{-\kappa''d_{k+1}(X)},\\
 r_k=\frac{\bar B}{\mathsf t_k}C_{\rm xt}^2
       \mathfrak b(\mathsf t_k)^2
       \{c_{\rm kp}^{\rm th}
          +c_{\rm rj}^{\rm th}e^{-(\check R_{k+1}-2)}
          +c_{\rm gf}\},\\
 c_{\rm gf}\le\tfrac32c_G(LM)^3,\qquad
 c_G=6L^6/r_1^2.
 \end{gathered}
\]
Its strip rows have bound
$2^{16}e^{-\check R_{k+1}}r_ke^{-\kappa_+d}$.
Under the common small-coupling ceiling
\[
 \log\gamma^{-2}\ge2\mathfrak p_0,\qquad
 \tfrac32\bar B\gamma^2C_{\rm xt}^2A_1^2
  (\log\gamma^{-2})^{2\mathfrak p_0}
  \{c_{\rm kp}^{\rm th}+c_{\rm rj}^{\rm th}
                                  +\tfrac32c_G(LM)^3\}\le1,
\]
the retained family and its rows together are bounded by
$e^{-\kappa_+d}$.
\end{lemma}
\begin{proof}
Localize the Schur form in \eqref{eq:qcd-collar-mismatch}
by its kept and rejected families. Its identity in the
external field gives the stated local expression.
Every fine cube used in the gauge-fixing sum is a source
cube of the local Gaussian triple. Its restricted triple
is in the class on which
$\widehat G_\square$ has row bound $c_G$, also on the
complex slices. A big cube has at most $3(LM)^3$
positive bonds; this gives the displayed $c_{\rm gf}$.
The kept and rejected estimates and
$|a'_{\rm st}|\le C_{\rm xt}\mathfrak b(\mathsf t_k)$
give $r_k$. All terms vanish without a stored collar
bond, hence are boundary terms.

The kept and rejected quadratic strip identities and
\cref{C-399-42} for $Q^{\rm gf}$ give the exact sum
$R(a')-R(a'_{\rm st})$ and the row bound. The strip
floor makes the multiplier at most $1/2$.
Finally $\mathfrak b(t)^2/t=
A_1^2t^{-1}(\log t)^{2\mathfrak p_0}$ decreases once
$\log t\ge2\mathfrak p_0$. Its value at $\gamma^{-2}$
and the stated ceiling bound $(3/2)r_k$ by one.
\end{proof}

\subsection{The actual frame and the exact step}

\begin{lemma}[The joined frame on a history]\label{C-399-32}
Suppose the history induction holds through level $k$,
with earlier fermion-origin sizes
$\mathsf{sz}_j^\natural\le16\varepsilon_\flat/17$.
Suppose also that the present size
$\mathsf{sz}_k^\natural\le\varepsilon_\flat$ and the
scalar-class, q-class and failed-test thresholds of the
common choice hold. Then the actual step is the joined
centered expansion of \cref{C-399-41}: it uses covariance
$x_k^{-1}C_k$, the calibration tests, the scalar classes
of \cref{C-399-46}, the tested q-origin of
\cref{C-399-33}, and the actual fermion potentials
contracted with the symmetric weakened covariance.
All big cubes are guards, including cubes outside
$\Omega_{k+1}$. Its output $E_{\rm lab}(X)$ obeys
\[
 \|E_{\rm lab}(X)\|\le\tfrac12e^{-\kappa''d_{k+1}(X)},\qquad
 \|(E_{\rm lab}(X))_{\ge2}\|
 \le K^{\rm lb}\mathsf{sz}_k^\natural
                         e^{-\kappa''d_{k+1}(X)}.
\]
It is defined on the admitted local pair, external-field
and stored-data domains and has the properties in
\cref{C-399-41}. The present-size hypothesis is verified
in \cref{C-399-21} before this step is applied.
\end{lemma}
\begin{proof}
In the scaled variable $\mathsf t_k^{1/2}B$, the integrated
covariance is $[\tau(A_1)_{II}]^{-1}$, the inverse block
of the reference form $\tau A_1$. The tested q-origin
supplies the remaining correction $(1-\tau)G_{II}$ to the
actual quadratic form, as in \cref{C-399-33}.
The positive real rescaling in \cref{C-399-44}
has the centered-frame bounds. The density tests
place the gauge fields in the local domains by
\cref{C-388-38}; the local vertices and contractions
there are defined by \cref{C-388-34,C-397-47}.
In particular the contraction parameters outside the
integration region have not been discarded.

The actual scalar list is admitted by \cref{C-399-47}.
Its two classes, the class with fermion legs and the
q-class give at most four ordinary origins. With
$y_k\ge Y_{\rm eff}^F(2\varepsilon_\flat/5)$ and the
q and failed-test ceilings, their amplitudes fit the
centered expansion's fixed envelope. Real calibration
changes the appropriate quadratic-growth bound by at
most $4/3$. Thus the admitted scalar contour estimate
applies to these members. The contraction properties
were checked after \cref{C-399-41}, so that lemma gives
the whole bound. Setting the fermion class to zero
leaves only scalar outputs. Marking that class therefore
gives the selective positive-degree bound.

At $k=0$ no old scalar array or old boundary family is
present. The all-field inverse $\Gamma_0^{-1}$ and the
first centered frame give the same construction and
the same enlarged constants. The q-only all-real
expansion is used solely for the comparison with
$\mathbf G_\tau$; the actual q-origin here keeps its
tests. Borel dependence and the restriction identities
follow before summation and persist by the common
absolute majorant.
\end{proof}

\begin{lemma}[Agreement on an interior domain]\label{C-399-45}
For $X\subset\Lambda_{k+1}$, the output of
\cref{C-399-32} is
\[
 E_{\rm lab}(X)=\mathbf G_\tau(X)+\widetilde E_k(X),
\]
where $\widetilde E_k$ is the small-field output on the
same torus. On nonwrapping domains, after the prescribed
transport it agrees with the infinite-lattice output.
\end{lemma}
\begin{proof}
On $X$ there are no external bonds or stored collar
coordinates and the conditional mean is zero. The
thresholds are the calibration thresholds. The
conditional structure, its parameter extension and its
rooted indices are therefore the one-lattice ones.
The scalar classes agree term by term: the boundary,
formed-factor and missing-weight variations vanish,
while the regular array is the same pointed array.
The fermion members, including their edge and wrapping
assignments, are restrictions of the same local
functions. Their weakened covariances agree by the
intrinsic restriction identity. Thus every conditional
integral and every parameter bracket is the same,
not just its support or bound. The small-field
definition subtracts the all-real q-family
$\mathbf G_\tau$; this gives the displayed identity.
For nonwrapping domains the local walks, tests and
vertices agree with their infinite-lattice counterparts,
which proves the last assertion.
\end{proof}

\begin{proposition}[The exact step with matter]\label{C-399-43}
Assume the admitted gauge coordinate identities, the
actual scalar and fermion member bounds of
\cref{C-399-32}, and the local strip identities above.
Then the gauge step with its fermion contraction is an
exact identity of coefficientwise measures. Its
conditional quadratic datum is
\[
 A_{\mathsf d}^F=\tau A_{\mathsf d}
       +(1-\tau)\chi_E^T(\widehat G_{\mathsf d})_{EE}\chi_E,
\]
and its scalar normalization changes by
\[
 c_t^F=-\tfrac32\log\tau\,\# I+\log\delta Z_k^F(1),
\]
where $\#I$ counts positive integrated gauge bonds.
The geometric coordinate maps are unchanged. The
fermion Gaussian is normalized inside the step and
its determinant is outside; all cluster and strip
families occur once.
\end{proposition}
\begin{proof}
The $II$ and $IE$ blocks of the displayed datum are
$\tau$ times the calibration blocks. Its $EE$ block
contains the separate collar gauge-fixing correction,
as required by \cref{C-399-33}. The integrated
normalization scales by $\tau^{-3\#I/2}$, because
the Lie algebra has three real coordinates on each
positive bond. The exact fermion Schur identity
adds $\delta Z_k^F$ and the normalized contraction.
The cluster exponential of \cref{C-399-41} evaluates
the resulting integral, and \cref{C-399-42,C-399-49}
move its external field to stored coordinates with
their exact differences. Together with the sum rules
of \cref{C-399-47,C-399-48}, these are precisely the
members of the transformed exponent.

The coordinate pushforward identity is applied to
each Grassmann coefficient. Split its real and
imaginary parts into positive and negative parts;
the local norm majorants give their integrability.
This justifies the four ordinary measure identities
and their recombination. The scalar fermion
determinant is continuous and positive at real
configurations, so the old coordinate exceptional
null sets remain null after multiplication by that
density. No positivity or Tonelli assertion is made
for a Grassmann-valued integral.
\end{proof}

\subsection{The store and boundary reentry}

\begin{lemma}[The stored quadratic factor]\label{C-399-17}
Let a Grassmann algebra have $N$ pairs of generators, both of
weight $h$, and let $A$ be their quadratic matrix. Then
\[
 \|e^{-\bar\Psi A\Psi}\|_h
 \le(1+h\|A\|^{1/2})^{2N}.
\]
Consequently the stored factor of a history is bounded by
$e^{c_{\rm store}\Sigma'(h)}$, where $c_{\rm store}$ may depend
on the fixed fibre weight and is independent of the history,
cutoff and volume.
\end{lemma}
\begin{proof}
The coefficient of degree $(p,p)$ is a minor of $A$.
Its absolute value is at most $\|A\|^p$, since the corresponding
matrix of exterior powers has norm at most $\|A\|^p$.
Thus the norm is at most
$\sum_p\binom Np^2h^{2p}\|A\|^p$, which is bounded by the
square of $\sum_p\binom Np h^p\|A\|^{p/2}$.
The all-field tower bound $\|\mathbb D_h\|\le c_{\rm tw}$ and
$N=\nu\Sigma'(h)$ prove the assertion. This estimates the
exponential directly; no bound of the quadratic form by its
operator norm times the number of places is required.
\end{proof}

The block Schur identity also holds with stored generators.
After integrating the places $A_j$, its determinant factor is
$\delta Z_j^F=\theta^{-\nu\#_{j+1}(\Omega_{j+1})}\det\Gamma_j$, its new
quadratic matrix is the Schur complement, and each old background
is the new background plus an image of the fluctuation $W$.
These identities are matrix identities before any estimate.

\begin{proposition}[Generated boundary terms]\label{C-399-15}
Under the induction and present-frame hypotheses of
\cref{C-399-32}, the output on a polymer meeting $Z_j$
is a generated boundary term. It is local, even and
jointly gauge invariant; it is holomorphic on the
admitted pair domains and real under the model's
conjugation. At creation its scalar part is bounded by
$B e^{-\kappa d_j(X)}$ and its part with fermion legs by
\[
 (1+\eta_{{\rm sr},j})K^{\rm lb}
 \mathsf{sz}_{j-1}^\natural e^{-\kappa_+d_j(X)},
 \qquad \eta_{{\rm sr},j}<1/8.
\]
Edge parts from pointed regular activities have the
corresponding $4K_M\bar v_j$ bound with their pointer
factor, where $\bar v_j\le b_{j-1}^{\rm blk}$ is the
small-field input size at their birth. The boundary
index sets obey
$\sum_{X\in\mathsf{Ix}^{\rm mk}_j}e^{-\kappa d_j(X)}
 \le C_{\mathsf B}\#_j(Z_j^\sim)$.
At a stored place of level $m<k$ use weight
$w_k(m,x)=(1+2^{-(k-m)})h_F$; at a formal background
leg of level $n<m$ in that layer use
$w_k(n,x)=(\frac12+2^{-(k-m)-1})h^{(m)}(n)$.
These weights decrease with $k$ and preserve the bounds
at all later steps that have been constructed.
\end{proposition}
\begin{proof}
Every output is assigned to a polymer containing the
closed enlargement of its integration support. If that
polymer is not contained in $\Lambda_j$, it meets
$Z_j^\sim$. The selective estimate of \cref{C-399-32}
and the strip estimate of \cref{C-399-42} give the
positive-degree bound. The normalizer, the q-only
Gaussian family and the collar correction are scalar;
their fixed bounds fit in $B$. The entropy sum through
a cube gives the index-set bound. Restriction of every
input to its filed support, followed by the component
identities of the conditional construction, gives
locality. Gauge covariance, conjugation and holomorphy
pass through the finite integrals and the absolutely
convergent sums.

For later reentry, retain the room in the weights.
At a row in layer $m$, split an old background into
$a+b+c$: the new formal background, the image of $W$,
and its change under the gauge fluctuation. For a
formal leg of level $n<m$ the available room is
\[
 h_b'=w_k(n,x)-w_{k+1}(n,x)
       =2^{-(k-m)-2}h^{(m)}(n).
\]
The new formal background is not expanded into all
stored generators: the whole row bound of
\cref{C-397-33} gives its admitted reading, while such
an expansion here would lose locality and its weight
reserve. Only $b$ and $c$ are expanded in walks.

A walk from layer $m<k$ to $A_k$, read at depth $k+1$,
crosses the intermediate layers and the band
$\Omega_k\setminus\Omega_{k+1}$. Since
$R_\flat/(3M_w)\ge8L$, it has at least $8L(k-m)$
steps. Half its available exponential weight pays
the length polynomials, and half gives
$e^{-4L(k-m)}\le L^{-4(k-m)}$. Consequently
\[
 \frac{h_b}{h_b'}\le4Lc_\psi(2L^{-4})^{k-m},
 \qquad c_\psi=2A_\kappa\Gamma_F/h_F .
\]
The fine shift into layer $m$ is bounded by
$c_hL^{-k}L^{-4(k-m)}\sup|B|$. Cauchy's estimate on
the link radius $\alpha_1^*L^{-m}$ gives the factor
$L^{-(k+1-m)}$ in addition to
$2C_{\mathcal A}\delta(\mathsf t_k)/\beta_0$.
Write $c=c_1+c_2$, where $c_1$ changes the background
kernel and $c_2$ reads the changed site factors at
its stored columns. A walk reaching a later layer
$m'$ loses one crossing band; its Cauchy factor
compensates this loss up to $L^3$. The complex row
bound and the room above give
\[
\begin{split}
 \frac{h_{c_1}}{h_b'}&\le
 \frac{2C_{\mathcal A}\delta(\mathsf t_k)}{\beta_0}
      8L^4A_\kappa C_H(2L^{-4})^{k-m},\\
 \frac{h_{c_2}}{h_b'}&\le
 \frac{2C_{\mathcal A}\delta(\mathsf t_k)}{\beta_0}
      6L^5A_\kappa C_{\mathfrak r}\beta_0
                                 (2L^{-4})^{k-m}.
\end{split}
\]
For the second estimate use
$|\mathfrak r_{m'}-1|
 \le C_{\mathfrak r}L^{-4(k-m')}\sup|B|$;
together with crossing it costs at most $L^4$
times the decay from $m$ to $k$.
A stored generator itself is split as
$\Psi+(\mathfrak r_m-1)\Psi$; its falling weight
bounds the second part by the ratio
$3C_{\mathfrak r}C_{\mathcal A}\delta(\mathsf t_k)
(2L^{-4})^{k-m}$.

In a made boundary term a nonstored leg also carries
the levelling function $u(x)$. The local shift bound is
$|u^{\pm1}-1|\le C_{\rm m}c_h
L^{-5(k-m)}\sup|B|$. On the doubled parameter contour,
its ratio to the room is at most
\[
 \frac{2C_{\mathcal A}\delta(\mathsf t_k)}{\beta_0}
 4C_{\rm m}c_h\beta_0(2L^{-5})^{k-m}.
\]
Here $u$ is evaluated at the same parameters as the
coefficient substitution, as in \cref{C-399-48}.
Coefficient variation itself is estimated by that
lemma, separately from these leg changes.

\emph{Combining the coefficient and leg changes.}
Fix a made term with base $X$ at level $j\le k$.
Let $g_s$ be its coefficient at the pinned substituted
triple, $u_s$ the levelling function, and $c_s$ the change
of its background legs, all formed at the same shift
parameters $s$ and at shift amplitude $\lambda=1$.
Thus $u_0=1$, $c_0=0$, and at $s=\mathbf1$ these are the
coefficient, levelling function and leg change of the step.
Stored generators of $h'$ are left in their places.
With $g_s[\ell]$ denoting substitution of the nonstored
legs $\ell$, put
\[
\begin{split}
 G(s)&=g_s[a],\\
 R^b(s)&=g_s[a+b]-g_s[a],\\
 R^c(s)&=g_s[a+b+c_s]-g_s[a+b],\\
 R^u(s)&=g_s[u_s(a+b+c_s)]-g_s[a+b+c_s].
\end{split}
\]
The change from the new formal legs to the old reading is
\[
 G(\mathbf1)-G(0)+R^b(\mathbf1)
                    +R^c(\mathbf1)+R^u(\mathbf1).
\]
Here $R^c(0)=R^u(0)=0$; $R^b(0)$ is the term with a
leg at $W$ and no dependence on $B$.

Expand only the parts $b,c_s$ into local terms and walks.
For a monomial and a choice $\omega$ of these parts, let
$N_\omega$ be the base $X^\uparrow$ enlarged by the cubes
met by their walks. For an edge term also include the
blocks read by its stored-leg site factors. The product
reads only those components of the shift that meet
$N_\omega$, and is consistent when its domain is enlarged.
Its brackets $[Y\setminus N_\omega\mid N_\omega]$ are
therefore assigned to $Y\supset N_\omega$. Their sum,
together with the value at the base, recovers the product
at $s=\mathbf1$. The bracket of $G$ gives its difference
between $s=\mathbf1$ and $s=0$.

The split above is used to estimate the terms. A separate
part of a made term can read the real factor of the
triple. For each choice $\omega$, its part from $G$, $R^b$
or $R^c$ is $g_s[\omega]$, and its part from $R^u$ is
$g_s[u_s\omega]-g_s[\omega]$. Their sum is
\[
 g_s[u_s\omega]
   =\bigl(\theta_{u_s}^{-1}g_s\bigr)[\omega]
   =\Pi_{j,X}(s,1;\mathfrak u)[\omega].
\]
All parts have the same roots $N_\omega$. Their brackets
at each $Y$ thus add to the bracket of $\Pi_{j,X}[\omega]$;
the value kept at the base has the same identity because
$u_0=1$. Lem.~\ref{C-399-48} and covariance of the chosen
legs imply that this sum is a function of the pair, with
the required tie and joint covariance. Summing these
identities at $Y$ proves the corresponding properties
of the member of the frame at $Y$.

For analyticity fix the real configuration $U$ and the
stored data on $Y$. On the normalized, tame slice, the
pinned substituted triple and the levelling function are
locally holomorphic in the form coordinates, $B$ and $s$.
The canonical substitution radius is at least two, so
the levelling function is also holomorphic for
$|\lambda|<2$. The parts $b,c_s$ are locally holomorphic
by their local and walk formulas: the moved configuration
stays in the admitted complex class for $\sup|B|\le\beta_0$.
The normalization preserves the coefficient and levelling
function. Composition with it therefore proves analyticity
on every tier. The walk sums converge uniformly on the
parameter domains, and their finite brackets retain this
property. This argument applies to the separate parts
used for estimates as well as to their recombination.
For an edge term the substituted triple is plain and
there is no $u_s$ or $R^u$; each part already reads the pair.

There is no factor equal to the number of legs.
Indeed, if $h_i=h_{a,i}+h'_{b,i}$ and
$h_{b,i}\le\rho_i h'_{b,i}$ with the split parts
within their rooms, the terms with first split
leg $i$ have coefficients
$t_i=h'_{b,i}\prod_{l<i}h_{a,l}\prod_{l>i}h_l$.
They telescope to
$\sum_i t_i\le\prod_i h_i$. Thus the sum of the
split terms is at most $\max_i\rho_i$ times the
original norm. This permits the walk carrying the
split leg to retain its own crossing factor.

\emph{Counting the domains of the brackets.}
For the parts $R^c,R^u$, put $p_l=h'_{b,l}/h_l$.
The nonsplit parts have weights at most
$(1-p_l/2)h_l$. Their telescoping sum is bounded
using
\[
 \sum_i p_i\prod_{l\ne i}(1-p_l/2)\le1,
 \qquad 0\le p_l\le1.
\]
Indeed, the left side is affine in each $p_l$;
at a vertex with $q$ nonzero coordinates it is
$q2^{1-q}\le1$, and at the zero vertex it is zero.
This gives respectively the largest ratio
$h_c/h_b'$ or $2h_u/h_b'$ times the coefficient norm.
The estimate for $R^b$ uses the preceding telescoping
with ratio $h_b/h_b'$. Each first split leg therefore
keeps its own ratio in all three parts.

This count must be made on one parameter contour for
each output polymer. Set
\[
 x=C_{\rm g}e^{(\kappa+2)c_0}e^{-(\bar\kappa-1)}
       \le e^{-8},\qquad
 r=1+e^{\bar\kappa-1}<e^{\bar\kappa}.
\]
The floor on $\bar\kappa$ gives these inequalities and
$r/(r-1)\le e^x$. For a fixed term and $Y$, let $R_N(s)$
be the sum of its parts from $R^b,R^c,R^u$ with roots
$N_\omega=N$. On the common torus
$\mathfrak T_r(Y)=\{s:|s_D|=r\text{ for every }D\subset Y\}$,
Cauchy's formula gives
\[
 R_N[Y\setminus N\mid N]
 =\int_{\mathfrak T_r(Y)}R_N(s)
     \prod_{D\in N}\frac{s_D}{s_D-1}
     \prod_{D\in Y\setminus N}\frac1{s_D-1}
     \prod_{D\subset Y}\frac{d\theta_D}{2\pi}.
\]
This also includes the base value when $N=Y$.
Consequently the norm of the part through the legs at $Y$
is at most
\[
 \sup_{s\in\mathfrak T_r(Y)}
 \sum_N e^{x\#N}e^{-(\bar\kappa-1)\#(Y\setminus N)}
                                      \|R_N(s)\|.
\]
The sum is evaluated at one $s$. This allows us to apply
the algebra norm and telescope the legs before taking
the supremum.

Every component of $Y\setminus N$ meets one of at most
$26\#N$ neighbouring cubes of $N$. There are at most
$C_{\rm g}^n$ connected sets of $n$ cubes through a given
cube, and adjoining a cube increases the size by at most
$c_0$. At any rate $0\le\kappa'\le\kappa+2$, the resulting
domain sum is therefore at most
\[
 (1-x)^{-26\#N}\le e^{27x\#N}.
\]
It follows that $e^{\kappa'd_k(Y)}$ times the preceding
norm is bounded by
\[
 \sup_{s\in\mathfrak T_r(Y)}
       \sum_N e^{28x\#N}e^{\kappa'd_k(N)}\|R_N(s)\|.
\]
At fixed $s$, decompose $g_s$ into products of unit
vectors in its Grassmann norm, substitute the chosen
parts of the legs, and use their walk bounds. Taking
the infimum over these decompositions gives the norm
of $g_s$, so the bound retains the ratio of the first
split leg. It introduces no separate supremum for each
root or monomial.

The remaining factor is paid by the same weight reserves.
Indeed,
\[
 \#N_\omega\le8\bigl(d_k(X^\uparrow)+1\bigr)
                  +8\sum_l(|\omega_l|+1),\qquad
 224x\le\frac1{13}.
\]
Thus
\[
 e^{28x\#N_\omega}
 \le e^{(d_k(X^\uparrow)+1)/13}
                    \prod_l e^{(|\omega_l|+1)/13}.
\]
The walk estimate holds with $e^{3|\omega|/2}$ in place
of $e^{|\omega|}$ and the same $A_\kappa$; it pays
$e^{|\omega_l|/13}$. The preceding bounds, with the
floor on $h_F$ and ceiling on the coupling, give
\[
 \frac{h_b+h_c+2h_u}{h_b'}\le\frac18,
 \qquad e^{1/13}(h_b+h_c+2h_u)\le\frac12h_b'.
\]
Here $h_u=C_{\rm m}c_hL^{-5(k-m)}\sup|B|\,w_k(n,x)$.
The factor two bounds its levelling factor on
$|\lambda|<2$. Hence every expanded leg other than
the first split leg retains its weight reserve.

Since $d_k(X^\uparrow)\le d_j(X)$, the base factor
$e^{d_k(X^\uparrow)/13}$ is paid by
$\kappa_+-\kappa\ge2c_E+1$. The count through a cube is
\[
 2048\sum_{m\ge0}
    \bigl(9216e^{-(2c_E+1)+1/13}\bigr)^m
       \le e^{-1/2}C_E,\qquad C_E=4096.
\]
With the edge pointer factor $(d_j(X)+3)^3$ the same
sum is at most $2048\cdot28\le e^{-1/2}94C_E$.
The two factors $e^{1/13}$ from the base and first
split leg fit in the spare $e^{1/2}$. For a mass edge
the base has size zero and $e^{2/13}<C_E$ suffices.
The coefficient variation is still estimated directly
by Lem.~\ref{C-399-48}, polymer by polymer at rate $r_V$.
The leg contributions use the count below, in the summed
norm at rate $\kappa-1\ge r_V$.

There are $L^{3(k-j)}$ lower cubes in a level-$k$
cube. A made term reaching layer $m>j$ has
$d_j(X)\ge8L^{m-j}-2$; an edge term has only
$7L-2$ in its first later layer. Use the uniform
lower bound $7L^{m-j}-2$. The complete count is
bounded by
$C_{\rm bd}^{\circ}
(c_\psi+2C_{\mathcal A}\delta(\mathsf t_k)/\beta_0)$,
where
\[
\begin{split}
 C_{\rm bd}^{\circ}={}&
 [9L^4A_\kappa C_H+(2+6L^5A_\kappa)C_{\mathfrak r}\beta_0
                         +4C_{\rm m}c_h\beta_0]\\
 &\times(1+376M^3)C_E
 \left(1+\sum_{d\ge0}(2/L)^d\right)
 \left(2+\sum_{e\ge1}L^{3e}
       e^{-\kappa(1-L^{-1})(7L^e-2)}\right),\\
 C_{\rm bd}={}&C_{\rm bd}^{\circ}
                    \max\{1,K_M^{-1},K^{\rm lb}/K_M\}.
\end{split}
\]
The polynomial $376M^3$ pays edge pointers, since
$\sum_{X\ni\square}(d_j(X)+3)^3
e^{-(\kappa_+-\kappa)d_j(X)}\le94C_E$.
All these constants precede $h_F$.
Put $\Sigma_k^\natural=\sup_{j<k}\mathsf{sz}_j^\natural$,
$\bar v_*=\sup_{j\le k}\bar v_j$, and
$\bar\eta=\sup_j\eta_{{\rm sr},j}<1/8$.
Reentry through explicit legs is bounded by
\begin{equation}\label{eq:qcd-boundary-contraction}
 C_{\rm bd}K_M(1+\bar\eta)
 \left(c_\psi+\frac{2C_{\mathcal A}\delta(\mathsf t_k)}{\beta_0}\right)
             \max\{\Sigma_k^\natural,\bar v_*\}.
\end{equation}
The mass-profile edges add
\[
 b_k^{\rm ms}
 =C_{\rm bd}\left(c_\psi+
       \frac{2C_{\mathcal A}\delta(\mathsf t_k)}{\beta_0}\right)
       3M^3w_0\nu C_H^2h_F^2\bar B_M\gamma^2.
\]
The regular terms require their zero moments and are
treated in the next proof. Once a new step has been
made, unchanged stored and formal legs carry the
smaller next weight, so their previously established
bounds persist.
\end{proof}

\begin{proposition}[Closing the matter induction on histories]
\label{C-399-21}
With the common numerical choices of \cref{C-396-12}, the actual
fermion-origin class at every history step satisfies
\[
 \mathsf{sz}_k^\natural\le\frac{16}{17}\varepsilon_\flat,
 \qquad
 \varepsilon_\flat=(64eL^3K_{\rm ap}^{\max})^{-1}.
\]
The class size is the supremum of its fibre norm on
the admitted tiers and external-field boxes, multiplied
by $e^{r_Vd_k(Y)}$, with the next weights at its
external and stored legs. The history step is defined
on these domains, has format \eqref{eq:qcd-history-format},
and obeys the exact identity of \cref{C-399-43}.
The bounds and locality of the regular and made
families persist through the finite induction.
\end{proposition}
\begin{proof}
We include earlier frame sizes among the induction
hypotheses. They are absent at $k=0$, where there
are no old activities or boundary terms and the
all-field inverse supplies the first step.
At a later step, complete the fermion square at
the new background. The determinant is the
normalizer of \cref{C-397-42}; the current generators
become $W$ plus their new background, and the
stored quadratic matrix becomes its exact Schur
complement. The remaining dependence on the gauge
fluctuation and $W$ forms the class whose size
we now estimate.

\emph{The terms that do not use a boundary.}
Put $\delta_k=\delta(\mathsf t_k)$. The small-field
vertex, the current mass member and the old regular
activities on one lattice give the bound
\[
\begin{split}
 b_k^{\rm blk}={}&C_Vh_F^2C_{\mathcal A}\delta_k
 +(c_\psi+C_GC_{\mathcal A}\delta_k)c_{\mathsf S}h_F^2
                                   \frac{2\bar B_M}{\mathsf t_k}\\
 &+(c_\psi+C_{\mathcal A}s_a+C_GC_{\mathcal A}\delta_k)
                C_\omega S_\omega C_2\phi(\mathsf t_k)^2.
\end{split}
\]
Here $s_a$ is the Cauchy shift factor in the
old-activity estimate; $S_\omega$ is its summed
scale weight. These are precisely the previous
small-field estimates read at the tier radius
$C_{\mathcal A}\delta_k$.
They control the current-region terms once the
old zero-moment cancellation is justified there.
The same vertex on stored layers and its
external collar dependence contribute at most
$2C_V^Fh_F^2C_{\mathcal A}\delta_k$.

\emph{Where an old coefficient can be evaluated.}
The tiers give a triple at level $t=k+1$. On
$X\subset\Lambda_\ell$, $j\le\ell\le k$, its
restriction to level $j$ has a fixed margin
$\tau_{\rm oc}$ inside the core radii, where
\[
 \tau_{\rm oc}
 =1-\frac{(1-\tau_{\rm m})\theta_{\rm oc}}{1+\beta_s},
 \qquad
 b_2=\frac{\tau_{\rm oc}\alpha_0}
             {18K^\sharp(1+c_\natural)}.
\]
For the local shift define
$|B|_X^{\rm sh}$ as the largest, over the bond
and difference-quotient places of $X$, of
$L^{-(m(x)-j)}$ times the exponentially weighted
supremum of $|B|$ in the shift estimate.
Then
\[
 |B|_X^{\rm sh}\le L^{-(\ell-j)}\sup|B|;
 \qquad
 |B|_X^{\rm sh}\le eL^{-160(k-m)}\sup|B|
\]
if all those places lie in layers at most $m<k$.
The restricted shift and its first difference and
source components scale with powers $1,2,3$,
respectively. Their sizes are thus at most
$|\mu|c_h|B|_X^{\rm sh}$.
When this is below $b_2$, the outer-substitution
estimate puts the shifted triple inside the
level-$j$ core. Its pair is
\[
 \left(e^{i\eta_k\mu\mathbf A}\mathbf U,\,
       L^{-3(t-j)}\mathbf J+
                         \mu L^{-3(k-j)}\mathbf K\right).
\]
Indeed the unshifted triple lies at radii
$(1-\tau_{\rm oc})\alpha_i$, and the
outer-substitution bound consumes less than
$\tau_{\rm oc}\alpha_i$. The admitted constants give
$128K^\sharp(1+c_\natural)c_{\rm os}\le1$ and
$128\tau_{\rm oc}>23>18$, so
$c_{\rm os}\alpha_0<b_2$. Hence the Cauchy disc
$|\mu|c_h\sup|B|\le L^{\ell-j}c_{\rm os}\alpha_0$
lies in the domain, including lag zero. This
verifies coefficient holomorphy and its uniform
bound at every shifted pair used below.

\emph{Zero moments in the current region.}
The cancellation is applied to local pieces,
not to an unchanged whole-volume background.
A piece $p$ consists of a polymer and one
local or walk term in each of its two legs;
let $Z_p$ contain those three supports.
Let $m_p$ be its trivial-background moment for
the empty history, and $m_p^h$ its moment along
the given history. A missing piece has moment
zero. If $Z_p\subset\Lambda_j^{\sim-1}$, all
pointed shares of the polymer are present and
both local background maps agree there; hence
$m_p^h=m_p$. Every other piece is a rim piece.
Thus, at a cell $c$, the moment left by the
history is the sum of $m_p^h-m_p$ over those
rim pieces. On a torus it also contains the
wrapping sum minus the nonfitting infinite-lattice
sum already treated in \cref{C-388-24}.

A different rule decides which terms may use
differences of legs. At lags $k-j\ge n_{\rm lag}$,
take a piece whole if $Z_p\not\subset\Omega_k$,
or if $Z_p\not\subset\Omega_{k+1}$ and
$d_j(Z_p)\ge8L^{k-j}-2$; otherwise use the
difference decomposition of \cref{C-388-24}.
The first rule, using $\Lambda_j^{\sim-1}$,
decides whether its number is unchanged. The
second, using the two current regions, decides
whether its local estimates have the needed
regular scale.

Every piece taken with differences is unchanged.
Its support and the local gauge cube used in
\cref{C-388-24} lie where the smoothed scale is
$L^k$. They lie within distance $8ML^k$ of
$\Omega_{k+1}$ and hence in $\Lambda_k$, since
the removed cell layer has width at least
$12ML^k$. The core and shifted-core checks
above therefore justify the old regular estimate
there. The H\"older difference of its background
legs supplies the same summable scale factor.
The terms with a split stored leg additionally
carry the factor below denoted $\Theta'_k$.

To check the algebra, sum only the pieces taken
with differences. Their constant terms are
$\bar\Phi(z_c)(\sum'_p m_p)\Phi(z_c)$.
The sum of all empty-history moments is zero
on the infinite lattice, or the wrapping-minus-
nonfitting sum on a torus. Subtracting the
moments of the pieces taken whole leaves exactly
the history rim moment above. Each such term
is filed on the level-$k$ hull of $Z_p$,
enlarged by the walks of its expanded legs.
That hull meets the complement of
$\Omega_{k+1}$, so no interior coefficient
acquires a dependence on the history.
Wrapping terms and the corresponding subtracted
empty-history terms are filed on the same
wrapping support and cancel there when both
are present.

For these whole terms, the two history legs
have relative sizes at most $L$ each, and at
most $L^{1-(k-j)}$ if $Z_p\subset\Omega_k$.
The empty-history comparison legs have size
$L^{-(k-j)}$. Count a term through every
cell of $Z_p$, not only its original filing
cell. A walk of length $n$ reaches at most
$(2n+5)^3$ relevant cubes; put
$n_{\rm wk}=\max_{u\ge0}(2u+5)^3e^{-u}$ and
$c_{\rm geo}=9n_{\rm wk}^2$.
The two walk counts pay these polynomials
with their exponential weights. The hull
weight leaves
$e^{-\kappa(1-L^{-1})d_j(Z_p)}$.
A piece leaving $\Omega_k$ has
$d_j(Z_p)\ge8L^{k+1-j}-2$; every other
whole piece has $d_j(Z_p)\ge8L^{k-j}-2$.
The factors $L^{3(k-j)}$ for cells and the
two leg sizes are consequently summable.
Their bound is
$\Theta_k C_\omega c_{\rm rim}C_2\phi(\mathsf t_k)^2$,
where
\[
 c_{\rm rim}=c_{\rm geo}(1+L^2)
       \sum_{d\ge1}L^{3d}
                      e^{-\kappa(1-L^{-1})(8L^d-2)}.
\]

\emph{Regular terms in stored layers.}
At a pointer in $\mathcal L_m$, retain the
formal lower-level legs with the weights of
that layer. For large $m-j$, pieces contained
in $\Omega_{m-1}\setminus\Omega_{m+2}$ use
the difference estimate at the local scale;
the remaining pieces are taken whole.
The former are in a regular core at some
$\check m\ge j+2$, so the shifted-domain
argument applies with $\ell=\check m$.
The lower lags are bounded directly by
their finite count. For degree at least
four, the extra leg weights supply the
power counting; for degree two the
H\"older difference is essential.

There is also a cell term in the zeroth
moment. Its coefficient $q_j(z)$ has
$\sum_{z\in c}q_j(z)=0$ and commutes with
colour. Transport the two fields inside
$c$ to $z_c$ by block links and their
inverses. Then
\[
\begin{split}
 \sum_{z\in c}\bar\Phi(z)q_j(z)\Phi(z)
 =\sum_{z\in c}\big[
  (\bar\Phi(z)-\bar\Phi(z_c))q_j(z)\Phi(z)
 +\bar\Phi(z_c)q_j(z)(\Phi(z)-\Phi(z_c))\big].
\end{split}
\]
Each term has one difference. The block-link
H\"older estimate applies directly; no
replacement by straight transport is made.
The two transports have length at most
$2M$ each. Their nonunitary parts cost
at most $\exp(6MC_Q\alpha_1L^{-(m-j)})
\le e^\kappa$ under the chosen radius
ceiling and $m-j\ge n_{\rm lag}\ge4$.
The term is filed with the cell and the
expanded leg and shift supports, hence
remains a boundary term.

For a whole piece based at a stored pointer,
reaching a later layer $m'>j$ forces
$d_j(Z_p)\ge7L^{m'-j}-2$; reaching
$\Omega_{k+1}$ gives the stronger
$8L^{k+1-j}-2$. Use the weighted walk
count with $e^{3|\omega|/2}$: one factor
$e^{|\omega|}$ pays its reachable cubes,
and the remainder pays crossing. This
still has constant $A_\kappa$, since
$e^{1/8}/(1-e^{-11/8})<1.517$ and the
local term fits its defining factor two.
The sum over pointers of one polymer
is bounded by its original insertion norm.
Thus the bound for these terms is
$\Theta_k C_\omega c_{\rm rim}^{\rm pt}
C_2\phi(\mathsf t_k)^2$, with
\[
 c_{\rm rim}^{\rm pt}=
 \frac{3L}{L-2}c_{\rm geo}(1+L^2)
 \left(1+\sum_{d\ge1}L^{3d}
                 e^{-\kappa(1-L^{-1})(7L^d-2)}\right).
\]
Wrapping and nonfitting pieces obey the same
count because their size is at least
$L^{k-j}-2$.

Finally, a pointed share whose pointer and
filing cell lie on opposite sides of
$\partial\Omega_{k+1}$ is taken whole.
In a level-$k$ cube at most
$6(t+1)L^{2(k-j)}$ lower cells lie within
distance $(t+1)M$ of those boundary faces.
The residual exponential and this count
give at most $39$ times the birth norm.
The walks of two legs contribute
$1+2n_{\rm wk}$. For higher degree, the
remaining legs pay their walk count;
at lag zero use
$2p^2(3/4+\Theta'_k/4)^{2p-1}<4.3$
under $\Theta'_k\le2\cdot10^{-3}$.
This gives
$39(1+2n_{\rm wk})\Theta_k S'_\omega
C_2\phi(\mathsf t_k)^2$, where
$S'_\omega=\sum_{d\ge0}L^{-d}(1+d\log L)^{2\mathfrak p_0}$.
These shares are filed with their whole supports,
so their domains meet $Z_{k+1}$.

To collect the estimates, set
\[
\begin{split}
 \Theta'_k={}&4Lc_\psi+
 \frac{2C_{\mathcal A}\delta_k}{\beta_0}
 [9L^4A_\kappa C_H+(2+6L^5A_\kappa)C_{\mathfrak r}\beta_0],\\
 \Theta_k={}&c_\psi+C_{\mathcal A}s_a+
                           C_GC_{\mathcal A}\delta_k+\Theta'_k.
\end{split}
\]
The stored-layer difference and finite-lag counts
are bounded by a constant $C_{\rm lay}$ fixed
before $h_F$. More explicitly, one can take
\[
 C_{\rm lay}
 =4M^3A_\kappa L^{1+\alpha}C(\alpha)C_H
       \max\{1,C_{\rm leg}^{\square},C_\omega\}\,c_{\rm ps},
 \quad
 c_{\rm ps}=(2A_\kappa C_H)^2c_{\rm geo}
                                  L^{3n_{\rm lag}+12}.
\]
The finite-lag factor pays those lags and the
extra local or walk terms; the H\"older factor
is needed only in the remaining degree-two
terms. The cell estimate above is included
in $C_{\rm leg}^{\square}$.
Since $L^{3(k-m)}(2L^{-4})^{k-m}=(2/L)^{k-m}$,
the stored-layer sum has the bound
\[
\begin{split}
 b_k^{\rm lay}={}&2C_V^Fh_F^2C_{\mathcal A}\delta_k\\
 &+\Theta'_k
 \left(C_\omega S_\omega+\frac{L}{L-2}C_{\rm lay}S'_\omega\right)
                                  C_2\phi(\mathsf t_k)^2\\
 &+\Theta_k\left[
 C_\omega(c_{\rm rim}+c_{\rm rim}^{\rm pt})
              +39(1+2n_{\rm wk})S'_\omega\right]
                                  C_2\phi(\mathsf t_k)^2\\
 &+\frac L{L-2}C_{\rm lay}\Theta'_k
                     w_0\nu C_H^2h_F^2L\bar B_M\gamma^2 .
\end{split}
\]
Together with the mass-edge bound $b_k^{\rm ms}$
from the preceding proof and $b_k^{\rm blk}$,
this bounds every regular term, vertex and
mass edge. Each decreases to zero with
$\gamma$ after the earlier constants and
$h_F$ have been fixed.

\emph{The two boundary contributions.}
The old made and edge terms reenter through
their explicit legs by
\eqref{eq:qcd-boundary-contraction}.
The floor $h_F\ge32C_{\rm bd}K_MA_\kappa\Gamma_F$
and ceiling
$32C_{\rm bd}K_MC_{\mathcal A}\delta_k\le\beta_0$
make this at most
$\frac9{64}\max\{\Sigma_k^\natural,\bar v_*\}$.
Their coefficient variation, \cref{C-399-48},
has bound
\[
 C^{\rm bd}\vartheta_{0,k}
  \left(\tfrac98K^{\rm lb}\Sigma_k^\natural
                       +108M^3K_M\bar v_*\right).
\]
The factor $108M^3$ is the edge pointer bound
$4M^3(d+3)^3e^{-(\kappa_+-\kappa)d}\le108M^3$.
The separate ceiling
$16C^{\rm bd}\vartheta_{0,k}
(\frac98K^{\rm lb}+108M^3K_M)\le1$
makes this at most
$\frac1{16}\max\{\Sigma_k^\natural,\bar v_*\}$.

Putting $b_k^{\rm tot}=b_k^{\rm blk}+b_k^{\rm lay}
+b_k^{\rm ms}$, we have proved
\[
 \mathsf{sz}_k^\natural
 \le\frac{13}{64}\max\{\Sigma_k^\natural,\bar v_*\}
                       +b_k^{\rm tot}.
\]
At the first step the boundary maximum is absent.
Let $\bar b=\sup_k b_k^{\rm tot}$. Since
$\bar v_j\le b_{j-1}^{\rm blk}\le\bar b$,
induction gives
$\mathsf{sz}_k^\natural\le64\bar b/51$.
The three ceilings
$b^{\rm blk}(\gamma),b^{\rm lay}(\gamma),
b^{\rm ms}(\gamma)\le\varepsilon_\flat/4$
therefore give
$64\bar b/51\le16\varepsilon_\flat/17$.
This proves the present-frame size from earlier
sizes before using \cref{C-399-32}.

The scalar entries satisfy
\cref{C-399-47}, including the boundary coefficient
variations just checked, and their q and tail
classes meet \cref{C-399-44}. The actual joined
frame is consequently admissible. Its interior
outputs agree by \cref{C-399-45}; extract their
vacuum, action and mass terms there alone,
and file all profile complements on their
rim supports. The other outputs, the
normalizer outside the interior, and the
collar correction are the boundary families
of \cref{C-399-15}. The coordinatewise identity
of \cref{C-399-43} proves the new format and
the exact step. The stored quadratic factor
is bounded by \cref{C-399-17}. This closes
the finite induction and preserves every
local domain used in it.
\end{proof}

\subsection{Vacuum constants and history masses}

Let $e_j$ be the full vacuum constant per site extracted by the empty
history at level $j$, and set $E_k=\sum_{j\le k}e_jn_j$.
In particular it includes both the gauge and fermion
normalizers and the value of the q Gaussian family.
The two gauge Gaussian conventions differ by
\[
 3(L^3-1)\sum_{j<k}n_{j+1}\log(\mathsf t_j/x_j).
\]
This number is retained in $E_k$. It is not a bounded error per
level-$k$ site.
Indeed a block fluctuation has $9(L^3-1)$ real
coordinates: $3(L^3-1)$ gauge coordinates and
$6(L^3-1)$ transverse coordinates. At the trivial
background the action Hessian vanishes on the former,
where the gauge-fixing form is $|\lambda|^2$.
Scaling the action by $\tau$ therefore multiplies
the determinant of the full form by
$\tau^{6(L^3-1)}$. Equivalently the covariance
normalizer contributes $\tau^{-9(L^3-1)/2}$ and
the q-Gaussian contributes $\tau^{3(L^3-1)/2}$.
Their product is $\tau^{-3(L^3-1)}$, which gives
the displayed vacuum correction. The collar term
$R$ vanishes at the trivial zero stored field.

\begin{lemma}[The vacuum constant of a label]\label{C-399-24}
There are label-dependent complex numbers $f_j(\mathcal S)$ such
that
\[
 -\mathcal E^F_{k,\mathcal S}=E_k-\sum_{j\le k}f_j(\mathcal S),
 \qquad
 |f_j(\mathcal S)|\le C_\flat^F(1+\log\mathsf t_{j-1})
             \#_j(Z_j^\sim).
\]
The empty-history finite-size remainder is bounded by a
constant $C_{\rm fs}$ independent of $K$ and $\ell$.
It is assigned to the regular scalar sum, with its opposite
subtracted from the extracted vacuum constant. Thus $E_k$
contains the infinite-lattice per-site numbers $e_jn_j$.
\end{lemma}
\begin{proof}
Define $f_j$ as the number extracted by the empty history on
the same torus minus the number extracted by the given history.
Write $\mathfrak f_k=\sum_{j\le k}\mathfrak f^{(j)}$
for the empty-history finite-size number. With the stated
assignment the extracted constant is
\[
 \sum_{j\le k}(e_jn_j+\mathfrak f^{(j)}-f_j)
                     -\mathfrak f_k=E_k-\sum_{j\le k}f_j.
\]
This proves the exact identity. The partial pointer and domain
sums defining $f_j$ need not be real; the conjugation identity
holds for the full summed energy. The difference has
three parts. The missing gauge integral costs its vacuum
number, bounded by $C(1+\log\mathsf t_{j-1})$, at each cube
of $Z_j$. In the fermion normalizer write the trace as a sum
of diagonal place contributions. At a place inside $\Omega_j$
its difference from the empty-history contribution requires a
walk reaching $\Omega_j^c$ and is at most
$Ce^{-c_f\operatorname{dist}_j(p,\Omega_j^c)}$.
The missing places outside $\Omega_j$ have a uniformly bounded
free fermion number. Both sums cost $C\#_j(Z_j^\sim)$.
Finally a missing regular pair through a pointer $z$ must meet
$Z_j$; the activity and entropy bounds give
$CE_0e^{-(\kappa+1-\log9216)
(\operatorname{dist}_j(z,Z_j)/M-2)}$.
Summation over $z$ again costs $C\#_j(Z_j^\sim)$.

A wrapping polymer or trace walk has size at least
$\mathrm{side}_j/(2M)-2$. Hence the finite-size number at that
level is at most
$CE_0\mathrm{side}_j^3
e^{-(\kappa+1-\log9216)(\mathrm{side}_j/(2M)-2)}$.
The sides run through distinct elements of $2L^{\mathbb N}$,
so the sum is finite. All these numbers depend on the label,
not on an evaluation configuration.
\end{proof}

We record the pure history estimate in the form used here.
Its hypotheses include the admitted calibration geometry,
the failed-test support and curvature conditions, the pivot
determinant identities and the retained quadratic lower bound
on each stored piece. Suppose a density, scalar or normed,
is bounded on its realized free coordinates by
\[
 e^{\operatorname{Re}(E-\sum_j f_j)+C_E\Sigma_{\mathcal L}-A(\mathsf{bw},U)}
 \prod_{(t,D)}|\mathsf u_{t,D}[Q_{t,D}]|,
\]
where $\mathsf{bw}\ge(1-a)\beta$, $0\le a<1$, and the
piece forms have the corresponding positive retained lower
bound. Here $\mathsf u_{t,D}$ and $Q_{t,D}$ are the weight and
retained quadratic form of a stored component $D$ created at level
$t$. The measure $d\mathscr M_h^{\rm fr}$ is the product of Haar
probability measures on the free configuration coordinates and
Lebesgue measures on the free Lie-algebra coordinates, including
the current field, as in \cite[Equation~(984)]{YMM1}.
The count $n_t^{\rm cr}$
is the number of cells created at that level; $g_t$ and $p_0(g_t)$
are its coupling and test threshold in the cited gauge construction.
We will also use its nonnegative history weight
\[
 \Theta^{(k)}(h)=w\,\mathrm{Dep}_k(h)+c_a\mathrm{Vol}_k(h),
 \qquad
 \mathrm{Dep}_k(h)=c_{\rm Dep}\sum_t p_0(g_t)^2n_t^{\rm cr}(h).
\]
Here $0\le w<1$, $c_a,C_{\rm vol}\ge0$ and
$0<c_{\rm Dep}\le c_*$ are fixed. The volume functional is
$\mathrm{Vol}_k(h)=C_{\rm vol}\sum_j\check R_j\#_j(Z_j(h))$,
where $\check R_j$ is the coarse-cell radius of the gauge history
at level $j$. Thus the coefficient of a created-cell deposit in
$\Theta^{(k)}$ is $c_\Theta=wc_{\rm Dep}<c_*$.
The weighted history sum uses the same component estimates,
with part of their decay reserved for this weight; see
\cite[Equations~(1455)--(1456) and their preceding definitions,
p.~592]{YMM1}.
The component majorant and its product integral
\cite[Props.~66.2 and 66.3]{YMM1} give
\begin{equation}\label{eq:qcd-history-interface}
 e^{-\operatorname{Re}E+\operatorname{Re}\sum_jf_j}
 \int\|g_h\|\,d\mathscr M_h^{\rm fr}
 \le e^{C_{\rm maj}(n_k+\Sigma_Z)}
 \prod_t e^{-(1-a)c_*p_0(g_t)^2n_t^{\rm cr}}.
\end{equation}
This is the cited component estimate, applied to the displayed
density: its proof uses the created-cell action allocation,
the piece pivot determinants, the positive piece-form bound
and Tonelli on the disjoint free coordinates. If also
$|f_j|\le C_\flat(1+\log\mathsf t_{j-1})\#_j(Z_j^\sim)$,
then for $2\mathfrak p_0>3r_c+2$ and sufficiently large final
threshold, the history summation theorem gives
\begin{equation}\label{eq:qcd-history-sum-interface}
 \sum_h e^{-\operatorname{Re}E}\int\|g_h\|
          \,d\mathscr M_h^{\rm fr}\le e^{(C_{\rm maj}+2)n_k}.
\end{equation}
The threshold is formed from
$((1-a)c_*,C_{\rm maj},C_\flat,\eth)$, $0<\eth<1$, before
the cutoff and volume. These are the full inputs to the pure
summation interface; the density estimate remains to be
verified for the present model.

\begin{theorem}[The upper history bound]\label{C-399-3}
Every history has the form \eqref{eq:qcd-history-format}.
Its normed density satisfies
\[
 \|g_h\|\le
 e^{-\beta A(U)+C_E^F\Sigma'(h)+E_k-\operatorname{Re}\sum_jf_j}
 \prod_{(t,D)}|\mathsf u_{t,D}[Q_{t,D}]|,
\]
and
\[
 e^{-E_k+\operatorname{Re}\sum_jf_j}\int\|g_h\|\,d\mathscr M_h^{\rm fr}
 \le e^{C_{\rm maj}^\dagger(n_k+\Sigma_Z)}
 \prod_t e^{-c_*p_0(g_t)^2n_t^{\rm cr}},
 \quad C_{\rm maj}^\dagger=48L^3C_E^F+C_*'.
\]
Its Berezin integral is a positive measure, and the sum of
their masses is the partition function.
\end{theorem}
\begin{proof}
The Schur identity, the local conditional expansion and the
boundary reentry just proved give the format by induction.
Gauge changes, coordinate pushforwards and multiplication by
test factors act coefficientwise in the finite-dimensional
Grassmann algebra. Joint gauge invariance gives one value on
each collared evaluation orbit. The exact bare determinant
is nonnegative, so each history's Berezin mass is nonnegative.

The degree-zero regular and boundary sums, together with the
coupling-profile shifts, are bounded by
$C_5^F\Sigma_{\mathcal L}$. To see the profile bound, a
level-$j$ cube contains $M^3L^{3j}$ fine sites, its regular
plaquette action is at most $C_{\rm reg}^2L^{-4j}$, and its
coefficient is $L^jb_j$; its cost is therefore a constant
per cube. A boundary domain has a cube in $Z_j^\sim$, so the
entropy bound and
$\sum_j\#_j(Z_j^\sim)\le C_Z\Sigma_{\mathcal L}$ give its
contribution. The wrapping constant is covered by
$C_{\rm fs}\Sigma_{\mathcal L}$.

For completeness the same count in all degrees is finite.
Each background leg from level $j$ to a layer $n$ costs at most
$CL^{1-(n-j)}$. Degree at least four gives net decay
$L^{-(n-j)}$ after the cube count, with fixed powers of $L$
in the constant. In degree two, sum the gauge-comparison
bounds over the walks at all rows of $X$, as in the proof
of \cref{C-388-24}. The possible first cubes number at most
$N_X=C_{\rm nb}M^3\#X$. For each leg the sum of its
weighted bounds is at most
\[
 2C_HL^{1-(n-j)}(1+N_XA_\kappa).
\]
The two factors $1+N_XA_\kappa$, the squared path length
and the complex transports cost at most a fixed multiple of
\[
 (1+8C_{\rm nb}M^3A_\kappa)^2
 2e^{\kappa/8}M^2\sup_{d\ge0}(1+d)^4e^{-7\kappa d/8}.
\]
The walk weights pay the remaining length polynomials.
With $L^{3(n-j)}$ pointers per level-$n$ cube, the scale
factors are
\[
 L^{3(n-j)}L^{2-2(n-j)}L^{2-2(n-j)}=L^4L^{-(n-j)}.
\]
These fixed factors enter the constants before $h_F$ is
chosen. At the trivial field group the kernel by its total
zero moment and use differences of legs; their H\"older
exponent $3/4$ gives a possible factor $L^{(n-j)/4}$.
The birth size is at most
$CL^{-(n-j)}(1+(n-j)\log L)^{2\mathfrak p_0}
\phi(x_{n-1})^2$, leaving a summable $L^{-3(n-j)/4}$.
A missing cancellation near a rim requires a polymer spanning
at least $RL^{n-j-1}$ and is covered by its exponential
decay. The mass profile costs
$w_0\nu C_H^2h_F^2L\bar B_M\Sigma_{\mathcal L}$.
Reading formal legs through the whole real background costs
at most one per leg by the enlarged row bound. Together with
Lem.~\ref{C-399-17} these estimates give a finite $C_E^F$.

The density estimate uses only degree zero and the stored
factor. Before the store is integrated, its scalar coefficient
is the positive product of the successive determinants.
Equivalently it is
$F_h=\det(D+m_0)/\det\mathbb D_h>0$; the Schur identities
give this formula, including at arbitrary gauge fields.
Taking degree zero of the format gives the scalar history
density for $e^{-\beta A}F_h$. Its exponent has the scalar
bound just established. Multiply its density by the
stored-factor norm bound. The same realized Borel coordinates
and piece forms therefore verify every hypothesis of
\eqref{eq:qcd-history-interface}, with loss zero and the
constant per cube at most $48C_E^F$.
This proves both displayed inequalities.
\end{proof}

\begin{theorem}[One vacuum energy in both bounds]\label{C-399-4}
With a sufficiently large final threshold,
\[
 e^{E_k-C_P(1+\log x_k)n_k}\le Z_{K,\ell}
 \le\sum_h\int\|g_h\|\,d\mathscr M_h^{\rm fr}
 \le e^{E_k+(C+2)n_k}.
\]
The constants are uniform in the cutoff and volume; $C_P$
may depend on the positive lower mass bound.
\end{theorem}
\begin{proof}
The upper bound is \eqref{eq:qcd-history-sum-interface}, with
$E=E_k$ and the label error from Lem.~\ref{C-399-24}.
For the lower bound retain the empty history at the last
level. Its mass is not estimated by the free determinant
alone: its remaining interacting Grassmann integral must
also be bounded below.

On the last level read the old activities at weight
$2A_\kappa h_*$ and the coarse generators at weight $2h_*$,
where $h_*$ is a coarse Gram weight depending on $m_-$.
Each leg read through its local part and walk terms costs
at most $A_\kappa$, which the activity weight pays.
Impose the final ceiling $\gamma\le\gamma(2A_\kappa h_*,5B_s)$.
Regroup their trivial-background degree-two terms at their
pointers, so that the actual zero-moment cancellation occurs
inside each regrouped polymer. The power counting above gives
\[
 \|P\|_{2h_*}\le C_*\phi(\mathsf t_{K-1})^2.
\]
Include the mass member $-M_{h_\emptyset}$ as well as the
bilinear part of $P$ in the perturbation $p$. The former has
operator norm at most $C_{\mathsf S_1}\bar B_M\gamma^2$;
impose $C_{\mathsf S_1}\bar B_M\gamma^2\le c_bm_-/8$.
The activity bilinear is small by the displayed bound after
the final floor is increased. Thus $\|p\|\le c_bm_-/2$. Hence
$G_*=D_K+p$ has real part at least $c_bm_-/2$, and its
inverse has a uniform exponentially decaying Gram bound.
Apply the finite Grassmann forest expansion to the remaining
terms of degree at least four, between fixed coarse cubes.
The interpolation matrix is positive semidefinite and
preserves the covariance's Gram bound; the derivative estimate
$\|\partial F\|_{h_*}\le h_*^{-1}\|F\|_{2h_*}$ and the
line bound $2e^{-\alpha_*\operatorname{dist}}$ give convergence
when $\|P\|_{2h_*}C(M_*,\nu)$ is small. Its connected sum
has modulus at most $\ell^3$. The original integral is
nonnegative, and its convergent logarithm is real by
conjugation. It is therefore at least
$(c_bm_-/2)^{\nu\ell^3}e^{-\ell^3}$.

Restrict the coarse links to the ball of radius $x_K^{-1/2}$.
Its action costs at most $24(B_3c_r)^2\ell^3$, and the Haar
ball estimate gives volume at least
$\exp[-(\frac92\log x_K+3\log7)\ell^3]$.
The scalar regular-family bound now gives the last-level
empty-history mass at least
$e^{E_K-C_P(1+\log x_K)\ell^3}$.

The positive exact history decomposition shows that the
empty-history mass at level $k$ dominates that at level $K$.
Finally $|e_j|\le C_e'(1+\log x_{j-1})$ and
$\sum_{j>k}n_j\le n_k/(L^3-1)$ give
$|E_K-E_k|\le C_e'(1+\log x_k)n_k/(L^3-1)$.
Enlarging $C_P$ proves the asserted lower bound with the
same $E_k$ as in the upper bound.
\end{proof}

\section{The source chain and local source energies}\label{sec:qcd-source-chain}

To control the correlations we insert sources before integrating any
fields.  A source for the action density changes the plaquette weight;
a source for a fermion bilinear changes the fine mass matrix.  The
second change also produces a relevant matrix mass at later scales.
We first control this induced mass and the energy extracted from each
local step.  We then prove that the same energy normalizes the sum of
all histories.  This distinction matters at complex sources, where a
history density need not be positive.

The resulting correlation estimate is Thm.~\ref{C-407-6} below.
Its constants are independent of both the cutoff and the torus side.
The estimate uses locality of the source energy and reflection
positivity: the former makes separated sources almost additive, and
the latter converts a whole-volume bound into a bound for one cell.
Throughout this section we use the ordered choices of constants
and the pure gauge construction verified in
Prop.~\ref{prop:qcd-joint-admission}.

The sources are $s=(w,J)$, with $w$ complex Lipschitz and $J$
a bounded field in
\[
 \mathbb M=\operatorname{End}(\mathbb C^2_{\rm s}\otimes
 \mathbb C^2_\tau\otimes\mathbb C^{N_f})\otimes1_{\rm c}.
\]
The physical bilinear sources lie in the specified span of
$\mathfrak B$, but the intermediate matrix source domain is
the entire displayed space. Put $\phi_j=\phi(x_j)$ and choose
$n_\pi\ge1$ with $(1+1/n_\pi)^{\mathfrak p_0}\le4/3$.
Set
\[
 \pi_j=(1+(K-j)/n_\pi)^{\mathfrak p_0},\quad
 \|w\|_{\sharp,j}=\max\{\|w\|_\infty,
                     40M\lambda_j\pi_j\operatorname{Lip}w\}.
\]
The fixed physical source norm is
$\|w\|_\sharp=\max\{\|w\|_\infty,40M\operatorname{Lip}w\}
=\|w\|_{\sharp,K}$.
In fine lattice coordinates let $B_k(x)$ be the ball of
radius $ML^k$ centered at $x$. The corresponding local norm, in physical coordinates, is
\[
 \|w\|_{\sharp,k;x}=\max\left\{
 \sup_{a_KB_k(x)}|w|,\,
 40M\lambda_k\pi_k\operatorname{Lip}(w|_{a_KB_k(x)})\right\}.
\]
The geometric estimates used below are
\begin{equation}\label{eq:qcd-source-geometry}
 \phi_j/\phi_{j+1}\le q_\phi=2L^{-1/2},\quad
 \sum_{i\le j}\phi_i\le\phi_j/(1-q_\phi),\quad
 \pi_j/\pi_{j+1}\le4/3.
\end{equation}
We write
$S_\uparrow=\sup_k\sum_{j\le k}\phi_j/\phi_k
\le(1-q_\phi)^{-1}$.
They follow from $|x_j-L^{i-j}x_i|\le\bar B$ and the
decreasing range of $\phi$, after increasing the final floor.

At pair radius $a$ define
\[
 \tau_j(w;a)=a x_j^{1/2}40M\lambda_j\operatorname{Lip}w/\mathsf b,
 \quad a_j^{\rm rad}=C_{\rm rad}\delta(\mathsf t_j),\quad
 \mathsf s_j(a)=C_{\mathcal A}\delta(\mathsf t_j)/a.
\]
For $i\le j$,
\begin{equation}\label{eq:qcd-source-radius-lag}
 \tau_i(w;a)\le\sqrt2L^{-(j-i)/2}\tau_j(w;a),\qquad
 \mathsf s_j(a)\le C_{\mathcal A}/C_{\rm rad}\quad(a\ge a_j^{\rm rad}).
\end{equation}
The first estimate uses $x_i\le2L^{j-i}x_j$ and
$\lambda_i=L^{-(j-i)}\lambda_j$.
Since $\widehat\phi(x)=x^{1/2}\phi(x)$ obeys
$\widehat\phi(x_j)\le c_\pi\widehat\phi(2\mathsf b)\pi_j$,
\[
 \tau_j(w;a_j^{\rm rad})\le
 8c_{\rm th}c_\pi C_{\rm rad}A_1\varrho
                       \widehat\phi(2\mathsf b)/\mathsf b
 \quad(\|w\|_{\sharp,j}<8\varrho).
\]
This is uniformly small at the final floor. At fixed radius
$\alpha_0$ it is uniformly small on the fixed ball
$\|w\|_\sharp<8\varrho$, where it is at most
$16\alpha_0\varrho\phi_j$. It need not be uniformly bounded
on the larger level-dependent ball at that fixed radius.

Fix a small operator radius $r_M>0$, put $r_J=r_M/L$, and let
\[
 \mathscr O_j=\{\|w\|_{\sharp,j}<4\varrho,
                    \lambda_j\|J\|_\infty<r_J\},\quad
 \Xi_j(s)=\lambda_j\left(
 \pi_j\operatorname{Lip}w/\varrho+L\|J\|_\infty/r_M\right).
\]
Then $\Xi_j<2$, $\mathscr O_j\subset\mathscr O_i$ for $i\le j$,
and $L\Xi_i\le(4/3)\Xi_{i+1}$.
The coordinates of a step are
$\mathfrak c=(w,\mathsf n,(\sigma_i)_{i\le j})$ with
$\mathsf n=-a_KJ$. Their domain is
\[
 \mathscr D_j=\{\|w\|_{\sharp,j}<4\varrho,
 L^{j+1}\|\mathsf n\|<r_M,
 \|\sigma_i\|<B_\sigma\phi_i\ (i\le j)\}.
\]
At pair radius $a$ replace the condition on $w$ by
$\|w\|_\infty<4\varrho$ and
$C_{\rm lin}\tau_j(w;a)<c_{\rm cl}/16$; call this
$\mathscr D_j(a)$. A doubled domain doubles each radius.
The final ceilings give
$2\mathscr D_j\subset2\mathscr D_j(a_j^{\rm rad})$.

On $2\mathscr D_j(a)$ the operator source satisfies
$L^{j+1}\|\mathsf n\|<2r_M$. The weighted inverse and walk
construction therefore applies through $\mathcal K_{j+1}$
and $\Gamma_j$, uniformly on the full matrix ball, by
Lem.~\ref{C-397-29}. In particular it gives holomorphic
localized covariances with their Gram bound and backgrounds
with the same localized row and H\"older estimates, at the
fixed enlarged constants. No positivity of
$m_0+\mathsf n$ is used on $\mathscr O_j$.

\subsection{The source step and its cell units}

For a matrix kernel $k$ define its symmetric zeroth moment at $y$ by
\[
 \mathrm{ord}_0[k](y)=\tfrac12\sum_{y'}[k(y,y')+k(y',y)].
\]
It is the entire filed cell unit, rather than each separate
matrix entry, whose total is canceled in the extraction.

\begin{proposition}[One step at the sources]\label{C-407-16}
Suppose the old data are defined on $2\mathscr D_i(a)$,
$i\le k$, where $a_k^{\rm rad}\le a\le\alpha_0$.
Assume, polymer by polymer and at each pointer, that they read
the pair only on its core, are jointly gauge invariant, and
are holomorphic in the pair and the independent source coordinates
on that core. Their scalar parts are strictly gauge invariant.
Assume also that the families, including their pointed scalar
parts, agree on common domains when the pair radius is reduced.
For the parts with legs, sum the norms of the unextracted
activity, the mass member and the individual moved forms;
assume this sum is at most $2C_w^{\rm s}\phi_{i-1}$.
After subtracting the explicit normalizers and the Gaussian
family, take the maximum of the unpointed scalar fluctuation
norm and all its pointed norms. Assume this maximum is at most
$2C_w^{\rm s}\phi_{i-1}+C'_{\rm lin}\tau_{i-1}(w;a)$.
These bounds hold at every admitted reduced radius.
Then the next source step is defined and has these bounds
at level $k+1$. Its data are holomorphic in the pair and
coordinates, and agree when restricted to a smaller
admitted pair radius.

Let $k_X(\mathfrak c)$ be its trivial-background degree-two
kernel before local extraction and let $\mathfrak s_{k+1}$
be the constant-field mass normalization. Its local increment is
\begin{equation}\label{eq:qcd-source-increment}
 \nu_{k+1}(y;\mathfrak c)=\mu_{k+1}
 -\mathfrak s_{k+1}^{-1}\sum_X
 \bigl(\mathrm{ord}_0[k_X(\mathfrak c)](y)
                 -\mathrm{ord}_0[k_X(0)](y)\bigr),
\end{equation}
where the sum runs over polymers of diameter less than a quarter
of the torus side. It is holomorphic and, for
$\mathfrak c\in\mathscr D_k$, $\mathfrak c'\in2\mathscr D_k$,
\begin{equation}\label{eq:qcd-increment-locality}
 |\nu_{k+1}(\mathfrak c)(y)-\nu_{k+1}(\mathfrak c')(y)|
 \le B_\nu\phi_{k+1}\|\mathfrak c-\mathfrak c'\|_{k,y},
\end{equation}
where
\[
 \|\mathfrak c-\mathfrak c'\|_{k,y}=
 \sup_xe^{-\kappa_\sigma|x-y|/L^{k+1}}
 \max\left\{\frac{\|w-w'\|_{\sharp,k;x}}\varrho,
 \frac{L^{k+1}|\mathsf n-\mathsf n'|(x)}{r_M},
 \max_{i\le k}\frac{|\sigma_i-\sigma_i'|(x)}{B_\sigma\phi_i}\right\}.
\]
Here $B_\nu=\alpha_\nu+\beta_\nu B_\sigma$, with
$\alpha_\nu$ fixed before $B_\sigma$ and $\beta_\nu\to0$
as $h_F\to\infty$.
\end{proposition}
\begin{proof}
The relative weight is $v_k=-L^{K-k}w/x_k$.
Put its centered quadratic in the q origin, as in
Lem.~\ref{C-399-32}, and keep its Gaussian family
$\mathbf G^v$ outside the small family. Its norm is at most
$C_G^v(|x_k-\mathsf t_k|/\mathsf t_k+\sup|v_k|)$.
The remaining gauge vertices have size
$7K_P\delta(\mathsf t_k)$.
The weight leaves a linear term. Subtract its constant
weight at a base plaquette $p_b$ for each bond $b$.
At a physical critical background this changes nothing,
since the derivative of the unweighted action along $b$
vanishes. Define its continuation by that difference
formula. The source Taylor estimate on a radius-$a$ pair
gains its curvature factor $a$. Let $\mathrm{sw}_k(v_k)$ be the
largest difference of $v_k$ between plaquettes with barycentres
at physical $\ell^\infty$ distance at most $M\lambda_k$.
Since $x_kv_k=-\lambda_k^{-1}w$, its swing bound is
$x_k\mathrm{sw}_k(v_k)\le\sqrt3M\operatorname{Lip}w$.
This bound and the polymer count give
$\|V^{\rm lin}\|\le C_{\rm lin}\tau_k(w;a)$.
This is the reason to retain the radius in the linear size.

Split the other potential into its scalar part $V_0$ and
its part $V_\psi$ with legs. Localize the old explicit
families at their fixed radii; this costs
$c_{\mathsf s}K_{\rm old}^{\rm s}\delta(\mathsf t_k)$.
Localizing the old scalar fluctuation families costs
\[
 K_\Lambda^{\rm s}\mathsf s_k(a)
 \sum_{i\le k}L^{-(k-i)}\epsilon_i^{{\rm s},0}(a).
\]
Choose $C_{\rm rad}\ge12\mathsf C_{\mathcal A}K_\Lambda^{\rm s}K_M^\bullet$.
Together with the linear term, the gauge vertices and the
final ceilings, this gives $\|V_0\|\le c_{\rm cl}/2$.
The source operator bounds and the old-leg estimate give
\begin{equation}\label{eq:qcd-source-leg-input}
\begin{split}
 \|V_\psi\|\le{}&C_V^{\rm s}h_F^2A_1\phi_k
 +c''_\psi c_{\mathsf S}h_F^2(\bar B_M+2B_\sigma)\phi_k\\
 &+c''_\psi\sum_{i\le k}\omega_{k+1-i}
                           \epsilon_i^{{\rm s},\psi}(a)
 +c''_\psi\Sigma_k^{\rm cs}.
\end{split}
\end{equation}
Here $\omega_n=C_\omega L^{n/4}$, and
\[
 c''_\psi=c_\psi+\mathsf C_{\mathcal A}(\alpha_0/a)\mathsf s_a
       +C_G\mathsf C_{\mathcal A}\delta_k\le(9/8)c'_\psi.
\]
For this bound choose
$C_{\rm rad}\ge4\mathsf C_{\mathcal A}c_h\alpha_0h_F/(a_{\rm an}A_\kappa\Gamma_F)$,
with the original value of $\Gamma_F$.
Thus the shift is small relative to every admitted radius.

The cell-sum term requires a more precise organization than the
identity for the total extracted kernel. Write
\[
 a_X(y)=\mathrm{ord}_0[k_X(\mathfrak c)](y)
              -\mathrm{ord}_0[k_X(0)](y)
       =\sum_p a_{X,p}(y).
\]
Here $p$ records a pair of local or walk terms in the two background
legs, and $Z_p$ is its actual supporting polymer. The increment
rule includes only $X$ of diameter less than one quarter of the
torus side. Although the total extracted kernel can be written as
$\sum_Xt_X+q_j+r_j$, with
\[
 t_X=k_X-\delta_{yy'}\mathrm{ord}_0[k_X],\qquad
 q_j=L^{-j}\mathsf S_1(\nu_j\psi^{\rm bg})
           -\mathfrak s_j\sum_z\bar\Psi(z)\nu_j(z)\Psi(z),
\]
this expression does not give the required local summands. We use
instead $u_X=k_X-\delta_{yy'}a_X$, filed at the cell of $X$;
on a single cube this includes the old zero-source local kernel
$L^{-j}\mu_j\mathsf S_1(1_\square\psi^{\rm bg})|_{\mathsf n=0}$.
Move $\bar\Psi(y)a_{X,p}(y)\Psi(y)$ to the cell containing $y$.
The accompanying local forms are
\[
 P_z^\mu=L^{-j}\mu_j
   \bigl[\mathsf S_1(1_z\psi^{\rm bg})
               -\mathsf S_1(1_z\psi^{\rm bg})|_{\mathsf n=0}\bigr],
 \qquad
 P_{z;X,p}=-\mathfrak s_j^{-1}L^{-j}
                 \mathsf S_1(a_{X,p}(z)1_z\psi^{\rm bg}).
\]
The first stands on the cube $z$ and the second on $Z_p$, pointed
at $z$. Together with the zero-source cube terms already included
in $u_X$, these forms sum to the extracted form
$L^{-j}\mathsf S_1(\nu_j\psi^{\rm bg})$. At zero source the moved
pieces and these forms vanish, leaving the old polymers. All
coefficients are holomorphic on their cores and commute with
colour; the summands are jointly gauge invariant. Thus the
old-leg estimate applies to each summand on its actual support.

Let $A_j(y,z)$ be the zeroth-moment map of
$L^{-j}\mathsf S_1(N1_z\psi^{\rm bg})$ at $\mathsf n=0$, and
let $A_j^{\mathsf n}$ denote its value at the matrix source.
The moment identity for each pair of walks uses the same lifted
paths as the local extraction. It gives, for a filed cell $c$,
\[
\begin{split}
 m_c-m_c(0)=\sum_{z\in c}\bigg[&
  \sum_y(A_j^{\mathsf n}-A_j)(y,z)\mu_j\\
 &+\sum_{X,p}\Bigl(1-\mathfrak s_j^{-1}
                 \sum_yA_j^{\mathsf n}(y,z)\Bigr)a_{X,p}(z)
   +\sum_{X,p}'a_{X,p}(z)\bigg].
\end{split}
\]
The prime denotes the polymers omitted by the diameter rule.
On the infinite lattice $\sum_yA_j(y,z)=\mathfrak s_j$;
the displayed expression then reduces to the change of the
background map applied to $\nu_j$, and vanishes at
$\mathsf n=0$. On a torus the remaining zero-source cell sum
is the difference of the filed wrapping-polymer terms.

File a change of a background map, with walk support $Y_q$,
on the level-$k$ hull of $c\cup Y_q$, $Z_p$, or
$Z_p\cup Y_q$, as appropriate. These hulls all contain a cube
of $c$; the outer legs enlarge them in the usual way. Measure
the coefficient at level $j$ with the weight
$e^{\kappa L^{-(k-j)}d_j(Y)}(1+d_j(Y))^3$, and denote the
result by $\|\cdot\|_{\rm nat}$. The hull count and the leg
estimate give a contribution at most
$C_{\rm hl}L^{k-j}h_F^2\sup_c\|m_c-m_c(0)\|_{\rm nat}$
per level-$k$ cube. The zero-moment units have the gain
$L^{(k-j)/4}$; degree at least four and field dependence
have the gain $L^{-(k-j)}$.

For completeness, the estimates for the moved pieces are
\[
 \sum_{p,y}\|a_{X,p}(y)\|
       e^{\kappa[d_j(Z_p)-d_j(X)]}
 \le C A_\kappa^2h_F^{-2}
              \bigl(\|E(X;\mathfrak c)\|+\|E(X;0)\|\bigr).
\]
The localized row sums give this inequality. Reaching the filed
cell from a cell $c'$ leaves the summable factor $e^{-r(c',c)}$.
Consequently, writing $n=k-j\ge1$ and
$\epsilon_j=\epsilon_j^{{\rm s},\psi}$,
\begin{equation}\label{eq:qcd-moved-piece-bounds}
\begin{gathered}
 \sum_{z\in c,X,p}\|a_{X,p}(z)\|
       e^{\kappa L^{-n}d_j(Z_p)}(1+d_j(Z_p))^3
       \le C_{\rm pw}h_F^{-2}\epsilon_j,
 \qquad
 \sum_{Z_p\supset\square,z}\|a_{X,p}(z)\|
       e^{\kappa d_j(Z_p)}\le C_{\rm pw}h_F^{-2}\epsilon_j.
\end{gathered}
\end{equation}
Here
$C_{\rm pw}=2e^{\kappa/2}S'_\kappa C C_{\rm wk}A_\kappa^2$,
where $S'_\kappa=\sup_d(1+d)^3e^{-5\kappa d/13}$.
For the omitted polymers the first bound gains
$e^{-c_t\ell/\lambda_j}$. Schwarz's lemma in the matrix source,
applied to the walks at rate $\kappa+2$, gives
\[
 \sum_q\left\|\sum_y(A_j^{\mathsf n}-A_j)_q(y,z)N\right\|
       e^{\kappa\Delta_q^+}(1+\Delta_q^+)^3
 \le C_{\rm bg}^{\rm w}\frac{L^j\|\mathsf n\|}{r_M}\|N\|,
 \quad
 C_{\rm bg}^{\rm w}=c_3^2\nu_PA_\kappa^2C_H^2,
 \quad c_3=\frac{27}{8e}.
\]
The torus error in $\mathfrak s_j^{-1}\sum_y A_j$ is at most
$\eta_{\min}^{-2}C_Ae^{-4c_t\ell/\lambda_j}$. Combining these
estimates yields
\[
\begin{split}
 \|m_c-m_c(0)\|_{\rm nat}\le{}&
 3M^3 C_{\rm bg}^{\rm w}\bar B_M
                  \frac{L^j\|\mathsf n\|}{r_M}\phi_j\\
 &+\left[\eta_{\min}^{-2}
       \left(C_{\rm bg}^{\rm w}\frac{L^j\|\mathsf n\|}{r_M}
                    +C_Ae^{-4c_t\ell/\lambda_j}\right)
                    +e^{-c_t\ell/\lambda_j}\right]
                         C_{\rm pw}h_F^{-2}\epsilon_j.
\end{split}
\]
At lag zero there is no spare decay for the polynomial factor.
Use the second moved-piece bound instead: the same per-cube
estimate holds with $C_{\rm hl}h_F^2$, and with the last
omitted-polymer tail replaced by $1$. In the source step itself
only lags at least $n_{\rm lag}$ enter $\Sigma_k^{\rm cs}$;
shorter lags are bounded as whole activities. The lag-zero
estimate will also be used for the regular sum on a history.

Now $L^{k-j}L^j\|\mathsf n\|/r_M<2/L$. For positive lags
the corresponding torus-tail factors are bounded by
$C_{\rm dx}$ and $C'_{\rm dx}$; for all lags use
$\max\{1,C'_{\rm dx}\}$. Summing proves
\begin{equation}\label{eq:qcd-source-cell-sum}
 \Sigma_k^{\rm cs}\le C_0^{\rm cs}h_F^2\phi_k+
 c_{\rm cs}\sum_{j\le k}\omega_{k+1-j}
                                  \epsilon_j^{{\rm s},\psi}.
\end{equation}
with
\[
 C_0^{\rm cs}=\frac{6C_{\rm hl}M^3C_{\rm bg}^{\rm w}\bar B_M}
                         {L(1-2L^{-1/2})},\qquad
 c_{\rm cs}=\frac{C_{\rm hl}C_{\rm pw}}{C_\omega}
 \left[\eta_{\min}^{-2}
       \left(\frac{2C_{\rm bg}^{\rm w}}L+C_AC_{\rm dx}\right)
          +\max\{1,C'_{\rm dx}\}\right].
\]
Neither a source-chain bound nor spatial regularity of the
coordinates $\sigma_i$ entered this estimate.

By the final ceiling \eqref{eq:qcd-source-leg-input} is at
most $c_{\rm cl}/8$. Multiply $V_\psi$ by a complex scalar.
The positive-degree output is zero at scalar zero, and its
cluster bound holds up to radius
$c_{\rm cl}/(2\|V_\psi\|)$. Schwarz's lemma gives
$\|\widetilde E_{\ge2}\|\le5K_M\|V_\psi\|$.
Taking out the moment and filing its moved pieces costs
$K_\mu^{\rm s}=K_\mu+\eta_{\min}^{-2}c_{\mathsf S}M^{-3}C_{\rm pw}$.
Put $\theta_E^{\rm s}=(K_\mu^{\rm s}/K_\mu)\theta_E
=\sqrt2K_\mu^{\rm s}K_Mc'_\psi$. Then
$5K_\mu^{\rm s}K_Mc''_\psi\le4\theta_E^{\rm s}$.
Consequently
\[
 \epsilon_{k+1}^{{\rm s},\psi}\le C_w^\psi\phi_k+
 4\theta_E^{\rm s}\sum_{i\le k}\omega_{k+1-i}
 [(1+c_{\rm cs})\epsilon_i^{{\rm s},\psi}
       +c_{\mathsf S}h_F^2(\bar B_M+2B_\sigma)\phi_i].
\]
The growing lag is summable against
$\phi_{i-1}/\phi_k\le CL^{-(k+1-i)/2}
(1+(k+1-i)\log L)^{\mathfrak p_0}$.
The prescribed floor on $h_F$ closes the leg size.

We next verify the asserted properties at each pointer. Each
member of the potential reads the independent coordinates on
its filing polymer: the gauge and linear vertices read $w$,
the mass member reads $\sigma_k$, and the local inverses read
$\mathsf n$. An old term is filed with the polymer of its
coefficient and all walks of its legs. Its local source
dependence is therefore retained. Local moments are likewise
filed with their actual pieces. The connected expansion then
reads only the coordinates on its union polymer.

For a place $p$, put $\dot V=(V_\psi)_p$, using the share
$n^{-1}N_p$ in degree $n\ge2$; one may also sum these shares
over the places of a site. This map preserves locality and
evenness, commutes with gauge transformations, and has norm
at most one on every polynomial. Define
\[
 H(X;\mathfrak c,t,\lambda)
   =\widetilde E[V_0+tV_\psi+\lambda\dot V](X),\qquad
 E^\bullet(X;\dot V)
   =\int_0^1\partial_\lambda H(X;\mathfrak c,t,0)\,dt.
\]
For $|t|<9/8$ and $|\lambda|<1$ the potential has size at most
\[
 \tfrac12c_{\rm cl}+(9/8+1)\tfrac18c_{\rm cl}
       =\tfrac{49}{64}c_{\rm cl}<c_{\rm cl}.
\]
Thus \cref{C-388-19,C-388-39} apply on an open set containing
the integration segment. The Cauchy formula for
$\partial_\lambda H$, followed by the uniformly bounded
integral in $t$, proves holomorphy in $\mathfrak c$ and the
pair. Both operations preserve the gauge identity of
\cref{C-388-36}. Taking degree zero gives strict scalar
invariance. The first-variation bound applies at every
pointer, with the norm of $V_\psi$ as its insertion bound.

For the restriction assertion let $a_k^{\rm rad}\le a'\le a$.
The coordinate domains satisfy
$2\mathscr D_k(a)\subset2\mathscr D_k(a')$.
The explicit old families are localized at their fixed radii
in both constructions. The old scalar fluctuation arrays,
including every pointed member, agree on the radius-$a'$
cores by the inductive hypothesis. Their bracket localization
commutes with this restriction, with the same slot and zero
counterterms. Hence their fluctuations agree at every admitted
$B$; at $k=0$ this array is empty. The two potentials and their
insertions therefore agree at the smaller radius, as do $H$
and $E^\bullet$.

The scalar terms pointed at bonds obey the same restriction
identity. The tests, quadratic origin and gauge vertices are
defined at the fixed radii. Their paths and insertions depend
only on their values at the given pair. The remaining scalar
origin is the localized old array just considered, together
with the linear source term. It has the same value on the
smaller core at every fluctuation, so its paths, insertions
and connected sums also agree there. Summing with the fixed
tent functions preserves this identity.

For degree zero let $D$ denote the new scalar activity
less its explicit Gaussian and normalizer families.
The scalar subtraction is exactly
\[
\begin{split}
 F_0^{(k+1)}(X;U;z;\mathfrak c)={}&D(X;U;z;\mathfrak c)
      -D(X;1;z;\mathfrak c)+D(X;1;z;0)\\
 &-1_{X=\{z\}}\widetilde e_{k+1}-\mathbf L_0(X;U;z).
\end{split}
\]
Thus $\|F_0\|\le2\|D\|+\|D(0)\|+
M^3|\widetilde e_{k+1}|+\|\mathbf L_0\|$, where
$M^3|\widetilde e_{k+1}|\le\|D(0)\|$ by its translation-cell
mean. The local action part $\mathbf L_0$ is the zero-source
one, formed at its fixed radius.
Every term in this subtraction is local and holomorphic
in the independent coordinates. The evaluations at the
trivial pair, which belongs to every radius, preserve
agreement on smaller cores; $\widetilde e_{k+1}$ and
$\mathbf L_0$ are independent of that radius. Thus the
pointed scalar families have all the properties required
in the next step's hypothesis.
There is no inverse-square radius loss from extracting a
new action coefficient. The scalar recursion is
\[
\begin{split}
 \epsilon_{k+1}^{{\rm s},0}\le{}&
 (C_w+C_w^\flat+C_w^\psi)\phi_k
 +\tfrac{17}8\theta_sh_F^2(\bar B_M+2B_\sigma)\phi_k\\
 &+2\theta_E\sum_{i\le k}\omega_{k+1-i}
 [(1+c_{\rm cs})\epsilon_i^{{\rm s},\psi}
       +c_{\mathsf S}h_F^2(\bar B_M+2B_\sigma)\phi_i]\\
 &+\tfrac12C'_{\rm lin}\tau_k(w;a)
 +2K_MK_\Lambda^{\rm s}\mathsf s_k(a)
                     \sum_{i\le k}L^{-(k-i)}\epsilon_i^{{\rm s},0}.
\end{split}
\]
The last sum contributes at most
$\frac14C_w^{\rm s}\phi_k+
\tfrac1{10}C'_{\rm lin}\tau_k$,
by \eqref{eq:qcd-source-radius-lag} and
$2K_MK_\Lambda^{\rm s}\mathsf s_k(a)\le1/5$.
The lag sum of old leg sizes costs at most
$\frac12C_w^{\rm s}\phi_k$. Choose
\[
 C_w^{\rm s}=C_w+C_w^\flat+C_w^\psi+
 3\theta_sh_F^2(\bar B_M+2B_\sigma)+
 4L^{1/2}C_{\rm pc}S_{1/4}\theta_E^{\rm s}c_{\mathsf S}h_F^2
                                  (\bar B_M+2B_\sigma).
\]
The scalar size closes, with the linear source term retained
at its actual radius-dependent size.

The moment rule \eqref{eq:qcd-source-increment} is a
holomorphic value at the trivial pair, independent of the
radius at which it was formed. For its locality estimate,
expand its background rows in walks. A term $f$ reads $w$
and the $\sigma_i$ on its polymer $X$, and $\mathsf n$ on
$X$ and its walk $\omega$. Extend the difference of $w$
from $X$ by the real and imaginary McShane extensions,
cut off at their supremum. Its level norm costs at most
$C(1+d(X))$ times the maximum normalized coordinate
difference on $X\cup\omega$. A Cauchy line from
$\mathfrak c$ inside the doubled domain gives
$|f(\mathfrak c')-f(\mathfrak c)|\le
2CM_0(1+d(X))t$, where $M_0$ is its term bound.
The maximum, rather than a sum over the coordinates, occurs.
Choose
$\kappa_\sigma\le\min\{\kappa/(4c_{\mathcal D}M),1/(24M_w)\}$.
Half the polymer decay and half the walk decay pay the
distance in \eqref{eq:qcd-increment-locality}; the rest
pay the polynomial and walk sums. This gives
$B_\nu=4c_{\rm inc}h_F^{-2}C_w^{\rm s}$.
Only the displayed terms in $C_w^{\rm s}$ involve
$B_\sigma$, and their coefficients are proportional to
$\theta_E$. Hence $B_\nu=\alpha_\nu+\beta_\nu B_\sigma$
with the claimed order of choices.
\end{proof}

\begin{lemma}[Comparison with the infinite lattice]\label{C-407-33}
At constant coordinates and the trivial gauge field, a source
step on the torus and on $\mathbb Z^3$ with the same old
families has, on polymers of diameter less than a quarter
of the side, activity difference bounded by
\[
 K_M e^{-2c_t\ell/\lambda_{k+1}}
 (\|V\|+e^{-\gamma_{P,k}})e^{-\kappa_+d(X)}.
\]
The corresponding moment difference costs at most
$C_{am}^\psi h_F^{-2}$ times this bound.
\end{lemma}
\begin{proof}
The translation-invariant fine inverse on the torus is
\[
 \sum_{a\in\mathbb Z^3}(-1)^{a_1+a_2+a_3}
       (\mathcal K_j^\infty)^{-1}(x,x'+\ell L^Ka).
\]
Exponential decay gives absolute convergence and verifies
both inverse identities. The same periodization holds for
the fluctuation covariance and background maps by their
Schur formulas. Between points closer than a quarter of
the side, nonzero images are at least three quarters of
the side away; the difference is at most
$Ce^{-4c_t\ell/\lambda_{k+1}}$ in the weighted kernel norms,
for $c_t\le\alpha_D/16$. Local inverses and supported
walks in such a polymer are the same on both lattices.
Only these image errors and the tails of background rows
remain. The activity is holomorphic on the operator balls
with the enlarged constants; the unperturbed data lie
inside them. Schwarz's lemma on the complex line between
the two sets of data bounds the output difference by its
uniform cluster bound times the relative kernel error.
Absorb the fixed factor and use the weaker displayed
exponential. The fibre coefficient estimate gives the
moment assertion.
\end{proof}

\begin{proposition}[The constant reference chain]\label{C-407-17}
For $|c|<8\varrho$ the infinite-lattice constant chain
$\sigma_k^{\rm c}(c)$ is a scalar, holomorphic in $c$ and
real for real $c$. It obeys
\[
 |s_k+\sigma_k^{\rm c}(c)|\le2\bar B_M/\mathsf t_k,
 \qquad |\sigma_k^{\rm c}(c)|
       \le B_0\phi_k|c|/\varrho.
\]
On the torus its constant-coordinate increment differs
from $\mu_{k+1}+\sigma_{k+1}^{\rm c}-L\sigma_k^{\rm c}$
by at most
$C_{\rm tor}\phi_{k+1}e^{-c_t\ell/\lambda_{k+1}}|c|/\varrho$.
The constants $B_0,C_{\rm tor}$ precede $B_\sigma$.
\end{proposition}
\begin{proof}
Put $x_i'=x_i-L^{K-i}c$. We first identify this splitting
with the holomorphic zero-source flow at bare label $\mathsf b-c$.
That flow extracts its own coupling increments $\tilde b_i$;
the source splitting retains $b_i$ and leaves
$(\tilde b_i-b_i)a_p$ in the scalar activities.
Compare the two on finite tori, with the same tests and
critical backgrounds. Assuming their mass increments agree
through level $i$, the two exponents at that level differ
only by a scalar constant and this redistribution of the
action. The next integral is the same. Its fermion-degree-two
parts therefore agree; different gauge Gaussian normalizers
have degree zero only. At the trivial pair take the zeroth
moments. The residual torus moments tend to zero by
\cref{C-407-33} and the wrapping-polymer tail. The nonzero
constant-field normalization then identifies the next mass
increments. This proves the identification by induction.
Holomorphy is taken from the first splitting. The second
contains action multiples and is not asserted to lie in
the normalized family space.

Here the complex-coupling flow is controlled by an induction
on the fixed disc $|c|<8\varrho$, with tests and calibrations
frozen. Impose $64\varrho\le r_x\gamma^{-2}$ and
$26C_w\phi(\gamma^{-2})\gamma^2\le r_x$.
The coupling recursion, whose fixed one-loop terms do not
depend on the coupling, then gives
\[
 |\tilde x_i-x_i|\le\tfrac23r_xx_i,
       \qquad |\tilde x_i|\ge\mathsf t_i/4.
\]
The exact one-loop telescope and its nonnegative slices give
$|G_i(\tilde x)|\le4B_F/\mathsf t_i$. The cubic remainder
sum is at most $1/(2\mathsf t_i)$ at the final ceiling.
Consequently
$|\tilde s_i|\le(4B_F+1/2)/\mathsf t_i
\le2\bar B_M/\mathsf t_i$. These bounds keep the next
step in its complex coupling and mass domains, closing
the induction and justifying composition of the holomorphic
steps throughout the disc.

The difference in coupling increments is
$|\tilde b_i-b_i|\le2C_w\phi(x_{i-1})$, whence
$|\tilde x_i-x_i'|\le4C_w\phi(x_i)$. Its change in the
one-loop coefficient is $O(\phi(x_i)/x_i^2)$, absorbed
by the cubic remainder after lowering $\gamma$.
The one-loop mass structure is
$\mu_j^{\rm 1L}=G_j-LG_{j-1}$. Therefore
\[
 \sigma_k^{\rm c}(c)=G_k(x')-G_k(x)+
 \sum_{j\le k}L^{k-j}
            (\mu_j^{\rm rem}(c)-\mu_j^{\rm rem}(0)).
\]
The nonnegative loop slices have total at most $B_F$.
Since $x_i\ge\mathsf t_i/2$ and $|x_i'|\ge\mathsf t_i/4$,
$|1/x_i'-1/x_i|\le8|c|/(\mathsf b\mathsf t_i)$.
Thus $|G_k(x')-G_k(x)|\le8B_F|c|/(\mathsf b\mathsf t_k)$.
Schwarz's lemma on the $8\varrho$ disc and the
cubic remainder give
\[
 \left|\sum_{j\le k}L^{k-j}
       (\mu_j^{\rm rem}(c)-\mu_j^{\rm rem}(0))\right|
 \le\frac{C_\rho S_3}{2\varrho}\phi_\sharp(x_k)^3|c|.
\]
The forward mass-flow bound on this disc gives the stated
window; increasing the final floor so that
$\phi_\sharp^3\le\phi$ gives $B_0=16B_F+C_\rho S_3$.
The identification with the constant-coupling flow gives
$\sigma_k^{\rm c}(c)=\widetilde s_k-s_k$ on $\mathbb Z^3$.
It is therefore a scalar number, holomorphic in $c$ and real
at real $c$, by the reality of that flow. This conclusion
does not assert the same reality for individual pointed
or polymer pieces.

For the torus comparison, match each polymer of diameter less
than a quarter of the side with its translate in $\mathbb Z^3$.
Define $\Delta_i$ by the level-$i$ family norm with
\[
 \sum_\sigma\|F^{\rm tor}_\sigma(Z)-F^\infty_\sigma(Z)\|
\]
in place of the norm of the activity at $Z$. The summands are
the new activity before extraction, each moved form $P_{z;X,p}$
on its own $Z_p=Z$, and the local mass member. The last member
has zero difference: both lattices use the same $\mu_i$ and
the same algebra of formal legs. Let $\Delta_i^0$ be the maximum
of the scalar difference norm and the corresponding difference
norms at every pointer, including the action terms. At constant
coordinates $\mathsf n=0$, the new cell-sum identity has
only the error in $\mathfrak s_i^{-1}\sum_y A_i(y,z)$ and
the omitted large polymers. Keep each summand on its $Z_p$.
The latter polymers cannot enter a new polymer of diameter
less than one quarter of the side. The former, together
with the residual zero-source torus moment, give a forcing
$r_i^{\rm c}$ bounded by
$C\phi_i e^{-4c_t\ell/\lambda_i}$, using the moved-piece
bound and the constant-chain sizes $\epsilon_i^{\rm c}
\le C\phi_{i-1}$. These constants precede $B_\sigma$.
Extraction is bounded in this norm. Indeed, if $Z_p$ has
diameter less than a quarter of the side, its walks do not
wrap around the torus. Their local inverses and moment maps
agree on both lattices. Thus the difference of the two forms
is the moment of the single difference family $E^{\rm dif}$
before extraction. The piece count in
\cref{eq:qcd-moved-piece-bounds} gives
\[
 \sum_{Z_p\supset\square}\sum_z
     \|a^{\rm dif}_{X,p}(z)\|e^{\kappa d_i(Z_p)}
       \le C_{\rm pw}h_F^{-2}\|E^{\rm dif}\|.
\]
Multiplying by the norm of a local mass form and using
$\mathfrak s_i\ge\eta_{\min}^2$ bounds the forms by
\[
 \eta_{\min}^{-2}c_{\mathsf S}M^{-3}C_{\rm pw}
 \|E^{\rm dif}\|.
\]
Including the unextracted term therefore
costs at most $K_\mu^{\rm s}$, the same factor used for the
source-step sizes. The derivative of the new activity with
legs in the old leg potential is at most $50K_M/9$ on the
admitted ball. Since
$\theta_E^{\rm s}=\sqrt2K_\mu^{\rm s}K_Mc'_\psi$,
its contribution is at most $4\theta_E^{\rm s}$ times the
old leg norm. Its derivative in the scalar potential has
an additional factor $\|\mathcal V_\psi\|/c_{\rm cl}$,
because the part with legs vanishes when $\mathcal V_\psi=0$.
These estimates, \cref{C-407-33} and the local moment rule give
\[
\begin{split}
 \Delta_{k+1}\le{}&C_{\rm img}K_Me^{-2c_t\ell/\lambda_{k+1}}
 \epsilon_{k+1}+4\theta_E^{\rm s}\sum_{i\le k}\omega_{k+1-i}\Delta_i\\
 &+\epsilon_\gamma\sum_{i\le k}L^{-(k-i)}\Delta_i^0
                  +C\sum_{i\le k}L^{k+1-i}|r_i^{\rm c}|,\\
 \Delta_{k+1}^0\le{}&C'K_Me^{-2c_t\ell/\lambda_{k+1}}
 \epsilon_{k+1}+C_{\rm dl}\sum_{i\le k}\omega_{k+1-i}\Delta_i\\
 &+\epsilon_\gamma\sum_{i\le k}L^{-(k-i)}\Delta_i^0
                  +C\sum_{i\le k}L^{k+1-i}|r_i^{\rm c}|.
\end{split}
\]
Here $\epsilon_{k+1}\le C\phi(\mathsf t_k)$, since the
linear term is zero at a constant source; all constants
precede $B_\sigma$, and $\epsilon_\gamma\to0$.
The lag sums are bounded by $C_\omega S_\omega$ and
$S_L=\sum_{n\ge1}L^{1-3n/2}(1+n\log L)^{\mathfrak p_0}$.
The $r_i^{\rm c}$ terms are summable by their doubly exponential
fall in $k-i$. Choose the earlier leg floor so that
$4\theta_E^{\rm s}C_\omega S_\omega\le2^{-1/2}$ and the final
ceiling so that
\[
 \epsilon_\gamma S_L\le
 \min\{1/2,(1-2^{-1/2})/(4C_{\rm dl}C_\omega S_\omega)\}.
\]
If $C'$ bounds the forcing terms, the two recursions close
with $D=4C'/(1-2^{-1/2})$ and
$D_0=2(C'+C_{\rm dl}C_\omega S_\omega D)$.
The large-diameter polymers omitted by the rule have the
same exponential bound. The increment difference is
holomorphic and zero at $c=0$; Schwarz's lemma supplies
its factor $|c|/\varrho$.
\end{proof}

Define the source chain by
\[
 \sigma_0=0,\quad
 \sigma_{k+1}=L\sigma_k+\nu_{k+1}(\mathfrak c_k(s))-\mu_{k+1},
 \quad \mathfrak c_k(s)=(w,-a_KJ,(\sigma_i(s))_{i\le k}).
\]
The increment is read as a constant on its coarse cell.
Write $\mathcal N_k=-\lambda_kJ+\sigma_k$ and
$\mathcal Z_k=1-L^{K-k}\partial_J\sigma_k|_0$ on constant
matrix sources.

\begin{lemma}[The source chain]\label{C-407-13}
The chain is formed and holomorphic on $\mathscr O_k$, stays
in $\mathscr D_k$, and satisfies
\[
 \|\sigma_k-\sigma_k^{\rm c}(w(\cdot))\|
 \le C_{\rm var}\phi_k\Xi_k+
 2C_{\rm tor}C_{\mathcal L}\phi_ke^{-c_t\ell/(2\lambda_k)}
                            \|w\|_\infty/\varrho.
\]
Moreover
\[
 \|\sigma_k(0,J)\|\le B_\nu S_\uparrow\phi_k
                  \lambda_k\|J\|/r_M,\qquad
 \|\mathcal Z_k-1\|\le B_\nu S_\uparrow\phi_k/r_M,
\]
and $\|\mathcal Z_{k+1}-\mathcal Z_k\|
\le B_\nu\phi_{k+1}/r_M$.
At every site the response is equivariant under the two
flavour factors:
\[
 V\mathcal Z_k(J)V^{-1}=\mathcal Z_k(VJV^{-1}),\qquad
 V=e_+\otimes V_1+e_-\otimes V_2,
 \quad V_1,V_2\in\mathrm U(N_f).
\]
For $S=\operatorname{supp}(s-s')$,
\[
 |\sigma_k(s)(x)-\sigma_k(s')(x)|
 \le C_t\phi_ke^{-\kappa_\sigma\operatorname{dist}(x,S)/L^k}.
\]
All statements hold uniformly in the volume and cutoff.
\end{lemma}
\begin{proof}
Put $q=(4/3)q_\phi<3/4$,
$C_{\rm tl}=\sup_{t\ge0}(1+t)e^{-\kappa_\sigma t}$,
$C_{\rm var}=2B_\nu(1+C_{\rm tl}+42M)/(1-\frac83L^{-1/2})$,
and $C_{\mathcal L}=\sum_{u\ge0}(2L^{1/2})^u
e^{-c_tL(L^u-1)}$.
Choose
\[
 B_\sigma\ge4B_0+3C_{\rm var}+8C_{\rm tor}C_{\mathcal L},
 \qquad B_\sigma\ge2C_{\mathcal L}B_\nu.
\]
These inequalities can be met: first make $\beta_\nu$ small
by raising $h_F$, then choose $B_\sigma$.

At a point $x_0$ compare with the reference chain at
$c=w(x_0)$. The weighted local difference in $w$ is bounded
by $(C_{\rm tl}+41M)\varrho\Xi_{k+1}$; that in the fine
mass by $r_M\Xi_{k+1}/L$. Cauchy's estimate for the constant
reference chain bounds its variation in $w$. If $D_k$
is the supremum of the resulting chain difference, the
increment locality estimate gives the induction
\[
 D_k\le C_{\rm var}\phi_k\Xi_k+T_k,\quad
 T_k=2C_{\rm tor}\sum_{j\le k}L^{k-j}\phi_j
 e^{-c_t\ell/(2\lambda_j)}\|w\|_\infty/\varrho.
\]
Indeed the coefficient of $C_{\rm var}\phi_{k+1}\Xi_{k+1}$
in the next step is at most
$q+B_\nu(C_{\rm tl}+41M)/C_{\rm var}
+4B_\nu/(3LB_\sigma)\le1$.
The remaining forcing terms are at most $T_{k+1}$ by
$2C_{\mathcal L}B_\nu\le B_\sigma$.
The defining series for $C_{\mathcal L}$ gives
$T_k\le2C_{\rm tor}C_{\mathcal L}\phi_k
e^{-c_t\ell/(2\lambda_k)}\|w\|_\infty/\varrho$.
The window follows from $\Xi_k<2$ and $\|w\|_\infty<4\varrho$;
holomorphy follows inductively by composition.

At $w=0$ the normalized matrix-source entry dominates the
feedback entries. The increment bound is
$B_\nu\phi_j\lambda_j\|J\|/r_M$.
Summing it with $L^{k-j}$ gives the first matrix bound.
Differentiating this holomorphic line at zero gives the
two response estimates.

For locality set $a_j(t)=\sup_{\operatorname{dist}(x,S)\ge t}
|\sigma_j(s)(x)-\sigma_j(s')(x)|$. The same local estimate
gives
\[
 a_j(t)\le La_{j-1}(t)+B_\nu\phi_j
 \left[C_{\rm loc}e^{-\kappa_\sigma t/L^j}
 +e^{\kappa_\sigma\sqrt3}\sup_{u\ge0}e^{-\kappa_\sigma u/L^j}
 \max_{i<j}\frac{a_i(t-u)}{B_\sigma\phi_i}\right],
\]
with $a_i(t-u)=2B_\sigma\phi_i$ if $u>t$ and
$C_{\rm loc}=8e^{\kappa_\sigma\sqrt3+\kappa/4}$.
For $t\le t_\flat L^j$ the claimed exponential follows
from the window by choosing $C_t\ge2B_\sigma
e^{\kappa_\sigma t_\flat}$. For larger $t$, choose
$2L^{1/2}e^{-\kappa_\sigma(L-1)t_\flat}\le1/4$.
The first term then costs at most one quarter of the
inductive bound; the forcing term costs one quarter after
raising $C_t$, and the feedback costs at most one half
since $e^{2\kappa_\sigma}B_\nu/B_\sigma\le1/2$.
This proves the exponential locality estimate.
\end{proof}

The flavour identity passes through the moment rule, the
averages and the covariances; differentiating it at zero
gives the displayed equivariance of $\mathcal Z_k$. At $J=0$
it implies, site by site,
\[
 \sigma_k(x;w,0)=n_k^0(x)+n_k^3(x)\tau_3,
 \qquad n_k^0(x),n_k^3(x)\in\operatorname{End}(\mathbb C^2_{\rm s}).
\]
These coefficients are spin operators, with no assertion of
pointwise reality. The scalar constant reference chain gives,
with $\Delta=\sigma_k(x;w,0)-\sigma_k^{\rm c}(w(x))$,
$n_k^3\tau_3=(\Delta-\tau_1\Delta\tau_1)/2$.
The source-chain estimate therefore bounds $\|n_k^3\|$ by
the variation term and the torus error. At constant $w$ only
the torus error remains; it vanishes for the constant chain
on $\mathbb Z^3$. The conjugation used for whole energies
and partition functions is
$s^\star=(\bar w,\gamma_4J^\dagger\gamma_4)$.
Its pseudoscalar signs must be retained, even for real
coefficients in the physical basis.

\begin{lemma}[Windows along arbitrary histories]\label{C-407-21}
For the chain obtained using only polymers inside the
integration regions of a history, let
$S_j=\bigcup_{i\le j}\widetilde Z_i$ and $R'_\flat=7LM$.
Then
\[
 |\sigma_j^h(s)(x)-\sigma_j(s)(x)|
 \le8B_\nu\phi_j e^{-\kappa_\sigma
          \operatorname{dist}(x,S_j)/(2L^j)}
 \quad(\operatorname{dist}(x,S_j)\ge R'_\flat L^j),
\]
and throughout $\Omega_j$ the difference is at most
$16B_\nu\phi_j$. Thus $\|\sigma_j^h\|\le2B_\sigma\phi_j$.
Replacing the full-chain coordinates by history coordinates
in the local source energies changes their level-$j$ sum
by at most $C'_{\mathfrak E}\#_j(\widetilde Z_j)$.
\end{lemma}
\begin{proof}
Omitted polymers reach a point at distance $t$ only with
$e^{-\kappa_\sigma t/L^j}$. The common polymers obey
the increment locality estimate, with the second
coordinate in its doubled domain. Hence the recurrence
for $a_j(t)$ in the preceding proof holds, now taking
$a_i(t-u)=3B_\sigma\phi_i$ below $R'_\flat L^i$.
The claimed bound follows with its exponent divided by
two: the old-level term costs at most one quarter if
$Le^{-\kappa_\sigma R'_\flat}\le1/4$; the bracket is at
most six times the proposed exponential if
\[
 (C_{\rm loc}+3e^{2\kappa_\sigma})
 e^{-\kappa_\sigma R'_\flat/4}
       +8e^{2\kappa_\sigma}B_\nu/B_\sigma\le6.
\]
These are fixed earlier floors. At a point of $\Omega_j$
the preceding source difference is already in the
exponentially small range by layer separation. One more
increment, each bounded by $B_\nu\phi_j/2$, gives
$9B_\nu\phi_j\le16B_\nu\phi_j$.
For the energy sum use its coordinate Cauchy estimate on
the doubled domain. Only the $\sigma_i$ differ, so no
Lipschitz extension polynomial is needed. Its decay
times the displayed distance exponential is summable
over cubes to $C'_{\mathfrak E}\#_j(\widetilde Z_j)$.
\end{proof}

\subsection{Source energies and normalized partition functions}

Let $\mathcal P$ be the closed physical source ball
$\|w\|_\sharp\le2\varrho$, $\|J\|_\infty\le2\varrho_J$,
with $J$ in the complex span of the physical basis $\mathfrak B$,
where $2\varrho_J\le c_Jm_-\le r_J/2$ and $c_J\le1/2$.
It lies strictly inside every $\mathscr O_j$.
Let $\mathfrak E_j(X;s)$ be the trivial-pair value of the
new scalar activity before its constant is extracted,
including its gauge and fermion normalizers and Gaussian
family. Define
\[
 e_j(X;s)=\mathfrak E_j(X;s)-\mathfrak E_j(X;0),\qquad
 E_k(s)=E_k+\sum_{j\le k}\sum_X e_j(X;s),\qquad
 Z_k^{\rm sub}[s]=e^{-E_k(s)}Z[s].
\]
No coordinate-dependent constant remains in the activities
after this extraction. The source-dependent action and
first-moment coefficients remain there; only the coupling
shift of zero sources is extracted.

\begin{proposition}[Local source energies]\label{C-407-23}
For $j\le K$, $e_j(X;\cdot)$ is holomorphic on
$\mathscr O_{j-1}$, zero at zero, and
\[
 |e_j(X;s)|\le C_{\mathfrak E}e^{-\kappa d_j(X)}.
\]
For $S=\operatorname{supp}(s-s')$,
\begin{equation}\label{eq:qcd-source-energy-tail}
 |e_j(X;s)-e_j(X;s')|\le
 C_{\mathfrak E}e^{-\kappa d_j(X)}
 e^{-\kappa_\sigma\operatorname{dist}(X,S)/L^j}.
\end{equation}
The distance here is in fine lattice units. Consequently,
for any fixed $0<\kappa_e<\kappa-c_E$,
\[
 \sup_{j,\square,s}\sum_{X\ni\square}
       e^{\kappa_ed_j(X)}|e_j(X;s)|\le C_{\rm locE}.
\]
Conjugation holds for the sum at each level:
\[
 \sum_{X\in\mathcal D_j}e_j(X;s^\star)
       =\overline{\sum_{X\in\mathcal D_j}e_j(X;s)}.
\]
Consequently $E_k(s^\star)=\overline{E_k(s)}$.  No
polymerwise conjugation identity is asserted.
\end{proposition}
\begin{proof}
The trivial pair belongs to every pair-radius domain.
The source step is holomorphic there, and composition
with the holomorphic chain proves holomorphy. Its scalar
fluctuation contribution is bounded by the source-step
size at $a_{j-1}^{\rm rad}$, uniformly on the level domain.
The gauge normalizer is unchanged and cancels in $e_j$.
The Gaussian-family difference is at most
$2e^{-\kappa_+d_j(X)}$ by its explicit fixed bound.

For the fermion normalizer use the exact derivative
\[
 \log\det\Gamma_{j-1}(\mathsf n)-\log\det\Gamma_{j-1}(0)
 =\int_0^1\operatorname{tr}\left[
 \Gamma_{j-1}(t\mathsf n)^{-1}
 (\bar{\mathcal H}\widehat{\mathsf n}\mathcal H)_{AA}
                       (t\mathsf n)\right]dt.
\]
Expand the inverse and both backgrounds in their source
walks and file the trace on their hull. The fibre trace
costs $8N_f$ a place. The weighted walk sums give
$C_N^{\mathsf n}N_fL^j\|\mathsf n\|e^{-(\kappa+1)d_j(X)}
\le2C_N^{\mathsf n}N_fr_Me^{-(\kappa+1)d_j(X)}$.
Together these bounds give the first assertion.

For the difference estimate use the coordinate Cauchy
argument in the proof of Prop.~\ref{C-407-16}.
If $X$ misses $S$, the source entries agree on $X$.
A background walk can see their difference only by
reaching $S$. After half its decay has been reserved
for summation, the remaining half gives the factor
$e^{-\operatorname{dist}(X,S)/(12\sqrt3M_wL^j)}$.
The chain entries have the exponential locality of
Lem.~\ref{C-407-13}. As
$\kappa_\sigma\le1/(24M_w)$, these give
\eqref{eq:qcd-source-energy-tail}. If $X$ meets $S$,
use twice the first bound. The entropy sum proves the
last displayed estimate. For conjugation, sum the new activities
and the localized normalizers before extracting their values.
Read on shell, this is the logarithm of the whole step integral.
At the trivial pair conjugation carries its scalar part at $s$
to the complex conjugate of that at $s^\star$. The trivial
block field passes every density test, so it belongs to the
background pairs on which this identity holds. Subtracting
the real value at zero sources gives the asserted identity
for the sum of the $e_j$. The assignment of the terms to
individual polymers need not be preserved by conjugation.
\end{proof}

\subsection{Complex densities and the bounds on stored layers}

The local construction gives holomorphic activities. To use them in a
history we must also prove an identity of measures on the history's
coordinates. We do this coefficient by coefficient in the even
Grassmann algebra. This keeps the coordinate changes and their null
sets separate from the later estimate of the density.

\begin{lemma}[A history density at complex sources]\label{C-407-34}
Assume the admitted pure gauge construction and majorant settings at
level $k$, and let $s\in\mathcal P$ and $h\in\mathcal H_k$.
Write ${\rm Gr}^{\rm ev}_h$ for the even algebra of the generators
retained by $h$. Suppose that, on the coordinate space
$\mathsf X_h^{{\rm co}\prime}$, the term of $h$ has the form
\[
 \hat\tau_h^{\rm F}[s]
 =\mathsf{Stay}^{\rm T}_h\chi_k(\Omega_k)
       e^{\boldsymbol\Phi_h^w}\exp A^{\rm F}_{k,h}\,
       \mathscr M_h .
\]
Here the displayed density is defined and Borel off one
$\mathscr M_h$-null set, is integrable in the fibre norm, and the
identity holds for every Grassmann coefficient on every Borel set.
Then there is a Borel, fibre-norm integrable density $g_h$ on the
free-coordinate space such that
\[
 \check\tau_h^{\rm F}[s]=g_h\mathscr M_h^{\rm fr}.
\]
There is a Borel set $S'\subset\mathsf v_h(\mathcal V_0[h])$,
outside which $\int\|g_h\|_{h_F}\,d\mathscr M_h^{\rm fr}=0$,
with the following property. For $\mathsf v\in S'$ and any
configuration $U\in\mathcal V_0[h]$ with $\mathsf v_h(U)=\mathsf v$,
put $\phi=\xi_h^{{\rm co}\prime}(U)$ and $V_n=\mathcal M^nU$.
Then the two exponents are defined at $\phi$ and
\[
 g_h(\mathsf v)=
 \prod_{(t,D)\in{\rm Pc}(h)}g_{t,D}(V_{t-1})
       e^{\boldsymbol\Phi_h^w(\phi)}
       \exp A^{\rm F}_{k,h}(\phi),\qquad g_{t,D}>0.
\]
In particular its fibre norm is bounded by the product of the scalar
moduli and the norm of the last exponential. If on $S'$ this bound is
at most $\kappa_h\bar u_h$, where $\bar u_h$ is the product of the
pure gauge piece majorants at loss zero, then
\[
 \int\|g_h\|_{h_F}\,d\mathscr M_h^{\rm fr}
 \le\kappa_h\exp\!\left(C'_\star\sum_{t\le k}\#_t(Z_t(h))\right).
\]
For $h_F\ge1$, Berezin integration is bounded by this norm integral
on every Borel set.
\end{lemma}
\begin{proof}
Choose the monomial basis $(e_I)$ of ${\rm Gr}^{\rm ev}_h$.
The coefficients $\varsigma_I$ of the bare sourced density after the
partial block kernels are bounded complex Borel functions of $U$.
Indeed the kernels are defined at every configuration, the bare
density is continuous on the compact configuration space, and the
partial Berezin integrals are finite sums of products. The pure gauge
measure identities therefore apply to each $\varsigma_I$.

Let $f$ be the density in the hypothesis. Pull its one exceptional
set back by $\Pi_h^{\rm le}\times{\rm id}$ and put
$g_h=f\circ(\Pi_h^{\rm le}\times{\rm id})$ off this pullback and
$g_h=0$ on it. The coordinate-change identity gives
$\check\tau_h[\varsigma_I]=(g_h)_I\mathscr M_h^{\rm fr}$ for
every $I$. Since the algebra is finite dimensional, coefficientwise
integrability and fibre-norm integrability are equivalent. This
justifies the same construction for all coefficients at once.

Write $\rho_0(U)=e^{-\beta A(U)}$ for the bare positive gauge
density. Its minimum on the compact configuration space is positive.
The source-free positive term has a Borel full-measure support
$S'_0\subset\mathsf v_h(\mathcal V_0[h])$. For its complement $R$,
the pure gauge domination of a term at a bounded complex density gives
\[
 |\check\tau_h[\varsigma_I]|(R)
 \le\frac{\sup|\varsigma_I|}{\min\rho_0}\,
                  \check\tau_h[\rho_0](R)=0.
\]
It follows that $\int_R|(g_h)_I|=0$ for every $I$, hence that the
fibre-norm integral vanishes there. Intersect $S'_0$ with the open
set of points in their tiers and remove the pulled-back null set;
call the result $S'$. The density vanishes outside the tiers because
the staying factor does. This proves the asserted support property.

At $\mathsf v\in S'$ the coordinate identities identify the pullback
with $\phi$. The staying factor is the displayed product of positive
piece factors, and $\chi_k(\Omega_k)=1$. Thus both exponents are
defined and the stated identity follows. An exponential in the even
algebra is invertible: writing $A=A_0+(A-A_0)$, its scalar factor is
$e^{A_0}$ and its remaining factor is a finite nilpotent exponential,
with inverse $e^{-(A-A_0)}$. Submultiplicativity now gives the norm
bound. The integral of $\bar u_h$ is the pure gauge majorant bound,
and the complement of $S'$ has zero norm integral, proving the last
estimate. Finally $|{\rm Ber}F|\le\|F\|_{h_F}$ and the linear
Berezin map commutes with the finite-dimensional integral.
\end{proof}

The next three estimates check that the source-dependent pieces fit
the same history geometry. The regular terms with fermion legs are
larger than at zero source, and the stored background maps have a
smaller margin than a generic doubled operator bound would give.
Both facts must be accounted for before closing the history induction.

\begin{lemma}[The sourced regular sum]\label{C-407-35}
Let $s\in\mathcal P$, $1\le k\le K$, $h\in\mathcal H_k$, and evaluate the
regular activities at the history coordinates $\mathfrak c^h$ and
at any evaluation configuration in the support of its tests. Then
\[
 \left\|\sum_{j=1}^k\sum_{(z,X)\in\mathsf{Ix}^{\rm rs}_j(\mathcal S)}
 E^{(j)}_{\ge2}(X;\mathsf U,\bar\Psi^{\rm bg},\Psi^{\rm bg};
                         z;\mathfrak c^h_{j-1})\right\|_{h_F}
 \le C^{\rm s}_{\rm leg}\Sigma_{\mathcal L}(h),
\]
and the sourced mass member satisfies
\[
 \|\mathsf M_h(s)\|_{h_F}
 \le w_0\nu C_H^2h_F^2L(\bar B_M+2B_\sigma)
                         \Sigma_{\mathcal L}(h).
\]
More precisely, for $n\ge1$, the contributions filed at a cube of
the layer $\mathcal L_n$ are bounded by
$C^{{\square},{\rm s}}_{\rm leg}\phi_n$ and
$w_0\nu C_H^2h_F^2L(\bar B_M+2B_\sigma)\phi_{n-1}$,
respectively. The mass profile vanishes on $\mathcal L_0$.
The constants are fixed before the lower mass bound
and the final coupling ceiling.
\end{lemma}
\begin{proof}
Lem.~\ref{C-407-21} places the coordinates in
$2\mathscr D_{j-1}(\alpha_0)$. There the part with legs has size
$2C_w^{\rm s}\phi_{j-1}$. Its background maps have the localized
row and H\"older bounds at the fine mass matrix, with the admitted
enlarged constants. Apply the regular-sum argument cell by cell,
using the radial gauge centered at the cell where a summand is
filed. Its center belongs to the summand's polymer, so the same
diameter bound controls the gauge field.

The degree-two term is split into its zero-moment units and the
remaining cell sums. The first part has the gain from subtracting
the constant leg. At a row of $\mathcal L_n$, the H\"older estimate
for a localized background row gives the same gain in the lag
$n-j$ as in the zero-source regular sum. The degree-at-least-four
and field-dependent parts use the corresponding power counting.
For the cell sums one uses the estimate proved in the source step,
\eqref{eq:qcd-source-cell-sum}; its coefficients are read at the
history coordinates, without assuming that the induced mass field
varies slowly. The moved pieces are filed at the hulls of their
polymers, while the forms taken out at a pointer are filed on the
polymer supporting that form.

The differences between these filings and the interior sums occur
at the rims. At lag $n-j$ their bound contains
\[
 L^{n-j}e^{-\kappa(7L^{n-j}-2)}.
\]
The zero-source torus cell sums have the additional bound
\[
 2C_\omega C_{\rm wr}L^{n-j}
                   e^{-\kappa(L^{n-j+1}-2)}.
\]
Both sequences are summable. Combining them with the weighted
piece count of the source step and its $\phi_{j-1}$ size gives
$C^{{\square},{\rm s}}_{\rm leg}\phi_n$ per layer cube.
Since $x_n\ge\tfrac12\gamma^{-2}$ and
$\gamma^{-2}\ge e^{18\mathfrak p_0}$, this is bounded by a
constant independent of the level. Summing the cubes gives the
first estimate.

Every generated mass profile is supported in $\Lambda_j\subset\Omega_1$,
$j\ge1$, and hence vanishes on $\mathcal L_0$.
On $\mathcal L_n$, $n\ge1$, the norm of the mass coefficient is at most
$L^{-(n-1)}(\bar B_M+2B_\sigma)\phi_{n-1}$. Count the fine
places in a layer cube and apply the row-sum bound to its two
background legs. The factors are precisely
$w_0\nu C_H^2h_F^2L$, giving the second estimate.
\end{proof}

\begin{lemma}[Scalar source activities and their scale sum]\label{C-407-40}
On $2\mathscr D_{j-1}(\alpha_0)$, the scalar activities obey,
polymer by polymer and pointer by pointer,
\[
 \sup_{X,z}e^{(\kappa+1)d_j(X)}
    \sup|E^{(j)}_0(X;\cdot;z;\mathfrak c)|
 \le E_0^{\rm s}:=E_0+2+K_M^\bullet c_{\rm cl},
 \qquad K_M^\bullet=K_M+C_EK^{\rm v}_{\mathbf M}.
\]
Consequently, at the coordinates of a history,
\[
 \begin{split}
 &\sum_{j\le k}\sum_{(z,X)\in\mathsf{Ix}^{\rm rs}_j(\mathcal S)}
 |E^{(j)}_0(X;\mathsf U;z;\mathfrak c^h_{j-1})
       -E^{(j)}_0(X;\mathbf1;z;\mathfrak c^h_{j-1})|\\
 &\qquad+\sum_{j\le k}|L^jb_jA(\varphi_j,\mathsf U)|
 \le C_5^{\rm s}\Sigma_{\mathcal L}(h),\qquad
 C_5^{\rm s}=C_{\rm m}^{\rm sc}
                 (E_0^{\rm s}+B_\sharp+M^{-3}\mathsf B).
 \end{split}
\]
The scalar edge-part bound also holds with $E_0^{\rm s}$ in
place of $E_0$. These two sourced constants are fixed before
the constant $C_E^{\rm F}$ used in the history estimate.
\end{lemma}
\begin{proof}
Decompose the scalar part into the gauge normalizer, the closed
fermion loops, the Gaussian family at the weight, the scalar
fluctuation family, and the local constant taken out. The gauge
normalizer and the local source-free subtractions are unchanged.
The closed loops at the fine mass matrix have their original
pointerwise bound: $L^{j-1}\|\mathsf n\|<2r_M/L\le\varrho_*$
puts their local walk expansion in its admitted domain. Their
source-dependent values at the trivial pair have already been
assigned to the source energy.

Point the Gaussian family along the segment of its quadratic
potential. The insertions are
$-(1-\tau)\hat g_b+\mathcal V_b^v$, with
$\mathcal V_b^v=-\tau T_2[v](z_b,z)$, and their sum is the
quadratic potential itself. Its admitted ceiling gives a bound
one at each pointer. Comparing it with its zero-source value
costs at most two.

For the scalar fluctuation family, insertions at fermion places
are bounded by the part with legs of the potential. The cluster
majorant is multilinear and nondecreasing in the vertex norms,
so these insertions obey the same bound as that part. The bond
insertions require the pointed scalar estimate of
Lem.~\ref{C-388-37}. The gauge and old-activity insertions cost
at most $9$ times their input norms, followed by the first-variation
constant $K_\Pi$ and the tent count $3L^3$. The resulting scalar
constant is $K_M^\bullet$. For the linear source term the tent
insertion is $\mathcal V^{\rm lin}(Y;P,h_zB)$, $0\le h_z\le1$;
the term's bound is a sum over its individual bonds and remains
$C_{\rm lin}\tau_{j-1}(w;\alpha_0)$. Schwarz's lemma for the
additional origin bounds its first variation by $K_\Pi$ times
this norm.

The scalar recursion of the source step therefore closes for
the pointed terms too, with bound
\[
 2C_w^{\rm s}\phi_{j-1}
      +C'_{\rm lin}\tau_{j-1}(w;\alpha_0).
\]
Here the old families in the recursive term are the whole
pointed families; no extra pointer sum is inserted. The
$C_{\rm rad}$ floor formed with $K_M^\bullet$ makes the
coefficient of that term at most $1/6$. The two displayed
source sizes are each at most $K_Mc_{\rm cl}/2$ by their
admission ceilings. The unchanged terms, the Gaussian cost
two, and this fluctuation bound give $E_0^{\rm s}$.

These scalar families are functions of the background pair,
holomorphic on its fixed-radius domain and gauge invariant
term by term. Their first curvature derivative at the trivial
pair vanishes. Thus their variation on a layer cube at lag
$n-j$ has the volume factor $L^{3(n-j)}$ and the quadratic
curvature gain $L^{-4(n-j)}$. The scalar sum-over-scales
estimate applies with $E_0^{\rm s}$, and $|b_j|\le B_\sharp$
is unchanged. This proves the stated $C_5^{\rm s}$ bound.
The scalar edge parts are the omitted pointed sums together
with the shares of the old action subtractions; applying their
original count to the pointerwise bound gives the final assertion.
\end{proof}

\begin{lemma}[The remaining weight of stored formal legs]\label{C-407-41}
Let $s\in\mathcal P$ and let a real evaluation configuration
satisfy the regular-background hypotheses for $h\in\mathcal H_k$.
Use the original background constant $C_H$, before its enlargement
for complex mass matrices, and set
\[
 c_{\mathsf M}=2600\nu_P/\alpha_D,\qquad
 C_{\rm bl}=\nu'\bigl(2+4/(L-1)\bigr),\qquad
 C_{\rm fl}=4C_{\rm bl}c_{\mathsf M}C_H.
\]
Assume $C_{\rm fl}r_M\le L$. For a row $x$ of level $i$ in a
stored layer $\mathcal L_m$, $i<m<k$, the background map at
$\mathsf n=-a_KJ$ satisfies
\[
 \sum_P|\mathcal H_{(i)}(x;P)|
 \le\bigl(\tfrac12+2^{-(k-m)-1}\bigr)C_HL^{-(m-i)}.
\]
The row twin has the same bound. Hence a formal leg retained at
the weight of Prop.~\ref{C-399-15} costs at most one when
it is evaluated in ${\rm Gr}_h$ at weight $h_F$.
\end{lemma}
\begin{proof}
At zero fine source the row sum is at most
$\tfrac12 C_HL^{-(m-i)}$. To estimate its change, put
$\mathsf s=\mathsf s_k$ and $d=d_h$. The fine background
column has the form $\psi_P=\mathcal K_k(\mathsf n)^{-1}g_P$
on $A_0$, where $g_P$ is independent of $\mathsf n$; it is
unchanged on $S_0$. The resolvent identity gives
\[
 \psi_P(\mathsf n)-\psi_P(0)
   =-\mathcal K_k(\mathsf n)^{-1}\hat{\mathsf n}\psi_P(0).
\]
The kernel bound \eqref{eq:qcd-fine-resolvent-kernel} and the
zero-source background row sum of \cref{C-397-33} imply
\[
 \sum_P|\psi_P(\mathsf n)(v)-\psi_P(0)(v)|
 \le2C_H^2\|\mathsf n\|\mathsf s(x_v)^{-1}
 \sum_{v'}(1+|x_v-x_{v'}|)^{-2}e^{-\alpha_Dd(x_v,x_{v'})/4}.
\]
The last sum is at most $c_{\mathsf M}\mathsf s(x_v)$.
Indeed, let $\rho=R\mathsf s(x_v)+1\le2R\mathsf s(x_v)$
and $t=|x_v-x_{v'}|_1$. The layered metric gives
$d\ge R\log(1+t/\rho)$. For $t\le\rho$ this is at least
$t/(4\mathsf s(x_v))$; for $t>\rho$ the exponential is
bounded by $(1+t/\rho)^{-p}$ with $p=\alpha_DR/4\ge8$.
There are at most $6(1+t)^2$ sites at distance $t$, each with
$\nu_P$ places. The two sums are bounded by
\[
 54\nu_P\left(\frac{32\mathsf s(x_v)}{\alpha_D}
              +\frac{16\mathsf s(x_v)}{\alpha_D}\right)
 \le c_{\mathsf M}\mathsf s(x_v).
\]

At level $i\ge1$ the closed Schur formula expresses the field
as $q_i\cdots q_1\psi+\theta_i^{-1}\mathsf e_i$. Only $\psi$
changes with $\mathsf n$. The composed averages contribute
$2\sqrt{48}L^i$, and the part of $\mathsf e_i$ depending on $\psi$
contributes $4\sqrt{48}L^i/(L-1)$. The factor $\sqrt{48}$
comes from $\iota^*$ in the per-chamber supremum norm.
Since $\sqrt{48}<\nu'=8$, their sum is at most
$C_{\rm bl}L^i$; the averages act trivially in flavour.
The change of the level
$i$ row is consequently at most
$2C_{\rm bl}c_{\mathsf M}C_H^2L^i\|\mathsf n\|$.

At a row in $\mathcal L_m$ the physical source bound gives
\[
 L^m\|\mathsf n\|=\lambda_m\|J\|_\infty
 \le L^{-(k-m)}r_M/(2L).
\]
The change is therefore at most
$\tfrac14L^{-(k-m)}C_HL^{-(m-i)}$, which is bounded by
$2^{-(k-m)-1}C_HL^{-(m-i)}$. Add the zero-source half.
The row twin is estimated in the same way. The stored weights
do not increase with the time of evaluation, so the ratio of
the resulting row weight to the weight at which a leg was
retained is at most one.
\end{proof}

\subsection{The history induction with complex sources}

At a complex source the scalar coefficient of a history is no
longer a positive measure. We therefore carry the identity of
measures and the estimates in all Grassmann degrees through the
induction. Lem.~\ref{C-407-34} then supplies the density and its
integrable norm majorant. The gauge tests and the stored variables
remain those of the unsourced construction.

\begin{definition}[The sourced inductive statement]\label{C-407-36}
For $s\in\mathcal P$, let $\mathsf I^{\rm F}(k)[s]$ mean that every
history through level $k$ has, coefficient by coefficient, the
measure representation
\[
 \mathsf{Stay}_h\,1_h^{\rm T}\chi_k(\Omega_k)
       e^{\boldsymbol\Phi_h^w}\exp A^{\rm F}_{k,h}\,\mathscr M_h,
\]
on a common full-measure set, with Borel density and one value at
all evaluation configurations of each point. The exponent has the
same terms and the same places as the unsourced exponent, evaluated
at the fine matrix mass and the coordinate chain of $h$. More
precisely, it has the following properties.
\begin{enumerate}
\item The exponent of the formed factors is
$\boldsymbol\Phi_h^w=\boldsymbol\Phi_h+\sum_{t,D}K^v_{t,D}$,
where the sum is over the kept pieces,
$|K^v_{t,D}|\le\gamma_{\rm kp}S_{t,D}/2$. Its quadratic bound
holds with $\gamma^\dagger=\gamma_{\rm kp}/2$, and its
formed-factor bound with $(\mathsf B_{\rm fr}^\flat+1,\kappa)$.
\item The action has running weight $\beta_h^F-a_K^{-1}w$.
Its difference from the fixed weight $\beta-a_K^{-1}w$ is
the source-independent layer term in
\eqref{eq:qcd-history-format}, already included in the scalar
bound. For history integration use
$\mathsf{bw}=\beta-a_K^{-1}\operatorname{Re}w
\ge(1-\mathsf a)\beta$, where the fixed loss
$\mathsf a\in(0,1/2]$ satisfies the source ceiling
$4\varrho\le\mathsf a(\gamma^{-2}-\bar B)$.
\item The stored quadratic form is $\mathbb D_h(\mathsf n)$.
Its row and column sums are bounded by $C_D^{\Sigma,\rm s}$,
and its exponential by $e^{\mathsf c_G\Sigma'(h)}$, with
$\max\{\mathsf c_{\rm tw},C_D^{\Sigma,\rm s}\}$ in
$\mathsf c_G$.
\item The mass profile is the sourced history profile of
Lem.~\ref{C-407-21}; it and the regular sum with legs have the
bounds of Lem.~\ref{C-407-35}.
\item The regular activities are those of
Prop.~\ref{C-407-16} at $\mathfrak c^h$. Extraction keeps
each moved form on its actual polymer $Z_p$, with its own
pointer. The pointed sum rule holds, the source-step sizes
hold, and the scalar entries and their edge parts satisfy
Lem.~\ref{C-407-40}.
\item The made boundary terms and the edge parts satisfy the
unsourced boundary estimates with $\mathsf B_\natural$ and
$E_0^{\rm s}$ in place of $\mathsf B$ and $E_0$. No reality
is required at a complex source. They include the rejected
halves, the reblocked linear term, the sourced strip rows and
the forms pointed on the rims. Their dependence on the history
has the coordinate tails of Lem.~\ref{C-407-21}. At step $j$
the interior leg size obeys
$\bar v_j^{\rm s}\le b^{\rm blk,s}\le c_{\rm cl}/8$,
and the frame leg size obeys
$\mathsf{sz}_j^\natural\le(64/51)b^{\rm tot,s}$.
The rim output with legs is bounded by
$K^{\rm lb}\mathsf{sz}_j^\natural e^{-\kappa''d_j(X)}$;
the sourced strip rows add at most
$\eta_j^{\rm sr}K^{\rm lb}\mathsf{sz}_j^\natural
e^{-\kappa_+d_j(X)}$, where $\eta_j^{\rm sr}<1/8$.
Formal legs in stored layers have the cost at most one of
Lem.~\ref{C-407-41}.
\item The vacuum constant of the label satisfies
\[
 -\mathcal E^{\rm F}_{k,\mathcal S}
   =E_k(s)-\sum_{n\le k}f_n(\mathcal S)
                    +\sum_{n\le k}g_n(\mathcal S;s),
 \qquad |g_n(\mathcal S;s)|\le C'_{\mathfrak E}
                                     \#_n\widetilde Z_n.
\]
Each $g_n$ is holomorphic and vanishes at zero source.
\item In all degrees the sum of the mass, regular and boundary
terms and the action counterterms is at most
$(C_E^{\rm F}-\mathsf c_G)\Sigma'(h)$. The constant includes
the scalar scale sum, scalar edge terms and all rim forms
specified above.
\end{enumerate}
At level zero this means the bare density with no activity,
formed factor or vacuum constant.
\end{definition}

The next lemma identifies the expansion used at one step. Its
estimates use only the bounds at previous steps. The estimates
for the newly generated boundary terms, recorded in the proof of
Prop.~\ref{C-407-22}, are one-step estimates under precisely
that hypothesis.

\begin{lemma}[The centered frame at a source]\label{C-407-38}
Use the common constant choices of \cref{C-396-12}, including
the source ceilings. Suppose
$\mathsf I^{\rm F}(k)[s]$ holds, $s\in\mathcal P$, and consider a
step along $h'=(h,\mathsf o)$ with $\Omega_1(h')\ne\emptyset$.
Set $t=k+1$, $\tau=x_k/\mathsf t_k$, $g_k=\mathsf t_k^{-1/2}$,
and $v'=\tau v_k=-w/\mathsf b$. The quadratic weighted term,
the linear term and the cross term, read at the integrated
fluctuation alone, form one centered class
\[
 \mathbb W^{\rm ce}=\mathbb W^{\rm q}
                         +\mathbb W^{\rm l}+\mathbb W^\times.
\]
They are analytic and Borel on the tiers, jointly gauge invariant,
local on their assigned polymers and complex linear in $w$;
conjugation holds at real $w$. At the rate
$r_V=(1-2\delta_{\rm sp})(\kappa-2)$, in the centered radial norm,
\[
\begin{split}
 \|\mathbb W^{\rm q}+\mathbb W^{\rm l}\|^{\rm rad}_{{\rm I}\star,r_V}
 &\le\mathsf k_\odot(\sigma_s+\varrho'_sg_k)
                         \le\tfrac12\mathsf m_\odot,\\
 \|\mathbb W^\times\|^{\rm rad}_{{\rm I}\star,r_V}
 &\le C_{\rm cr}^{\rm V}\mathsf C_{\mathcal A}A_1\varrho
           \mathsf t_k^{-1}(\log\mathsf t_k)^{\mathfrak p_0}
                         \le\tfrac14c_M^{\rm O},
\end{split}
\]
where $\sigma_s=2\varrho/(\mathsf b-\bar B)$,
$\varrho'_s=\sqrt3\varrho/20$, and
$\mathsf m_\odot\le c_M^{\rm O}=w_0^{\rm O}\varepsilon^{\rm O}$.
At the pair $(\mathsf t_k,x_k)$ their centered size is at most
$c_M^{\rm O}\le w_0^{\rm O}$.

The remaining classes are of sort $\Sigma_{\rm v}$. Their sizes
on the tiers satisfy
\[
 \mathsf{sz}^{\rm rg}\le\varepsilon_\flat,\quad
 \mathsf{sz}^{\rm bd}\le\tfrac12\varepsilon_\flat,\quad
 \mathsf{sz}^{\rm q}\le\varepsilon_\flat,\quad
 \mathsf{sz}^\natural\le\tfrac{64}{51}b^{\rm tot,s}
                         \le\tfrac{16}{17}\varepsilon_\flat.
\]
At the first step the boundary class is absent and the regular
class together with the additional scalar matter terms has size
at most $3\varepsilon_\flat/2$. The centered expansion exists,
is holomorphic on the interior of $\mathcal P$, and has
\[
 \varepsilon_{\mathsf w}^{\mathfrak F+}L^3
       \le\frac5{32\mathsf K_{
       ap}^+},\qquad
 \mathsf b_{\rm cs}^{\mathfrak F+}\le\frac9{32}.
\]
Its cluster sums have the locality, invariance, tie and
measurability properties of the joined expansion. On a domain
$X\not\subset\Lambda_t$ its output after subtracting the
weighted Gaussian part is at most
$(9/16)e^{-\kappa''d_t(X)}$; its part with legs is at most
$K^{\rm lb}\mathsf{sz}^\natural e^{-\kappa''d_t(X)}$.
Inside $\Lambda_t$ the expansion agrees term by term with
Prop.~\ref{C-407-16}, including its pointed outputs.
\end{lemma}
\begin{proof}
The weighted quadratic and linear terms are the centered
Taylor members of the weighted action. The subtraction at a
base plaquette removes the constant weight in the linear and
cross terms. The bounds for the cross term use its conditional
mean chains, detailed below. Each chain is assigned to the hull
of its footprint and reads the stored field only at its final
bond. Thus it is local, and its assignment is independent of
the fields. The Taylor terms and these chains are analytic,
Borel and gauge covariant, so their assigned sums have the
stated properties. The source ceilings give the two displayed
norm bounds. Their sum is at most $3c_M^{\rm O}/4$ in the
calibration frame. On passage to the pair the quadratic-growth
factor is at most $\max\{1,\tau^{-1}\}\le4/3$, proving the
centered-class bound.

The regular scalar class is the ordinary regular class with
the pointed entries at $\mathfrak c^h$, the action shares removed
with the same $b_j$, and the counterterms zero. The plaquette
weight made from the history is unchanged: it reads the regions
and the numbers $b_j$, not the external source. The boundary
scalar class includes the sourced made terms and formed factors.
One must apply their estimates origin by origin. The box theorem
requiring real entries cannot itself be applied to complex
entries. Its constituent estimates require only analyticity,
invariance, the tie and the bounds on entries, made terms and
formed factors. These hold here with
$(E_0^{\rm s},\mathsf B_\natural,\mathsf B_{\rm fr}^\flat+1)$,
by Lem.~\ref{C-407-40} and the inductive statement. In
particular an edge part has at most
$M^3\#_{\pi_j}X\le8M^3e^{d_j(X)}$ pointers. Its scalar part
is therefore at most $8M^3E_0^{\rm s}e^{-\kappa d_j(X)}$,
which is covered by $\mathsf B_\natural$.

The two scalar class estimates, at the calibration couplings,
are
\[
 \max\{1,\tau\}C_{\rm eff}
  \bigl(\mathsf s_k^\sharp[E_0^{\rm s}]
                   +\vartheta_{0,k}g_k^2y_k^{2q}\bigr),
 \qquad
 C_{\rm eff}(\mathsf B_\natural+
                     \mathsf B_{\rm fr}^\flat+1)\vartheta_{0,k}.
\]
Since $\max\{1,\tau\}\le5/4$ and
$\mathsf s_k^\sharp[E]$ increases with $E$, the doubled-box
ceiling makes each at most $\varepsilon_\flat/2$. The weighted
Taylor terms of order at least three belong to the regular
class. Their norm is bounded by
\[
 \tfrac12\max\{\mathsf K^0(r_P^\circ)^{-3},
                         \mathsf K_0r_0^{-3}\}
 (\mathsf C_{\mathcal A}+\theta_{\rm sl})^3 A_1^3
       \gamma(1+2\log\gamma^{-1})^{3\mathfrak p_0}
       \le\tfrac12\varepsilon_\flat.
\]
This follows from the third-order Taylor remainder on the
original tiers and $\sup|v'|\le1/8$. The first step has the
same estimate on its block-local domains. This proves the
regular and boundary bounds. The quadratic class keeps its
unsourced bound under the pair comparison. The leg class is
bounded by the one-step recurrence below: it reads earlier
frame sizes only, and gives $(64/51)b^{\rm tot,s}$ at the
new step.

The boundary members at a complex source have the same
analyticity on the tiers as in the proof of
Prop.~\ref{C-399-15}. Their coefficients are
holomorphic on their filing slices; the background
maps at $m_0+\widehat{\mathsf n}$ have the bounds and
holomorphy of Lem.~\ref{C-397-29} and the localized
operator construction; and the shift,
substituted triples and levelling functions read no
source. The normalization and composition argument
therefore applies unchanged. In particular it uses
the admitted slices for a made boundary coefficient.
It does not require an extension of that coefficient
to every complex link configuration.

The coefficient and leg pieces of a made term are
recombined at each polymer and each choice of leg parts
as the bracket of $\Pi_{j,X}[\omega]$. This gives the
pair dependence of the member of the frame. Edge
pieces have that dependence individually. Their
products are Borel by the coefficient, background and
shift formulas, and uniform convergence preserves
Borel dependence in the walk sums. Thus the source
frame has the pair, analyticity and measurability
properties required by the joined expansion, as well
as the sizes just proved.

The noncentered sizes sum to less than $6\varepsilon_\flat$.
Using $\varepsilon_\flat=(64eL^3\mathsf K_{\rm ap}^{\max})^{-1}$
and $\varepsilon^{\rm O}=(16eL^3\mathsf K_{\rm ap}^{+})^{-1}$,
the amplitude is at most
$3/(32\mathsf K_{\rm ap}^+)+1/(16\mathsf K_{\rm ap}^+)
=5/(32\mathsf K_{\rm ap}^+)$. The large-field contribution
is at most $1/(8\mathsf K_{\rm ap}^{\max})$, so the cluster
bound is $9/32$. The canonical source-step hypotheses and the
pair comparison give the remaining hypotheses of the centered
expansion, including those of its first-step version.
Lem.~\ref{C-399-41} applies this expansion in the even
Grassmann algebra. Its majorant reads the sizes just obtained;
the positive Gaussian, its tests and its admissible centers
are gauge objects. It supplies the cluster properties and the
bound $9/32$ for both the full frame and its weighted Gaussian
subframe. Subtracting them gives $9/16$. Marking the class
with legs gives the stated $K^{\rm lb}$ bound.

On an interior domain the conditional mean is zero, no external
field is read, the boundary class vanishes, and the cross term
has no polymer. The centered quadratic and the ordinary
quadratic class combine into the weighted quadratic origin;
the linear member joins the scalar potential. The origin-joining
identity of Prop.~\ref{C-388-34} and the interior
identification of Lem.~\ref{C-399-45} give the asserted
termwise equality, also at every pointer.
\end{proof}

\begin{proposition}[Induction at a complex source]\label{C-407-37}
With the common constant choices of \cref{C-396-12},
$\mathsf I^{\rm F}(0)[s]$ holds for
every $s\in\mathcal P$, and
$\mathsf I^{\rm F}(k)[s]\Longrightarrow
\mathsf I^{\rm F}(k+1)[s]$ for $0\le k<K$.
\end{proposition}
\begin{proof}
At level zero the exponent is the weighted Wilson action and
the fermion form at $m_0-a_K\widehat J$. Its coefficients are
polynomials in the links, and its even exponential is a Borel
density. If $\Omega_1(h')$ is empty, every later region is empty
and no variable is integrated. The term is simply the image of
the tested bare density. All profiles vanish, there are no
regular or boundary terms, and the label constant cancels the
regular vacuum number. The stored-form bound and the bound on
that number prove the assertion in this case.

Otherwise fix $h'=(h,\mathsf o)$. The gauge change of variables,
tests and gauge-fixing factors are unchanged. The exact Schur
identities at the matrix mass give the partial fermion integral.
The background change and completion of the square use the
sourced operators, with their doubled bounds. Stored generators
are kept at their old places; the other legs acquire the
canonical levelling factors. The boundary coefficient estimate
of Lem.~\ref{C-399-48} applies with $\mathsf B_\natural$,
since it uses analyticity and bounds without a reality condition.
The normalizer is $\delta Z_k^{\rm F}$ at the matrix source.

Lem.~\ref{C-407-38} gives the joined centered expansion, with
the first-step version when $k=0$. The first fermion covariance
is defined at every configuration by its positive real-part
bound, so no regularity hypothesis is added at that step.
Inside $\Lambda_{k+1}$ the history coordinates lie in the
doubled source window, and the expansion is the source step of
Prop.~\ref{C-407-16}. Take out its vacuum number, action
coefficient, source energy, and local mass increment. The
increment uses only pieces with $Z_p\subset\Lambda_{k+1}$.
Its forms stay on these $Z_p$: interior pointers enter the
regular sum, rim pointers enter the edge parts. Reinsert the
mass through the new profile; its complement remains on the
rim. At each pointer the extracted form, its reinserted part
and its remainder sum to zero.

We spell out the measure identity because it is independent of
positivity at a complex source. Transport of the weighted
Wilson action over the bare plaquettes gives the sourced
change-of-variables identity. The Taylor decomposition at the
conditional mean is exactly the sum of the weighted quadratic,
linear and cross terms, the term on stored fluctuations, and
the terms of order at least three. Split the stored quadratic
term into its kept and rejected halves. The product of this
splitting identity and the joined expansion cancels the
preintegral value at the conditional mean. The scalar linear
term evaluated at the stored field leaves the integral as an
exponential and is reblocked into the new boundary term. The
sum rule for the sourced strip rows is the one in
Lem.~\ref{C-399-42}. These are the sourced counterparts of
the identities in Prop.~\ref{C-399-43}.

All operations are coefficientwise in the finite-dimensional
even algebra. In particular the conditioning integral is
followed by the Gaussian Berezin integral with the symmetrized
sourced decoupling covariance of Prop.~\ref{C-397-40}.
Its conjugation maps the source to $\gamma_4J^\dagger\gamma_4$,
and its charge conjugation to $J^{\mathsf c}$; it is not a
positivity assertion at a fixed complex source. For an
interchange of integrals, first use Tonelli for the fibre norm,
then Fubini for each coefficient. The needed integrability and
common null set are supplied by Lem.~\ref{C-407-34}. Wherever
the scalar argument uses a nonvanishing exponential, use the
invertibility of an even exponential. Its scalar coefficient
is itself a nonvanishing complex exponential.

The transported source exponent, the canonical shifts and the
two halves are Borel. Activity and boundary coefficients are
locally uniform limits of finite sums of Borel coefficients
and localized covariances. The coordinate fields $\sigma_n^h$
depend on the source, not on the configuration. The new density
is therefore Borel. Evaluation configurations of a point
differ by a gauge transformation. The weighted action is
invariant plaquette by plaquette, the halves are covariant,
and the matrix mass and increments commute with colour. This
proves the one-value assertion.

The remaining task is quantitative. The kept-half and store
bounds are proved below; Lem.~\ref{C-407-35} bounds the mass
and regular terms with legs; Lem.~\ref{C-407-40} bounds the
scalar entries and their scale sum; and
Lem.~\ref{C-407-41} pays the stored formal legs. The one-step
boundary recurrence below uses only preceding frame sizes and
preserves $(64/51)b^{\rm tot,s}$. The rim outputs of
Lem.~\ref{C-407-38}, the sourced strip bounds and the pointed
extraction just described give the boundary clause. Finally
the local energy comparison of Lem.~\ref{C-407-21} gives the
label errors. These estimates give each clause of
Def.~\ref{C-407-36} at level $k+1$ and complete the
induction.
\end{proof}

\begin{proposition}[All histories at the sources]\label{C-407-22}
For $s\in\mathcal P$ the sourced history measures are
complex measures and their sum is the exact bare sourced
measure after the block maps. Uniformly for $k\le K$,
\[
 \sum_h e^{\Theta^{(k)}(h)-\operatorname{Re}E_k(s)}
       \int\|g_h(s)\|\,d\mathscr M_h^{\rm fr}
 \le e^{(C+2)n_k}.
\]
Thus $Z_k^{\rm sub}$ is holomorphic on the interior of
$\mathcal P$, continuous on it, and
\[
 |Z_k^{\rm sub}[s]|\le e^{(C+2)n_k},\qquad
 Z_K^{\rm sub}[0]\ge e^{-C_P(1+\log x_K)\ell^3}.
\]
The factorization is $Z[s]=e^{E_k(s)}Z_k^{\rm sub}[s]$,
and $Z_k^{\rm sub}[s^\star]=\overline{Z_k^{\rm sub}[s]}$.
\end{proposition}
\begin{proof}
Prop.~\ref{C-407-37} gives the complex measure identity
and the full-degree form of every history. The history chain
remains in the
doubled coordinate window by Lem.~\ref{C-407-21}.
The physical source ball lies in the fixed-radius
domains, so all local norms needed by the histories
are the fixed-radius norms of the source step.

We verify the additional scalar payments. Set
$a=\frac12(1-c_\Theta/c_*)\in(0,1/2]$ and impose
$4\varrho\le a(\gamma^{-2}-\bar B)$.
Then the real part of the sourced Wilson weight is at
least $(1-a)$ times the unsourced weight.
The source quadratic on stored fields is decomposed by
the pure weighted-form construction into kept
and rejected halves. Its kept part obeys
$|K_{t,D}^v|\le C_v^{\rm kp}\sup|v|S_{t,D}$.
Impose
$4C_v^{\rm kp}\varrho\le\gamma_{\rm kp}(\gamma^{-2}-\bar B)$.
The retained real quadratic lower bound then remains
at least $\gamma_{\rm kp}S_{t,D}/2$. Its rejected part
is a bounded scalar boundary family. This application
uses only the supremum of the weight, the original
positive retained forms and the original tiers.
More precisely, at later steps the local source forms
obey $|q_Y(P;B,B)|\le k_\upsilon e^{-r_\upsilon d(Y)}
\sup|B|^2$, with $k_\upsilon$ proportional to $\sup|v|$.
Use the original fixed-real-factor analytic slices and
their graded-to-layered restrictions. With
$\beta_\upsilon=8C_Ee^{6\delta_\upsilon M}k_\upsilon$,
$r_\upsilon-\delta_\upsilon M\ge c_E$,
$W_t\ge LM$ and
$\delta_\upsilon LM/2\ge342(\kappa'+1)$,
\cite[Thm.~AG and Prop.~79.3]{YMM1} give
\[
 \|\widehat\Upsilon\|\le C_\Upsilon\beta_\upsilon,
 \quad \operatorname{osc}\widehat\varphi_X
 \le b_{\rm fr}^{\rm s}\mathfrak b_k^2e^{-\kappa'd(X)},
 \quad b_{\rm fr}^{\rm s}
 =c_\upsilon^2B_f\beta_\upsilon K_{\rm ch}^\sharp
 C_{\rm xt}^2e^{-\kappa'}.
\]
Choose the source ceiling so that $b_{\rm fr}^{\rm s}\le1/4$.
The first step has the direct block-local quadratic bound
and the analogous coefficient at most one. These are
retained formed factors, with their original quadratic
tier factor, and not small fluctuation activities.
Their later Cauchy variation is precisely the variation
of the original formed-factor family with this enlarged
fixed coefficient. The source forms meet these hypotheses
by their quadratic Taylor bound and original kernel walk
rates; their norms are linear in $\sup|v|$.

The linear source term evaluated on the stored field
and mean leaves a collar term of size at most
\[
 C_{\rm lin}\max\{\alpha_0,12L^6\}\varrho
 (1+2c_{\mathsf M}^{\rm w})C_{\rm xt}A_1
 \mathsf t_j^{-1/2}(\log\mathsf t_j)^{\mathfrak p_0}.
\]
It is below $c_{\rm cl}/16$ at the final floor. For the
source cross term use its weight-difference continuation,
\[
 X^v=-\tau\sum_b\sum_{Y\ni b}
                  F_{Y,2}^{v-v(p_b)}(\widetilde z_b,F_*).
\]
Its constant-weight part is zero at the physical center
and is not continued. Expand $m=-K_1a'$ by the
conditional-mean resolvent and leg chains. A chain
$(b'',\jmath,b_1,S,b')$ reads the pair on its footprint
and $a'$ at a stored bond $b'$. File it at $Y$ joined
to its footprint hull, which meets the collar.
If its tree length is $l$, that hull contains at most
$125L^3+250L^2l/M$ fine cubes. The chain exponential
therefore pays its enlargement and polymer weight under
the existing walk-rate floors. The swing
$x_j\sup_Y|v-v(p_b)|\le2\#_{\mathcal P_j}Y\sqrt3M\operatorname{Lip}w$
and the stored tier bound then give total size
\[
 C_{\rm cr}\mathsf C_{\mathcal A}A_1\varrho
       \mathsf t_j^{-1}(\log\mathsf t_j)^{\mathfrak p_0}
 \le c_{\rm cl}/16.
\]
The rejected half has the separate bound
\[
 C\mathfrak b(\mathsf t_{t-1})^2
 e^{-(\check R_t-2)}e^{-\kappa''d_t(X)}.
\]
Since $\check R_t\ge(\log\mathsf t_t)^{r_c}$ and
$r_c\ge2$, the threshold-square times this suppression
is uniformly bounded. Its entropy sum therefore costs
a fixed number per rim cube. This, rather than the
tier-weighted retained oscillation bound alone, pays
the rejected halves. The retained exponents keep their
original formed-factor norm with the original tiers.
The linear term left outside the integral is reblocked from
$X$ to its coarse hull $X^{\rm up}$. The inequalities
$d_{k+1}(X^{\rm up})\le L^{-1}d_k(X)$ and
$\lambda_V\ge L^{-1}\kappa''+2c_E$ give
\[
 e^{-\lambda_Vd_k(X)}
 \le e^{-\kappa''d_{k+1}(X^{\rm up})}e^{-c_Ed_k(X)}.
\]
The hull count is at most $(L+2)^3S_{\rm s}$. The reblocked
family is therefore bounded by
\[
 K_{\rm rb}c_{\rm cl}e^{-\kappa''d_{k+1}(X')}/16,
 \qquad K_{\rm rb}=(L+2)^3S_{\rm s}e^{4\lambda_1}.
\]
It reads the pair and stored fluctuation on $X'$ alone and
vanishes unless $X'$ meets the collar.

Stripping produces additional boundary rows. Those from the
two halves have their quadratic bound times
$C_{\rm fl}^{\rm tr}e^{-\check R_t}$, uniformly bounded because
$\mathfrak b(\mathsf t_{t-1})^2e^{-\check R_t}
\le e^{-2}C_{\mathfrak b}$. Apply Lem.~\ref{C-399-42} to the
other two families, the rim cluster sums and the reblocked
linear term. They are local, jointly gauge invariant and
analytic on its complex box; their bounds there are respectively
$9/32$ and $K_{\rm rb}c_{\rm cl}/16$, at rate $\kappa''$.
The rows sum to the difference between the field before the
strip and the stored field and vanish away from the strip.
Their extra factor is
\[
 \eta_t^{\rm sr}=2^{14}e^{-\check R_t}
       \le2^{14}e^{-324}<\frac1{16}.
\]
In positive degree the rim family and its rows are thus bounded
by $(9/8)K^{\rm lb}\mathsf{sz}_t^\natural$.
All scalar additions just listed are counted in
$\mathsf B_\natural$ and $C_E^{\rm F}$.

The leg recurrence uses only preceding frame sizes. Write
$\Theta_k^{\rm s}=c''_\psi+\Theta_k^{\prime\rm s}$ for the
split-leg factor at new and stored generators, respectively;
$\Theta_k^{\prime\rm s}\le2\cdot10^{-3}$. At lags below
$n_{\rm lag}$ keep each whole pointed activity. At longer
lags use the units, moved pieces and forms constructed above.
A piece agrees with its one-lattice counterpart only if its
actual polymer lies in the regular core. At a cell of a
stored layer $\mathcal L_m$, take differences only for an
unchanged piece contained in $\Omega_{m-1}\setminus\Omega_{m+2}$.
Every other piece is taken whole.

The difference pieces have the cell-sum bound of
Lem.~\ref{C-407-35}. A whole piece reaching layer $m>j$ has
$d_j(Z)\ge7L^{m-j}-2$, so its cells and the unused decay are
paid by
\[
 c_{\rm rim}^{\rm s}=c_{\rm geo}(1+L^2)
       \sum_{d\ge1}L^{3d}e^{-\kappa(1-L^{-1})(7L^d-2)}.
\]
The two families of moved pieces cost $1+2C_{\rm pw}$.
For a split stored leg the factor $(2L^{-4})^{k-m}$ combines
with the cell count to give $(2/L)^{k-m}$, whose sum is at
most $L/(L-2)$. For degree $2p\ge4$ each unsplit leg keeps
at most three quarters of its weight; telescoping and the
walk count cost at most
\[
 2p(1+2pn_{\rm wk}^2)
      (3/4+\Theta_k^{\prime\rm s}/4)^{2p-1}
       \le8.6n_{\rm wk}^2\le c_{\rm geo}.
\]
The torus cell sums are filed on wrapping polymers with
$d_j(Z)\ge L^{k-j}-2$ and use the convergent bound
$\sum_{d\ge1}L^{3d}e^{-\kappa(1-L^{-1})(L^d-2)}\le1$.
Formal stored legs cost at most one by Lem.~\ref{C-407-41}.

These estimates give the stored-layer share
\[
\begin{split}
 b_k^{\rm sl,s}={}&C_{\rm lay}^{\rm s}\frac L{L-2}
 \left[4Lc_\psi^{\rm s}
 +\frac{2\mathsf C_{\mathcal A}\delta(\mathsf t_k)}{\beta_0}
  \bigl(9L^4A_\kappa C_H+
           (2+6L^5A_\kappa)C_{\mathfrak r}\beta_0\bigr)\right]\\
 &\hspace{7mm}\cdot
 \bigl(1+w_0\nu C_H^2h_F^2L(\bar B_M+2B_\sigma)\bigr)\phi_k.
\end{split}
\]
Besides the interior share $b_k^{\rm blk,s}$, the new terms
have the shares
\[
\begin{split}
 b_k^{\rm lay,s}={}&2C_V^{\rm F,s}h_F^2\mathsf C_{\mathcal A}
                         \delta(\mathsf t_k)\\
 &+\Theta_k^{\prime\rm s}
   \bigl[(1+c_{\rm cs})2C_{\rm pc}S_{1/4}C_w^{\rm s}
                      +C_0^{\rm cs}h_F^2\bigr]\phi_k\\
 &+\Theta_k^{\rm s}\bigl(1+C_\omega(1+2C_{\rm pw})\bigr)
                     c_{\rm rim}^{\rm s}2C_w^{\rm s}\phi_k
       +\Theta_k^{\rm s}2C_{\rm pc}S_{1/4}C_w^{\rm s}\phi_k\\
 &+\frac L{L-2}\Theta_k^{\prime\rm s}C_\omega(1+2C_{\rm pw})
   \bigl(2c_{\rm geo}(1+L^2)+c_{\rm rim}^{\rm s}\bigr)
                     2C_w^{\rm s}\phi_k+b_k^{\rm sl,s},\\
 b_k^{\rm ms,s}={}&C_{\rm bd}^{\rm s}
       \left(c_\psi^{\rm s}
             +\frac{2\mathsf C_{\mathcal A}\delta(\mathsf t_k)}{\beta_0}\right)
  \bigg[3M^3w_0\nu C_H^2h_F^2
       \left(\bar B_M\gamma^2+\frac{B_\nu}6
                                   \phi(\tfrac12\gamma^{-2})\right)\\
 &\hspace{25mm}
     +2\eta_{\min}^{-2}c_{\mathsf S}M^{-3}C_{\rm pw}C_w^{\rm s}
                                   \phi(\tfrac12\gamma^{-2})\bigg].
\end{split}
\]
The mass share includes the rim complement with coefficient
$L^{-j}[\varphi_j\mu_j-(1-\varphi_j)(\nu_j^h-\mu_j)]$,
bounded by $L^{-j}(|\mu_j|+B_\nu\phi_j/2)$, and the forms
$P_{z;X,p}$ on their actual $Z_p$, of summed norm at most
$2\eta_{\min}^{-2}c_{\mathsf S}M^{-3}C_{\rm pw}C_w^{\rm s}
\phi_{j-1}$. Their coefficients read no configuration; their
reentry is through their legs alone.

For the forms this uses a loss of one unit of decay.
Their known bound sums $\|P_{z;X,p}\|e^{\kappa d_j(Z_p)}$
through a cube. Apply the preceding bracket count at
rate $\kappa-1$, with roots containing the hull
$Z_p^\uparrow$. After the walk factors and the weight
room have been used, its remaining factor is
\[
 e^{(\kappa-1+1/13)d_k(Z_p^\uparrow)}e^{2/13}
   =e^{(\kappa-12/13)d_k(Z_p^\uparrow)}e^{2/13}.
\]
At $j=k$ it is bounded by $e^{\kappa d_j(Z_p)}e^{2/13}$.
At $j<k$, use $d_k(Z_p^\uparrow)\le L^{-(k-j)}d_j(Z_p)$
and
\[
 (\kappa-12/13)L^{-1}+\kappa(1-L^{-1})\le\kappa.
\]
The weight of the form therefore also pays the reach
factor used in the boundary count. Finally
$e^{2/13}<C_E$, already included in $C_{\rm bd}^{\rm s}$.
This proves the displayed contribution of the forms
in the summed norm at rate $\kappa-1\ge r_V$.
The made and ordinary edge terms instead have the
spare rate $\kappa_+$ and use their earlier domain count.
Their coefficient and leg contributions remain
$1/16$ and $9/64$ times the boundary maximum,
respectively; they are recombined before asserting
pair dependence of the frame.

Set $\Sigma_{k+1}^\natural=\sup_{1\le j\le k}
\mathsf{sz}_j^\natural$ and
$\bar v^{\rm s}=\sup_{1\le j\le k}\bar v_j^{\rm s}$, both
zero at $k=0$. The old boundary terms entering through legs
contribute at most $\frac18\frac98
\max\{\Sigma_{k+1}^\natural,\bar v^{\rm s}\}$; their
coefficient variations contribute at most $\frac1{16}$ times
the same maximum, by Lem.~\ref{C-399-48} and the coefficient
ceiling. Thus
\[
 \mathsf{sz}_{k+1}^\natural
 \le\frac{13}{64}\max\{\Sigma_{k+1}^\natural,\bar v^{\rm s}\}
       +b_k^{\rm tot,s},\qquad
 b_k^{\rm tot,s}=b_k^{\rm blk,s}+b_k^{\rm lay,s}+b_k^{\rm ms,s}.
\]
Every share is bounded by its value at $\gamma$, and
$\bar v^{\rm s}\le b^{\rm blk,s}\le b^{\rm tot,s}$.
Since $(13/64)(64/51)+1=64/51$, induction preserves
$(64/51)b^{\rm tot,s}$. Each of the three share ceilings is
$\varepsilon_\flat/4$, giving the bound
$16\varepsilon_\flat/17$ for the leg class. This verifies
the one-step estimate used in Lem.~\ref{C-407-38}.

The complex stored tower has uniformly bounded row and
column sums by its closed Schur form and the weighted
inverse at $\mathsf n$; its fine-store entries remain
of range one. Lem.~\ref{C-399-17} therefore bounds its
exponential with a fixed enlarged constant. The stored
mass profile is the convex interpolation of
$L^{-n}(s_n+\sigma_n^h)$ and
$L^{-(n-1)}(s_{n-1}+\sigma_{n-1}^h)$; its norm per
layer cube is bounded by the source window. The source
cell-unit estimate supplies the needed power counting.
For the component integral keep the actual allocation of the
piece stocks. At zero sources its majorant retains
$e^{-\gamma_{\rm kp}S_{t,D}/4}e^{-P^G_{t,D}/2}$ and
the created cells are paid by
\[
 \sum_{t,D}\left(P^W_{t,D}+\tfrac12P^G_{t,D}
                      +\tfrac14\gamma_{\rm kp}S_{t,D}\right)
       \ge c_*N(h),\qquad
 N(h)=\sum_t p_0(g_t)^2 n_t^{\rm cr}(h).
\]
Here $P^W,P^G$ are the nonnegative Wilson and gauge-fixing
stocks, and $\sum P^W\le\beta A(U)$. The source leaves
$e^{-\gamma_{\rm kp}S_{t,D}/2}$ after the kept halves:
one quarter remains in the same majorant and the other
quarter, together with half the gauge-fixing stock, pays
the created cells. The Wilson stock is at least
$(1-a)P^W$. Since $1-a\le1$, the payment is at least
$(1-a)c_*N(h)$, with the same integrable piece majorant
as at zero sources. The enlarged bounded terms proved
above only change its fixed cube constant. Consequently
the integrated component bound is
\eqref{eq:qcd-history-interface} with rate $(1-a)c_*$
and $E_k(s)$ in place of $E_k$. This argument does not
require the stronger pointwise loss-$a$ property of
the retained forms. No positivity of the complex density
is asserted.

The extracted constants along $h$ differ from $E_k(s)$
by the missing source polymers, boundary numbers and
the coordinate-chain difference of
Lem.~\ref{C-407-21}. Each costs at most
$C'_{\mathfrak E}\#_j(\widetilde Z_j)$ and so has the
same label-error form as Lem.~\ref{C-399-24}.
The component estimate has created-cell rate
$(1-a)c_*>c_\Theta$. Thus the weight $e^{\Theta^{(k)}(h)}$
is paid by the same created-cell estimate. Apply the pure history summation interface
with this rate and error constant
$C_\flat^F+C'_{\mathfrak E}$, choosing its threshold
before the cutoff. This proves the norm sum.
Berezin integration has norm at most one at $h_F\ge1$,
so it bounds $Z_k^{\rm sub}$. Its lower bound at zero
is Thm.~\ref{C-399-4}. Holomorphy follows from that
of the finite sourced partition function and $E_k(s)$.
\end{proof}

\begin{proposition}[Sources bounded at an intermediate scale]
\label{C-407-43}
Use the common assignment of \cref{C-396-12}, including the
positive-mass ceilings for the lower bound in \cref{C-399-4}.
For $1\le k\le K$ put
\[
 \mathcal P_k^\#=
 \{(w,J):\ \|w\|_\sharp<2\varrho,
                     \ \lambda_k\|J\|_\infty\le r_J/2\}.
\]
The part of $\mathcal P$ with $\|w\|_\sharp<2\varrho$
lies in $\mathcal P_K^\#$, and
$\mathcal P_K^\#\subset\mathcal P_k^\#\subset\mathcal P_j^\#$
for $1\le j\le k$.
The history induction of \cref{C-407-37}, with the format of
\cref{C-407-36}, holds through level $k$ for every
$s\in\mathcal P_k^\#$, with the same constants. In particular,
\[
 \sum_{h\in\mathcal H_k}e^{\Theta^{(k)}(h)}
    \int\|g_h(s)\|\,d\mathscr M_h^{\rm fr}
 \le\exp[\operatorname{Re}E_k(s)+(C+2)n_k].
\]
The factorization and bounds are
\begin{equation}\label{eq:qcd-level-source-partition}
\begin{gathered}
 Z[s]=e^{E_k(s)}Z_k^{\rm sub}[s],\qquad
 |Z_k^{\rm sub}[s]|\le e^{(C+2)n_k},\qquad
 Z_k^{\rm sub}[0]\ge e^{-C_P(1+\log x_k)n_k},\\
 \frac{|Z[s]|}{Z[0]}
 \le\exp\!\left[\operatorname{Re}(E_k(s)-E_k(0))
                +(C+2+C_P(1+\log x_k))n_k\right].
\end{gathered}
\end{equation}
The function $Z_k^{\rm sub}$ is holomorphic on the interior,
continuous on $\mathcal P_k^\#$, and respects source
conjugation. The constant $C$ is independent of the positive
mass floor; $C_P$ has the dependence stated in \cref{C-399-4}.
\end{proposition}
\begin{proof}
The physical source ball has
$\|J\|_\infty\le2\varrho_J\le c_Jm_-\le r_J/2$.
Since $\lambda_j\le\lambda_k$ for $j\le k$, the inclusions
follow. Moreover $\mathcal P_k^\#\subset
\mathscr O_j\cap\mathscr O_j(\alpha_0)$ for $0\le j\le k$.
Thus the local source chains, the history chains and the energies
are defined in their fixed-radius coordinate domains. To extend
the history bound we must also check the operators and stored
terms; the existence of the local chain alone is insufficient.

Write $\mathsf n=-a_KJ$. The source ceilings give
\begin{equation}\label{eq:qcd-level-source-sizes}
\begin{gathered}
 L^j\|\mathsf n\|\le L^{-(k-j)}r_J/2\quad(j\le k),
 \qquad L^{k+1}\|\mathsf n\|\le r_M/2,\\
 \lambda_j\|J\|_\infty\le r_J/2
       =r_M/(2L)\le\varrho_*/4\le\varrho_\sharp/4,
 \qquad m_0+\|\mathsf n\|\le1.
\end{gathered}
\end{equation}
Here $r_J=r_M/L$, $r_M\le L\varrho_*/2$, and
$\varrho_*\le\varrho_\sharp$ are the already chosen radii.
We check each use of $J$ in the induction up to a level $t\le k$.

At a step from $t<k$, the first bound in
\eqref{eq:qcd-level-source-sizes} puts the fine mass matrix
in the source ball of \cref{C-397-29}. The tower, backgrounds,
decoupled covariances and normalizer therefore have their
usual doubled bounds. This includes the first step, where
$a_K\|J\|_\infty\le L^{-k}r_J/2$. Restricting $J$ to a
polymer preserves this ball, so the local conditional
contraction and measure identities also apply. The conjugate
source remains in the same class because
$\|\gamma_4J^\dagger\gamma_4\|=\|J\|$.

For the sum of old terms on a layer of level $j\le t$,
$L^j\|\mathsf n\|\le r_M/(2L)$ is inside the cell-unit
window of \cref{C-407-16}. For a stored layer $m<t$ the
stronger lag estimate
\[
 L^m\|\mathsf n\|\le L^{-(t-m)}r_M/(2L)
\]
retains the gain used in the stored-leg estimate
\cref{C-407-41}. Indeed its change of row sums is at most
$\tfrac14L^{-(t-m)}C_HL^{-(m-i)}$, hence at most
$2^{-(t-m)-1}C_HL^{-(m-i)}$, under the chosen ceiling
$C_{\rm fl}r_M\le L$. The sums over layers consequently use
the same constants as on $\mathcal P$.

The stored tower has the same row and column bounds as well.
On its fine layer the entries are bounded by
$\mathsf c_{\rm disc}+m_0+\|\mathsf n\|
\le\mathsf c_{\rm disc}+1$; on the other layers use the
source operator bounds just checked. The history mass fields
are those of \cref{C-407-21}, with
$\|\sigma_j^h\|\le2B_\sigma\phi_j$. Their layer coefficients
are bounded by
$L^{-(j-1)}(\bar B_M+2B_\sigma)\phi_{j-1}$, as in
the original sourced format.

The normalizer walks and the label errors use the same
coordinate balls. They give the same holomorphic vacuum
constants and the error $C'_{\mathfrak E}\#_j(\widetilde Z_j)$.
Every remaining source estimate reads only $w$: its supremum,
its swing, the kept and rejected quadratic halves, and the
loss of the real action weight. The condition
$\|w\|_\sharp<2\varrho$ is unchanged. The bare sourced
density is entire and has continuous coefficients on the
compact link space, so the measure operations remain valid.

These checks prove the format by induction on $t$, at the
fixed class $\mathcal P_k^\#$. The proof of \cref{C-407-22}
now gives the history bound with the same created-cell rate
and layer constants. Berezin integration gives the upper
partition bound. The lower bound at zero is
\cref{C-399-4}; division proves the last inequality.
Finally $Z_k^{\rm sub}=e^{-E_k}Z$, with $E_k$ holomorphic on
the open set $\mathscr O_k$ containing $\mathcal P_k^\#$.
This proves continuity, holomorphy and conjugation.
\end{proof}

\subsection{Separated energies and the cell bound}

For a physical set $S$ let $\#_j^+(S)$ count the level-$j$
$M$-cubes within distance $M\lambda_j$ of $S$, and put
\[
 \mathscr C_k(S;\delta)=C_{\mathscr C}
       \sum_{j\le k}\#_j^+(S)e^{-\kappa_E\delta/\lambda_j}.
\]
The following estimates are uniform in $K$:
\begin{equation}\label{eq:qcd-source-count}
 \mathscr C_K(T_\ell;\delta)\le C\ell^3(1+\delta^{-3}),
 \qquad
 \mathscr C_K(S;\delta)\le
 C[1+\log_L^+(1/\delta)+(R/\delta)^3]
\end{equation}
if $S$ lies in a ball of radius $R$.
Indeed $\#_j^+(S)\le
\min\{n_j/M^3,27(1+\operatorname{diam}S/(M\lambda_j))^3\}$.
The sums follow by splitting at $\lambda_j=\delta$;
the fine tail is bounded by
$\sum_{u\ge0}L^{3u}e^{-\kappa_E L^u}<\infty$.

\begin{proposition}[Separated source energies]\label{C-407-25}
Let $s_1,\ldots,s_N$ be sources in the ball satisfying
\[
 \operatorname{dist}(\operatorname{supp}s_i,\operatorname{supp}s_l)
 \ge\delta>0\qquad(i\ne l),
\]
and suppose their sum is in the ball. Then
\[
 \left|E_K(\sum_i s_i)-\sum_iE_K(s_i)+(N-1)E_K(0)\right|
 \le\sum_i\mathscr C_K(\operatorname{supp}s_i;\delta).
\]
For a holomorphic source polydisc $s(z)=\sum_i z_is_i$,
with radii $\rho_i$, if $S_1,S_2$ are at distance at least
$d$ and the supports are pairwise disjoint, then
\[
 |\partial_1\cdots\partial_n E_K(s(z))|_{z=0}
 \le2^{n-1}\prod_i\rho_i^{-1}\mathscr C_K(S_1;d).
\]
\end{proposition}
\begin{proof}
For each polymer choose a nearest source support. Its
difference from the sum of the single-source energies
is bounded by the tails from all other supports. Split
the polymer decay and use
$\operatorname{diam}X\le c_{\mathcal D}M\lambda_j(1+d_j(X))$.
At least half the separation from another support must
be paid by that diameter or its distance to $X$.
Taking $\kappa_E\le\min\{\kappa_\sigma/8,
\kappa/(4c_{\mathcal D}M)\}$ leaves the factor
$e^{-\kappa_E\delta/\lambda_j}$ and a summable polymer
weight. Packing the separated supports on fine levels
and using the crude cube count on coarse levels gives
the first estimate.

For the second apply
$\Delta_n=\prod_i(1-P_i)$, where $P_i$ sets $z_i=0$.
Its norm is at most $2^{n-1}$ times any one single-source
difference norm. For each $X$,
$\max\{\operatorname{dist}(X,S_1),\operatorname{dist}(X,S_2)\}
\ge(d-\operatorname{diam}X)/2$.
Splitting its exponential weights gives
$Ce^{-\kappa_Ed/\lambda_j}e^{-\kappa d_j(X)/2}
e^{-\kappa_\sigma\operatorname{dist}(X,S_1)/(2\lambda_j)}$.
Sum through the nearest cube to $S_1$ and apply Cauchy's
estimate in the $n$ variables. This is the stated bound.
\end{proof}

\begin{proposition}[A source in one cell]\label{C-407-30}
Let a board of physical unit cells have its walls through
links, and let $s\in\mathcal P$ be supported in one cell
at distance at least $\delta$ from its walls. For the
$\mathsf N=\ell^3$ products $r$ of wall reflections, let
$r_{\rm lin}s$ be its linear source image, including the
spin lift on $J$, and set
\[
 \lambda^{\mathfrak b}(s)=\mathsf N^{-1}
        \sum_r(E_K(r_{\rm lin}s)-E_K(0)),\qquad
 c_V=C+2+C_P[1+\log(L(g^{-2}+B_\sharp))].
\]
Then $\lambda^{\mathfrak b}$ is holomorphic, and
\begin{equation}\label{eq:qcd-cell-source-bound}
 |Z[s]/Z[0]|\le
 e^{\operatorname{Re}\lambda^{\mathfrak b}(s)+c_V+
                 \mathscr C_K(\operatorname{supp}s;\delta)}.
\end{equation}
For sources in distinct cells the corresponding moment
bound is the product of these one-cell bounds.
\end{proposition}
\begin{proof}
The source exponential is even and lies in its cell
algebra. Reflection maps it to the reflected source
exponential: in the physical basis
$\gamma(\hat n)\Gamma^\dagger\gamma(\hat n)=\Gamma$.
Its fold is thus the exponential of the disjointly
reflected source. The supremum and Lipschitz source
norms are preserved, since the source vanishes in a
strip along each wall. The coordinate-link reflection
positivity and its chessboard inequality apply.

For the folded source, use Prop.~\ref{C-407-22}
and the zero-source floor in Thm.~\ref{C-399-4}.
Take the $\mathsf N$th root. Its normalization cost is
$e^{c_V}$. The separated-energy estimate replaces
$\mathsf N^{-1}(E_K(s^{\rm fold})-E_K(0))$ by the
average of its individual reflected energies at cost
$\mathscr C_K(S;\delta)$. Odd reflection words use the
conjugated source, and the source conjugation changes
the energy to its complex conjugate; therefore its real
part is the real part of the holomorphic linear-image
average displayed above. This gives the cell seminorm
bound. The multiple-cell chessboard inequality proves
the final assertion, and also bounds a single-cell
expectation by its seminorm.
\end{proof}

\begin{proposition}[A base source and a choice of stopping scale]
\label{C-407-44}
Use the source-energy hypotheses of \cref{C-407-23} and the
common constant assignment of \cref{C-396-12}.
\begin{enumerate}[label=\textup{(\roman*)}]
\item Let $s_0,s_0+s_i,s_0+\sum_i s_i\in\mathscr O_k$,
where the supports $S_i$ of the increments $s_i$ are pairwise
at distance at least $\delta>0$. Then
\[
 \left|E_k(s_0+\sum_i s_i)-\sum_iE_k(s_0+s_i)
                         +(n-1)E_k(s_0)\right|
 \le\sum_i\mathscr C_k(S_i;\delta).
\]
For $n\ge2$, suppose instead that the supports are pairwise
disjoint, $\operatorname{dist}(S_1,S_2)\ge d>0$, and
$s_0+\sum_i z_is_i\in\mathscr O_k$ on the polydisc of
radii $\rho_i$. Then
\[
 \left|\partial_1\cdots\partial_n
 E_k(s_0+\sum_i z_is_i)\big|_{z=0}\right|
 \le2^{n-1}\prod_i\rho_i^{-1}\mathscr C_k(S_1;d).
\]
\item Suppose $s_0=(0,J_0)$, $s_i=(0,J_i)$,
$\|J_0\|_\infty\le r_J/2$ and
$\|J_i\|_\infty\le\varepsilon\le r_J/8$, with the
supports $S_i$ pairwise at distance at least $\delta$.
The first bound of \textup{(i)} at $k=K$ holds with
\[
 \mathscr C^\flat(S;\delta)
  =C_{\mathscr C}\sum_{j=1}^K\#_j^+(S)
                e^{-\kappa_E\delta/\lambda_j}\omega_j,
 \quad
 \omega_j=\begin{cases}
 1,&\lambda_j\le\delta,\\
 \lambda_{j-1}/4,&\lambda_j>\delta.
 \end{cases}
\]
For $S$ in a ball of radius $R_S$,
\begin{equation}\label{eq:qcd-bilinear-base-additivity}
 \mathscr C^\flat(S;\delta)
       \le C_{\mathscr C}^\flat[1+(R_S/\delta)^3],
\end{equation}
uniformly in the cutoff and volume.
\item Let a coordinate-link board have cells of physical
side $l=qa_K$, where $q\ge3$ is odd and $\ell/l$ is an
even integer. Write $\mathsf N=(\ell/l)^3$ and let $s$ be
supported in one cell $\square$, at positive distance from
its walls. For the source exponential $G_s$, let $s^{\rm ch}$
be the source obtained by folding into all cells. Then
\[
 \|G_s\|_{\rm ch}^{\mathsf N}=Z[s^{\rm ch}]/Z[0]\ge0,
 \qquad
 |Z[\sum_\square s_\square]/Z[0]|
       \le\prod_\square\|G_{s_\square}\|_{\rm ch}.
\]
Under the partition-bound hypotheses of \cref{C-407-43},
$s\in\mathcal P_k^\#$ implies $s^{\rm ch}\in\mathcal P_k^\#$
and
\[
 \|G_s\|_{\rm ch}\le
 \exp\!\left[\frac{\operatorname{Re}(E_k(s^{\rm ch})-E_k(0))}
                        {\mathsf N}
      +(C+2+C_P(1+\log x_k))\frac{|\square|}{\lambda_k^3}\right].
\]
If the distance from the walls is at least $\delta$, put
$\lambda_k^{\mathfrak b}(s)=\mathsf N^{-1}
\sum_r(E_k(r_{\rm lin}s)-E_k(0))$. Then also
\begin{equation}\label{eq:qcd-intermediate-cell-bound}
 \|G_s\|_{\rm ch}\le
 \exp\!\left[\operatorname{Re}\lambda_k^{\mathfrak b}(s)
 +(C+2+C_P(1+\log x_k))\frac{|\square|}{\lambda_k^3}
                  +\mathscr C_k(S;\delta)\right].
\end{equation}
\end{enumerate}
The constant $C_{\mathscr C}^\flat$ depends only on
$C_{\mathscr C},\kappa_E,M,L$ and is chosen before the
positive mass floor and the final small-coupling threshold.
\end{proposition}
\begin{proof}
For \textup{(i)}, apply the proof of \cref{C-407-25} with
$s_0$ added to every argument. More explicitly, the contribution
of a polymer is
\[
 \Delta_j(X)=e_j(X;s_0+\sum_i s_i)-\sum_i e_j(X;s_0+s_i)
                                      +(n-1)e_j(X;s_0).
\]
Choose a nearest support $S_{i(X)}$ and write this as
\[
 [e_j(X;s_0+\sum_i s_i)-e_j(X;s_0+s_{i(X)})]
 -\sum_{i\ne i(X)}[e_j(X;s_0+s_i)-e_j(X;s_0)].
\]
Each difference is bounded by \eqref{eq:qcd-source-energy-tail}
using the supports other than $S_{i(X)}$. The same fine-scale
packing gives $e^{-\kappa_E\delta/\lambda_j}$; at coarser
scales use the differences from $s_0$ and the cube count.
The source-independent constant in that proof therefore
remains valid. For mixed derivatives apply
$\prod_i(1-P_i)$ to $e_j(X;s_0+\sum_i z_is_i)$.
Each $1-P_i$ changes the source only on $S_i$.
The distance splitting in \cref{C-407-25}, followed by Cauchy's
estimate on smaller polydiscs and a radius limit, proves the
second assertion.

For \textup{(ii)}, the sum of the increments has norm at most
$\varepsilon$, since their supports are disjoint. The fine
levels $\lambda_j\le\delta$ have the bound just proved.
For a coarser level let $J$ be one increment or their sum and put
\[
 F(\tau)=e_j(X;(0,J_0+\tau J))-e_j(X;(0,J_0)).
\]
The domain $\mathscr O_{j-1}$ contains this family for
\[
 |\tau|<T_j=\frac{r_J/\lambda_{j-1}-\|J_0\|_\infty}{\varepsilon}
              \ge\frac4{\lambda_{j-1}}.
\]
Here $F(0)=0$. Its bound from
\eqref{eq:qcd-source-energy-tail} retains the distance from
$X$ to the support of $J$ throughout the disc. Schwarz's
lemma gives that bound at $\tau=1$ multiplied by
$T_j^{-1}\le\lambda_{j-1}/4$. The coarse-level sum in the
proof of \textup{(i)} therefore acquires exactly $\omega_j$.

To sum it, write $u=K-j$, so $\lambda_j=L^{-u}$, and use
$\#_j^+(S)\le64(1+2R_SL^u/M)^3$. On the levels
$\lambda_j>\delta$, which occur only if $\delta<1$, the sum
is at most $C_{\mathscr C}$ times
\[
 432\sum_{u\ge0}L^{-u-1}
 +432R_S^3\sum_{L^u<1/\delta}L^{2u-1}
 \le36+40(R_S/\delta)^3.
\]
The fine levels are bounded by the exponentially convergent
sum used in \eqref{eq:qcd-source-count}. This proves
\eqref{eq:qcd-bilinear-base-additivity}; the logarithmic
count of coarse levels has been replaced by a geometric sum.

For \textup{(iii)}, $G_s$ is even and belongs to its cell
algebra. Its reflections multiply to $G_{s^{\rm ch}}$.
The reflection-positive chessboard argument in the proof of
\cref{C-407-30} uses only the even number of slabs in each
direction, so gives the displayed identity and product bound
on this board. The reflected sources have the same supremum
and Lipschitz norms; because they vanish near the walls, the
norm of their disjoint sum is their maximum. Source conjugation
and the spin lifts preserve the operator norm of $J$.
Thus folding preserves $\mathcal P_k^\#$.

Apply \eqref{eq:qcd-level-source-partition} to $s^{\rm ch}$
and take the $\mathsf N$th root. Since
$n_k/\mathsf N=|\square|/\lambda_k^3$, this gives the first
cell estimate. If the support lies $\delta$ inside the cell,
its distinct reflected images are separated by $2\delta$.
Part \textup{(i)} replaces their total energy by the sum of
their individual energies. The cube count under a reflection
costs at most $27$, already included in $C_{\mathscr C}$.
Finally conjugate images have the same real energy as their
linear images. This proves \eqref{eq:qcd-intermediate-cell-bound}.
\end{proof}

\subsection{Uniform separated cumulants}

The insertion norms are
\[
 \|f\|_0=\max\{\|f\|_\infty,40M\operatorname{Lip}f\},
 \qquad \|f\|_\Gamma=\|f\|_\infty.
\]
Set
\[
 \varrho_0=\varrho,\qquad \varrho_\Gamma=\varrho_J,\qquad
 \varrho_{\min}=\min\{\varrho,\varrho_J\}.
\]
The truncated function is the mixed derivative of
$\log Z[\sum_i z_is_i]$ at zero. Only its germ at zero is
used: $Z[0]>0$ ensures that this germ exists.

\begin{theorem}[Separated cumulant bounds]\label{C-407-6}
Let $n\ge2$, let the supports of $f_i$ be pairwise
disjoint with $S_1,S_2$ at distance at least $d>0$, and
let each support lie in a ball of radius $R$.
Uniformly in the cutoff and volume,
\[
 |S_n^T((\alpha_i,f_i)_{i=1}^n)|
 \le (1+d^{-1})^3 B(R)^n(n!)^{5+p/2}
                         \prod_i\|f_i\|_{\alpha_i},
\]
where
\[
 B(R)=C_Be^{c_V}(1+R)^3/\varrho_{\min},\qquad
 p=432C_{\mathscr C}/\log L.
\]
Without the support-radius restriction one also has the
torus bound
$C_n(d)e^{nc_V\ell^3}\prod_i\|f_i\|_{\alpha_i}$.
For the finitely many cutoffs with $L^K<16n$, a direct bound is
\[
 [C_B'(1+N_f/m_-)(1+R)^3]^n(n!)^5
                                      \prod_i\|f_i\|_\infty.
\]
\end{theorem}
\begin{proof}
For any holomorphic counterterm $\Lambda(s)$, define
\[
 G_\Lambda(z)=e^{-\Lambda(s(z))+\Lambda(0)}Z[s(z)]/Z[0],
 \quad G_\Lambda(0)=1.
\]
The exact differentiation formula is
\[
 S_n^T=\partial_{[n]}\Lambda+
 \sum_\pi(-1)^{|\pi|-1}(|\pi|-1)!
                         \prod_{B\in\pi}\partial_BG_\Lambda.
\]
The sum of its absolute combinatorial coefficients is
at most $2^nn!$: order the blocks and then order their
elements increasingly, and bound the resulting records
by permutations with cuts. This formula needs no
zero-free complex source neighborhood beyond the germ.

First take $\Lambda=E_K$. On the polydisc with radii
$2\varrho_{\alpha_i}/\|f_i\|_{\alpha_i}$, disjoint
supports keep the total source in $\mathcal P$.
The subtracted-partition upper bound and zero-source
floor give $|G_\Lambda|\le e^{c_V\ell^3}$.
Cauchy's estimate bounds every block derivative.
Prop.~\ref{C-407-25} and
\eqref{eq:qcd-source-count} bound the mixed energy
derivative by $2^{n-1}\prod_i\rho_i^{-1}
C\ell^3(1+d^{-3})$. These bounds prove the torus assertion.

Suppose now $L^K\ge16n$. Set $\eta=(16n)^{-1}$.
Use the product hat partition of unity at mesh $\eta$.
Each function splits into at most $(2R/\eta+3)^3$
terms supported in cubes of side $2\eta$, and each
piece has its insertion norm bounded by
$(1+80M/\eta)\|f_i\|_{\alpha_i}$.
Fix one piece from each insertion. In each coordinate
choose walls of unit cells at half-integral fine-grid
positions, leaving distance $\eta$ to all these supports.
A support forbids at most $4\eta L^K+1$ positions;
there are $n$ supports and $L^K\ge16n$, so fewer than
$L^K$ positions are forbidden. This chooses a valid
through-link board for every allowed torus, without
a lower bound involving $nR$ on its side.

Group the pieces by their cells and let $\Lambda$ be
the sum of the holomorphic cell energies
$\lambda^{\mathfrak b}$ for these groups. On a
coordinate face indexed by $B$, the multiple-cell bound
gives
\[
 |G_\Lambda(z_B)|\le
 \exp\left[|B|c_V+\sum_{i\in B}
                      \mathscr C_K(S_i;\eta)\right].
\]
Since each $S_i$ has diameter at most a fixed multiple
of $\eta$, the refined count gives
$\mathscr C_K(S_i;\eta)\le
216C_{\mathscr C}(1+\log_L(1/\eta))+C$.
Thus the partition derivative sum costs at most
$C^ne^{nc_V}(16n)^{np/2}2^nn!\prod_i\rho_i^{-1}$.
The mixed derivative of $\Lambda$ is zero unless all
pieces are in one cell. In that case the reflected
energy bound of Prop.~\ref{C-407-25} gives
$C(1+d^{-1})^3\prod_i\rho_i^{-1}$ times its fixed
combinatorial factor.
Finally sum over the piece choices. Their number and
norm enlargement cost at most
$(32nR+3)^{3n}(1+1280Mn)^n$.
Using $n^n\le e^nn!$, these four powers of $n$ and
the partition factor give $(n!)^5$, and the energy
count gives the additional $(n!)^{p/2}$.
Absorb fixed constants into $B(R)$.

For completeness the coarse cutoffs require no
renormalization estimate. At every configuration
$\|(D+m_0)^{-1}\|\le1/m_0$. The determinant inequality
$|\det(1+T)|\le e^{\|T\|_1}$ therefore gives
\[
 |Z[s]|\le Z[0]\exp\left[
 2a_K^{-1}\sum_p|w(p)|+
 \frac{384N_f}{m_-}\sum_x\|J(x)\|\right].
\]
A support in a radius-$R$ ball contains at most
$81a_K^{-3}(1+R)^3$ plaquettes and
$27a_K^{-3}(1+R)^3$ sites. Choose the independent
source radii as the reciprocals of these costs.
The moment bound is then $e^n$, and the partition
formula costs at most $2^nn!$ times their product.
Since $a_K^{-1}<16n$, the action cost is bounded by
$C(16n)^4(1+R)^3\|f\|_\infty$ and the bilinear cost
by $C(N_f/m_-)(16n)^3(1+R)^3\|f\|_\infty$.
The same inequality $n^n\le e^nn!$ proves the last
displayed estimate. These cases cover every cutoff.
\end{proof}

\subsection{A common choice of constants}

We now choose the constants used in the preceding estimates.
The order matters: a bound for the induced mass is not an
input to the constants that prove that bound. Throughout,
the group is $\mathrm{SU}(2)$, the chamber parameter is $r=1/2$,
and $N_f\ge1$ counts the four-component flavours, each formed
from the two conjugate two-component blocks.

The gauge induction requires numerical bounds and also conditions
on its radii, lattices and auxiliary objects. We first verify these conditions for the constants
chosen in \cite[Prop.~69.1]{YMM1}. The subsequent numerical construction
will preserve such a choice while imposing the matter bounds.

\begin{lemma}[Admission of the pure gauge induction]
\label{def:qcd-gauge-admission}
Choose the gauge constants in the order and with the defining formulas
of \cite[Prop.~69.1 and its proof]{YMM1}. Fix the history mollifier:
it is smooth, nonnegative, has integral one, is supported in the
supremum-norm ball of radius $1/4$, and is invariant under signed
coordinate permutations. Use the critical backgrounds, coordinate
maps, retained quadratic forms and pivot logarithms of that
construction, with their specified branches.

Let $x_0,\ldots,x_K$ be an admissible coupling sequence as in
\cref{sec:qcd-gauge-input}, with $x_j\ge\gamma^{-2}$, and suppose
$MR(x_K)\le L^m$. At every step and every history the gauge
construction has its required domains and numerical bounds. Its
local auxiliary maps are holomorphic on their prescribed common
balls and vanish at zero datum. Its retained exponents are defined
on the full prescribed layered spaces, have finite oscillations,
and agree exactly on extending a history. Consequently the canonical
finite gauge induction applies. All assertions persist after
decreasing $\gamma$.
\end{lemma}

\begin{proof}
We first pass from the common numerical choice to a coupling
sequence and a torus. We then check the two conditions involving
actual functions: the auxiliary background maps and the retained
exponents. This separates conditions needed to define a step from
bounds on its output, which are proved by induction.

\emph{The common numerical choice.}
The proof of \cite[Prop.~69.1]{YMM1} fixes the exponents
$r_c\ge2$, $\mathfrak p_0\ge6r_c+1$ and
$q\ge41\mathfrak p_0+1$, followed by the reserve fractions,
$L$, the response parameters, the decay rates, the geometric
lengths and the analytic radii. The constants in a defining formula
are evaluated after its arguments have been fixed. The potential
bound $E_0$, the increment bound $B_\sharp$ and the history bounds
are then fixed. The threshold $\gamma$ is chosen last. In particular
none of these choices uses $K$, a particular coupling sequence or
the torus.

Here is why the order permits the last choice. Write
$F_p(x)=x^{-1/2}(\log x)^p$. After the earlier constants have
been fixed, the remaining small quantities are bounded by finite
sums of constant multiples of
\[
 F_p(x),\qquad (\log x)^{-(q-\mathfrak p_0)},\qquad
 x^a(\log x)^p e^{-c(\log x)^{r_c}},\qquad
 e^{-c(\log x)^{2\mathfrak p_0}},
 \quad c>0.
\]
The exponents and coefficients in each occurrence are fixed.
These expressions tend to zero as $x\to\infty$; the third uses
$r_c\ge2$. Each required upper bound can therefore be imposed on
the whole interval $x\ge\gamma^{-2}$. The lower bounds on
$\log\gamma^{-2}$ are finite. Bounded windows, such as those for
the reserve fractions, the normalizer radius and the boundary
rate, are made nonempty at their earlier stages of the choice;
they are not repaired by the last choice of $\gamma$.
This is the numerical construction proved in the cited proposition.

For use in the next proposition, the same order allows an additional
floor or ceiling only at a stage where its dependencies have already
been fixed and its permitted range remains nonempty. All later
constants defined by formulas must then be evaluated at the new
earlier values. Thus a larger rate or a smaller radius does not
mean retaining numerical values obtained from the previous choice.
The next proposition checks these requirements for the matter
conditions explicitly.

\emph{Couplings, radii and room on the torus.}
Put $y_j=\log x_j$, $\check R_j=R(x_j)$,
$\epsilon_j=A_0F_{\mathfrak p_0}(x_j)$,
$\delta_j=A_1F_{\mathfrak p_0}(x_j)$ and
$\mathfrak b_j=A_1y_j^{\mathfrak p_0}$. The choice of $\gamma$
places $y_j$ above all the monotonicity thresholds. Consequently
\[
 \epsilon_j\le\epsilon(\gamma^{-2}),\qquad
 \delta_j\le\delta(\gamma^{-2}),\qquad
 y_j\ge\log\gamma^{-2},\qquad
 \mathfrak b_j\ge A_1(\log\gamma^{-2})^{\mathfrak p_0}\ge1.
\]
In particular every linear smallness condition for the background,
the shift, the averaging logarithm and the coordinate charts holds
at every level. Conditions on the relative change of coupling
follow from the same bound $|b_j|\le B_\sharp$ used in choosing
the uniform regularity constants.

The running radii are
\[
 \alpha_{\iota,j}=C_\iota F_q(x_j),\qquad \iota=0,1.
\]
At either stage containing level $j$ one uses these same values.
Truncating a coupling sequence therefore restricts the table
without changing any previous radius. Since $\log\gamma^{-2}>2q$,
$F_q$ is decreasing on the required interval. The fixed ceiling
conditions include
\[
 \begin{gathered}
 (1+\beta_s)C_0F_q(\gamma^{-2})\le\min\{\alpha_0,\alpha_2\},
 \qquad
 (1+\beta_s)C_1F_q(\gamma^{-2})\le\alpha_1,\\
 \mathsf o L^3M B^{(\rm cg)}C_0F_q(\gamma^{-2})\le\alpha_0,
 \qquad a_\flat\le\alpha_0,\qquad \mathsf o L^3\ge3.
 \end{gathered}
\]
Here $\mathsf o$ and $B^{(\rm cg)}$ are respectively the fixed
local-gauge enlargement and its gauge bound in the layered-space
construction of \cite{YMM1}. These inequalities place each
layered-space reading in its fixed analytic class. They bound the
real and complex local potentials, the current, and the reduced
towers required by that class; the tower conditions are part of
the layered-space definition. They are not inferred from curvature
alone. 

The adjacent-scale radius inequality is
$\check R_j\le L^{K-j}\check R_K$. Thus
\[
 M\check R_j\le L^{m+K-j}.
\]
The lower bounds $\check R_j\ge\max\{L^2,\mathsf oL+1\}$ and
$M>6$ imply
\[
 L^2M\le L^{m+K-j},\qquad
 \mathsf oLM+6<L^{m+K-j},\qquad m'<m+K-j.
\]
The grain, mesh and containing boxes fit at every level, including
the last. The same applies to every truncated history. These are
the torus-room inequalities following \cite[Prop.~69.1]{YMM1}.
The larger local charts used elsewhere in this paper require the
stronger inequality $LMR(x_K)\le L^m$; this is supplied by
$\ell\in\mathfrak L(g)$ in \cref{sec:qcd-gauge-input}, rather
than by the weaker room assumption of this lemma.

\emph{Gaussian integration and boundary terms.}
The historical Gaussian uses the covariance and conditional mean
of \cite[Prop.~57.5]{YMM1}. Denote its fixed collar factor by
$\Theta_{\rm col}$. Write its external datum as $a=a'+a''$, where $a'$
is real and belongs to the prescribed tier: it obeys
$|a'|<\mathfrak b_k\min\{C_{\rm xt},(4\Theta_{\rm col})^{-1}\}$
on the collar, $|a'|<\theta_c\mathfrak b_k$ on the rim,
and $|a'|<C_{\rm xt}\mathfrak b_k$ elsewhere.
The complex increment $a''$ vanishes on the rim
and satisfies $|a''|\le\epsilon_a\mathfrak b_k$
elsewhere. Both its real and imaginary parts enter
the conditional-mean estimate.
The chosen collar width and rim width control
the collar and rim contributions to the conditional mean. The
remaining contribution has the factor
$e^{-\delta_M(W_t-2)/2}$ and is made small by the running
separation. The complex covariance radius controls the imaginary
part of the conditional mean. The common choice imposes both
the mean bound and the normalizer positivity bound before the
Gaussian expansion is applied. In particular the bounds concern
the full external datum, not just its restriction to a polymer.

To spell out the quantitative reserve, let $\mathcal K_{\rm ap}$
be the fixed cluster prefactor and let $S_{\rm rg}(x)$,
$S_{\rm bd}(x)$ and $S_0(x)$ denote the three potential bounds
in the proof of \cite[Prop.~69.1]{YMM1}: regular and boundary
potentials at a later step, and the potential at the first step.
Their explicit formulas there have the first two types of decay
displayed above. The Gaussian large-field penalty is
\[
 \Gamma_P(x)=C_{\rm mn}^{-1}(LM)^{-3}
                \theta_c^2A_1^2(\log x)^{2\mathfrak p_0}.
\]
The common choice gives, uniformly for $x\ge\gamma^{-2}$,
\[
 \begin{gathered}
 S_{\rm rg}(x),S_{\rm bd}(x)
                   \le(12eL^3\mathcal K_{\rm ap})^{-1},\qquad
 S_0(x)\le(6eL^3\mathcal K_{\rm ap})^{-1},\\
 (LM)^3e^{-\Gamma_P(x)}\le(6\mathcal K_{\rm ap})^{-1}.
 \end{gathered}
\]
The cluster estimate is therefore at most $1/3$: at a later
step its two terms are bounded by
$\mathcal K_{\rm ap}eL^3(S_{\rm rg}+S_{\rm bd})\le1/6$
and $\mathcal K_{\rm ap}(LM)^3e^{-\Gamma_P}\le1/6$.
At the first step the first term is instead
$\mathcal K_{\rm ap}eL^3S_0\le1/6$. This separately checks
the initial step, where no old potential family is available.

For a step $t=k+1$, the rejected forms and the strip expansion
have the joint bound
\[
 \mathsf b_{\rm rj}+\mathsf b_{\rm sr}
 \le e^{-\check R_t}(c_3+c_4\mathfrak b_k^2),
\]
where $c_3,c_4$ are the fixed constants evaluated in the same
proof. They are fixed before $\gamma$. The flow bound gives
$x_k\le(3L/2)x_t$. With $y=\log x_t\ge\log(2L)$,
\[
 \mathfrak b_k^2\le A_1^2(2y)^{2\mathfrak p_0},\qquad
 \check R_t\ge y^{r_c}\ge y,\qquad
 (2y)^{2\mathfrak p_0}e^{-y/2}
                     \le(8\mathfrak p_0/e)^{2\mathfrak p_0}.
\]
Consequently the preceding bound is at most
\[
 \gamma\bigl[c_3+c_4A_1^2
                   (8\mathfrak p_0/e)^{2\mathfrak p_0}\bigr]
 \le\tfrac13.
\]
Together with the cluster bound this leaves the boundary
reserve $1-2/3$. The additional smallness conditions for the
strip coordinate changes are imposed in the same numerical
choice; their tails are a fixed constant times
$e^{-\check R_t}$, compared with the already specified
coordinate radius. Hence they hold for every history on the
given sequence. No constant is selected separately for a
particular background or history.

\emph{The auxiliary background maps.}
An octet $\square$ of level $n$ is the union of eight level
$(n-1)$ cubes meeting a vertex. Put $Z=\square^{\sim5}$ and
let $\mathbf B_j^\square$ be its reduced determining set at
truncation $1\le j\le n$. The canonical maps are
\[
 H_{\square,j}(B)=\mathcal H_{\mathbf B_j^\square,1}(B),
 \qquad \sup|B|<r_{\rm B}.
\]
Here $B$ assigns a complex Lie-algebra element to each
determining row. By \cite[Lem.~51.6]{YMM1}, this map is
the holomorphic potential in the normalized mean-Landau
gauge, based at the identity configuration. Its exponential,
together with the separate normalizing gauge, represents
the critical background over the datum $e^{iB}$.
The map vanishes at $B=0$. Its radius $r_{\rm B}>0$ and
derivative bounds are independent of height and volume.
Uniqueness for the mean-Landau condition and the node
normalization makes these maps commute with adjunction,
constant $\mathrm{SU}(2)$ gauge rotations and translations of lattice charts,
as checked before \cite[Prop.~51.14]{YMM1}.

The radii used to substitute an old interaction are the
absolute system at height $k$:
\[
 a_{\iota,j}=\alpha_\iota\quad(1\le j\le k),\qquad
 a_{\iota,k+1}=\frac{\alpha_\iota}{1+\beta_s},
 \qquad \iota=0,1.
\]
Its output radii are $\bar a_\iota=
(1+\beta_s)a_{\iota,k+1}=\alpha_\iota$, and
$\max_\iota\bar a_\iota/a_{\iota,j}=1$ for $j\le k$.
The common choice includes
\[
 \alpha_0+\alpha_1\le a^\star,\qquad
 c_{\rm sh}M\alpha_0<\alpha_1,\qquad
 \alpha_1^2\le c_\natural\alpha_0,\qquad
 \frac{(1+\beta_s)\delta(\gamma^{-2})}{\alpha_0}
                                                 \le s^\star.
\]
These are precisely the shape and shift-size conditions
of \cite[Def.~51.12]{YMM1}. The slow-growth comparisons
in that definition hold for the displayed absolute
system. The choices $a^\star\le a^\star_{\rm oc}$ and
$c_{\rm sh}\ge c^{\rm oc}_{\rm sh}$ are made when the
fixed analytic radii are chosen. This radius system has
a different purpose from the running radii of the
layered spaces above.

For $k\ge2$, the octet construction uses
\[
 b_\iota^{\rm oc}=\theta_{\rm oc}a_{\iota,k+1},\qquad
 a_{\iota,k}^{\rm oc}=\frac{b_\iota^{\rm oc}}{1+\beta_s},
 \qquad
 a_{\iota,j}^{\rm oc}=\min\{a_{\iota,j},b_\iota^{\rm oc}\}
 \quad(j<k).
\]
Here $0<\beta_s\le1/5$ and
$(3/2)\theta_{\rm oc}^2<1$. The proof preceding
\cite[Prop.~51.14]{YMM1} checks the same shape and
slow-growth inequalities for this system. Its shift-size
ceiling is $(1+\beta_s)s^\star/\theta_{\rm oc}$, the
ceiling used in that substitution.

Let a pair on $X'=(\square^{\sim5})^\uparrow$ belong to
$\mathcal U_k^\circ(X';b_0^{\rm oc},b_1^{\rm oc})$.
In particular, its restriction to $Z$ belongs to the
window domain: it has the local potential and curvature
bounds, and all the reduced real and complex towers
exist with the required bounds at every truncation.
Put the averaged datum $V$ in the comb gauge used in
\cite[Prop.~51.11]{YMM1}. With $\eta_j=L^{-j}$ in the
rescaled level-$j$ chart, set
\[
 \begin{split}
  B_V&=i^{-1}\log V,\qquad
  A_{\square,k}=H_{\square,k}(B_V),\\
  Q_{\square,j}&=i^{-1}\log M_{\mathbf B_j^\square}
                          (e^{i\eta_kA_{\square,k}}),\qquad
  A_{\square,j}=H_{\square,j}(Q_{\square,j}).
 \end{split}
\]
The reduced averaging map in this formula uses its
specified logarithm branch.

We check the domains of these compositions. In
\cite[Prop.~51.11]{YMM1}, take the top window radii
$\widetilde a_\iota=b_\iota^{\rm oc}/(1+2\beta_s)$ and
the lower radii $a_\iota=a_{\iota,j}^{\rm oc}$, with
$n=k$, $\tau=1$ and $B'=0$. Its hypotheses are
\[
 \widetilde a_0+\widetilde a_1\le a^\flat,\qquad
 a_0+a_1\le a^\flat,\qquad
 2B_3^2K_{\rm tw}Ma_0<a_1,\qquad
 \widetilde a_0\le t_{k-j}a_0.
\]
The constants $B_3,K_{\rm tw}$ and the radius $a^\flat$
are fixed in that proposition; the radius is chosen
after $M$ to satisfy the logarithm, averaging and
response ceilings. The octet radius system gives
these inequalities. The comparison sequence is
$t_\nu=(6/5)\sqrt{1+\nu}$. For
$\Theta_s=\sup_{\nu\ge1}t_\nu L^{-s\nu}$, the
remaining hypotheses are
\[
 \begin{gathered}
 \Theta_1\le1,\qquad
 (1+10\beta_s)\Theta_2+\beta_s/4+1/1000\le1/8,\\
 (1+3\beta_s)\Theta_3+5\beta_s\Theta_2
                    +\beta_s/8+1/500\le1/16.
 \end{gathered}
\]
Their suprema occur at $\nu=1$. The proof of
\cite[Prop.~51.14]{YMM1} evaluates the last two left
sides at $L=13$, $\beta_s=1/5$ as less than $0.082$
and $0.039$, respectively. Thus these hypotheses
hold uniformly in $k-j$.

The substitution proposition now proves that $B_V$
and $Q_{\square,j}$ lie inside the prescribed row
balls. It gives, on $\square^{\sim3}$,
\[
 |A_{\square,j}|,|\nabla A_{\square,j}|<a_1/2,
 \qquad |\partial e^{i\eta_jA_{\square,j}}-I|
                                      <a_0\eta_j^2/8,
 \qquad |J^{(j)}(e^{i\eta_jA_{\square,j}})|<a_0/16.
\]
Each required inner-polymer restriction therefore
belongs to the stated analytic core. The restricted
real background is identified by criticality and
uniqueness: variations preserving the restricted
datum extend by the original field, because the
boundary bonds are fixed determining bonds.
The identities extend to the complex logarithmic
domains by holomorphy. The resulting maps and their
domains are the canonical maps at every truncation.
At heights $k\le1$ no octet substitution is required.

\emph{Retained factors and their domains.}
There are two possible retained exponents for a pair $(t,X)$:
the kept quadratic form and the sum of external pivot
logarithms. Write $t=l+1$. Their arguments are the pair
read in level-$l$ units and the level-$l$ stored field,
extended by zero on the other bonds of $X$.
The radius ceilings above put this pair in
$\mathscr O(X;\alpha_0,\alpha_1,\alpha_2)$ on the whole
layered space, including its stored-data tiers.
This class membership is the conclusion of
\cite[Prop.~57.4]{YMM1}. Here $\mathfrak M$ is the fixed
power of $L$ specifying the cube partition of the layered
regions. The local hypotheses are part of the layered-space
definition: at each layer the
curvature, current, complex potential and first
difference obey the graded radii; real local gauges
exist on the prescribed cubes. Also, every cube made of
level-$n$ $\mathfrak M$-cubes, contained in the domain and
the level-$n$ region, and of side at most $12\mathfrak M$
in level-$n$ units, has a containing embedded cube in the
domain with a real gauge of index $n$ and bound $a_\flat$.
The history gaps put the tiny cubes about fluctuation
bonds inside the preceding layer, and the torus-room
inequalities make them embedded. The radius tests
for the cited proposition are exactly the three
running-radius ceiling inequalities above, together
with $a_\flat\le\alpha_0$. They hold uniformly at
every level. Passing from the level-$t$ triple to
the level-$(t-1)$ pair rescales $A'$ by $L^{-1}$
and $\mathbf J$ by $L^{-3}$; hence the underlying
configuration and its current in bare units are
unchanged. The proposition then supplies
$\mathscr O(X;\alpha_0,\alpha_1,\alpha_2)$ on the
full layered domain.

The fixed rate and pivot conditions are
\[
 \delta_\star^{\rm pl}LM\ge C_2(\kappa+1),\qquad
 L\max\{\alpha_0,\alpha_1\}\le a_{\rm nf}.
\]
The common choice satisfies these, with the stronger rate
reserve used in subsequent locality estimates.

The construction preceding \cite[Prop.~62.3]{YMM1} checks the pivot
condition even for a coarse bond crossing the last layer.
The two averaging blocks then lie in a cube of side
$2\mathfrak M$ in the last two layers. Their real and
complex potentials are bounded by $L\alpha_0$ and
$L\alpha_1$. The pivot estimate at
$a=L\max\{\alpha_0,\alpha_1\}\le a_{\rm nf}$ therefore
applies to this bond as well. It gives
$|z_c-6\log L|\le C_za$ on the specified logarithm branch.
Thus neither crossing bonds nor the last layer require
restricting the domain of a factor.

The kept quadratic form is holomorphic in the pair, and
on the stored-data tier it obeys
\[
 |\varphi^{\rm kp}_{t,X}|
 \le\tfrac12K_{\rm ch}C_{\rm xt}^2\mathfrak b_l^2
                         e^{-\kappa(d_t(X)+1)}.
\]
For the other exponent there are at most $6M^3$ pivot
bonds in $X$. Its domain is one cube or two adjacent
cubes, so $d_t(X)=0$, and
\[
 |\varphi^{\rm el}_{t,X}-6N_X\log L|
                         \le6M^3C_za_{\rm nf},
 \qquad N_X\le6M^3.
\]
Define
\[
 \operatorname{osc}\varphi
   =\sup_{\mathfrak q}\sup_{P,P'}
       |\varphi(P;\mathfrak q)-\varphi(P';\mathfrak q)|,
\]
where $\mathfrak q$ ranges over the allowed stored-data tiers
and $P,P'$ over the corresponding full layered space.
At each fixed $\mathfrak q$, subtract the values at $P$ and $P'$. The first estimate is used
twice; in the second, the same constant $6N_X\log L$
cancels. Since $\mathfrak b_l\ge1$, this gives
\[
 \sum_{\nu\in\{{\rm kp},{\rm el}\}}
        \operatorname{osc}\varphi^\nu_{t,X}
 \le B_{\rm fr}^{\flat}\mathfrak b_l^2e^{-\kappa d_t(X)},
 \qquad
 B_{\rm fr}^{\flat}
       =K_{\rm ch}C_{\rm xt}^2+12M^3C_za_{\rm nf},
\]
where an absent exponent contributes zero. This is a
bound on actual finite oscillations over the specified
domains. Each factor is Borel measurable and holomorphic
on fixed-real slices; gauge transformations and conjugation
preserve its defining formula and branch.

Keep the name $(\nu,t,X)$ attached to each factor.
Its pair reading and stored-data tier use only the history
through step $t$, and its radii use only $x_0,\ldots,x_t$.
Extending a history changes none of these. Thus names,
domains and exponents are retained exactly. There are at
most two names for each $(t,X)$ and finitely many domains
on a finite lattice, so the family is finite. This also
shows why the argument does not require an inductive bound
on the previously collected interaction to define these
factors.

\emph{The finite induction.}
The actual-step hypotheses are now available: the regular
backgrounds and their coordinate neighborhoods; the common
radii and stored-data tiers; the Gaussian and strip bounds;
and the auxiliary maps and retained factors on their full
domains. Start with Wilson's positive finite-volume
measure. Given the integrable density at one level,
\cite[Prop.~69.3]{YMM1} supplies the next density and its
exact integration identity. Its proof makes the two Borel
changes of variables, carries the kept quadratic forms
and pivot factors with their fixed names, and applies the
Gaussian and strip expansions on precisely the domains
checked above. Its abstract measure step also supplies
integrability and the exceptional null sets; these are
not assumptions on every point in a chart image.

If the initial small region of a step is empty, all its
integration coordinates and local families are empty.
The coordinate changes are then zero-dimensional, and
the same identity is evaluation at the retained field.
This case is included in the cited step proposition.
The first nonempty step uses $S_0$, as checked above.
At each subsequent step the numerical bounds give the
same hypotheses for the new regular and boundary terms.
The agreement of radii, maps and retained exponents under
truncation makes these successive constructions one
finite induction, as in \cite[Thm.~AC]{YMM1}.
Finally, all coupling-dependent conditions were imposed
uniformly for $x\ge\gamma^{-2}$. Decreasing $\gamma$
therefore preserves every assertion on the resulting
class of coupling sequences and tori.
\end{proof}

\begin{proposition}[A common choice and the finite induction]
\label{C-396-12}\label{prop:qcd-joint-admission}
Fix $b_F>0$, $0<m_-\le m_+\le1/2$ and the action-source
radius $\varrho>0$. The additional numerical conditions in the
local, history and source estimates can be satisfied in the
ordered choice of constants, independently of the cutoff and
volume. At the resulting complete assignment the hypotheses
of the pure gauge induction hold for both the calibration
and running sequences. There are $B_\sharp>0$ and
$0<\gamma\le1/2$ with the following property. Put
$\bar B=LB_\sharp/(L-1)$. For $0<g\le\gamma$, an admissible
torus side $\ell\in\mathfrak L(g)$, $K\ge1$, and
\[
 \mathsf b\in\bar I(g)
 =\bigl[g^{-2}+\bar B,\,L(g^{-2}+B_\sharp)-\bar B\bigr],
 \qquad \mathsf m\in[m_-,m_+],
\]
the finite flow and every history term are defined, with
the asserted domains, locality, bounds and exact identities.
This includes the complex source step on its stated open
coordinate domain, the common vacuum energy in the two
partition bounds, and the local source-energy and cumulant
bounds. The choice remains admissible, with these conclusions,
after decreasing $\gamma$.
\end{proposition}

\begin{proof}
We first choose the constants in the order of the gauge
construction and verify the additional matter inequalities.
We then verify the hypotheses of the gauge induction and
apply it to the matter objects constructed above. The pure
gauge theorem supplies the regular-pair cores, historical
Gaussian triples and gauge coordinate maps. The fermion
potentials satisfy their own local and history estimates.

\emph{The choices before the fibre weight.}
Choose the pure fixed fractions and integer exponents first:
$r_c\ge2$, $\mathfrak p_0\ge6r_c+1$, and the larger auxiliary
exponent required in the pure table. Choose $\delta_{\rm sp}$,
$\beta_s$, the Gaussian reserves and the halo data in their
open windows. The mollifier in the history profiles is also
fixed here. Choose an odd $L\ge13$ above the pure and matter
floors. The requirement on $b_F$ is only positivity.

Form the free and background operator constants next.
The complex gauge and matrix-source radii $\alpha_1^*$ and
$\varrho_\sharp$ precede the supremum inverse constant
$C_{\rm sup}$. The latter precedes the supremum walk constant
$C_q^{\rm sup}$ and the least power $M_w\ge L$ of $L$ such that
\[
 C_{\rm nb}\max(C_q,C_q^{\rm sup})/M_w\le e^{-3}.
\]
Thus $q_{\rm sup}=C_q^{\rm sup}/M_w$ satisfies
$C_{\rm nb}q_{\rm sup}\le e^{-3}$.
The row-sum constants $C_\Sigma$ for the background maps
and $C_D^\Sigma$ for the stored tower are fixed before
$C_H$; neither depends on $M$ or $N_f$.
Form also $C_D^{\Sigma,{\rm s}}$ from
the same row-sum argument with $2C_D$, decay rate $\alpha_D/2$,
and the fine-operator bound $\mathsf c_{\rm disc}+2$.
This sourced tower constant also precedes the fibre weight.
Here $C_H^\circ$ includes the preliminary background-kernel
bound of \cref{C-397-40}, before summing over all stored
places, and the parametrix and fine-resolvent constants in
\cref{C-397-33}. These use the already fixed local regularity
constants and radii, including its curvature and absorption
ceilings. The row-sum enlargement is
$(1+C_\Sigma L^{3/2})C_H^\circ$. Choose
\[
 C_H\ge\max\{6C_S(c_{\rm hop}+b_F),\,8LC_{\rm sup},\,
 2(1+C_\Sigma L^{3/2})C_H^\circ,\,
 C_D(b_F+C_D^\Sigma)+1,\,2+2C_{\rm nb}C'_H\}.
\]
The single-walk constants $C'_H$ and $C'(\alpha)$ are formed
before $C_H$ from the local inverse and source-map bounds.
For $0<\alpha<1$ choose also
\[
 C(\alpha)\ge\max\{LC_H/(1-\alpha),\,
                         5L^2+4C_{\rm nb}C'(\alpha)\}.
\]
The radius $\varrho_*=(2C_DC_H^2)^{-1}$ then lies in the
earlier source ball and below $(8L^2C_S)^{-1}$.
All inverse and walk estimates are formed on that ball;
increasing $C_H$ only decreases the radius subsequently used.
The free ellipticity and finite four-resolvent calculation
give the early coupling constant $c_2^F$ before $M$ and
$\kappa$. The constants $B_F$ and $\bar B_M=3B_F+1$ of the
mass sum rule have the same earlier placement. In particular,
the coupling bound does not use a late normalizer bound.

The centered conditional construction requires a finite
enlargement of the pure rate floors. Keep the pure Gaussian
map
\[
 N_{\rm gs}=64B_{\rm nm}^5\gamma_{\rm pos}^{-2}+c_r+2,
 \qquad
 \bar c_P^+=C_{\rm mn}^3\lambda_r^{-1}(N_{\rm gs}+1).
\]
These are fixed before the rates. With
$\kappa''=(1+4\beta_s)\kappa+1+c_E$ and
$\bar r_H^+=\kappa''+c_E+C_{\rm cs}(\bar c_P^++2)$, require
\[
\begin{split}
 L&\ge\frac{2C_{\rm cs}(1+4\beta_s)}{1-2\delta_{\rm sp}},\\
 \tfrac12(1-2\delta_{\rm sp})\kappa
 &\ge2(1-2\delta_{\rm sp})+2c_E
    +\frac{C_{\rm cs}}L[1+2c_E+C_{\rm cs}(\bar c_P^++2)],\\
 \bar\kappa&\ge1+C_{\rm cs}\bar r_H^++\log C_{\rm cs},
 \qquad
 \bar\kappa\ge(\kappa+2)c_0+\log C_g+9.
\end{split}
\]
Choose $\kappa$ and then $\bar\kappa$ above these and the
original entropy and strip floors. The geometry constant
$c_0$ is independent of both rates and of $M$.

Choose the remaining pure geometric lengths in their order
and choose the power $M=L^{m'}$ last among the lengths.
Let $C_{\rm rc}$ be the geometric constant for recounting localized
mass-member walks through cubes: the extra norm factor is at most
$1+C_{\rm rc}M_w/M$. Indeed a walk reaching beyond distance $t$
from its starting cube face has $|\omega|+1\ge t/(CM_w)$ and
meets at most $C(1+|\omega|)^3$ cubes; the available
$e^{|\omega|}$ walk weight pays this count.
Besides the pure requirements, impose
\[
\begin{gathered}
 8LM\ge\max\{R_0,1+16/\delta_*,2\alpha_D^{-1}\log L,LM_w\},
 \quad M\ge3M_w,\\
 c_fM\ge C_{\rm dc}\max\{\bar\kappa,\kappa+2\},\qquad
 \delta_FM\ge\kappa+2,\qquad M\ge2C_{\rm rc}M_w,\\
 (LM)^3\ge C_{\rm mn}^{-1}\bar c_P^+,\\
 \delta_\star^{\rm pl}M\ge500L^2(3\kappa+2).
\end{gathered}
\]
The last inequality is the conditional-mean-chain
floor. The pure kept/rejected quadratic-form floors are
included at this stage. Each has the form $M\ge C$,
$cM\ge C+p\log M$, or $CM^pe^{-cM}\le\epsilon$ with
fixed positive $c,\epsilon$. Hence they hold for all
sufficiently large powers of $L$.
The closed enlargement has reach $D_{\rm ex}=M$ in the
lower-level units, and $M/L$ in the upper-level units.
It need not contain an additional layer of big cubes.
Thus the condition with $\min(M,D_{\rm ex})$, read in
lower-level units, is exactly the displayed floor on $M$.

The source-tail conditions are compatible with this order.
Set
\[
 \kappa_\sigma=\min\{\kappa/(4c_{\mathcal D}M),
                   1/(24M_w),1/4\},\quad R_\flat'=7LM,
 \quad C_{\rm loc}=8e^{\sqrt3\kappa_\sigma+\kappa/4}.
\]
Require $\kappa_\sigma R_\flat'\ge\log(4L)$ and
$\kappa_\sigma R_\flat'\ge\kappa+9$.
For the first branch these are floors on $\kappa$, since
$7L>4c_{\mathcal D}$ and $c_{\mathcal D}\le6$.
For the second and third branches they are floors on $M$.
Indeed $c_f=(3LM_w)^{-1}$ and $C_{\rm dc}\ge16c_0/L$
already imply $M\ge48c_0M_w(\kappa+2)$.
Consequently
\[
 Le^{-\kappa_\sigma R_\flat'}\le\tfrac14,
 \qquad
 (C_{\rm loc}+3e^{2\kappa_\sigma})
                   e^{-\kappa_\sigma R_\flat'/4}\le2.
\]

Now reduce the pure analytic radii in their original order.
Include $\alpha_0\le\epsilon_*L^{-2}$,
$2L\alpha_1\le\alpha_1^*$,
$6C_Q\alpha_1ML^{-4}\le\kappa$ and $114M^3\alpha_0^2\le1$.
Retain also the tree-path ceilings
$2C_t\alpha_1M\le\kappa$ and $18C_t\alpha_1M_w\le1$.
Here $C_t$ is the fixed geometric tree-path constant. The
last two requirements are ceilings on the radius chosen
after $M$, not upper bounds on that earlier length.
Include also the conditional mean/error
conditions
\[
 \epsilon_{\rm dc}\le(LM)^{-3},\quad
 \theta_m+\epsilon_m\le\vartheta_*<1,\quad
 (LM)^3\epsilon_m^2\le1,\quad
 1+\vartheta_*<C_{\mathcal A},\quad
 \vartheta_*<\theta_{\rm sl}.
\]
The pure proof makes the exponentially small error by the
preceding length choices and the remaining errors by these
radius choices. Its mean bound uses the actual historical
external tiers, not an independent mean parameter.

Before $E_0$ form $A_\kappa=2C_{\rm nb}e^{(\kappa+2)c_0}$,
the finite small-lag count used in reading old legs, the
normalizer constant $C_N$, and the cluster constant $K_M$.
They use only the already chosen geometry, rates, operator
bounds and analytic radii. They depend on neither $E_0$ nor
$B_\sharp$ nor $h_F$. The fixed potential ceiling
$c_{\rm cl}$ is chosen independently of the array bounds;
$K_M$ does not depend on that choice.
The degree-zero pointing estimate also gives
$K_M^\bullet=K_M+C_EK_M^{\rm v}$ at this stage.

Three demands on determined pure constants must be imposed
on their chosen inputs. The maps themselves are retained.
First form the finite matter boundary bound
\[
 B^F=24(LM)^3E_{\rm nf}+C_NN_f+3,
\]
where $E_{\rm nf}$ is the pure first-step bound.
This number depends on neither $K_M$, $K^{\rm lb}$, $E_0$
nor $B_\sharp$. The reserve $3$ bounds the sum
$1/2+1+1/16+1$: the history output, the pure rejected
terms, the matter strip row and the retained Gaussian mismatch.
Their bounds follow from the later frame and strip smallness
conditions and \cref{C-399-42,C-399-49}. Choose $E_0$ above
its pure floors, with $8M^3E_0\ge B^F$, and with the full
matter-array floor
\[
 E_0\ge E_{\mathbf N}+C(C_N+M^3)N_f
 +C(1+L^{-1})(c_{2,\rm loop}+N_fc_2^F)
 +CK_{\mathbf G}+C_n+C C_{am}r_\alpha^{-2}C_NN_f+2.
\]
Every quantity on its right was fixed earlier.
The original map $B\ge\max(1,B_{\mathbf T},8M^3E_0)$ then
gives $B\ge B^F$. Second impose
$\varepsilon_{\mathcal W}\alpha_0\le2a_Q$; the unchanged map
\[
 e_{\rm cov}=\min\{1,c_2^{\rm SO},
                         \varepsilon_{\mathcal W}\alpha_0/2\}
\]
gives $e_{\rm cov}\le a_Q$.
Third choose $A_0$ after $A_1$ with
$4\pi A_0\ge c_{\rm sh}A_1$.
Since $C_{\rm xt}\ge\pi(1+c_sA_0/A_1)$ with $c_s>2$,
the original map $C_{\mathcal A}=\max(3,2C_{\rm xt})$
gives $C_{\mathcal A}\ge c_{\rm sh}$.
These are respectively a floor on $E_0$, a ceiling on
$\alpha_0$, and a floor on $A_0$. They introduce no return
to a determined constant.

Choose $B_\sharp$ after $E_0$. Include both the pure floor
$B_\sharp\ge G_\sharp=C_{am}r_\alpha^{-2}E_0$ and
\[
 B_\sharp\ge(1+L^{-1})(c_{2,\rm loop}+N_fc_2^F)+1.
\]
Form $C_{\mathfrak r},C_\delta,C_K,C_G$ after $E_0$ and
$B_\sharp$, and then choose the vertex radius $\beta_0$ with
\[
 C_G\beta_0\le\tfrac12,\qquad
 c_h\beta_0\le\min\{\epsilon_P,\alpha_1/2,a_{\rm an}/2\},
 \qquad 2C_P\beta_0\le\min\{r_P^\circ,\alpha_1/2\}.
\]
The last inequality supplies the separate pivot estimate at
the first step; the positive-level shift bound alone does
not supply it.
Form now the remaining old-reading and boundary-counting
constants. Together with the already fixed $A_\kappa$,
$K_M$ and $\Gamma_F$, the constants $K_\mu$, $C_{\rm bd}$,
$C_\omega$ and $S_\omega$ precede the fibre weight.
Their definitions contain no fibre weight. The cancellation of
$h^2$ between a mass-member norm and its moment extraction
is the reason $K_\mu$ is independent of $h$.
For the source estimates form the weighted piece constant
$C_{\rm pw}$ and the enlarged cell-unit constant $C_{(N)}$
before $C_\omega$, and then $c_{\rm cs}$ and $C_{\rm pc}$.
Their bounds divide out the same two leg weights $h^2$;
none depends on the subsequent choice of $h_F$.

\emph{The fibre weight and source constants.}
Choose $h_F$ sufficiently large. Besides
$h_F\ge2A_\kappa\Gamma_FL/(L-1)$ and
$h_F\ge32C_{\rm bd}K_MA_\kappa\Gamma_F$ and
\[
 h_F\ge\frac94 C_{\rm bd}K_MA_\kappa\Gamma_F
                      c_{\rm cl}/\varepsilon_\flat,
\]
impose the flat and differentiated contraction floors. Their factor is
\[
 \theta_E(h)=\sqrt2K_\mu K_M
 \left[\frac{2A_\kappa\Gamma_F}h+
 \left(\frac{c_h}{a_{\rm an}}+C_G\right)
                         \frac{c_{\rm cl}}{4C_Vh^2}\right]
 \longrightarrow0.
\]
In particular choose $2\sqrt2\theta_EC_\omega S_\omega\le1/2$.
For the source extraction, first form
\[
 K_\mu^{\rm s}=K_\mu+c_b^{-2}c_{\mathsf S}M^{-3}C_{\rm pw},
 \qquad \theta_E^{\rm s}(h)=(K_\mu^{\rm s}/K_\mu)\theta_E(h).
\]
The weighted piece constant $C_{\rm pw}$ and all the factors
in $K_\mu^{\rm s}$ precede $h_F$. The source history uses
$\Gamma_F^{\rm s}=2\Gamma_F$ and the boundary constant
$C_{\rm bd}^{\rm s}$ formed with the doubled operator bounds.
Require
\[
 h_F\ge32C_{\rm bd}^{\rm s}K_MA_\kappa\Gamma_F^{\rm s},
 \qquad 2\sqrt2\theta_E^{\rm s}C_\omega S_\omega\le\tfrac12.
\]
Put $D=1-8/(3\sqrt L)>0$ and $P=1+C_{\rm tl}+42M$.
The increment calculation has
$B_\nu=\alpha_\nu+\beta_\nu B_\sigma$, where the coefficient
$\beta_\nu$ is a fixed multiple of $\theta_E(h_F)$.
Require, at this stage,
\[
 4\theta_E^{\rm s} C_{\rm pc}S_{1/4}(1+c_{\rm cs})\le\tfrac12,
 \qquad 6P\beta_\nu\le D/2,
 \qquad 2C_{\mathcal L}\beta_\nu\le\tfrac12.
\]
The factors multiplying $\theta_E$ in these three conditions
are already fixed. In particular the cell-sum argument
does not use $B_\nu$ or $B_\sigma$. These are genuine floors
on $h_F$.

Immediately afterwards choose $C_{\rm rad}$. All its floors
are explicit and have fixed right sides:
the auxiliary threshold
\[
 \delta_{\rm ax}=\min\{1,s^*\alpha_0/(1+\beta_s),
              \varepsilon_{\mathfrak h}/C_{\mathcal A},
              \delta(e^{2\mathfrak p_0})\}
\]
is formed from constants preceding $h_F$.
\[
\begin{split}
 C_{\rm rad}\ge\max\{& (1+\beta_s)/s^*,\,
 12C_{\mathcal A}K_\Lambda^{\rm s}K_M^\bullet,\,
 4c_h C_{\mathcal A}\alpha_0h_F/(a_{\rm an}A_\kappa\Gamma_F),\\
 &C_{\rm sf}\sqrt L A_0/(\varepsilon_{\mathcal W}A_1),\,
 512c_{\rm ins}(1+c_{\rm cl}^{-1})C_{\mathcal A}K_\Lambda^{\rm s}K_M,\\
 &\sqrt2L(2+C_{am}r_\alpha^{-2})K_M^\bullet K_\Lambda C_{\mathcal A}
                                  /(\theta_E(h_F)C_\omega),\\
 &16(1+s^*K_\Lambda^{\rm s}K_M)(1+\alpha_0/\delta_{\rm ax})^2\}.
\end{split}
\]
The last two floors come from radius localization and the
later insertion estimate. These floors are independent of $C_2$, $C_3$, $B_\sigma$
or a final small-coupling threshold.

Form the zero-source constants in their proved order:
$r_x,C_{\rm sl},C_r,C_{(2)},C_0,C_1,C_\chi,c_s,\theta_s$,
then $C_2,C_T,C_3$, then $C_\rho',C_\rho,C_s$, and then
$C_w,C_b,C_e,B_z,C_\mu$. Here $C_3$ is the cubic mass
remainder constant; later source remainder constants are
different constants. The separate mass-member estimates
ensure that none of these depends on the eventual source-window
radius. Next form $S_3,C_w^\psi,B_0,C_{\rm tor}$.
The constant-source comparison uses only zero-source
constants and therefore also precedes $B_\sigma$.

The two apparent return inequalities for $B_\sigma$ are
\[
 B_\sigma\ge4B_0+3C_{\rm var}+8C_{\rm tor}C_{\mathcal L},
 \qquad B_\sigma\ge2C_{\mathcal L}B_\nu,
 \qquad C_{\rm var}=2PB_\nu/D.
\]
Substitute $B_\nu=\alpha_\nu+\beta_\nu B_\sigma$.
The preceding floors make each return coefficient at most
$1/2$. For example the explicit finite choice
\[
 B_\sigma\ge\max\left\{
 2\left(4B_0+8C_{\rm tor}C_{\mathcal L}
                         +\frac{6P}D\alpha_\nu\right),
 4C_{\mathcal L}\alpha_\nu\right\}
\]
satisfies both. Now form $C_w^{\rm s},B_\nu,C_{\rm var}$,
$C_{\rm tor}'$, $C_\sigma$ and $C_\nu$, and finally
\[
 B_s=\bar B_M+\max\{5C_\sigma,2B_\sigma\},\qquad A=5B_s.
\]
The general sourced steps use the open set on which their
potential has size below $c_{\rm cl}$. They do not require
their order-$\phi$ leg families to belong to the smaller
balls about the normalized zero-source families. The cubic
mass decomposition is used only in the constant-coordinate
window, with $C_\rho$ fixed before $B_\sigma$.
All local-energy, stored-source, history-majorant and later
insertion constants are now finite and can be formed.

Allow the remaining constants to depend on the fixed mass interval,
and choose the physical bilinear-source radius at this point.
More explicitly, put
$S_D=\sum_{y\in\mathbb Z^3}e^{-\alpha_D|y|/2}<\infty$.
The matrix-mass radius, already fixed from the early operator
constants, is
\[
 r_M=\min\left\{1,\frac1{4C_{\mathcal K}},\frac1{4C_DS_D},
 \frac1{8e^{3\alpha_D/2}b_F^2C_{\mathcal K}^2C_DS_D},
 \frac{L\varrho_*}2,\frac L{C_{\rm fl}}\right\}>0,
\]
where
\[
 c_{\mathsf M}=2600\nu_P/\alpha_D,\quad
 C_{\rm bl}=\nu'\left(2+\frac4{L-1}\right),\quad
 C_{\rm fl}=4C_{\rm bl}c_{\mathsf M}C_H.
\]
Here $\nu'=8$, $\nu_P$ is the number of fine places over a
site, and $C_H$ is the original background constant.
Thus the matrix-mass change costs at most one on a stored
formal leg. The radius is chosen from earlier constants;
the later auxiliary fibre norm does not change it.
Fix $0<c_J\le\min(1/2,\varrho_*,r_M/L)$ and
$0<2\varrho_J\le c_Jm_-$. Thus all physical sources lie
inside the full matrix-source operator ball. Put
$h_*=\max\{h_F,\sqrt{2/(c_bm_-)}\}$.
For the final Berezin bound recompute the norm constants
at $H=2A_\kappa h_*$. This changes only that auxiliary norm
estimate and its threshold. The flow and source-chain
constants just fixed remain those at $h_F$.

\emph{The final threshold.}
Choose $\gamma$ last. Every remaining requirement is a
positive upper bound on a function tending to zero with
$\gamma$, or a finite lower bound on $\log\gamma^{-2}$.
This includes the original pure requirements, the background
regularity bound $\bar C_{\rm reg}\epsilon(\gamma^{-2})\le\epsilon_*$,
the relative-coefficient bound $4\bar B\gamma^2\le\sigma_*$,
and $\gamma^{-2}\ge e^{18\mathfrak p_0}$.
For the actual density-test fields and covariance classes require
\[
\begin{gathered}
 2\epsilon(\gamma^{-2})\le\varepsilon_{\mathcal W}\alpha_0,
 \qquad 2\epsilon(\gamma^{-2})\le\varepsilon_{\mathcal W}\mathsf a_H,
 \qquad \delta(\gamma^{-2})\le\mathsf d_H,\\
 C_H^{\mathcal W}\delta(\gamma^{-2})
                      \le(1-L^{-2})e_{\rm cov}.
\end{gathered}
\]
The last inequality closes the step's estimate
$\mathsf m_k\le L^{-2}\mathsf m_{k+1}
+C_H^{\mathcal W}\delta(\mathsf t_k)$ at the same covariance
radius. The constant $C_H^{\mathcal W}$ belongs to the gauge
sum rule and is distinct from the fermion background constant.

The actual history Gaussian is used at the positive real
ratio $\tau=x_j/\mathsf t_j$, with $|\tau-1|\le1/4$.
Its construction and retained collar term are those of
\cref{C-399-44,C-399-43}. Impose the explicit frame thresholds
and the three $\varepsilon_\flat/4$ bounds in
\cref{C-399-47,C-399-48,C-399-21}. The three-share
recursion of \cref{C-399-21} has factor $64/51$,
so these three shares give
$\frac{64}{51}\frac34\varepsilon_\flat
=\frac{16}{17}\varepsilon_\flat<\varepsilon_\flat$.
In particular include the boundary-sum ceiling
\[
 16C^{\rm bd}\frac{A_1}{C_0}
 (\log(\tfrac12\gamma^{-2}))^{-(q-\mathfrak p_0)}
 \left(\tfrac98K^{\rm lb}+108M^3K_M\right)\le1,
 \qquad q>\mathfrak p_0.
\]
All quantities except $\gamma$ have been fixed.
The retained Gaussian mismatch requires, in addition,
\[
 \frac32\bar B\gamma^2 C_{\rm xt}^2 A_1^2
 (\log\gamma^{-2})^{2\mathfrak p_0}
 \left(c_{\rm kp}^{\rm th}+c_{\rm rj}^{\rm th}
                   +\frac32c_{\mathsf G}(LM)^3\right)\le1.
\]
The coefficients here are the fixed kept, rejected and gauge
quadratic-form constants. This is the domainwise bound of
\cref{C-399-49}, including the strip rows.
The effective scalar-frame threshold is
$\log\gamma^{-2}\ge Y_{\rm eff}^{\rm F}(2\varepsilon_\flat/5)$,
formed with $(E_0,B,B_{\rm fr}^{\flat})$ as in
\cref{C-399-47}; the q-class and failed-test bounds of
\cref{C-399-44} are imposed at the same stage.

The source scalar arrays use the larger budget
\[
 E_0^{\rm s}=E_0+2+K_M^\bullet c_{\rm cl},\qquad
 K_M^\bullet=K_M+C_EK_M^{\rm v}.
\]
Their boundary budget is
\[
 B_\natural=B+8M^3(E_0^{\rm s}-E_0)
          +C_v^{\rm rj}+C_v^{\rm sr}+2+K_{\rm rb}c_{\rm cl}/16.
\]
Use the source effective-frame threshold
$\log\gamma^{-2}\ge Y_{\rm eff}(\varepsilon_\flat/(2+2\sqrt2))$
formed at the doubled box
$\sqrt2(E_0^{\rm s},B_\natural,B_{\rm fr}^{\flat}+1)$
of \cref{C-407-38}, and the three source shares, including
the formal-leg bound of \cref{C-407-41}. These are
requirements on the actual complex source step; its identity
does not follow from positivity at zero sources.
Among the source ceilings are
\[
\begin{gathered}
 C_{\rm rad}\delta(\gamma^{-2})\le\alpha_0,\quad
 16\varrho\le r_v(\gamma^{-2}-\bar B),\quad
 64\varrho\le r_x\gamma^{-2},\\
 32C_{\rm bd}^{\rm s}K_M C_{\mathcal A}
                         \delta(\gamma^{-2})\le\beta_0.
\end{gathered}
\]
The source class with fermion legs is bounded by a fixed
constant times $\phi(\gamma^{-2}/2)$. The linear collar
term is bounded by a constant times
$\gamma(1+2\log\gamma^{-1})^{\mathfrak p_0}$, and the
cross term by a constant times
$\gamma^2(1+2\log\gamma^{-1})^{\mathfrak p_0}$.
The kept quadratic term leaves half the retained positive
floor; the rejected term is estimated using
$\mathfrak b(t)^2e^{-(\check R-2)}\le C_{\mathfrak b}$.
The precise coefficients and thresholds were fixed in the
source-history proof. Each additional condition is therefore
satisfied by reducing $\gamma$.
The constant-coordinate torus comparison also requires
\[
 C_\psi\phi(\tfrac12\gamma^{-2})\le c_{\rm cl}/20,\qquad
 \epsilon_\gamma S_L\le
 \min\left\{\tfrac12,
       \frac{1-2^{-1/2}}{4C_{\rm dl}C_\omega S_\omega}\right\}.
\]
Its constants are formed before $B_\sigma$, and
$\epsilon_\gamma\to0$ as $\gamma\to0$.

For the common lower bound impose
\[
 C_{\mathsf S_1}\bar B_M\gamma^2\le c_bm_-/8,
 \qquad \phi(\gamma^{-2})\le c_*(m_-),
 \qquad \gamma\le\gamma(2A_\kappa h_*,5B_s).
\]
The first pays the running mass member; the second pays
the activity bilinear and the last interacting forest
expansion. The threshold $\gamma(h,A)$ is obtained by
taking the infimum of its positive continuous ceilings
on the compact rectangle $[h_F,h]\times[1,A]$.
It is positive and decreases in $h,A$. Thus this last
choice implies all earlier uses at $h_F$ and radius $5B_s$.

Finally choose the history-sum threshold with the actual
matter and source majorants. In the unmarked application
take $a=1/2$ and rate $c_*/2$. More generally, a fixed
deposit mark with created-cell coefficient $c_\Theta<c_*$
uses $a=(1-c_\Theta/c_*)/2$ and rate
$c=(c_*+c_\Theta)/2>c_\Theta$. Its deposit dials can be
written as $(w',c_{\rm Dep}')=(c_\Theta/c,c)$, so that
$w'<1$ and $c_{\rm Dep}'\le c$. Apply the actual thresholds
of \cite[Prop.~67.6]{YMM1}, with these dials, the enlarged majorant
and the logarithmic label-error constant. Their hypotheses
hold since $2\mathfrak p_0>3r_c+2$ and the remaining
threshold is finite. The quarter-stock allocation in
Prop.~\ref{C-407-22} gives this integrated rate; it
does not assert a stronger pointwise loss-$a$ property.
The later insertion ceilings are also scalar envelopes
tending to zero. A finite minimum of all these positive
thresholds completes the common choice.

\emph{Application of the gauge induction.}
The pure gauge constants are chosen in the order of
\cite[Prop.~69.1 and its proof, pp.~595--599]{YMM1}.
At each stage above, the additional matter conditions impose
floors, positive ceilings or nonempty windows on freely chosen
constants, using only quantities already fixed. Every new constant
defined by a formula is placed after its arguments. The defining
formulas for the pure gauge constants are retained. Requirements
on earlier constants are imposed at their original stages, before
the subsequent constants are evaluated.

In particular, the condition $B\ge B^F$ is obtained by requiring
$8M^3E_0\ge B^F$ in the choice of $E_0$. The condition
$e_{\rm cov}\le a_Q$ is obtained by imposing
$\varepsilon_{\mathcal W}\alpha_0\le2a_Q$ on $\alpha_0$.
The condition $C_{\mathcal A}\ge c_{\rm sh}$ is obtained from
the floor on $A_0$ checked above. In each case the original
formula gives the required value; no constant defined by a
formula is chosen independently of its arguments.

Choose the constants successively with these additional bounds.
Omitting the matter-only constants then leaves a choice satisfying
the pure gauge conditions, since their order, formulas and
inequalities have been preserved. Finally, all conditions on
$\gamma$ hold throughout a sufficiently small interval
$(0,\gamma_*]$. For the pure conditions this follows from their
uniform imposition for $x\ge\gamma_*^{-2}$ in the proof of
\cite[Prop.~69.1]{YMM1}; for the additional conditions it follows
from the positive thresholds and vanishing upper bounds verified
above. Taking their finite minimum completes the choice and shows
that decreasing $\gamma$ preserves it.

\Cref{def:qcd-gauge-admission} therefore supplies all
conditions of the pure gauge induction at this one restricted
assignment, with the canonical data and substitutions of that
lemma. Both coupling sequences used here are in its
range. The calibration sequence
$\mathsf t_j=L^{K-j}\mathsf b$ has zero increments.
At each stage of the running construction, the increments
already obtained satisfy $|b_j|\le B_\sharp$. For
$\mathsf b\in\bar I(g)$ the recurrence
$x_{j+1}=L^{-1}x_j-b_{j+1}$ gives
\[
 \mathsf t_j\ge g^{-2}+\bar B,\qquad
 |x_j-\mathsf t_j|\le\bar B,\qquad x_j\ge g^{-2}.
\]
The allowed torus sides supply the required chart room for
both sequences, as in the proof of \cref{C-388-38}.
The gauge induction consequently applies to each finite
prefix of either sequence. It is quantified over all such
admissible sequences; its use does not presume that a full
matter flow has already been constructed.

\emph{The matter and source steps.}
The gauge result supplies the charts, auxiliary systems and
coordinate identities. We now verify its applications to
the actual fermion and source coefficients using the
estimates proved above.

The fermion covariance and its local normalizer are constructed
on their intrinsic domains in \cref{C-397-40,C-397-47}.
The density-test and sum-rule applicability is
\cref{C-388-38}; the small-field induction and mass estimates
give the regular activities before any large-field history
is used. The histories then use the algebra-valued expansion,
the positive real Gaussian rescaling and the agreement on
interior domains of \cref{C-399-41,C-399-44,C-399-45}.
The degree-zero arrays and the reentering boundary terms obey
\cref{C-399-47,C-399-48}. Consequently \cref{C-399-43}
applies at every finite step, as an identity of measures in
the even Grassmann algebra.

For complex sources, \cref{C-407-34} gives the coefficientwise
measure identity on one common full-measure set, and
\cref{C-407-37} closes the source induction with the scalar
and formal-leg estimates just used. The common vacuum
assignment is preserved by these identities. The final
history sum and lower bound therefore give
\cref{C-407-22,C-407-23} and the uniform cumulant estimates.
This proves the asserted finite construction.
\end{proof}

The range of bare labels is a full tail, despite its
description by reference intervals. Decrease $\gamma$ so
that $\gamma^{-2}\ge2\bar B$. For any
$\mathsf b\ge\gamma^{-2}+\bar B$, take
$g=(\mathsf b-\bar B)^{-1/2}\le\gamma$.
Its lower endpoint is $\mathsf b$, while its upper endpoint
is $L\mathsf b-2\bar B\ge\mathsf b$.
Consequently
\[
 \bigcup_{0<g\le\gamma}\bar I(g)
        =[\gamma^{-2}+\bar B,\infty).
\]
For each such $g$ the admissible sides are cofinal.
Thus the volume-first conclusions apply to every
sufficiently large fixed $\mathsf b$, with the stated
positive mass interval. Using the closed intervals is
essential at the displayed lower endpoint.

\section{The short-distance two-point functions}
\label{sec:qcd-two-point-laws}

We keep the physical coupling $\mathsf b$ fixed and the mass in
$[m_-,m_+]\subset(0,1/2]$. Write
$\lambda_j=L^{j-K}$, $\widehat\phi(t)=(1+\log t)^{\mathfrak p_0}$
and $\phi(t)=t^{-1/2}\widehat\phi(t)$.
Only two-point functions of the action density, or of two bilinears,
are considered in this section.

For $0<\vartheta<1$ and $d>0$, a test pair consists of nonnegative
$f,h\in C^1_c$ supported in balls of radius $d$, with
\[
 d\le |x-y|\le10d\quad(x\in\operatorname{supp}f,\ y\in\operatorname{supp}h),
 \quad \int f\ge\vartheta d^3\|f\|_\infty,
 \quad \operatorname{Lip}f\le\frac{\|f\|_\infty}{\vartheta d},
\]
and the same last two conditions for $h$. We also require the pair to lie in a common
ball of radius $5d$. Put
\[
 \mathcal K_q(f,h)=\iint f(x)h(y)|x-y|^{-q}\,dx\,dy.
\]
Then
\begin{equation}\label{eq:qcd-template-lower}
 10^{-q}\vartheta^2d^{6-q}\|f\|_\infty\|h\|_\infty
 \le\mathcal K_q(f,h)
 \le(4\pi/3)^2d^{6-q}\|f\|_\infty\|h\|_\infty.
\end{equation}

\begin{proposition}[The two leading kernels]\label{C-408-27}
Every volume-first continuum limit constructed below satisfies, on
these test pairs,
\[
 S_2((0,f),(0,h))=\frac9{16\pi^2}\mathsf b^{-2}
                  \mathcal K_6(f,h)(1+\epsilon_0),
\]
\[
 S_2((\Gamma,f),(\Gamma',h))=\frac1{2\pi^2}
                  \mathcal K_4(f,h)(\delta_{\Gamma\Gamma'}+\epsilon_B).
\]
For each fixed $\vartheta$, the errors are bounded by a function
$\epsilon(d)$ tending to zero as $d\downarrow0$. The bound is
uniform in the permitted states and in the mass interval. In
particular, each displayed field has a nonzero two-point function.
\end{proposition}

\subsection{Reduction to a local source energy}

Let $E_K(s)$ be the source energy, and let $e_j(X;s)$ be its
level-$j$ terms, so that $E_K(s)-E_K(0)=\sum_{j,X}e_j(X;s)$.
The local source bound gives, for sources supported in the ball
of radius $1/4$ about a board-cell centre,
\begin{equation}\label{eq:qcd-cell-holomorphic-bound}
 \Psi(s)=e^{-\Lambda(s)}\frac{Z(s)}{Z(0)},\qquad
 \Psi(0)=1,\qquad |\Psi(s)|\le B,
 \quad
 \Lambda(s)=\frac1{\mathsf N}\sum_{r}\{E_K(rs)-E_K(0)\}.
\end{equation}
Here $r$ runs over the $\mathsf N$ board reflections, and $B$
is independent of the cutoff and the side. The source domain is
the fixed domain of the preceding section. Write
$\|F\|_\circ=\|f\|_\sharp/\varrho$ for $F=(0,f)$ and
$\|F\|_\circ=\|f\|_\infty/\varrho_J$ for $F=(\Gamma,f)$.

\begin{lemma}[Packing with signs]\label{C-408-20}
For a test pair with $8a_K\le d\le1/200$, its truncated two-point function
differs from a mean of mixed derivatives of $E_K$, over translates
of the pair and their board reflections, by at most
\[
 C d^3\|F\|_\circ\|H\|_\circ.
\]
The packing has $N\ge(96d)^{-3}$ translates per cell. Its union
meets at most $C N\mathsf N(d/\lambda)^3$ cubes of side
$\lambda>d$.
\end{lemma}

\begin{proof}
For separated sources $F,H$, the source norm of $z_1F+z_2H$
is the maximum of the two norms. This is immediate for the supremum
norm. For the Lipschitz norm, put
\[
 L_*=\max\{\operatorname{Lip}(z_1f),\operatorname{Lip}(z_2h)\}.
\]
If $x,y$ lie in different supports,
let the segment leave the first at $p$ and enter the second at $q$.
The sum of the two function values is bounded by
$L_*(|x-p|+|q-y|)$, hence by
$L_*|x-y|$. Thus disjoint translates with arbitrary
signs retain the same norms.

On the open bidisc $|z_i|<2/\|F_i\|_\circ$, the function in
\eqref{eq:qcd-cell-holomorphic-bound} is bounded by $B$.
Apply Schwarz's lemma on each complex line to $(\Psi-1)/(B+1)$.
On the bidisc smaller by $1/(4B)$, $|\Psi-1|\le1/2$.
Its logarithm therefore exists there. Apply Cauchy's estimate
on smaller circles and let their radii increase to obtain
\[
 |\partial_1\partial_2\log\Psi(0)|
                  \le4B^2\|F\|_\circ\|H\|_\circ.
\]
The derivative is the bilinear form $S^T_{2,K}-\partial_1\partial_2\Lambda$.
Choose translates $F_t,H_t$ of the pair inside the cell ball,
with the balls of radius $5d$ mutually at distance at least $d$.
A lattice grid with mesh between $12d$ and $13d$, in a centred
cube of side $1/4$, supplies at least $(96d)^{-3}$ translates.
The cube and the balls of radius $5d$ lie in the cell ball, since
$\sqrt3/8+5d<1/4$. In each coordinate the number of grid points
is at least $(52d)^{-1}-1\ge(96d)^{-1}$.
Average the preceding bilinear form on
$\sum_t\varepsilon_tF_t,\sum_t\varepsilon_tH_t$ over independent
signs. Only equal indices survive, so the mean over the $N$ pairs
has bound $4B^2N^{-1}\|F\|_\circ\|H\|_\circ$.
Lattice translation invariance makes the truncated correlation the
same for every pair. The mean of $\partial_1\partial_2\Lambda$
remains, since $E_K$ itself depends on the board. This proves the
claim. The same grid, repeated by board reflections, meets at most
$C N\mathsf N(d/\lambda)^3$ cubes of side $\lambda>d$:
each cell has a fixed-volume occupied region and $N\ge c d^{-3}$.
\end{proof}

\subsection{The Gaussian and fermion parts of the energy}

At level $j$ use the source domain
\begin{equation}\label{eq:qcd-level-source-domain}
 \|w\|_\infty,\ 40M\lambda_j\operatorname{Lip}w
       <\frac{\varrho}{\widehat\phi(x_j)},\qquad
 \lambda_j\|J\|_\infty<R_J,
\end{equation}
where $R_J=\min\{r_M,\varrho_S,\varrho'_S/L\}$ is below the
matrix-mass and free-operator source radii. The steps through
level $j$ are made on one constant pair-radius system
$a=C_{\rm rad}\delta(\mathsf t_{j-1})$. At earlier levels $i\le j$ the
linear gauge-source term has size $t_i^{\rm lin}$, with
\[
 t_i^{\rm lin}\le\sqrt2L^{-(j-i)/2}T,\qquad T=t_j^{\rm lin}<c_{\rm cl}/16.
\]
This term need not be $O(\phi(x_i))$. Its Gaussian contribution
must be retained.

The comparison with a Gaussian expansion involves two changes:
removing the tests, and changing the sort assigned to a potential
that is defined at every real fluctuation. The latter change is
not an identity of individual cluster terms. The following bound
accounts for it.

\begin{lemma}[Changing the sort of an all-fluctuation family]\label{C-408-34}
Assume the local-step and cluster-expansion hypotheses of
\cref{C-388-34,C-388-19}, with a region $R$, its decoupled
Gaussian structure, finite threshold $\mathfrak b$ and all bond
parameters equal to one. Let a base system have only origins of
sort $\Sigma_{\rm q}$. Let $m\ge1$ families $\mathbb W_r$ be
defined at every real fluctuation and satisfy
$\|\mathbb W_r\|_{\star,\lambda_V}\le w_\star$. If the base
has at most $4-m$ origins, denote by $E_{\rm v}$ and $E_{\rm q}$
the cluster sums after adjoining these families with the
indicated sort. Then, for $X\subset R$,
\[
 |E_{\rm v}(X)-E_{\rm q}(X)|
       \le2K_P^\Pi p_\star e^{-\lambda_{\rm out}d_t(X)}.
\]
If the base has at most $3-m$ origins, an additional family
$\dot{\mathbb W}$ defined at every real fluctuation, with
$0<\|\dot{\mathbb W}\|_{\star,\lambda_V}<\infty$, satisfies
\[
 |E'_{\rm v}[\dot{\mathbb W}](X)
             -E'_{\rm q}[\dot{\mathbb W}](X)|
 \le2K_P^\Pi p_\star
       \frac{\|\dot{\mathbb W}\|_{\star,\lambda_V}}{w_\star}
                           e^{-\lambda_{\rm out}d_t(X)}.
\]
The constants are those of the cluster expansion; in its
small-field application $w_\star=c_{\rm cl}$ and $p_\star$
is its large-field bond weight after the bond count.
\end{lemma}
\begin{proof}
For a set $\mathbf D$ of vertices let $Y^{\mathbb W}(\mathbf D)$
be the union of the domains belonging to the added families.
An index of the system of sort $\Sigma_{\rm v}$ has its
large-field bonds outside $\mathfrak X\langle
Y^{\mathbb W}(\mathbf D)\rangle$ and has the inner test on
that set. Expand this test as
\[
 \mathsf{in}_{\mathfrak e'}=
  \prod_{b\in\mathfrak e'^+}(1-1_{\{|B(b)|\ge\mathfrak b\}})
  =\sum_{\mathfrak e''\subset\mathfrak e'}
                (-1)^{\#\mathfrak e''^+}\mathsf{out}_{\mathfrak e''}.
\]
This is legitimate at every real fluctuation because the
potentials are defined there. Conditioning and the decoupling
bracket are linear. The terms with $\mathfrak e''=\emptyset$
agree index by index with the terms of the system of sort
$\Sigma_{\rm q}$ whose large-field bonds avoid the added
vertices. In the difference, every remaining term therefore
has a marked large-field bond lying in an added vertex.

For the terms already present in the system of sort
$\Sigma_{\rm q}$, the marked-bond estimate gives
$K_P^\Pi p_\star e^{-\lambda_{\rm out}d_t(X)}$.
The terms created by the test expansion have the same Gaussian
majorant: each extra bond carries its usual weight $p$.
Give such a bond the region $b^\sharp$ and the enlarged weight
$\widehat w(b)=pe^{4\lambda_1}$. The map from an expanded
index to its vertices, old bonds and extra bonds is injective:
the vertices determine $Y^{\mathbb W}$, and membership in
its bond set distinguishes the extra bonds from the old ones.
The old regions already cover the cluster; the additional
bond regions only increase the covering sum. At a coarse cube
there are at most $6((L+2)M)^3$ such oriented bonds. Their
enlarged weights sum to at most $p_\star$, while all element
weights sum to at most twice the number of coarse cubes.
Thus the unmarked activity majorant is unchanged, and its
marked version has one extra factor $p_\star$:
\[
 \bar h'(\gamma)\le e^{8c_w-4\lambda_1}
                            e^{-\lambda_2d_t(\gamma)},\qquad
 \bar z'(\gamma)\le p_\star e^{8c_w-4\lambda_1}
                            e^{-\lambda_2d_t(\gamma)}.
\]
The cluster count gives the same bound
$K_P^\Pi p_\star e^{-\lambda_{\rm out}d_t(X)}$ for these
terms and proves absolute convergence of the expanded sum.
Adding the two bounds proves the first assertion.

For an insertion adjoin the extra origin
$\epsilon\dot{\mathbb W}$. Both systems are holomorphic for
$|\epsilon|<w_\star/\|\dot{\mathbb W}\|_{\star,\lambda_V}$.
Apply the first assertion on each smaller circle and then
Cauchy's estimate at zero. Letting the radius increase to
the stated one gives the insertion bound.
\end{proof}

\begin{lemma}[Energy decomposition]\label{C-408-21}
On \eqref{eq:qcd-level-source-domain},
\[
 e_j(X;w,J)=e_j^{\rm G}(X;w)+e_j^{\rm F}(X;J)
                 +e_j^{\rm H}(X;w,J)+e_j^{\rm I}(X;w,J).
\]
The sums of the first two terms satisfy the following identities
of logarithmic germs at zero sources:
\[
 \sum_{j,X}e_j^{\rm G}
 =\log\frac{Z^{\rm G}[w]}{Z^{\rm G}[0]}
       -\log\frac{Z^{\rm G}_{\rm last}[w]}{Z^{\rm G}_{\rm last}[0]},
\]
\[
 \sum_{j,X}e_j^{\rm F}
 =\log\frac{\det(D[1]+m_0-a_K\hat J)}{\det(D[1]+m_0)}
       -\log\frac{\det\mathbb D_K(1;J)}{\det\mathbb D_K(1;0)}.
\]
Here $Z^{\rm G}$ is the linearized gauge Gaussian integral, with
the constant last-level gauge mode fixed. The fermion identity also defines the
analytic continuation of the quotient of determinants by the product
of the invertible fluctuation determinants. Separate global logarithms
of its numerator and denominator are not required on the enlarged
operator-source ball; both are invertible at zero sources. The term $e_j^{\rm H}$
vanishes to fourth order in $w$, uniformly in $J$. Moreover
\[
 |e_j^{\rm I}(X;s)|\le C_I\phi(x_j)e^{-\kappa d_j(X)},
\]
and changing a source outside $X$ gives the same bound multiplied by
$\exp[-\kappa_\sigma\operatorname{dist}(X,\operatorname{supp}(s-s'))/\lambda_j]$.
\end{lemma}

\begin{proof}
We give the decomposition of the scalar activities, since it is
needed to keep the error small on the domains in
\eqref{eq:qcd-level-source-domain}. Let $\mathcal L_i$ be the
linear source term plus the linear variation obtained by localizing
only the fluctuation-dependent $b$-parts of the earlier Gaussian
activities; let $q_i$ be the quadratic variation of those same parts.
The pair-dependent $c$-parts are included in $V_i^{\rm rest}$ below:
their absolute localization has the small shift factor
$\delta(\mathsf t_{i-1})$. Both $\mathcal L_i$ and $q_i$ are given at all
fluctuations. In the untested Gaussian expansion form
the following families at a pair $P$. Let
$\mathcal V_i^{v'}(P)$ be the relative-weight quadratic origin
of the source step at $\tau=1$, with
$v'_{i-1}=-L^{K-i+1}w/\mathsf t_{i-1}$.
Write $E^\circ_{P,i}[\mathcal W](X)$ for the untested Gaussian
cluster logarithm with origins $\mathcal V_i^{v'}(P),\mathcal W$,
less its value with only the first origin. All logarithms use
the branch at zero sources. Define
\[
 \mathcal E_i(\eta)(X;P)
 =E^\circ_{P,i}[\eta\mathcal L_i(P)+q_i(P)](X),\qquad
 F_i^{\rm G}=\mathcal E_i(0)+\tfrac12\mathcal E_i''(0).
\]
More explicitly put
\[
 c_i(X;P)=E^\circ_{P,i}[q_i(P)](X),\qquad
 b_i(X;P;\mathcal L,\mathcal L')=
 \left.\tfrac12\partial_{\eta_1}\partial_{\eta_2}
 E^\circ_{P,i}[\eta_1\mathcal L+\eta_2\mathcal L'+q_i(P)](X)
 \right|_{\eta_1=\eta_2=0}.
\]
Then $F_i^{\rm G}(X;P)=c_i(X;P)+
b_i(X;P;\mathcal L_i(P),\mathcal L_i(P))$.
These are the c- and b-parts used in the recursive localization.
The origin containing the weighted quadratic form is held fixed.
These families are formed at the pair itself. Cauchy's estimate
in the two linear parameters gives
\[
 |c_i(X;P)|\le K'_M\|q_i(P)\|e^{-\kappa_+d_i(X)},\qquad
 |b_i(X;P;\mathcal L,\mathcal L')|
  \le8K'_Mc_{\rm cl}^{-1}\|\mathcal L\|\|\mathcal L'\|
                                  e^{-\kappa_+d_i(X)}.
\]
The bilinear coefficient is bounded for linear terms of any
size, using discs of radii $c_{\rm cl}/(4\|\mathcal L\|)$
and $c_{\rm cl}/(4\|\mathcal L'\|)$.

The Gaussian model uses the appropriate jets at the trivial
pair. Constant gauge transformations fix that pair and act on
the fluctuation by the adjoint representation, which has no
invariant vector. Hence $\mathcal L_i(1)=0$. Let
$\mathcal L_i^{\rm G}(A;B)$ be the first background-field jet
of $\mathcal L_i$ at the pair of $e^{iA}$; it is bilinear in
$A,B$. Put $Q_i^{\rm G}(B)=q_i(1;B)$. These jets are
independent of the localization radii. At the trivial pair set
\[
 \mathcal E_i^{\rm G}(\eta)(X;A)
  =E_i^\circ[\eta\mathcal L_i^{\rm G}(A)+Q_i^{\rm G}](X),
 \quad
 F_i^{\rm G,lin}(X;A)
   =\mathcal E_i^{\rm G}(0)(X;A)
               +\tfrac12\partial_\eta^2\mathcal E_i^{\rm G}(0)(X;A).
\]
Reflection of the fluctuation makes the expansion even in
$\eta$. The family $F_i^{\rm G,lin}$ is a polynomial of
degree at most two in $A$, with no linear part. At the trivial
pair $F_i^{\rm G}(X;1)=c_i(X;1)=F_i^{\rm G,lin}(X;0)$.
The pair-dependent bilinear part has no terms of background
orders zero or one; its quadratic jet is that of
$F_i^{\rm G,lin}$. The higher jets of the pair-dependent
constant part belong to the remainder, not to the Gaussian
model.

At the trivial pair the decoupled structure is complete on
the torus. The sum over $X$ is therefore the logarithm of the
whole Gaussian integral,
\[
 \sum_X\mathcal E_i^{\rm G}(\eta)(X;A)
   =\log\int d\mu_i^{\rm G}(B)
        e^{\eta\langle l_i(A),B\rangle+Q_i^{\rm G}(B)},
\]
where $\mu_i^{\rm G}$ is the normalized Gaussian at the
relative weight. This logarithm is quadratic in $\eta$.
Thus the local terms of orders at least four have zero sum,
although they need not vanish polymer by polymer. Completing
the square gives
\[
 \sum_XF_i^{\rm G,lin}(X;A)
   =\tfrac12\langle l_i(A),\mathcal G_i l_i(A)\rangle
                           +\log\mathfrak Z_i(w),
\]
where $\mathcal G_i$ and $\mathfrak Z_i$ are the covariance
and mass after adding $Q_i^{\rm G}$. Take the localization
sum rule in background/fluctuation bidegrees $(0,2)$ and
$(1,1)$. It identifies $Q_i^{\rm G}$ and $l_i$ with the old
quadratic tree terms transported through the linearized
background maps. These terms include the cross terms between
different slices. Together with the logarithmic mass of the
relative-weight Gaussian origin they give exact successive
Gaussian integrations. Their telescoping is the Gaussian
identity in the statement, with the constant block mean fixed
to remove the harmonic zero mode.

The constructions are made at the pair, before evaluation at the
trivial pair. Use the bond tents $h_z(b)$ of
\cite[Equation~(760)]{YMM1}, so that $\sum_zh_z(b)=1$.
For $\mathcal W_\eta=\eta\mathcal L_i(P)+q_i(P)$ put
\[
 \mathcal W_{\eta,t}(Y;B)=\mathcal W_\eta(Y;tB),\qquad
 \dot{\mathcal W}_{\eta,t,b}(Y;B)=
 \left.\partial_{t_b}\mathcal W_\eta(Y;(t_cB(c))_c)\right|_{t_c=t}.
\]
Let $\dot{\mathcal W}_{\eta,t,z}=\sum_bh_z(b)\dot{\mathcal W}_{\eta,t,b}$.
Keeping the first origin fixed, the pointed family is
\[
 F_{i,\eta}^\bullet(X,z)=\int_0^1
 E'_{P,i}[\mathcal W_{\eta,t};\dot{\mathcal W}_{\eta,t,z}](X)\,dt,
\]
where the prime differentiates the second origin. At $t=0$ that
origin is zero, so summing over $z$ gives $\mathcal E_i(\eta)$.
Its orders zero and two in $\eta$ give the pointing of $F_i^{\rm G}$.
Thus localization commutes with the split, pointer by pointer.
Polarization and $|h_zB|\le|B|$ give the insertion bound
$|\eta|\|\mathcal L_i(P)\|+2\|q_i(P)\|$. It follows that
\[
 |c_i(X;P;z)|\le2K'_M\|q_i(P)\|e^{-\kappa_+d_i(X)},\qquad
 |F_i^{\rm b}(X;P;z)|\le6K'_Mc_{\rm cl}^{-1}
                  \|\mathcal L_i(P)\|^2e^{-\kappa_+d_i(X)}.
\]
These coefficient bounds hold for linear terms of any size;
the pointed bilinear part vanishes at the trivial pair.
The Gaussian activity has a constant part and a part bilinear in
$\mathcal L_i$. On radius $a$ their bounds are respectively
$C K'_M t_i^{\rm lin} e^{-\kappa_+d(X)}$ and
$C K'_M c_{\rm cl}^{-1}(t_i^{\rm lin})^2e^{-\kappa_+d(X)}$.
On the fixed radius the latter has the additional factor
$(\alpha_0/a)^2$. The scalar localization has lag weight
$L^{-(i-1-i')}$ and factor at most
$C\,C_{\mathcal A}/C_{\rm rad}$. Since
\[
 \sum_{i'<i}L^{-(i-1-i')}t_{i'}^{\rm lin}\le\tfrac12t_i^{\rm lin},
\]
the two radius floors in the common constants give
$\|\mathcal L_i\|\le17t_i^{\rm lin}/16$ and
$\|q_i\|\le t_i^{\rm lin}/16$ on radius $a$.
On the fixed radii the bounds are
$\|\mathcal L_i\|\le2(\alpha_0/a)t_i^{\rm lin}$ and the
same bound for $q_i$. The bilinear families, summed and
pointed, are thus at most
$10K'_Mc_{\rm cl}^{-1}(t_i^{\rm lin})^2e^{-\kappa_+d_i(X)}$
on radius $a$ and at most
$32K'_Mc_{\rm cl}^{-1}(\alpha_0/a)^2(t_i^{\rm lin})^2
e^{-\kappa_+d_i(X)}$ on the fixed radii. The constant-part
variation from the trivial pair is at most
$K'_Mt_i^{\rm lin}e^{-\kappa_+d_i(X)}/8$, or twice this
at a pointer. This proves inductively the definedness of
every Gaussian family used at the next level.

For the higher fluctuation Taylor terms use the absolute localization
on a fixed auxiliary disc of radius
$\mathsf C_{\mathcal A}\delta_{\rm ax}$, where
\[
 \delta_{\rm ax}=\min\{1,\mathsf s^\star\alpha_0/(1+\beta_s),
       \varepsilon_{\mathfrak h}/\mathsf C_{\mathcal A},
       \delta(e^{2\mathfrak p_0})\}=\delta(\mathsf t_{\rm ax})>0
\]
is fixed before $h_F$ and $\gamma$. It is the same localization on its smaller step
disc. Its order-$n$ coefficient costs
$\rho_i^n$, $\rho_i=\delta(\mathsf t_{i-1})/\delta_{\rm ax}\le1/2$.
The cubic remainder of the bilinear Gaussian activity is therefore
bounded by
\[
 C\rho_i^3(\alpha_0/a)^2Tt_i^{\rm lin}\le C\phi(x_{i-1}).
\]
The constant Gaussian parts have the same bound after localization,
by its factor $C\delta(\mathsf t_{i-1})/\alpha_0$.
The gauge vertices, explicit closed-loop variations, the mass
response and all fermion activities also contribute at most
$C\phi(x_{i-1})$. For closed loops use the array bound of
\cref{C-397-42}; their value and first gauge variation vanish
at the flat field at every $J$. An axial gauge on a box containing
$X$ and second-order Schwarz give an extra squared curvature
threshold; the spare exponential rate pays the box-size polynomial.
No through-cube norm at rate $\kappa$ is inferred from a single
spare unit of array decay.

Inductively split the old scalar family, at every pointer, as
$F_i^{\rm G}+F_{i,\ge4}+F_{i,\rm h}+F_{i,\rm r}$, with
the Gaussian family taken less its trivial-pair value.
Put $K''=16c_{\rm ins}K_M(1+c_{\rm cl}^{-1})$ and
$N_{\rm h}=K_M(1+c_{\rm cl})/2$. The inductive bounds are
\[
 \|F_{i,\rm r}\|_{i;a}\le C_{\rm r}^{\rm lev}\phi(x_{i-1}),
 \qquad \|F_{i,\ge4}\|_{i;a}+\|F_{i,\rm h}\|_{i;a}
                                                   \le N_{\rm h}.
\]
We carry locality and holomorphy in the independent
coordinates through this induction. Write
$\mathcal O_j^{\rm lev}$ for the domain
\eqref{eq:qcd-level-source-domain} and put
\[
 \mathfrak C_j=
 \left\{\mathfrak c=(w,-a_KJ,(\sigma_k)_{k<j}):
 (w,J)\in\mathcal O_j^{\rm lev},\quad
 \|\sigma_k\|<2B_\sigma\phi(x_k)\ (k<j)\right\}.
\]
Here every $\sigma_k$ is a field on the bare sites with
values in $\mathbb M$, and its norm is the supremum over
those sites. On the finite torus this is an open set
in a finite-dimensional complex vector space. At every
pointer and fixed point of a core of radius $a$, each
of $F_i^{\rm G}$, its constant and bilinear parts,
$F_{i,\ge4}$, $F_{i,\rm h}$ and $F_{i,\rm r}$ is
holomorphic in $\mathfrak c\in\mathfrak C_j$ and reads
these coordinates in its polymer alone. The Gaussian
families have this property also on the fixed pair
radii. They and $F_{i,\ge4}$ read only $w$.
Each family is a local, gauge-invariant function of
the pair, holomorphic on the cores of radius $a$.
The bounds above hold throughout $\mathfrak C_j$.

The first two non-Gaussian families are holomorphic and fourth
order in $w$, uniformly in $J$ and the independent coordinates
$\sigma_k$. The scalar subtraction includes the zero-source
family $\mathbf d_i=D(0)|_1-\widetilde e_i$ at each pointer
and the zero-source action part $\mathbf L_0$:
$F_0^{(i)}=\widetilde E_0(1,1)+\mathbf d_i-\mathbf L_0$.
Here $\widetilde E_0$ denotes degree zero less its trivial-pair
value. In particular no new action coefficient is extracted
on the shrinking radius.
Write the next potential as
$\Gamma_i+zV_i^{\rm rest}+\xi V_i^{\rm h}$,
$\Gamma_i=\mathcal L_i+q_i$.
Its logarithm is holomorphic for
$|z|\le c_{\rm cl}/(4\|V_i^{\rm rest}\|)$ and
$|\xi|\le c_{\rm cl}/(4\|V_i^{\rm h}\|)$.
Schwarz's lemma gives changes bounded by
$K''\|V_i^{\rm rest}\|$ and $K''\|V_i^{\rm h}\|$.
This estimate holds at every pointer. The terms from the tests
and the quadratic origin do not depend on $z,\xi$ and cancel
in these differences; the insertion terms have the uniform
bound $2K'_Mc_{\rm ins}c_{\rm cl}e^{-\kappa_+d_i(X)}$.
At $z=\xi=0$, expand in $\eta\mathcal L_i+q_i$.
The terms of order at least four are fourth order in $w$ and
have norm at most $N_{\rm h}/3$ by Cauchy on $|\eta|=2$.
The terms of orders zero and two differ from $F_i^{\rm G}$
only by the tests, the origin sort and $|x_{i-1}/\mathsf t_{i-1}-1|$.
Lem.~\ref{C-408-34} bounds the change of origin sort,
including the pointing insertions. Its all-fluctuation
hypothesis holds for the quadratic and linear families above,
and there are at most four origins including an insertion.
Removing the tests gives the same marked large-field bound
$CK_Me^{-cb_*^2}$.
Finally Schwarz's lemma in the coupling disc gives
$C K_M|x_{i-1}/\mathsf t_{i-1}-1|$. Both errors are
$O(\phi(x_{i-1}))$.

More explicitly, with $\bar v$ the supremum of the pointing
insertion norm, the four changes at each pointer cost
$K_{\mathbf M}^{\rm v}p$ for the tests and the quadratic-origin
pointing term, $2K_P^\Pi p_\star\bar v/c_{\rm cl}$ for the
sort, $K_P^\Pi p_\star\bar v/c_{\rm cl}$ for removal of tests,
and $2K_\Pi\bar v\bar B/(\sigma_\star\mathsf t_{i-1})$
for the coupling. The value at the trivial pair adds a factor
two; taking the orders zero and two on $|\eta|=2$ adds a factor
$5/4$. Since $\bar v\le9t_i^{\rm lin}/4<c_{\rm cl}/3$,
$K_P^\Pi p_\star\le K'_Mp$, and
$p\le\phi(\mathsf t_{i-1})^4$, these terms give
$C_1^{\rm lev}\phi(x_{i-1})$ in the pointed family norm.
The constant $C_1^{\rm lev}$ does not depend on $C_{\rm rad}$.

Consequently, if $\varepsilon_i^{\rm r}=\|F_{i,\rm r}\|$,
\[
 \varepsilon_i^{\rm r}\le C_0^{{\rm lev}\prime}\phi(x_{i-1})
       +\frac{13}{12}K''K_\Lambda^{\rm s}
          \frac{\mathsf C_{\mathcal A}}{C_{\rm rad}}
                              \max_{i'<i}\varepsilon_{i'}^{\rm r},
 \qquad
 C_0^{{\rm lev}\prime}=K''C_0^{\rm lev}+C_1^{\rm lev}
                      +C_w+\theta_sh_F^2(\bar B_M+2B_\sigma).
\]
The floor $C_{\rm rad}\ge512c_{\rm ins}(1+c_{\rm cl}^{-1})
\mathsf C_{\mathcal A}K_\Lambda^{\rm s}K_M$ makes the last
coefficient at most $(13/12)(16/512)<1/24$. It is imposed
before $C_0^{\rm lev}$ and $C_0^{{\rm lev}\prime}$ are formed.
Since the earlier $x_{i'-1}$ are larger and $\phi$ decreases,
$C_{\rm r}^{\rm lev}=2C_0^{{\rm lev}\prime}$ closes the
induction. The same coefficient bounds $F_{i,\rm h}$ by
$N_{\rm h}/24$; together with $N_{\rm h}/3$ this closes the
higher-order bound. Moreover
$F_{i,\rm h}=\int_0^1\partial_\xi\widetilde E_0(1,\xi)\,d\xi$
is fourth order in $w$, because the derivative is linear in
$V_i^{\rm h}$. With $C_{\rm rest}=C_0^{\rm lev}+C_{\rm r}^{\rm lev}$,
the ceiling $8C_{\rm rest}\phi(\gamma^{-2}/2)\le c_{\rm cl}$
makes every $z$ disc have radius at least two.

The additional assertions in the induction follow
alongside these bounds. By Prop.~\ref{C-407-16},
the explicit members of the step and its fermion
covariances are holomorphic in the independent
coordinates. A walk reads $\mathsf n$ on the same set
on which it reads the configuration. Every expanded
leg is assigned to a polymer containing that walk;
the mass-coordinate terms read the bare sites of
their assigned polymers. Thus a potential member
at $Y$ reads the coordinates in $Y$ alone.

For a localized earlier scalar family, the induction
gives a holomorphic family of bare arrays with zero
counterterms. The parameter version of scalar
localization preserves holomorphy at a fixed pair
and fluctuation. Its restriction to $Y$ uses only
the entries of that array on polymers contained in
$Y$. Two coordinate vectors agreeing on $Y$ give
the same restricted array, and hence the same
localized member there. Uniform bounds on the
field polydisc give joint holomorphy in the
fluctuation and coordinates. Taking the first two
Taylor terms in the fluctuation preserves these
properties.

The cluster expansion at $X$ reads potential
members with $Y^\sharp\subset X$ and the fermion
covariance decoupled to cubes in $X$. Those
covariances read $\mathsf n$ in $X$ alone. Hence
the expansion is holomorphic in
$(\mathfrak c,z,\xi)$ and in $(\mathfrak c,\eta)$
on the discs used above, and reads the coordinates
in $X$ alone. This also holds at every pointer:
the test and quadratic-origin terms read only
$w$ locally, and the bond and fermion insertions
are derivatives of the same local holomorphic
families. The estimates for an inserted origin
justify their parameter integrals.

Differences in $z,\xi$, Cauchy integrals on
$|\eta|=2$, and evaluation at the trivial pair
preserve coordinate holomorphy and locality.
They therefore preserve these properties in the
definitions of all four families. The gauge
identities hold before these linear operations;
their scalar parts are invariant and functions
of the pair. This closes the additional clauses
of the induction. At step $j$ the same argument
applies to the new activity, since its potential
uses only the previously constructed families.

At the trivial pair $\mathcal L_j=0$: a linear gauge functional
there would be an invariant adjoint colour vector. The scalar
constant, together with the logarithmic mass of the fixed
relative-weight origin, gives $e_j^{\rm G}$; the free fermion normalizer
gives $e_j^{\rm F}$, the fourth-order families give $e_j^{\rm H}$,
and the remainder gives $e_j^{\rm I}$. The Gaussian sums telescope
by exact Gaussian integration. In this identity the linear term is
$\mathcal L_i^{\rm G}$, the first background-gauge jet of
$\mathcal L_i$, and the quadratic term is $Q_i^{\rm G}=q_i(1)$;
the earlier tree terms are generated by $\mathcal L_i^{\rm G}$.
The pair-dependent parts have already been assigned to the small
remainder. The fermion logarithmic germs telescope by
\cref{C-397-12}. This proves both sum identities.

\emph{The tail.}
Let $s,s'\in\mathcal O_j^{\rm lev}$,
$S=\operatorname{supp}(s-s')$, and
$D=\operatorname{dist}(X,S)$. If $X$ meets $S$,
the two bounds for $e_j^{\rm I}$ suffice, after
doubling their constant. Otherwise $w=w'$ and
$J=J'$ on $X$. The source-chain bound gives,
at every bare site of $X$ and for $k<j$,
\[
 |\sigma_k(s)-\sigma_k(s')|
 \le C_{\rm t}\phi(x_k)e^{-\kappa_\sigma D/\lambda_j},
\]
since $\lambda_k\le\lambda_j$ and both sources
belong to $\mathscr O_k$.

We first record what this implies for a local
holomorphic function $f$ of the independent
coordinates, bounded on $\mathfrak C_j$ by
$M_0e^{-\kappa d_j(X)}$. Set
\[
 \delta_k=(\sigma_k(s)-\sigma_k(s'))1_X,
 \qquad
 \theta=\frac{C_{\rm t}}{B_\sigma}
                   e^{-\kappa_\sigma D/\lambda_j}.
\]
Keep $w',-a_KJ'$ fixed and replace every
$\sigma_k(s')$ by $\sigma_k(s')+\zeta\delta_k$.
This affine curve lies in $\mathfrak C_j$ for
$|\zeta|<1/\theta$, because
\[
 \|\sigma_k(s')+\zeta\delta_k\|
 \le B_\sigma\phi(x_k)(1+|\zeta|\theta)
             <2B_\sigma\phi(x_k).
\]
When $\theta<1$, its value at $\zeta=1$ agrees
with the chain of $s$ on $X$. Locality identifies
the values of $f$ there. Schwarz's lemma applied
to the difference from $\zeta=0$ gives
$2\theta M_0e^{-\kappa d_j(X)}$. When
$\theta\ge1$, the bounds at the two chain
points give the same inequality. If $\theta=0$
the coordinate vectors agree on $X$ and the
difference vanishes. Consequently
\[
 |f(\mathfrak c_{j-1}(s))-f(\mathfrak c_{j-1}(s'))|
 \le C_{\rm tail}M_0e^{-\kappa d_j(X)}
                         e^{-\kappa_\sigma D/\lambda_j},
 \qquad C_{\rm tail}=2\max\{1,C_{\rm t}/B_\sigma\}.
\]

Apply this estimate at each $z$ on the disc
\[
 |z|\le R=\frac{c_{\rm cl}}
                   {4C_{\rm rest}\phi(x_{j-1})}\ge2
\]
to $\mathcal E(z;\mathfrak c)$, the degree-zero
part at the trivial pair of the step-$j$ activity
$\widetilde E(z,0)(X)$, before its scalar value
is subtracted. The preceding coordinate induction
applies on all of $\mathfrak C_j$, not only at
the coordinate vectors of physical sources.
The potential is within $c_{\rm cl}$ of the
quadratic origin, so the cluster bound is
\[
 |\mathcal E(z;\mathfrak c)|
       \le K_M(c_{\rm cl}+1)e^{-\kappa d_j(X)}.
\]
It follows that
$\Delta(z)=\mathcal E(z;\mathfrak c_{j-1}(s))
-\mathcal E(z;\mathfrak c_{j-1}(s'))$ satisfies
\[
 |\Delta(z)|\le
 C_{\rm tail}K_M(c_{\rm cl}+1)e^{-\kappa d_j(X)}
                              e^{-\kappa_\sigma D/\lambda_j}.
\]
At $z=0$ the activity has only the Gaussian
scalar potential and quadratic origin. These
read $w$ in $X$ alone, so $\Delta(0)=0$.

The definition of the decomposition gives
\[
 e_j^{\rm I}(X;s)-e_j^{\rm I}(X;s')=\Delta(1).
\]
Indeed, the gauge normalizer, weighted Gaussian
part and $e_j^{\rm G}$ read $w$ in $X$ alone
and have equal values at the two sources.
The free fermion normalizer cancels with
$e_j^{\rm F}$. The change obtained by turning
on $\xi$ is precisely $e_j^{\rm H}$ and is
removed in $e_j^{\rm I}$. A second application
of Schwarz's lemma, now in $z$, therefore gives
\[
\begin{split}
 |e_j^{\rm I}(X;s)-e_j^{\rm I}(X;s')|
 &\le R^{-1}\sup_{|z|\le R}|\Delta(z)|\\
 &\le4C_{\rm rest}C_{\rm tail}K_M(1+c_{\rm cl}^{-1})
       \phi(x_{j-1})e^{-\kappa d_j(X)}
                         e^{-\kappa_\sigma D/\lambda_j}.
\end{split}
\]
Take $C_I$ to be the larger of this constant
and twice the preceding bound constant for
$e_j^{\rm I}$. Since $\phi(x_{j-1})\le\phi(x_j)$,
this proves the asserted tail. The constants
$C_{\rm tail}$ and $C_I$ are fixed before $m_-$.
\end{proof}

\subsection{The free templates and the sum over scales}

\begin{lemma}[The templates]\label{C-408-25}
For a test pair with $8a_K\le d$, $20d\le\ell$ and
$10m_+d\le1$,
the Gaussian sum of mixed derivatives is
$\frac9{16\pi^2}\mathsf b^{-2}\mathcal K_6$ with error at most
\[
 C\mathsf b^{-2}\left[\frac{a_K}{\vartheta d}
                   +(d/\ell)^3+d^3\right]\|f\|_\infty\|h\|_\infty.
\]
The free fermion sum is
$\delta_{\Gamma\Gamma'}(2\pi^2)^{-1}\mathcal K_4$ with error at most
\[
 C\left[\left(\frac{a_K}{\vartheta d}+(m_+d)^2
                       +(d/\ell)^2\right)d^2
       +d^6m_-^{-2}+d^4m_-^{-1}\right]
                          \|f\|_\infty\|h\|_\infty.
\]
\end{lemma}

\begin{proof}
The flat constrained-curvature identity
\cite[Prop.~120.3]{YMM1} identifies the free covariance
per colour with the closed-curvature projection. Its canonical
step hypotheses hold by \cref{C-396-12}. Away from coincidence its continuum kernel is the
dipole matrix
$\mathbb K_3(x)=(3xx^{\mathsf T}-|x|^2I)/(4\pi|x|^5)$.
The squared entries sum to $3/(8\pi^2|x|^6)$.
The free case of \cite[Prop.~121.2]{YMM1}, including
its lattice-kernel comparison and Riemann-sum estimate, gives error $C[a_K/(\vartheta d)+(d/\ell)^3]$ on the pair.
Wick's rule multiplies this square by $\dim\mathfrak{su}(2)/2=3/2$,
giving $9/(16\pi^2)$.
For the scale-one remainder use instead the all-distance bound
\cite[Prop.~120.5]{YMM1}:
\[
 |C^{\rm bg}_{K,s}(p,q)|\le C(a_K/a_s)^3.
\]
At $a_s=1$ it is $Ca_K^3$. Thus the square of this kernel
and its cross terms with the free kernel are bounded, after
the mesh factors, by $C(1+|x-y|^{-3})$. Integrating over
the test pair gives a contribution at most
$C\mathsf b^{-2}d^3\|f\|_\infty\|h\|_\infty$.

For the fermions put $G_0=(D[1]+m_0)^{-1}$.
The free symbol's light Schur complement is
$m_0+i\gamma\cdot p+O(|p|^2)$, with the differentiated bounds
proved above. A smooth dyadic partition and integration by parts
give, for $n\ne0$,
\[
 \iota^*G_0\iota(n)=\mathscr G_{m_0}(n)+O(|n|^{-3}),\qquad
 |\Pi_1G_0\iota(n)|+|\iota^*G_0\Pi_1(n)|\le C|n|^{-3},
 \quad |\Pi_1G_0\Pi_1(n)|\le C|n|^{-4},
\]
where
\[
 \mathscr G_m(n)=\frac{\gamma(n)(1+m|n|)e^{-m|n|}}{4\pi|n|^3}
                          +\frac{m e^{-m|n|}}{4\pi|n|}.
\]
Indeed the light error and mixed symbols have degree zero, and
the nonanalytic heavy term has degree one; their dyadic Fourier
sums have respectively the displayed third and fourth powers.
The analytic heavy inverse has exponential decay. These bounds
are uniform in $0\le m_0\le1/2$.

The light propagator between individual chambers has factor
$48^{-1}$, which cancels the $48^2$ chamber pairs in the loop.
For the reflection-real basis,
\[
 -\operatorname{tr}[\Gamma\gamma(u)\Gamma'\gamma(-u)]
                         =8\delta_{\Gamma\Gamma'}|u|^2.
\]
This follows from $\operatorname{tr}(t_jt_{j'})=\delta_{jj'}$,
$\tau_3\tau_a\tau_3=\epsilon_a\tau_a$ and $c_a^2\epsilon_a=1$.
The scalar-vector cross term has zero spin trace. The remaining
mass correction is $O((m_+d)^2)$, and a heavy factor costs
$O(a_K/d)$. Thus
$-a_K^2\operatorname{Tr}(\hat f\Gamma G_0\hat h\Gamma'G_0)$
has leading coefficient $8/(16\pi^2)=1/(2\pi^2)$.
The Riemann-sum error is $C a_K/(\vartheta d)$ relative to
$d^2\|f\|_\infty\|h\|_\infty$.
The antiperiodic image formula and positive $m_-$ give
the additional $C(d/\ell)^2$ term.

Finally $\|\mathbb D_K(1;0)^{-1}\|\le C/m_-$.
Factor the first source derivative into its two background rows.
Their squared row sums are $O(a_K^2)$, so its trace norm is
at most $Ca_K^3\sum_x|f(x)|\le Cd^3\|f\|_\infty$.
For the second derivative use the resolvent between the separated
supports; its $C a_K^2|x-y|^{-2}$ bound gives trace norm at most
$Cd^4\|f\|_\infty\|h\|_\infty$.
The two terms in the second logarithmic derivative are therefore
bounded by $C(d^6m_-^{-2}+d^4m_-^{-1})$.
This uses trace norms of the localized factors, with no global
volume or rank factor. It is the last-level subtraction in the
free fermion sum.
\end{proof}

\begin{lemma}[The interaction error]\label{C-408-26}
For the packing in \cref{C-408-20}, the mean of the mixed derivatives
of $\sum_{j,X}e_j^{\rm I}$ is bounded, at bilinear indices, by
\[
 C R_J^{-2}\phi(\mathsf b/d)d^2\|f\|_\infty\|h\|_\infty,
\]
and at action-density indices by
\[
 C\left(\frac{40M}{\vartheta\varrho}\right)^2
 \phi(\mathsf b/d)\widehat\phi(\mathsf b/d)^2
                       \|f\|_\infty\|h\|_\infty.
\]
\end{lemma}

\begin{proof}
On \eqref{eq:qcd-level-source-domain} the two source radii are
\[
 r_{\Gamma,j}=\frac{R_J}{\lambda_j\|f\|_\infty},\qquad
 r_{0,j}=\frac{\varrho}{\widehat\phi(x_j)
                  \max(\|f\|_\infty,40M\lambda_j\operatorname{Lip}f)}.
\]
Cauchy's estimate gives
$C\phi(x_j)e^{-\kappa d_j(X)}r_{\alpha,j}^{-1}r_{\alpha',j}^{-1}$.
If $\lambda_j\le d$, a mixed derivative must connect the two
supports. The tail in \cref{C-408-21} gives a factor
$e^{-c d/\lambda_j}$, and the number of nearby cells is
$C(d/\lambda_j)^3$.
If $\lambda_j>d$, apply the same independent signs to the packed
sources. The packing meets at most
$C N\mathsf N(d/\lambda_j)^3$ cells; after taking the mean its
contribution is again bounded by $C(d/\lambda_j)^3$ times the
single-cell Cauchy bound. This volume factor is essential.

At bilinear indices sum
\[
 C\sum_j\phi(x_j)\lambda_j^2(d/\lambda_j)^3
 \begin{cases}e^{-cd/\lambda_j},&\lambda_j\le d,\\1,&\lambda_j>d.
 \end{cases}
\]
Since $x_j\asymp\mathsf b/\lambda_j$, both geometric sums are
led by $\lambda_j\asymp d$ and are bounded by
$C\phi(\mathsf b/d)d^2$.
For the action density replace $\lambda_j^2/R_J^2$ by
\[
 \varrho^{-2}\widehat\phi(x_j)^2
                    \max(1,40M\lambda_j/(\vartheta d))^2.
\]
Below $d$ the exponential cuts the sum. Above $d$ its terms
are bounded by $C\phi(x_j)\widehat\phi(x_j)^2d/\lambda_j$,
again led by $\lambda_j\asymp d$. The polynomial logarithms
are paid by the geometric ratios. This proves both bounds
uniformly in the number of levels.
\end{proof}

\begin{proof}[Proof of \cref{C-408-27}]
The fourth-order term in \cref{C-408-21} has zero mixed derivative
at both kinds of indices. The Gaussian part does not read $J$,
and the free fermion part does not read $w$. Combine
\cref{C-408-20,C-408-25,C-408-26} and divide by
\eqref{eq:qcd-template-lower}. On the lattice the additional
relative cutoff error is $C a_K/(\vartheta^3d)$.
On each volume sequence $20d\le\ell$ eventually.
First take the volume limit, then the cutoff subsequence at each
fixed test pair. These cutoff and image errors vanish.
The remaining bounds tend to zero with $d$, since
$\phi(\mathsf b/d)\widehat\phi(\mathsf b/d)^2\to0$;
the packing remainder is $O(d)$ after division by either
template. Taking $d$ small enough gives a positive lower bound
at every equal pair of indices. This proves nonvanishing.
\end{proof}

\section{Limits taken volume first}\label{sec:qcd-continuum}

At cutoff $K$ choose a limit $\nu_K$ of the torus gauge
marginals as their physical sides tend to infinity through
$\ell=2L^m$ satisfying the room condition above. Such limits
exist by \cref{C-389-10}. Their local expectations are
$\omega_{\nu_K}(F)=\int\langle F\rangle_U\,\nu_K(dU)$.
Let $S^T_{n,K,\nu_K}$ denote the cumulants of the physical
observables defined in the introduction. Bounds that are uniform
in the torus side pass to these expectations, and hence to
their cumulants, by \cref{C-389-10}.

\begin{proposition}[Compactness on separated supports]\label{C-389-25}
For any sequence $K_j\to\infty$ and any choices $\nu_{K_j}$,
there is a common subsequence on which all the cumulants of
order at least two converge on every separated tuple of
$C_c^1(\mathbb R^3)$ functions. The limits $S_n$ satisfy
the bound of \cref{C-407-6}. At fixed indices they are
distributions of locally finite order off the coincidence
diagonals. They obey the permutation, translation, cubic,
flavour and conjugation identities of the lattice model.
\end{proposition}
\begin{proof}
Fix the indices, number of entries and pairwise separated compact
support sets, each in a ball of radius $R$. The bound of
\cref{C-407-6} is uniform in $K$, the volume and the positions
of those balls. It gives equicontinuity in the product of the
insertion norms. Choose the support sets to be finite unions
of closed rational cubes and choose a countable dense set of
$C^1$ functions in each of them. These spaces are separable
both for the $C^1$ norm used by the action density and for the
supremum norm used by a bilinear. A diagonal subsequence makes
all the resulting cumulants converge. The uniform multilinear
bound extends convergence to every separated tuple: its supports
fit inside pairwise separated sets of the indicated countable
kind. The same bound passes to the limit. The kernel theorem
on products of separated open sets gives distributions there,
with agreement on overlaps by uniqueness.

The exact finite-lattice identities pass to the limit.
For a real translation $v$, choose $v_j\in a_{K_j}\mathbb Z^3$
converging to $v$. Translated test functions converge in their
insertion norms; for the $C^1$ norm this uses uniform continuity
of their first derivatives. Their supports remain in fixed
separated neighborhoods. Equicontinuity and exact lattice
translation invariance give invariance under $v$.
\end{proof}

Let $\sigma$ be the centered functional of
\eqref{eq:qcd-centered-moments}. It is the limit of the
corresponding centered lattice products.
It is useful to keep centering in this formula: no convergence
of the uncentered one-point functions is needed.

\begin{proposition}[Positivity at coordinate and diagonal planes]
\label{C-389-27}
The functional $\sigma$ is reflection positive at every coordinate
plane and every plane with normal $e_\mu\pm e_\nu$.
Reflection acts geometrically on the reflection-real indices.
\end{proposition}
\begin{proof}
The plane $P_0=\{x_1=x_2\}$ is a mirror of every lattice
$a_K\mathbb Z^3$. For a finite combination $F=\sum_i z_iX_i$
whose supports lie strictly on one side, the centered lattice
products $F_K$ lie in its even half algebra for all large $K$.
Centering commutes with reflection, since
$\omega_{\nu_K}(\Theta_{P_0}G)=
\overline{\omega_{\nu_K}(G)}$. Therefore
\[
 0\le\omega_{\nu_K}(F_K\Theta_{P_0}F_K)
   =\sum_{i,l}\bar z_i z_l\,
       \sigma_K(\theta_{P_0}X_i\cdot X_l).
\]
Here the lattice positivity is \cref{C-386-13}; each
concatenation is separated across the plane. The partition
formula and convergence of the cumulants pass this inequality
to $\sigma$. Translations and cubic covariance give all the
diagonal planes stated. The same argument at $x_3=0$ gives
the coordinate planes.
\end{proof}

\section{Rotations and reconstruction for the centered fields}

Let $\mathcal I=\{0\}\cup\mathfrak B$ be the finite
reflection-real index set, with signs $\pi_\alpha\in\{1,-1\}$.
For a tuple $X=((\alpha_i,f_i))_{i=1}^n$ of real
$C_c^1(\mathbb R^3)$ functions with pairwise disjoint supports,
write $\pi(X)=\prod_i\pi_{\alpha_i}$.
The geometric action is
$rX=((\alpha_i,f_i\circ r^{-1}))_i$.
Let $\sigma$ be normalized, multilinear and permutation
symmetric on these tuples, with
\[
 \sigma(\varnothing)=1,\qquad
 \overline{\sigma(X)}=\pi(X)\sigma(X).
\]
It extends complex linearly to the entries and to finite
formal combinations of tuples wherever the indicated
concatenation is separated.  If $F=\sum_i\zeta_iX_i$ and
$G=\sum_j\xi_jY_j$ lie in the positive side of a plane,
\[
 \sigma(\theta\bar F\cdot G)=
       \sum_{i,j}\bar\zeta_i\xi_j
                      \sigma(\theta X_i\cdot Y_j).
\]
Reflection acts geometrically and introduces no index sign
in this basis.  The sign in improper Euclidean covariance
is a different operation.

For the centered cumulants of the preceding sections,
$\sigma$ is given by \eqref{eq:qcd-centered-moments}.
Thus every singleton has value zero.  This is a functional
of even observables; it is not assumed to be the moment
functional of a positive joint probability law for the
bilinears.

\subsection{The analytic hypotheses}

Set $\|f\|_\sharp=\max\{\|f\|_\infty,
40M\|\nabla f\|_\infty\}$,
$\|f\|_0=\|f\|_\sharp$ and
$\|f\|_\Gamma=\|f\|_\infty$ for $\Gamma\in\mathfrak B$.
We use the following assumptions on $\sigma$.
\begin{enumerate}
\item[(R0)] $\sigma((\alpha,f))=0$ for every single entry.
\item[(R1)] It is translation invariant and satisfies
$\sigma(wX)=\pi(X)^{\delta(w)}\sigma(X)$ for $w\in W_3$,
where $\delta(w)=1$ for an improper map and zero otherwise.
\item[(R2)] For each $n\ge2$, $d,R>0$ there is
$C_n(d;R)$ such that
$|\sigma(X)|\le C_n(d;R)\prod_i\|f_i\|_{\alpha_i}$
whenever the supports are pairwise at distance at least
$d$ and each lies in a ball of radius $R$, independently
of the positions of those balls.
\item[(R3)] $\sigma(\theta\bar F\cdot F)\ge0$ for
$\theta(x_1,x_2,x_3)=(x_1,x_2,-x_3)$ and every finite
combination supported in $x_3>0$.
\item[(R4)] In the reflection Hilbert space at this plane,
the joint energy and first tangential momentum satisfy
$H\ge|P_1|$.
\item[(R5)] There are $0<c_{\mathrm I}\le1/16$,
$C_{\mathrm I}<\infty$ and $p_{\mathrm I}\ge0$ such that,
for $0<\eta\le1$, $|y_3|\le\eta/8$, real
$h\in C_c^1(B(y,c_{\mathrm I}\eta))$ and combinations
$F,G$ supported in $x_3>\eta/4$,
\[
 |\sigma(\theta\bar F\cdot(\alpha,h)\cdot G)|
 \le C_{\mathrm I}\eta^{-p_{\mathrm I}}\|h\|_\sharp
 \|[\tau_{-\eta e_3/4}F]\|
 \|[\tau_{-\eta e_3/4}G]\|.
\]
\item[(R6)] For each fixed $\alpha$, its two-point form
is reflection positive also at
$\Sigma_D=\{x_1+x_3=0\}$: for real $f_i$ supported on
one open side,
\[
 \sum_{i,j}\bar\zeta_i\zeta_j
 \sigma((\alpha,f_i\circ\theta_D),(\alpha,f_j))\ge0.
\]
\item[(R7)] In coordinates $y=(x_3,x_1)$, a label is
$\lambda=(\alpha,\upsilon)$, $\upsilon\in C_c^1(\mathbb R)$
real.  Let $\mathsf{Prof}_t$ consist of $C_c^1$ profiles
supported in $[t,2t]\times[-t,t]$, with
$|a|\le3$ and $\operatorname{Lip}a\le12/t$.
There is $a_\Gamma<4$ such that, for every label,
\[
 0\le\sigma((\alpha,(a\otimes\upsilon)\circ\theta),
                         (\alpha,a\otimes\upsilon))
       \le C_\lambda t^{-a_\Gamma},
 \qquad a\in\mathsf{Prof}_t,\quad0<t\le1.
\]
\end{enumerate}
The Hilbert space in (R4)--(R5) is the completion of the
positive-time combinations modulo the null space of
$\sigma(\theta\bar F\cdot G)$, and $\Omega=[\varnothing]$.
Assumptions (R1)--(R3) make its translations
$T_t[F]=[\tau_{te_3}F]$ and
$\mathcal U(v)[F]=[\tau_vF]$, $v\perp e_3$, respectively
a strongly continuous self-adjoint contraction semigroup
$e^{-tH}$ and a commuting unitary representation.
Indeed reflection and translation make $T_t$ symmetric,
and repeated Cauchy--Schwarz gives
\[
 \|T_tx\|\le\|x\|^{1-2^{-j}}
                    \|T_{2^jt}x\|^{2^{-j}}.
\]
The position-independent bound in (R2) bounds the limsup
of the second factor by one. Null vectors are preserved by the
same inequality, and moving-entry continuity follows
from (R2).  A countable dense set of $C^1$ tuples proves
separability.  In particular $e^{-tH}$ is injective.

\begin{proposition}[The two-plane spectral inequality]\label{C-389-28}
Suppose the translation-invariant separated functional has the
position bounds and reflection positivity at $e_3^\perp$ and
at $(e_3+\epsilon e_k)^\perp$, $k=1,2$, $\epsilon=\pm1$.
Then its joint spectral measure in the $e_3$ picture is
supported on $E\ge|q_k|$.
\end{proposition}
\begin{proof}
Put $e_\pm=(e_3\pm\epsilon e_k)/\sqrt2$ and
$W=\{x_3>|x_k|\}$. For a combination $F$ supported in $W$,
let $\mu_F$ be its joint spectral measure carried to the
right half-plane by $w=(E-i\epsilon q_k)/\sqrt2$. Then
\[
 g(s,s')=\sigma(\theta\bar F\cdot
                  \tau_{se_++s'e_-}F)
        =\int e^{-sw-s'\bar w}\,d\mu_F(w),\qquad s,s'\ge0.
\]
Let $\theta_+$ be reflection in $e_+^\perp$ and put
$F_1=\theta_+\theta F$. Both $F_1$ and $F$ lie on the
positive side of this plane. In its reflection picture,
\[
 g(s,s')=\langle[F_1]_+,T_s^+\mathcal U^+(s'e_-)[F]_+\rangle_+.
\]
For $\varphi\in L^1(0,\infty)$ and its Laplace transform
$q(z)=\int_0^\infty e^{-zs}\varphi(s)\,ds$, Fubini and
the spectral theorem in these two pictures give
\[
 \int|q(w)|^2\,d\mu_F(w)
 \le\|[F_1]_+\|\,\|[F]_+\|
        \sup_{u\ge0}|q(u)|\sup_{y\in\mathbb R}|q(iy)|.
\]
Take $\varphi(s)=s^{a-1}e^{-\eta s}/\Gamma(a)$, where
$a=\varepsilon+i\kappa$, $\varepsilon,\eta>0$.
Thus $q(z)=(z+\eta)^{-a}$. Its two suprema are bounded by
$\eta^{-\varepsilon}$ and
$\eta^{-\varepsilon}e^{|\kappa|\pi/2}$. Fatou, first as
$\varepsilon\downarrow0$ and then as $\eta\downarrow0$, gives
\[
 \int_{w\ne0}e^{2\kappa\arg w}\,d\mu_F(w)
       \le\|[F_1]_+\|\,\|[F]_+\|e^{|\kappa|\pi/2}.
\]
Letting $\kappa$ tend to $+\infty$ and $-\infty$ excludes
every closed sector with $|\arg w|>\pi/4$. This says
$E\ge|q_k|$ on the measure of $[F]$.
Every positive-time tuple can be translated into $W$ by
$\tau_{re_3}$ for large $r$. The bad spectral projection
$Q=1_{\{E<|q_k|\}}(H,P)$ then satisfies
$e^{-rH}Q[X]=Q[\tau_{re_3}X]=0$. The semigroup is
injective, so $Q[X]=0$ on a dense set and hence $Q=0$.
\end{proof}

\begin{theorem}[Finite-family rotation criterion]\label{C-389-30}
A functional with (R0)--(R7) is covariant under every
Euclidean isometry:
\[
 \sigma(rX)=\pi(X)^{\delta(r)}\sigma(X).
\]
It is reflection positive at every plane and its joint
energy--momentum spectrum is contained in
$\{(E,q):E\ge|q|\}$.
\end{theorem}

\begin{proof}
The analytic point-plate and tube constructions used in
this proof are the general functional results of
\cite[Props. 97.1, 97.2 and the kernel
construction preceding 98.2]{YMM1}.  Their inputs
are precisely the reflection pictures, the cone
$H\ge|P_\nu|$, the local insertion bound (R5), and
continuity (R2).  They act linearly on each inserted
test function.  We give the finite-label argument,
including the step in which the sign rule matters.

\emph{Pictures and point plates.}
Transporting reflection positivity by $w\in W_3$
multiplies the coefficient of a tuple by the real sign
$\pi(X)^{\delta(w)}$ and does not change positivity.
Thus each coordinate plane has a reflection picture.
Axis permutations carry (R4) to every ordered pair of
distinct axes.  Spectral calculus consequently gives
the contraction semigroup
\[
 T(z)=\int e^{-z_3E+i(z_1q_1+z_2q_2)}\,d\mathscr E,
 \quad \operatorname{Re}z_3\ge
               |\operatorname{Im}z_1|+|\operatorname{Im}z_2|,
\]
holomorphic in its open tube.
For its planar restriction write
$d_{\mathrm{Tb}}(z)=\operatorname{Re}z_3-
|\operatorname{Im}z_1|$; for the three-dimensional
point construction write
$d_3(z)=\operatorname{Re}z_3-
|\operatorname{Im}z_1|-|\operatorname{Im}z_2|$.
The slab operator
$B_l(\alpha,h)[G]=[(\alpha,h)\cdot\tau_{le_3}G]$
is defined for $\operatorname{supp}h\subset\{0<x_3<l\}$
and is bounded by (R5), with
\[
 B_l(\alpha,h)^*=B_l(\alpha,\overline{h\circ\theta_{l/2}}).
\]
This follows by reflecting, translating by $-le_3$ and
using permutation symmetry in its matrix elements.
There is no index sign in this adjoint identity.
For a profile supported in $|x_3|<d$ put
$B_d^{\mathrm c}(\alpha,h)=
B_{2d}(\alpha,\tau_{de_3}h)$.

Let $\Upsilon_\lambda$ be the length of the interval
hull of $\operatorname{supp}\upsilon$ and put
$\wp=p_{\mathrm I}+4$.  The point-plate construction
gives a holomorphic insertion $B_\lambda(p,q)$ on
the planar tubes with
\[
 \|B_\lambda(p,q)\|\le
 C_{\mathrm{plt}}(1+\Upsilon_\lambda)
 \|\upsilon\|_{C^1}
 \min\{1,d_{\mathrm{Tb}}(p),d_{\mathrm{Tb}}(q)\}^{-\wp}.
\]
It is useful to recall its mechanism. Choose a nonnegative radial
$j\in C_c^1(\mathbb R^2)$ of integral one, supported in $|y|\le1/2$,
and put $j_\delta(y)=\delta^{-2}j(y/\delta)$.
Mollify a point plate by $j_\delta$ and put
$B^\delta(p,q)=T(p-\delta e_3)
B_\delta^{\mathrm c}(j_\delta\otimes\lambda)T(q-\delta e_3)$.
The convolution identity recovers every ordinary slab
profile.  A holomorphic mean-value formula on a disc
of radius proportional to the tube distance, followed
by (R5), bounds these operators independently of
$\delta$.  Cauchy's estimate and symmetric convolution
make them Cauchy in norm at positive real times;
the shift law extends their limit throughout the tube.
This proof uses only the displayed norm of the fixed
label, so it applies separately to each $\alpha$.

Products give kernels
$K_n(Z;\lambda_1,\ldots,\lambda_n)$, holomorphic on
the tubes for every time order and representing
$\sigma$ on time-ordered plates.  The quarter turn
in the $(x_3,x_1)$ plane is proper and belongs to
$W_3$, so the kernels patch in four frames.
For two ordered lists the inner product of their
vectors is the kernel with the first list reflected
and reversed.  This is the labelled Gram identity;
no relation between two distinct labels has been used.

\emph{Harmonics and their top coefficient.}
For a labelled distinct-point configuration $Z$,
$\zeta\mapsto K_n(R_\zeta Z)$ is entire, periodic
with period $\pi/2$, and of growth
$O(e^{\wp n|\operatorname{Im}\zeta|})$.
The four-frame continuation and the point-plate
bound give these statements as in the cited kernel
construction.  The Fourier series is therefore finite:
\[
 K_n(R_{is}Z;\underline\lambda)
 =\sum_{4|m|\le\wp n}c_m(Z;\underline\lambda)e^{-4ms}.
\]
Shifting the real part of the angle gives the exact
rotation phase of $c_m$, also for rotations about any
centre.  Locally keeping an admissible frame fixed
makes each coefficient real analytic in the distinct
points.  The sign rule gives
\[
 \overline{c_m(Z;\underline\lambda)}
                =\pi(\underline\lambda)c_{-m}(Z;\underline\lambda).
\]
Consequently nonzero extreme harmonics occur in both
directions, even for mixed scalar and pseudoscalar labels.

Let $\bar m(n)$ be the largest $|m|$ that is nonzero
for some $n$-point configuration and some labels;
put $\bar m(n)=0$ if all these harmonics vanish.
In particular $\bar m(0)=\bar m(1)=0$.
Apply Cauchy--Schwarz to the labelled Gram identity
at $R_{is}$ and let $s\to+\infty$.  Comparing leading
exponents gives
\[
 2\bar m(a+b)\le\bar m(2a)+\bar m(2b).
\]
For $2a$ points multiply each finite Gram matrix by
$e^{-4\bar m(2a)s}$ and take its limit.  The result,
$t_a(\theta S\cup T)$, is positive Hermitian.
The rotation phase transports its Schwarz inequality
to any line $L$:
\[
 |t_a(\theta_LS\cup T)|^2\le
 |t_a(\theta_LS\cup S)|\,
 |t_a(\theta_LT\cup T)|.
\]
If some times coincide, perturb them by distinct
small amounts and use continuity.  This preserves
the same top coefficient.

\emph{Polygons and halving.}
Let $\mathcal N_a(S)$ mean that $t_a$ vanishes on
every labelled $2a$-point configuration containing
$S$.  If $S\subset\theta_LS'\cup S'$, $S'$ lies
on one side of $L$ and $\#S'\le a$, then
$\mathcal N_a(S)$ implies both
$\mathcal N_a(S')$ and $\mathcal N_a(\theta_LS')$.
For a configuration $X\supset S'$ on that side the
diagonal value $t_a(\theta_LX\cup X)$ vanishes.
The last Schwarz inequality makes every corresponding
mixed value vanish.  These mixed configurations form
a nonempty open set in the free point positions,
with arbitrary fixed labels.  Their distinct-point
domain is connected, since its removed affine
coincidence sets have codimension two.  The
real-analytic identity theorem proves the assertion
for all positions.

If $a$ did not divide $2\bar m(2a)$, a regular
$2a$-gon with one fixed label at all vertices would
have $t_a=0$: rotation by $\pi/a$ fixes it but
multiplies its top harmonic by a nontrivial phase.
Reflect an arc across the perpendicular bisector of
its next side.  Its first $\lceil j/2\rceil$ vertices
and their reflection contain an arc of $j$ vertices.
Repeated halving from the full polygon to a singleton
therefore gives $\mathcal N_a$ for that singleton.
Every point is a vertex of such a polygon and the
label was arbitrary.  This would make $t_a$ identically
zero, a contradiction.  Thus $a\mid2\bar m(2a)$.

For completeness put $\nu(a)=\bar m(2a)$.
The preceding convexity makes its increments an
increasing bounded integer sequence, since
$0\le\nu(a)\le\wp a/2$.  They are eventually equal
to an integer $d$, and $\nu(a)=da-c$ eventually,
where $c\ge0$ is the sum of their deficits from $d$.
The divisibility forces every sufficiently large
$a$ to divide $2c$, hence $c=0$.  All increments
equal $\nu(1)$, so
\[
 \bar m(2a)=a\bar m(2),\qquad
                  \bar m(n)\le\tfrac n2\bar m(2).
\]

\emph{Equal-index and mixed two-point functions.}
For one index $\alpha$ its two-point form is real,
since $\pi_\alpha^2=1$, and invariant under cubic
isometries.  Its one-field reflection space maps
isometrically into the full space by
$[f]\mapsto[(\alpha,f)]$, intertwining translations
and energy.  The two-point continuation from the
same general construction has spectral measure
$\rho_\lambda$ on the cone and kernel
\[
 k_{\lambda,r}(\zeta)=\int
 e^{-r(E\cos\zeta-iq\sin\zeta)}\,d\rho_\lambda(E,q).
\]
Assumption (R7) gives its growth at most
$e^{(4+a_\Gamma)|\operatorname{Im}\zeta|}$.
Since $4+a_\Gamma<8$ and its period is $\pi/2$,
$k_{\lambda,r}=c_0(r)+2c_1(r)\cos4\zeta$.
At $\zeta=is$ its positive cone measure gives
$k_{\lambda,r}(is)\ge0$; divide by $\cosh4s$ and
let $s\to\infty$ to obtain $c_1\ge0$.
For the opposite sign let $\rho_D$ be the positive
spectral measure of the two-point form turned by
$\pi/4$, supplied by (R6).  Its cone support is
not assumed.  For the boosts
$b_s(E,q)=(E\cosh s+q\sinh s,q\cosh s+E\sinh s)$,
the three-harmonic identity and uniqueness of the
Laplace--Fourier representation give
\[
 (\cosh4-1)\rho_D+(b_1)_*\rho_\lambda+
       (b_{-1})_*\rho_\lambda
                       =(\cosh4+1)\rho_\lambda.
\]
This is the positive-measure calculation in
\cite[Prop. 96.4]{YMM1}; the representation
uniqueness used there is Proposition 96.1.
In particular $\rho_D$ is bounded by a constant
times $\rho_\lambda$ and is supported on the cone.
Its continued kernel is
$c_0(r)-2c_1(r)\cosh4s\ge0$; the same limit gives
$c_1\le0$.  Thus the nonconstant
harmonics vanish whenever both labels are equal.

It is essential to compare mixed labels by the
top-harmonic Schwarz inequality.  Real polarization
between a scalar and a pseudoscalar would lose their
purely imaginary, reflection-odd correlation.
If $\bar m(2)>0$, choose a line separating two
arbitrary points $z,y'$ with labels $\lambda,\lambda'$.
Set $S=\{(\theta_Lz,\lambda)\}$ and
$T=\{(y',\lambda')\}$ on its same side.  Then
\[
 |t_1((z,\lambda),(y',\lambda'))|^2
 \le|t_1((z,\lambda),(\theta_Lz,\lambda))|
       |t_1((\theta_Ly',\lambda'),(y',\lambda'))|=0.
\]
This contradicts the definition of the top harmonic.
Hence $\bar m(2)=0$ and therefore $\bar m(n)=0$
for every $n$.

\emph{Separated tests and the group.}
The zero harmonic proves planar rotation invariance
first for plates with disjoint planar projections.
Convolution by product mollifiers and finite Riemann
sums of their translates extend it to all $C^1$
entries with those projections; (R2) controls the
moving-entry limits.  To remove the restriction,
partition in $x_2$ into intervals much shorter than
the separation.  Every intersecting pair of planar
projections then has an $x_2$ gap with a level outside
all interval hulls.  Translate all entries above
that level by $se_1$.  The correlation before and
after rotation is holomorphic on a common strip in
$s$, by the $e_2$ reflection picture and the cone
bound.  For large real $s$ the number of intersecting
pairs decreases.  Induction and the identity theorem
give equality at $s=0$.  This is the explicit
separated-test argument of
\cite[Prop. 95.4]{YMM1}; it does not require
positivity of a joint classical law.

We have obtained every rotation about $e_2$.
Conjugation with $W_3$ gives rotations about each
axis; these generate $\mathrm{SO}(3)$.  Together
with $-1\in W_3$ and the sign character this gives
the asserted $\mathrm O(3)$ covariance.  Proper
isometries transport coordinate reflection positivity
to every plane.  Spatial rotations act unitarily
in the $e_3$ picture and rotate $(P_1,P_2)$.
Their translates of $H\ge|P_1|$ imply
$E\ge\sup_\phi|q_1\cos\phi+q_2\sin\phi|=|q|$
on the joint spectrum.
\end{proof}

\subsection{Midpoint insertions and shrinking plates on the lattice}

We give the quantitative step needed for (R5) and (R7).
Its analytic inputs are the localized source energies
and the subtracted partition bound proved in the
source section.  In the notation used here they are:
\begin{enumerate}
\item $Z[s]/Z[0]=e^{\mathcal E(s)}
 Z_K^{\mathrm{sub}}[s]/Z_K^{\mathrm{sub}}[0]$,
$\mathcal E(s)=\sum_{j,X}e_j(X;s)$, with each
term holomorphic for the source radii $2\varrho$
and $2\varrho_J$ and vanishing at zero source; the summed energy
has the conjugation identity of \cref{C-407-23};
\item at physical cube side $\mathscr W_j=ML^ja_K$,
write $\hat d_j(X)=d_j(X)$ for the polymer tree size;
then the local energy majorants obey
\[
 |e_j(X;s)|\le\varepsilon_j(X),\qquad
 \sup_\square\sum_{X\supset\square}
       \varepsilon_j(X)e^{\kappa_e\hat d_j(X)}
                                      \le C_{\mathrm{loc},E};
\]
\item changing the source only at distance $D$
from $X$ changes this term by at most
$2\varepsilon_j(X)e^{-\kappa_eD/\mathscr W_j}$;
\item uniformly in the source domain,
\[
 \frac{|Z_K^{\mathrm{sub}}[s]|}{Z_K^{\mathrm{sub}}[0]}
 \le\exp(C_0\ell^3),\quad
 C_0=C+2+C_P(1+\log[L(g^{-2}+B_\sharp)]).
\]
\end{enumerate}
These estimates follow from \cref{C-407-22,C-407-23,C-396-12}.  The arguments below do
not assume that a fermion bilinear has a positive
classical law.

Fix a board with cells
$\square_k=o+lk+[-l/2,l/2]^3$, $l=L^{-i}$ and
$o\in a_K\mathbb Z^3$, whose walls are lattice
link mirrors.  On a compatible torus let
$\mathsf N=(\ell/l)^3$.  For an even element $D$
of one cell, let $\operatorname{Fold}D$ be its
product of images in all cells, using the
antilinear reflection for orientation-reversing
images, and put
\[
 \|D\|_{\mathrm{ch},\ell}
       =\omega_\ell(\operatorname{Fold}D)^{1/\mathsf N}.
\]
This root is nonnegative.  The relations of the
reflection group hold on even elements: parallel
reflections compose to a translation, orthogonal
reflections differ on a generator by minus one
and hence agree on even products, and a period
acts trivially on even products despite the
antiperiodic generator sign.

Here is the algebraic chessboard argument.  In
one direction view the torus as a cyclic list
of $2N$ slabs.  Reflection positivity at a wall
and its antipodal wall bounds the squared modulus of
an array expectation by the product of the
two half-reflected array expectations.  For
positive constant-array roots, maximize the absolute value of the
expectation divided by the product of these
roots over the finite set of arrays with the
given entries.  Half-reflection preserves
maximality.  A maximizing array with the
longest constant cyclic string must be
constant: reflection at an endpoint doubles
the string if its length is at most $N$,
and fills the cycle otherwise.  The maximum
is therefore one.  For a zero root perturb
each entry by a real $\epsilon$; its folded
expectation is a nonnegative polynomial
with leading coefficient one, so choose
$\epsilon\to0$ off its finitely many roots.
Repeat the argument in the three directions.
It proves
\[
 |\omega_\ell(\prod_kD_k)|
                 \le\prod_k\|D_k\|_{\mathrm{ch},\ell}.
\]
The same argument proves fine-to-coarse
domination, and invariance of the root under
the board reflections and a cell-centre
reflection.  This is the finite reflection
argument of \cite[Prop. 91.7]{YMM1},
now in the commuting even cell algebra.

For a fixed-cutoff periodic state $\omega_\nu$
obtained along sides $\ell_r\to\infty$, the
bound passes to the limit.  It also gives
the multiple-reflection inequality
\[
 |\omega_\nu(D\,\tau_{le_3}F\,\Theta_cF')|
 \le\liminf_r\|D\|_{\mathrm{ch},\ell_r}
                        \|F\|_{\nu,c}\|F'\|_{\nu,c},
\]
for even $F,F'$ above the lower wall $c$ and
$D$ in the intervening cell.  To prove it
without a bounded-operator assumption, put
$\widehat DX=D\tau_{le_3}X$ and
$D^\sharp=\Theta_{c+l/2}D$.
Translation invariance and the reflection
relations give $\widehat D^\sharp$ as its
formal adjoint for the positive reflection
form.  The alternating products yield
$s_n=\langle F,(\widehat D^\sharp\widehat D)^nF\rangle$
with $s_n^2\le s_{n-1}s_{n+1}$, hence
$s_1^n\le s_0^{n-1}s_n$.
Choose a coarser board containing the finite
dependence region of $F$, and take $n$ along
multiples of its side ratio.  Fine-to-coarse
chessboard domination gives
$s_n\le C_F^2(\liminf_r\|D\|_{\mathrm{ch},\ell_r})^{2n}$.
The constant $C_F$ is finite at fixed cutoff:
if $F=\sum_\Psi f_\Psi\Psi$, Gram's inequality
and $\|D_U^{-1}\|\le m_0^{-1}$ bound it by
$\sum_\Psi\|f_\Psi\|_\infty m_0^{-\deg\Psi/2}$.
Its $n$th root disappears as $n\to\infty$.
This proves the bound for $\widehat DF$;
one final Schwarz inequality gives the
displayed multiple-reflection estimate.
It holds along every subsequence of sides.

\begin{lemma}[A cell normalizer]\label{C-389-17}
Let $p=1$ or $2$, and let $f_q$ have disjoint
supports in the middle half of one cell,
with indices $\alpha_q$.  Put
\[
 D_z=\exp\sum_{q=1}^pz_q\Phi_{\alpha_q}(f_q),
 \qquad r_*=\min_q2\varrho_{\alpha_q}/\|f_q\|_{\alpha_q}.
\]
Along a subsequence of sides there is a
holomorphic scalar $\mathcal E^\circ(z)$ on
$|z_q|<r_*$ such that
\[
 |e^{-\mathcal E^\circ(z)}|
 \liminf_r\|D_z\|_{\mathrm{ch},\ell_r}\le Q_{\mathrm{mid}}(l),
 \quad Q_{\mathrm{mid}}(l)=
       e^{2C_{\mathrm{loc},E}c_{\mathrm{geo}}+C_0l^3},
\]
where
\[
 c_{\mathrm{geo}}=(1-L^{-3})^{-1}+
 \sup_{1<u_0\le L}\sum_{n\ge0}(u_0L^n)^3
       e^{-\kappa_e[u_0L^n/4-2]_+}<\infty.
\]
In particular $\bar Q_{\mathrm{mid}}=Q_{\mathrm{mid}}(1)$
is uniform for $0<l\le1$.
\end{lemma}

\begin{proof}
Reflect the cell source into each board cell,
conjugating $z$ on orientation-reversing
copies.  Their sum $s_z^{\mathrm{ch}}$ is
inside the source domain for $|z_q|<r_*$:
disjoint middle supports preserve the maximum
supremum and Lipschitz bounds, including for
the matrix indices of operator norm at most
one.  The reflection-real basis gives
\[
 \|D_z\|_{\mathrm{ch},\ell}^{\mathsf N}
                         =Z[s_z^{\mathrm{ch}}]/Z[0].
\]
Give each polymer a root cell by the centre
of its first cube in a fixed ordering.
Let $s_z^k$ be the source in root cell $k$
alone and define
\[
 \mathcal E_\ell^\circ(z)=\mathsf N^{-1}
 \sum_k\sum_{j,X:\,k(X)=k}\widetilde e_j(X;s_z^k),
\]
where $\widetilde e=e$ for a proper cell
image and $\widetilde e=\bar e$ for an
improper one.  In the latter case the
argument has $\bar z$, so the conjugated
composition is holomorphic in $z$.
The summable local-energy majorants bound
this function uniformly in $\ell$ at fixed
$K$, and
\[
 \operatorname{Re}\mathcal E(s_z^{\mathrm{ch}})
 =\mathsf N\operatorname{Re}\mathcal E_\ell^\circ(z)
 +\sum_{j,X}\operatorname{Re}
       [e_j(X;s_z^{\mathrm{ch}})-e_j(X;s_z^{k(X)})].
\]
For $\mathscr W_j\ge l$ the last differences
sum to at most
$2\mathsf NC_{\mathrm{loc},E}(l/\mathscr W_j)^3$.
For $\mathscr W_j<l$, a different cell's
middle is at distance at least $l/4$ from
the root cell.  Locality and the polymer
diameter bound give, after the size weight
is paid, a contribution per root cell at
most
$2C_{\mathrm{loc},E}(l/\mathscr W_j)^3
e^{-\kappa_e[l/(4\mathscr W_j)-2]_+}$.
The scales form a geometric sequence, so
their total divided by $\mathsf N$ is
at most $2C_{\mathrm{loc},E}c_{\mathrm{geo}}$.
The subtracted partition bound adds $C_0l^3$.
This proves the asserted bound for
$\mathcal E_\ell^\circ$.
Montel's theorem at fixed cutoff gives a
locally uniform subsequential holomorphic
limit, with all derivatives.  Taking the
liminf of the bound proves the lemma.
The subsequence may depend on the insertion;
the limiting state itself is fixed.
\end{proof}

\begin{proposition}[Centered midpoint insertion]\label{C-389-18}
For real $f$ in the middle of a board cell
of side $l\le1$,
\[
 |\omega_\nu([\Phi_\alpha(f)-\omega_\nu\Phi_\alpha(f)]
                   \tau_{le_3}F\,\Theta_cF')|
 \le\frac{\bar Q_{\mathrm{mid}}}{\varrho_\alpha}
        \|f\|_\alpha\|F\|_{\nu,c}\|F'\|_{\nu,c}.
\]
For a complex $f$, twice this constant is valid.
\end{proposition}

\begin{proof}
Take $p=1$ in the normalizer lemma and apply
multiple reflections along its subsequence.
The holomorphic function
$\omega_\nu(D_z\tau_{le_3}F\Theta_cF')
e^{-\mathcal E^\circ(z)}$ has modulus at
most $\bar Q_{\mathrm{mid}}\|F\|\|F'\|$.
Cauchy's estimate on radius
$r_*=2\varrho_\alpha/\|f\|_\alpha$
bounds its derivative at zero.  With
$F=F'=1$ it also bounds
$|\omega_\nu\Phi_\alpha(f)-
\mathcal E^{\circ\prime}(0)|$
by $\bar Q_{\mathrm{mid}}/r_*$.
The coefficient of $\mathcal E^{\circ\prime}(0)$
is bounded by $\|F\|\|F'\|$ using the
same multiple-reflection estimate with
$D=1$.  Adding the two derivatives gives
$2\bar Q_{\mathrm{mid}}/r_*=
\bar Q_{\mathrm{mid}}\|f\|_\alpha/\varrho_\alpha$.
Split a complex function into real and
imaginary parts to obtain its stated bound.
\end{proof}

\begin{proposition}[Shrinking plates]\label{C-389-19}
Let $0<t\le1$ and let real $f$ be supported
in $|x_1|\le t$, $|x_2-y_2|\le\Upsilon/2$,
$t\le x_3\le2t$.  Let $l$ be the least
power of $L$ at least
$4\max\{2t,\Upsilon/2\}+4a_K$.
Then
\[
 |S_2^T[\nu]((\alpha,f\circ\theta),(\alpha,f))|
 \le\frac{\|f\|_\alpha^2}{\varrho_\alpha^2}
 \left[\log2\,(Q_{\mathrm{mid}}(l)+1)^2+
 \frac{C_{\mathrm{loc},E}C_{\mathrm{dep}}
                         (1+\Upsilon+m')}{4t}\right],
\]
where $M=L^{m'}$ and $C_{\mathrm{dep}}$ depends
only on $L,\kappa_e$.  The side $l$ is at
most the least power $\bar l(\Upsilon)$ of
$L$ exceeding $4\max\{2,\Upsilon/2\}+4$.
\end{proposition}

\begin{proof}
Choose a nearest lattice board centre to
$(0,y_2,0)$; both reflected supports lie
in the middle of its cell.  Use the
normalizer lemma with $p=2$ and put
$\Psi(z)=\omega_\nu(D_z)e^{-\mathcal E^\circ(z)}$.
Then $\Psi(0)=1$ and $|\Psi|\le Q_{\mathrm{mid}}(l)$.
Schwarz's lemma on complex lines makes
$|\Psi-1|\le1/2$ on the bidisc of radius
$r_*/[2(Q_{\mathrm{mid}}+1)]$.
There $|\log\Psi|\le\log2$, so
\[
 |\partial_{12}\log\Psi(0)|
 \le4\log2\,(Q_{\mathrm{mid}}+1)^2/r_*^2.
\]
The second cumulant is this derivative
plus $\partial_{12}\mathcal E^\circ(0)$.

For an energy term of scale $\mathscr W_j$,
let $D_1,D_2$ be its distances from the
two plate supports.  Locality and Cauchy's
estimate give its derivative a bound
$2r_*^{-2}\varepsilon_j(X)
e^{-\kappa_e\max(D_1,D_2)/\mathscr W_j}$.
The factor two is included in $C_{\rm dep}$ below.
The support separation gives
$D_1+D_2\ge2t-(\hat d_j(X)+4)\mathscr W_j$,
and
$\max(D_1,D_2)\ge D_2/2+(D_1+D_2)/4$.
For a cube $\square\subset X$ also
$D_2\ge\operatorname{dist}(\square,
\operatorname{supp}f)-(\hat d_j(X)+4)\mathscr W_j$.
The source-energy size weight pays the
resulting $3\kappa_e\hat d_j(X)/4$; the fixed factor
$e^{3\kappa_e}$ is absorbed in $c(\kappa_e)$ below.
Counting cubes near the plate yields
\[
 |\partial_{12}\mathcal E_\ell^\circ(0)|
 \le\frac{C_{\mathrm{loc},E}c(\kappa_e)}{r_*^2}
 \sum_{j\le K}(2u_j+3)
       ((\Upsilon/t)u_j+3)(u_j+3)
                         e^{-\kappa_eu_j/2},
 \qquad u_j=t/\mathscr W_j.
\]
For $u_j\ge1$ the geometric sequence is
summable after its cubic polynomial.
For $u_j<1$ its terms are at most
$20(\Upsilon/\mathscr W_j+3)$.
Their first sum is at most
$\Upsilon/[t(1-L^{-1})]$, and their number
is at most $\log_L(M/t)+1\le(m'+2)/t$.
Thus the last sum is bounded by
$C_{\mathrm{dep}}(1+\Upsilon+m')/t$.
Pass to the normalizer subsequence limit,
use $r_*=2\varrho_\alpha/\|f\|_\alpha$,
and add the logarithm derivative bound.
The side bound follows directly from
$t\le1$ and $a_K\le1$; it permits boards
larger than one for large fixed width.
\end{proof}

\subsection{Checking the hypotheses on a volume-first limit}

For the chamber construction, the earlier compactness
argument gives (R0)--(R3), translation continuity and
the sign rule.  The exact diagonal lattice reflection
inequality gives (R6) on the same subsequence.  The
independent two-plane argument gives (R4).  The
remaining quantitative inputs are the midpoint
insertion estimate and the shrinking-plate estimate;
they are needed here with their full lattice
uniformities in the cutoff and board size.

For the real midpoint insertion estimate we may use
the weaker constant
$2\bar Q_{\mathrm{mid}}/\varrho_\alpha$,
$\varrho_0=\varrho$ and
$\varrho_\Gamma=\varrho_J$.
It implies (R5) with
\[
 c_{\mathrm I}=1/(40L),\qquad p_{\mathrm I}=0,
 \qquad C_{\mathrm I}=
            2\bar Q_{\mathrm{mid}}/\min\{\varrho,\varrho_J\}.
\]
Indeed choose $l=L^{-i}$ with
$\eta/(5L)<l\le\eta/5$, and a nearest board centre
$y^*\in L^{-(i+2)}\mathbb Z^3$.  Its lower wall is
$c=y_3^*-l/2$.  The support of $h$ lies in the
middle cube since
$c_{\mathrm I}\eta+l/(2L^2)<l/4$, and
\[
 -0.226\eta\le c<c+l\le0.226\eta<\eta/4.
\]
Write $F'=\tau_{2ce_3}F$ and
$G''=\tau_{-le_3}G$.  The lattice midpoint
estimate applies to
$\theta_c\bar F'\cdot(\alpha,h)\cdot\tau_{le_3}G''$.
Pass to the fixed-subsequence limit.  Its two
reflection norms become
$\|[\tau_{ce_3}F]\|$ and
$\|[\tau_{-(c+l)e_3}G]\|$.
Both are bounded by the norms in (R5), using
the positive-time contractions with times
$c+\eta/4$ and $\eta/4-c-l$.
The midpoint constant is uniform for $l\le1$,
so no power of $\eta^{-1}$ is lost.  We use the
factor two in the lattice insertion bound.

For a label $\lambda=(\alpha,\upsilon)$ and
$a\in\mathsf{Prof}_t$, the plate function
$f=a\otimes\upsilon$ obeys
\[
 \|f\|_\infty\le3\|\upsilon\|_\infty,
 \quad \operatorname{Lip}f\le15\|\upsilon\|_{C^1}/t,
 \quad \|f\|_\alpha\le600M\|\upsilon\|_{C^1}/t.
\]
The shrinking-plate estimate on the lattices,
uniform at a board side bounded in terms of
$\Upsilon_\lambda$ alone, then gives
$\sigma(\theta(\alpha,f)\cdot(\alpha,f))
\le C_\lambda t^{-3}$ in the limit.
This is (R7) with $a_\Gamma=3<4$.
Thus (R0)--(R7) hold and the rotation criterion applies.
No estimate at coincident supports
is used in this passage.

\subsection{Point bounds and a common field domain}

\begin{proposition}[An $L^1$ point bound]\label{C-389-33}
If (R0)--(R7) hold with $p_{\mathrm I}=0$, then there
is $C_{\mathrm{pt}}<\infty$ such that, for
$0<\eta\le1$, real $h\in C_c^1(\{|x_3|<\eta/2\})$
and combinations $F,F'$ supported in $x_3>\eta$,
\[
 |\sigma(\theta\bar F\cdot(\alpha,h)\cdot F')|
 \le16C_{\mathrm{pt}}\eta^{-4}\|h\|_{L^1}
 \|[\tau_{-\eta e_3}F]\|
 \|[\tau_{-\eta e_3}F']\|.
\]
The same constant works for the finite set of indices.
\end{proposition}

\begin{proof}
The point-field construction in three dimensions in
\cite[Prop. 97.2]{YMM1} gives a norm-holomorphic
family $B_\alpha(z,z')$ with the shift and adjoint laws
\[
 T(\kappa)B_\alpha(z,z')T(\kappa')
                   =B_\alpha(\kappa+z,z'+\kappa'),
 \quad B_\alpha(z,z')^*=B_\alpha(z'^*,z^*),
\]
and
$\|B_\alpha(z,z')\|\le C_{\mathrm{pt}}
\min\{1,d_3(z),d_3(z')\}^{-4}$.
Here $z^*=(-\bar z_1,-\bar z_2,\bar z_3)$.
Its proof uses a three-dimensional mollifier and a
mean-value formula on a polydisc of radius
$r=c_{\mathrm I}\eta/8$.  The smeared profile has
$\|\gamma\|_\sharp\le40Mm_3r^{-7}$, where $m_3$ is a fixed
constant of the chosen mollifier, and the
mean-value volume costs $(2r)^3$; (R5) therefore
costs $C\eta^{-4}$.  Symmetric convolution and
Cauchy's estimate make these mollified operators
converge in norm at positive real times.  The shift
law gives their unique continuation and the slab
identity
\[
 B_d^{\mathrm c}(\alpha,h)=
       \int h(x)B_\alpha(de_3+x,de_3-x)\,dx.
\]
Take $d=\eta$. For $|x_3|<\eta/2$ the two real tube distances are
at least $\eta/2$.  Integrate the operator norm
against $|h|$ and apply it to the two shifted
vectors.  The factor $2^4$ proves the estimate.
All steps are linear in the label, and the finite
index set has a common insertion constant.
\end{proof}

\begin{theorem}[Finite-family reconstruction]\label{C-389-34}
Let $\sigma$ satisfy the preceding hypotheses with
$p_{\mathrm I}=0$.  Its reflection Hilbert space
carries Hermitian operator-valued tempered
distributions $\varphi_\alpha$, a strongly continuous
unitary representation of the proper orthochronous
Poincar\'e group, and a common dense invariant
polynomial domain $\mathcal D_0$ containing $\Omega$.
The vacuum is cyclic and invariant, translations
have spectrum in the closed forward cone, the
fields are scalar under the connected Lorentz
group, and
\[
 [\varphi_\alpha(f),\varphi_\beta(g)]=0
 \quad\hbox{on }\mathcal D_0
\]
for spacelike separated test supports, for every
pair of indices.  Their Euclidean continuations
agree with $\sigma$ on all separated $C_c^1$ tuples.
No uniqueness of the translation-invariant vector
is asserted.
\end{theorem}

\begin{proof}
We use the constructive general theorem
\cite[Thm. AO]{YMM1} in its (Rec) form.
Its exact four hypotheses are full Euclidean
covariance, the separable reflection picture with
strongly continuous translations, the circular
spectral cone, and an $L^1$ insertion estimate on
$|x_3|<c_0\eta$ between positive-time vectors,
with polynomial factor $C_\wp\eta^{-\wp}$.
We have verified them with
$(C_\wp,\wp,c_0)=(16C_{\mathrm{pt}},4,1/2)$.
The real scalar assumption in that theorem is
used for the slab adjoint.  Here exactly the same
adjoint identity holds separately for every
reflection-real index; a sign under an improper
map is not a sign in this identity.

More explicitly, let
$\mathcal D_\infty=\bigcap_{j\ge0}\operatorname{dom}H^j$
and fix an integer $l>4$.  The point bound and its
shift law give a point form $q_\alpha$ whose
regularization
\[
 A_{\alpha,l}=(1+H)^{-l}q_\alpha(1+H)^{-l}
\]
is bounded.  This follows by inserting the Laplace
integral for each $(1+H)^{-l}$; near zero its norm
integrand is bounded by
$s^{l-1}s'^{l-1}\min(s,s')^{-4}$, which is
integrable for $l>4$.
The Minkowski field form is
\[
 q_{\alpha,f}(\psi,\psi')=
 \int f(x)q_\alpha(U(-x)\psi,U(-x)\psi')\,dx.
\]
It is bounded by
$\|A_{\alpha,l}\|\|f\|_1
\|(1+H)^l\psi\|\|(1+H)^l\psi'\|$.
Differentiate its translated matrix element and
integrate by parts in time to move energy powers
from the first vector to the second.  Riesz
representation and the spectral domain criterion
then give
\[
 \|(1+H)^j\varphi_\alpha(f)\psi\|
 \le\|A_{\alpha,l}\|
 \sum_{k=0}^{l+j}\binom{l+j}{k}
       \|\partial_t^kf\|_1
                         \|(1+H)^{2l+j}\psi\|.
\]
This proves a common invariant domain
$\mathcal D_\infty$ and tempered matrix elements
for every index, independently of any factorial
bound in the number of fields.  Their adjoint
relations give
$\varphi_\alpha(f)^*\supset\varphi_\alpha(\bar f)$.

The tube kernels with their labels are the same
holomorphic functions used above.  Their tempered
boundary values define all mixed vacuum
distributions.  Proper rotations continue to
boosts without a sign; the labelled analytic
continuation therefore constructs their common
unitary Lorentz representation.  Simultaneous
permutation symmetry of points and indices makes
the two boundary values agree at spacelike
separation, proving mixed locality by the same
tilted-tube argument as in the general theorem.
The span of all finite labelled field products
on $\Omega$ gives $\mathcal D_0$; it is invariant
and dense by the Euclidean recovery and the
density of separated positive-time tuples.
The construction identifies its Euclidean values
with $\sigma$.  Vacuum uniqueness in the general
theorem needs $\ker H=\mathbb C\Omega$, an
additional assumption not used here.
\end{proof}

If the finite index space also carries the flavour action
of the construction, define it on bilinear tuple classes by
\[
 \mathsf V(V)[((\Gamma_i,f_i))_i]
              =[((V\Gamma_iV^{-1},f_i))_i].
\]
Action-density entries are fixed, and the formula extends
multilinearly in the indices. Its matrices in the
reflection-real basis are real orthogonal and it
commutes with the semigroup, translations and proper
rotations.  It consequently commutes with the
Poincar\'e representation and carries
$\varphi_\Gamma$ to $\varphi_{V\Gamma V^{-1}}$,
fixing $\varphi_0$.  For the reflection $p$ of the
first spatial coordinate,
$\mathsf P[X]=\pi(X)[pX]$ is unitary and
\[
 \mathsf P\varphi_\alpha(f)\mathsf P^{-1}
          =\pi_\alpha\varphi_\alpha(f\circ p).
\]
It conjugates the Poincar\'e action by $p$ and
exchanges the two flavour factors.  These assertions
follow already for slab insertions and pass through
the same point and boundary limits.

The reconstruction fixes full tempered Minkowski
vacuum distributions. Its Euclidean kernels agree with
the limiting correlations at distinct points. The extension
to Schwartz tests flat on every coincidence diagonal uses
\cite[Prop. 108.4 and Section 93]{YMM1}, applied to each
fixed string of labels. With $\wp=4$, its point bound is
$|s_n(x)|\le C^n(2n^2)^{4n}
\min(1,\min_{i\ne j}|x_i-x_j|)^{-4n}$.
Taylor's formula for a flat test cancels this singular factor;
Schwartz decay then makes its product with $s_n$ integrable.
The tensor-density result identifies this integral with the
Schwinger functional on that test space. This extension
asserts neither convergence of lattice contact terms
nor uniqueness of an extension across those diagonals.
The separated-moment bound with $(n!)^\varkappa$ is
compatible with that identification; it is not a
substitute for the polynomial insertion estimate used
to construct the fields.

\begin{proof}[Proof of \Cref{C-396-28}]
Choose the constants by \cref{C-396-12}, which verifies
the gauge induction for both coupling sequences, and put
$\mathsf b_0=\gamma^{-2}+\bar B$. For every
$\mathsf b\ge\mathsf b_0$, the choice
$g=(\mathsf b-\bar B)^{-1/2}\le\gamma$ places
$\mathsf b$ in the interval of the pure gauge construction;
the permitted physical sides contain all sufficiently large
$2L^m$. The determinant, inverse and infinite-volume lattice
states are supplied by \cref{C-391-16,C-389-10,C-386-13}.
The block construction, forward flow and
history bounds give the source estimates of
\cref{C-407-22,C-407-6}. These bounds are uniform in the
torus side and pass to every $\nu_K$.

Prop.~\ref{C-389-25} gives a common spacing subsequence
on every separated tuple, together with clause \textup{(i)}
and the lattice symmetries. Prop.~\ref{C-389-27}
gives coordinate and diagonal reflection positivity on
that same limit. The reflection-picture construction and
\cref{C-389-28} give $H\ge|P_1|,|P_2|$.
The source energies and common vacuum energy imply the
midpoint and plate bounds in \cref{C-389-18,C-389-19}.
Their transfer to the continuum verifies (R5) with
$p_{\mathrm I}=0$, and (R7) with exponent $3<4$;
the other properties (R0)--(R7) have just been proved.
Thus \cref{C-389-30} gives full Euclidean covariance,
all-plane reflection positivity and the circular spectral
cone. This proves clause \textup{(ii)}.

The point bound \cref{C-389-33} verifies the final
reconstruction input. The finite-family reconstruction
\cref{C-389-34} then gives the common domain, proper
Poincar\'e action, Hermiticity and pairwise locality.
The exact flavour and parity identities pass through
the same construction. Finally, the cutoff errors in
\cref{C-408-27} vanish at each fixed test pair after
the volume limit. Its positive leading short-distance
terms give a nonzero two-point function for every index.
These are clauses \textup{(iii)} and \textup{(iv)}.
\end{proof}

\appendix
\section{Index of notation and conventions}
\label{app:qcd-notation}

The entries follow the order of the construction. The last column
gives the definition or the passage immediately preceding the cited result.

\begingroup
\small
\renewcommand{\arraystretch}{1.12}
\begin{longtable}{@{}p{0.24\textwidth}@{\hspace{0.025\textwidth}}p{0.735\textwidth}@{}}
\textbf{Notation} & \textbf{Meaning and location}\\
\hline
\endfirsthead
\textbf{Notation} & \textbf{Meaning and location (continued)}\\
\hline
\endhead
\endfoot
\endlastfoot
$W_3$, $v=(x,g)$, $v\cdot s$
& The signed permutation group, a chamber at the lattice site $x$,
and its four neighbours; \eqref{C-391-6}.\\
$N_f$, $\nu=8N_f$, $\gamma_\mu$, $\gamma_4$
& The number of four-component flavours, the dimension of the
spin--internal--colour--flavour fibre at one chamber, and the spin
matrices; \cref{sec:qcd-lattice-foundation,sec:qcd-block-integration}.\\
$D_U$, $D[U]$, $m_0$
& $D_U$ includes the bare mass; $D[U]=D_U-m_0$ is massless.
The physical mass is $\mathsf m=a_K^{-1}m_0$;
\eqref{C-391-10} and \cref{sec:qcd-block-integration}.\\
$a_p(U)$, $A(U)$, $\beta$
& The plaquette energy $1-\tfrac12\operatorname{tr}U_p$, its sum
over positive plaquettes, and its coefficient $\beta=L^K\mathsf b$;
\cref{sec:qcd-lattice-foundation}.\\
$\mathfrak B$, $\Phi_{\alpha,K}$, $S^T_{n,K,\nu_K}$
& The reflection-real meson basis, the action and meson insertions,
and their joint cumulants. The index $0$ denotes the action density;
\eqref{eq:qcd-physical-observables}. The centered moments $\sigma$
are defined in \eqref{eq:qcd-centered-moments}.\\
$L$, $M$, $a_K$, $\lambda_j$, $n_j$
& The blocking factor, the polymer-block side in level-$j$ units,
the spacing $a_K=L^{-K}$, the spacing $\lambda_j=L^{j-K}$,
and the site count $n_j=\ell^3\lambda_j^{-3}$ on $T_\ell$;
\cref{sec:qcd-gauge-input}.\\
$\mathsf t_j$, $x_j$, $b_j$
& The calibration $\mathsf t_j=L^{K-j}\mathsf b$, the running
action coefficient, and its extracted increment:
$x_{j+1}=L^{-1}x_j-b_{j+1}$; \cref{sec:qcd-gauge-input}.\\
$\mathfrak m_j$, $s_j$, $\mu_j$
& The rescaled bare mass $\mathfrak m_j=\lambda_j\mathsf m$,
the accumulated correction, and its increment
$s_{j+1}=Ls_j+\mu_{j+1}$, $s_0=0$;
\cref{sec:qcd-gauge-input,C-396-33}.\\
$R(t)$, $\epsilon(t)$, $\delta(t)$, $\mathfrak b(t)$
& The range and gauge-field thresholds. The threshold for the
rescaled fluctuation is $\mathfrak b(t)=t^{1/2}\delta(t)$;
\cref{sec:qcd-gauge-input}.\\
$\phi(t)$, $\widehat\phi(t)$, $\phi_\sharp(t)$, $\psi_\sharp(t)$
& The functions controlling logarithmic losses. They are
$t^{-1/2}(1+\log t)^{\mathfrak p_0}$,
$(1+\log t)^{\mathfrak p_0}$,
$t^{-1/2}(1+\log t)^{3\mathfrak p_0}$, and
$t^{-1/2}(1+\log t)^{9\mathfrak p_0}$, respectively;
\cref{sec:qcd-gauge-input}.\\
$\mathcal M$, $U_j(V)$
& The gauge block average and the critical gauge orbit on its
fibre $\mathcal M^jU=V$; \cref{sec:qcd-gauge-input}.\\
$(\mathbf U,\mathbf J)$, $(U,A',\mathbf J)$,
$\mathcal U_j^\circ(X;b_0,b_1)$
& A complex pair, a representative with
$\mathbf U=e^{iL^{-j}A'}U$ and $U$ real, and its analytic core.
The core includes bounds on the required block towers;
\cref{sec:qcd-gauge-input}.\\
$\mathcal P_j$, $\mathcal D_j$, $d_j(X)$, $X^\sharp$
& The partition into $M$-blocks, its nonempty wall-connected
polymers, their block-tree length, and the prescribed enlargement
used in localization; \cref{sec:qcd-gauge-input}.\\
$\iota$, $\Pi_0$, $\Pi_1$
& The normalized embedding of a site spinor into chamber constants,
the projection onto these constants, and its complement;
\cref{sec:qcd-block-integration}.\\
$Q(V)$, $R_k$, $q_{j+1}$
& The normalized site average, its iterated chamber average
$R_k=Q(\mathcal M^{k-1}U)\cdots Q(U)\iota^*$, and
$q_1=L^{-1/2}Q(U)\iota^*$,
$q_{j+1}=L^{-1/2}Q(\mathcal M^jU)$ for $j\ge1$;
\cref{C-387-24,C-387-30}.\\
$\theta=b_F$, $\theta_k$, $\mathfrak a_k$
& The block penalty and the coefficients obtained by composing
its Gaussian kernels. Here
$\theta_k=b_F(1-L^{-1})/(1-L^{-k})$ and
$\mathfrak a_k=L^{-k}\theta_k$; \cref{C-387-30}.\\
$\mathbb D_j$, $\Gamma_j$, $\Psi^{\rm bg}$
& The quadratic form after $j$ integrations, the form on the
variables integrated at the next step, and the background
substitution obtained by completing the square.
The fluctuation covariance is $\Gamma_j^{-1}$;
the definitions precede \cref{C-397-12}.\\
$\mathcal K_i$, $\mathsf W_i$
& The operator on the original chamber variables with all the
level-$i$ block penalties added, and the sum of those penalties;
\cref{sec:qcd-regular-operators}.\\
$\mathsf L_i$, $\mathsf s_i$, $d_{h,i}$
& The layer scale, its Lipschitz minorant, and the path distance
with bond length $1/\max(\mathsf s_i(x),\mathsf s_i(x'))$;
\cref{sec:qcd-regular-operators}.\\
$\mathcal H_{(n)}$, $S_j^{\rm sym}$
& The background map from the stored and current variables to
level $n$, and the symmetrized localized covariance;
\cref{C-397-40,C-397-33} and
\eqref{eq:qcd-symmetric-covariance}.\\
$\|F\|_h$, $h_F$, $\Gamma_F$
& The Grassmann fibre norm with a weight for each generator,
the fixed fermion weight, and the covariance allowance used
in Gaussian contraction; \cref{C-387-18,C-397-40} and the
definitions preceding \cref{C-388-14}.\\
$\|F\|_{k;a}$, $h^{(k)}(n)$
& The level-$k$ family norm at pair radius $a$, with weight
$e^{\kappa d_k(X)}$ in the polymer sum, and the leg weight
$C_HL^{-(k-n)}h_F$; the definitions preceding \cref{C-388-14}.\\
$\|\mathbb V\|_{\star,r}$, $\|V\|_r$
& The origin norm allowing quadratic fluctuation growth, and
the norm that sums local potentials over polymers containing
one specified cube; \eqref{eq:qcd-origin-norm} and
\cref{C-388-19}.\\
$h$, $\mathcal S(h)$, $\Omega_j$, $\Lambda_j$, $Z_j$
& A history, its discrete label, its nested regions, its integration
regions, and $Z_j=\Lambda_j^c$. The label is separate from
the free coordinates at which that history is evaluated;
\cref{sec:qcd-histories}.\\
$\mathcal L_j$, $\mathfrak P_h$, $\varphi_j$
& The layers $\Omega_j\setminus\Omega_{j+1}$, with last layer
$\Omega_k$; their fermion places; and the fixed history cutoffs.
The fermions on layers $j<k$ are stored;
\cref{sec:qcd-histories}.\\
$s=(w,J)$, $\mathsf n$, $\mathfrak c_k(s)$
& The physical action and matrix sources, the bare matrix
$\mathsf n=-a_KJ$, and the step coordinates
$(w,\mathsf n,(\sigma_i)_{i\le k})$;
\cref{sec:qcd-source-chain,C-407-13}.\\
$\|w\|_{\sharp,j}$, $\|w\|_\sharp$, $\pi_j$
& The norm
$\max(\|w\|_\infty,40M\lambda_j\pi_j\operatorname{Lip}w)$,
its final-level value, and
$\pi_j=(1+(K-j)/n_\pi)^{\mathfrak p_0}$;
\cref{sec:qcd-source-chain}.\\
$\mathscr O_j$, $\mathscr D_j$, $\mathscr D_j(a)$
& The physical source domain, the domain of independent step
coordinates, and the latter domain at pair radius $a$.
A doubled domain doubles each radius;
\cref{sec:qcd-source-chain}.\\
$\sigma_k$, $\nu_{k+1}$, $\mathcal N_k$, $\mathcal Z_k$
& The source correction and its extracted increment,
$\mathcal N_k=-\lambda_kJ+\sigma_k$, and the response
$\mathcal Z_k=1-L^{K-k}\partial_J\sigma_k|_0$;
the definitions preceding \cref{C-407-13}.\\
$E_k(s)$, $Z_k^{\rm sub}[s]$
& The extracted vacuum energy and the remaining partition
function, with $Z[s]=e^{E_k(s)}Z_k^{\rm sub}[s]$;
\cref{C-407-22}.\\
\end{longtable}
\endgroup

\subsection*{Conventions}

The barred and unbarred fermions are independent Grassmann
generators. We normalize Berezin integration by
$\mathrm{Ber}\,e^{-\bar\Psi A\Psi}=\det A$.
At complex gauge fields, $q^+$ denotes the holomorphic continuation
of the adjoint at real fields. The barred background substitution
is the row twin given before \cref{C-397-12}; it is not obtained
by taking the adjoint of the unbarred substitution.

Operator formulas use lattice coordinates unless a physical
scale is displayed. The physical torus has side $\ell$.
Gauge links are periodic, and spinors are antiperiodic in each
coordinate. In the volume-first construction, one first takes
torus limits at fixed $K$ and then takes a subsequence as
$K\to\infty$; see \cref{sec:qcd-continuum}.

\subsection*{Terms used in the history expansion}

\emph{Made and formed.}
A made boundary entry is a term generated at an earlier step
and retained near a boundary. A formed factor is one of the
retained factors of the calibration construction. Their
scalar contributions are listed separately in
\cref{C-399-46}; the corresponding historical families are
those of \cite[Prop.~58.26]{YMM1}.

\emph{Tie.}
This is agreement at real triples representing the same pair,
with the same stored data. Its preservation under substitution
is proved in \cite[Prop.~58.18]{YMM1}. For a term with
external fermion legs, the coefficient and the legs must be
transformed together, as in \cref{C-399-48}.

\emph{Tier.}
This is the prescribed range of a stored external field. Its
bound depends on whether a bond lies in the collar, on the rim,
or elsewhere; see \cite[Section~57.7]{YMM1}.
The superscript $\mathrm{rad}$ on a historical family norm
restricts the stored coordinates to these ranges, as specified
in \cite[Section~58.15]{YMM1}.

\emph{Guard.}
For a fine region $Z_\circ$, $\operatorname{Gd}(Z_\circ)$ is
the union of coarse cubes in $Z_\circ^\sharp$ meeting the
integration region. Localization parameters are fixed at one
on this set and vary outside it; see
\cite[Section~57.3]{YMM1}.

\emph{Pinned substitution and levelling function.}
A pinned substitution applies the shift and then the prescribed
local gauge transformation $u$. This transformation removes
the endpoint factors of the response and preserves the stored
transports at the specified pinnable tops. It is the levelling
function of \cite[Section~58.6]{YMM1}; its use in an earlier
term is described in \cite[Section~58.9]{YMM1}.
With fermion legs, the inverse action of the same $u$ is
included in their reading; see \cref{C-399-48}.

\begin{thebibliography}{9}
\bibitem{AbdesselamRivasseau1994}
A.~Abdesselam and V.~Rivasseau,
\emph{Trees, forests and jungles: a botanical garden for cluster expansions},
in \emph{Constructive Physics}, Lecture Notes in Physics \textbf{446},
Springer, 1995, pp.~7--36; arXiv:hep-th/9409094, Theorem~III.1.
\bibitem{Etingof}
P.~Etingof, \emph{Lie Groups and Lie Algebras}, MIT lecture notes,
\href{https://ocw.mit.edu/courses/18-755-lie-groups-and-lie-algebras-ii-spring-2024/mit18_755_s24_lec_full.pdf}{18.755, Spring 2024}, Theorem~37.5.
\bibitem{YMM1}
\emph{Scalar fields from three-dimensional lattice Yang--Mills theory},
companion manuscript, corrected edition of 26 October 2026.
The accompanying \href{../YM3/Revisions/YMM1-2026-10-26.pdf}{YMM1.pdf} is the edition cited here.
\end{thebibliography}
\end{document}
